\documentclass[11pt]{article}

\usepackage[T1]{fontenc}
\usepackage[utf8]{inputenc}
\usepackage[english]{babel}
\usepackage[margin=1in]{geometry}
\usepackage{lmodern}
\usepackage{microtype}
\usepackage{mathtools}
\usepackage{amsmath,amssymb,amsthm}
\usepackage{array}
\usepackage{booktabs}
\usepackage{tabularx}
\usepackage{longtable}
\usepackage{float}
\usepackage{graphicx}
\usepackage{pdflscape}
\usepackage{caption}
\usepackage{needspace}
\usepackage{enumitem}
\usepackage{fancyhdr}
\usepackage[numbers]{natbib}
\usepackage{doi}
\usepackage{xurl}
\usepackage{hyperref}
\usepackage{aliascnt}
\usepackage{cleveref}
\usepackage{orcidlink}
\usepackage{titling}

% Foundations IV — shared manuscript macros.
% Stable, cross-section macros. Add macros incrementally as new
% vocabulary stabilizes; never redefine locally inside a section.
% Governed by the shared notation and terminology discipline.

%--- axis labels (the 8 primitive-necessity axes of the F-IV catalog) -------
\newcommand{\axisAccess}{Access}
\newcommand{\axisClosure}{Closure}
\newcommand{\axisStability}{Stability}
\newcommand{\axisTransport}{Transport}
\newcommand{\axisStatus}{Status}
\newcommand{\axisRecords}{Records}
\newcommand{\axisSymmetry}{Symmetry}
\newcommand{\axisInterfaces}{Interfaces}
\newcommand{\axisSelfReference}{Self-reference}

%--- chapter groupings (the 6 drafting chapters that group the 8 axes) ------
\newcommand{\chapAccess}{Access}
\newcommand{\chapStability}{Stability}
\newcommand{\chapTransport}{Transport}
\newcommand{\chapSymmetry}{Symmetry}
\newcommand{\chapStatusRecordsCoherence}{Status, Records, Coherence}
\newcommand{\chapSelfReferenceLimits}{Self-Reference and Limits}
\newcommand{\chapClosureRealityBoundary}{Closure--Reality Boundary}

%--- catalog identity --------------------------------------------------------
\newcommand{\FIV}{Foundations~IV}
\newcommand{\FoundI}{Foundations~I}
\newcommand{\FoundII}{Foundations~II}
\newcommand{\FoundIII}{Foundations~III}
\newcommand{\NumLaws}{52}
\newcommand{\NumMechanizeNow}{46}
\newcommand{\NumAnchor}{2}
\newcommand{\NumPartial}{2}
\newcommand{\NumObligation}{2}
\newcommand{\NumAxes}{8}
\newcommand{\NumChapters}{6}

%--- quotient discipline (shared Transport / Access / Stability) ------------
\newcommand{\currQ}{\pi^{\mathrm{cur}}}  % current quotient map
\newcommand{\predM}{\pi^{\mathrm{pred}}} % predictive quotient map
\newcommand{\splitObs}{\Delta}         % split-pair obstruction
\newcommand{\quotMap}{\pi}             % quotient map
\newcommand{\accessQ}{q_I}             % access quotient
\newcommand{\obsSig}{\sigma_I}         % observation signature
\newcommand{\closureForm}{\mathsf{cl}} % formed-closure operator
\newcommand{\BirdInt}{\mathsf{BirdInt}}% interaction judgment

%--- substrate-anchor notation (Papers 1, 3 -- present in inherited Lean) ---
% Paper 1 (Adequacy Residuals) — for F6/F8 anchor paragraphs.
% Note: \Xi is a built-in Greek letter; we use it directly without
% redefinition. The adequacy residual is written $\Xi(D \mid L)$.
\newcommand{\Kdom}{K^{D}}              % dominator
\newcommand{\Kloss}{K^{L}}             % loss
\newcommand{\Schur}{\mathsf{Schur}}    % Schur identity
% Paper 3 (AOR) — for F53 anchor paragraph
\newcommand{\ClSB}{\mathrm{Cl}_{\mathrm{SB}}}  % reflective closure operator
\newcommand{\etaSB}{\eta_{\mathrm{SB}}} % AOR unit
\newcommand{\XiSB}{\Xi_{\mathrm{SB}}}   % AOR residual

%--- substrate-anchor notation (Papers 2, 6 -- substrate_pending in Lean) ----
% Paper 2 (SDTC) / Paper 6 (RH) — for F14 paragraph
\newcommand{\Jinv}{J}                  % involution
\newcommand{\Fixinv}{\mathrm{Fix}(\Jinv)}  % involution fixity set
\newcommand{\trAX}{\operatorname{tr}}  % trace functional

%--- normal-form template (uniform across all 52 laws) ----------------------
% Use these environments to wrap the catalog laws in body prose.
% Each F-IV law follows the 6-part template: Setup → Theorem → Proof
% spine → Case enumeration → BirdInt status → Nonclaims.
\newcommand{\NormalForm}[1]{\textbf{Normal form #1}}
\newcommand{\SetupBlock}{\textit{Setup.}}
\newcommand{\ProofSpine}{\textit{Proof spine.}}
\newcommand{\CaseEnum}{\textit{Cases.}}
\newcommand{\BirdIntStatus}{\textit{Interaction status.}}
\newcommand{\Nonclaims}{\textit{Nonclaims.}}

%--- status tags (display form for paper text) ------------------------------
% Body prose uses these textual descriptions; the literal Lean tag
% strings (`# obligation`, `# substrate_pending`) appear only in
% app:formalization per feedback_academic_register.md.
\newcommand{\statTheorem}{Theorem}
\newcommand{\statPartial}{Partial}
\newcommand{\statObligation}{obligation (axiom, opacity-guarded)}
\newcommand{\statAnchor}{Anchor}
\newcommand{\statSubstratePending}{substrate-pending (axiom-applied)}

%--- anchor cross-reference shorthands --------------------------------------
% These are paper-prose hooks for the substrate papers. They expand
% to descriptive title phrases, not the in-house ``Paper N'' series
% numbering. Bibtex keys live in paper/references.bib.
\newcommand{\PaperAR}{the adequacy-residuals paper}
\newcommand{\PaperSDTC}{the self-dual trace confinement paper}
\newcommand{\PaperAOR}{the audited-operational-realisability paper}
\newcommand{\PaperNeedleKiller}{the needle-killer paper}
\newcommand{\PaperNS}{the Navier--Stokes paper}
\newcommand{\PaperRH}{the Riemann Hypothesis paper}
\newcommand{\PaperPvNP}{the P versus NP paper}
\newcommand{\PaperCSL}{the CSL paper}

%--- opacity-guard predicate names (display form in body prose) -------------
% In body text the predicate is named descriptively; the literal Lean
% predicate name appears only in app:formalization.
\newcommand{\GuardPvNP}{saturated-SAT source predicate}
\newcommand{\GuardCSL}{CSL diagnosis source predicate}
\newcommand{\GuardDuality}{duality-fixity substrate-anchor predicate}
\newcommand{\GuardSSB}{future branch-selection formation source predicate}

%--- canonical theorem-name shortcuts (for cross-references in body) --------
% These render the canonical paper-prose names so cross-references are
% mechanically enforced. Add as needed during drafting.
\newcommand{\thmDescentRepair}{Descent--Repair Normal Form}
\newcommand{\thmHolonomyMemory}{Holonomy-Memory Repair Normal Form}
\newcommand{\thmLocalGlobalObstruction}{Local--Global Obstruction Normal Form}
\newcommand{\thmRecombinationComparison}{Recombination Comparison}
\newcommand{\thmNoNeedles}{No-Needles Stability}
\newcommand{\thmSufficiencyClosure}{Sufficiency Closure}
\newcommand{\thmStrictExtension}{Strict Extension}
\newcommand{\thmNoUnstatusedResidual}{No Unstatused Residual}
\newcommand{\thmQuotientality}{Quotientality}
\newcommand{\thmAdequacy}{Adequacy / No-Overread}
\newcommand{\thmNoFreeDistinction}{No-Free-Distinction}
\newcommand{\thmHiddenness}{Hiddenness Normal Form}
\newcommand{\thmCSLFormedClosure}{CSL Formed-Closure Law}
\newcommand{\thmLayerMultiplicity}{Layer Multiplicity}
\newcommand{\thmDualityFixity}{Duality Fixity}
\newcommand{\thmSelectionObstruction}{Selection Obstruction}
\newcommand{\thmSpontaneousSSB}{Spontaneous Symmetry Breaking}
\newcommand{\thmAnomalySymmetry}{Anomaly / Symmetry Obstruction}
\newcommand{\thmDualityEquivalence}{Duality Equivalence}
\newcommand{\thmFactRecordFormation}{Fact / Record Formation}
\newcommand{\thmHorizonSufficiency}{Horizon Sufficiency}
\newcommand{\thmClosureRealityBoundary}{Closure--Reality Boundary}
\newcommand{\thmClosureCompletenessBoundary}{Closure Completeness Boundary}
% Additional thm* shortcuts are added only when body prose needs them.

% F-code of a catalog row, printed as a link to its numbered result.
\makeatletter
\newcommand{\Fref}[1]{\@ifundefined{r@thm:#1}{\@ifundefined{r@lem:#1}{\@ifundefined{r@prop:#1}{#1}{\hyperref[prop:#1]{#1}}}{\hyperref[lem:#1]{#1}}}{\hyperref[thm:#1]{#1}}}
\makeatother


% Flush pending floats (as placeins' \FloatBarrier): a page break only
% when floats are still waiting to be placed.
\makeatletter
\newcommand{\FloatBarrier}{\par\begingroup\let\@elt\relax
  \edef\@tempa{\@currlist\@botlist\@deferlist\@dbldeferlist}\def\@tempb{}%
  \ifx\@tempa\@tempb\else\clearpage\fi\endgroup}
\makeatother
\newtheorem{theorem}{Theorem}[section]
% Environments share the theorem counter through aliased counters, so that
% cleveref prints the environment's own name.
\newaliascnt{lemma}{theorem}
\newtheorem{lemma}[lemma]{Lemma}
\aliascntresetthe{lemma}
\newaliascnt{proposition}{theorem}
\newtheorem{proposition}[proposition]{Proposition}
\aliascntresetthe{proposition}
\newaliascnt{corollary}{theorem}
\newtheorem{corollary}[corollary]{Corollary}
\aliascntresetthe{corollary}
\theoremstyle{definition}
\newaliascnt{definition}{theorem}
\newtheorem{definition}[definition]{Definition}
\aliascntresetthe{definition}
\newaliascnt{remark}{theorem}
\newtheorem{remark}[remark]{Remark}
\aliascntresetthe{remark}
\floatstyle{ruled}
\newfloat{algorithm}{tbp}{loa}
\floatname{algorithm}{Algorithm}

\urlstyle{same}
\captionsetup{
  font=small,
  labelfont=bf,
  labelsep=period
}
\captionsetup[table]{skip=6pt,position=top}
\captionsetup[figure]{skip=6pt}
\captionsetup[algorithm]{skip=6pt,position=top}
\crefname{algorithm}{algorithm}{algorithms}
\Crefname{algorithm}{Algorithm}{Algorithms}
\crefname{theorem}{theorem}{theorems}
\Crefname{theorem}{Theorem}{Theorems}
\crefname{lemma}{lemma}{lemmas}
\Crefname{lemma}{Lemma}{Lemmas}
\crefname{proposition}{proposition}{propositions}
\Crefname{proposition}{Proposition}{Propositions}
\crefname{corollary}{corollary}{corollaries}
\Crefname{corollary}{Corollary}{Corollaries}
\crefname{definition}{definition}{definitions}
\Crefname{definition}{Definition}{Definitions}
\crefname{remark}{remark}{remarks}
\Crefname{remark}{Remark}{Remarks}
\setlist[enumerate]{leftmargin=*, itemsep=2pt, topsep=4pt}
\setlength{\droptitle}{-3.2em}
\pretitle{\begin{center}\LARGE\bfseries}
\posttitle{\par\end{center}\vspace{0.5em}}
\preauthor{\begin{center}\normalsize}
\postauthor{\par\end{center}\vspace{-0.4em}}
\predate{\begin{center}\small}
\postdate{\par\end{center}\vspace{0.6em}}
\setlength{\emergencystretch}{2em}
\clubpenalty=10000
\widowpenalty=10000
\displaywidowpenalty=10000
\interfootnotelinepenalty=10000
\allowdisplaybreaks[2]
\fancypagestyle{firstpage}{\fancyhf{}
  \fancyfoot[C]{\scriptsize Automorph Inc., Wilmington, DE, USA \quad \texttt{ioannis@automorph.io} \quad \textcopyright\ 2026 \quad CC-BY 4.0 \quad \doi{10.5281/zenodo.23087000}}
  \renewcommand{\headrulewidth}{0pt}
  \renewcommand{\footrulewidth}{0pt}
}

\title{Six Birds Foundations IV:\\
A Catalog of Layer-Agnostic Structural Laws}
\author{Ioannis Tsiokos\,\orcidlink{0009-0009-7659-5964}}
\date{Version 3: October 1, 2026\\[2pt] \footnotesize (Version 1: June 17, 2026)}
\hypersetup{
  pdftitle={Six Birds Foundations IV: A Catalog of Layer-Agnostic Structural Laws},
  pdfauthor={Ioannis Tsiokos},
  hidelinks
}

\begin{document}

\sloppy
\hbadness=10000

\maketitle
\thispagestyle{firstpage}

\begin{abstract}
Foundations IV states and proves 52 structural laws about layers of
description.  A layer is a set of states read through a quotient, which
fixes the distinctions the layer can make.  Each law is a theorem with
its own hypotheses, stated without reference to what the layer is made
of, so that it applies to numbers, physical systems, cells, markets,
languages or programs whenever its hypotheses hold.  Most laws are
instances of one criterion, descent: a readout on states descends
through a quotient exactly when it is constant on the quotient's
classes, and when it is not, a pair of states that the quotient merges
and the readout separates witnesses the failure; the repair either
refines the quotient or coarsens the readout.  Most laws are
elementary consequences of this criterion, and the contribution is the
catalog itself: uniform statements with explicit hypotheses, proofs and
scope.  The laws give descent-style counterparts of notions such as
access and hidden predictive structure, stability and residuals,
transport and local-to-global gluing, symmetry and its breaking,
records, objectivity, probability, causality, composability,
information loss, complementarity, entropy, limits and completeness,
and the boundary between formal closure and physical realisability.  A
few laws use other arguments: a Loewner-order inequality (No-Needles
Stability), a trace squeeze (Duality Fixity, whose domination
hypothesis is shown to hold exactly when the confined readout
vanishes almost everywhere), a connectedness argument for phase boundaries, and finite
counterexamples showing that, for abstract realisability predicates
stable under reducts, formal closure and realisability need not imply
each other in either direction.  All laws
are proved here with two qualifications: the Hiddenness Normal Form is
conditional on an explicit source hypothesis, motivated by but not
verified for the saturated-SAT construction of an earlier paper, and
one equivalence of the Closure--Reality Boundary is imported from the
paper on audited operational realisability.  Special cases and analogies
from mathematics, the sciences and engineering follow each law and are
marked as such.  The laws are formalized in Lean~4 with Mathlib; for Duality
Fixity the formalization treats an abstract ledger rather than
Hilbert-space measures.
The catalog is not exhaustive.
\end{abstract}

\medskip
\noindent\textbf{Keywords.} quotient descent; factorization; split-pair
obstruction; local-to-global gluing; structural laws; layer-agnostic
statement; Six Birds Theory; emergence calculus.

\tableofcontents

\section{Introduction}\label{sec:intro}

Let \(q:X\to Q\) be a quotient map, let \(F:X\to Y\) be a map one
wants to transport through that quotient, and let \(r:Y\to R\) be the
target identification under which the transported output is to be
read.  The elementary question is whether the composite \(rF\)
depends only on the \(q\)-class of the input.  Equivalently, does
there exist a map \(\overline F:Q\to R\) with
\[
  \overline F q = rF?
\]
If \(x\) and \(x'\) have the same image in \(Q\) but \(rF(x)\ne
rF(x')\), then no such \(\overline F\) can exist.  Conversely, when
\(q\) is surjective, the absence of such pairs is exactly the
condition for descent.  There are two canonical repairs.
One may refine the source quotient so that the offending pair is no
longer identified, or one may coarsen the target by quotienting \(R\)
by the smallest equivalence relation generated by the identifications
\[
  rF(x)\sim rF(x')
  \qquad\text{whenever}\qquad
  q(x)=q(x').
\]
The latter repair forgets precisely the distinction that the source
quotient cannot support.  The obstruction is thus not merely a
failure: it is a diagnostic object, and the minimal repairs are
determined by the same data that produced it.

\paragraph{A running example.} Take \(X=\{1,2,3,4\}\) and
\(Q=\{a,b,c\}\) with \(q(1)=q(2)=a\), \(q(3)=b\), \(q(4)=c\); take
\(Y=R=\{u,v,w\}\) with \(r\) the identity, and \(F(1)=u\),
\(F(2)=v\), \(F(3)=F(4)=w\).  The pair \((1,2)\) is split:
\(q(1)=q(2)\) but \(F(1)\ne F(2)\), so \(F\) does not descend through
\(q\).  The source repair separates \(1\) and \(2\); the target repair
identifies \(u\) and \(v\).  Each chapter opening below returns to this
example briefly.

The example uses no operator theory, no regularity hypothesis and no
category larger than the quotient construction itself.  Yet the same
shape appears in more elaborate settings.  A local family may agree on
overlaps but fail to come from a global object.  A symmetric
construction may be unable to choose a branch without adding
symmetry-breaking data.  A formally closed record may exist without
being physically realised.  A finite observation may determine all
declared future probes, or fail to, precisely because some future probe
separates states that the present quotient cannot see.  In each case
there is a setup, an obstruction, a repair or boundary, and a scope set
by the hypotheses.

\paragraph{What the paper does.}
\FIV\ is a catalog of such patterns.  It introduces no new foundation;
it states, as theorems with explicit hypotheses, \NumLaws\ structural
laws that recur across the Six Birds series, and proves them.  Each law
is presented in the same order: its setup, the theorem, the proof, a
case enumeration when the statement branches, and, where useful, a
paragraph saying what the law does not claim
(\Cref{def:normal-form}).  The uniform presentation lets a reader ask
the same questions of every entry: what are the objects, what is the
obstruction, what is the repair or classification, and what lies
outside the claim?  It is not meant to suggest that the laws are one
theorem.

Most laws are elementary.  Many are the descent criterion applied to
particular data, and several carry much of their content in their
hypotheses: the domination bounds of Duality Fixity (\Fref{F14}) exist
exactly when the readout they confine vanishes almost everywhere, and part (i) of the
Hiddenness Normal Form (\Fref{F13a}) assumes four of its five
conclusions.  The value of the catalog lies in stating each pattern
once, with its hypotheses and its scope visible, so that it can be
recognized and checked in a new setting.

A \emph{law} here is not a physical law or a principle detached from
hypotheses.  It is a named theorem stated at the level of the data in
its setup.  Some laws are quotient-factorization statements, some are
finite classifications, and some are boundary statements showing that a
desirable implication fails without stronger data.  The names of the
laws, such as Probability, Causality or Entropy, say what kind of data
the maps carry; they do not announce new theorems about probability,
causality or entropy.

Most laws apply one criterion to different data: a readout descends
through a quotient exactly when its split-pair obstruction, the set of
pairs that the quotient merges and the readout separates, is empty
(\Cref{def:split-pair}).  A few use arguments of other kinds: a
Loewner-order inequality (No-Needles Stability, \Fref{F6}), a trace
squeeze (Duality Fixity, \Fref{F14}), fiber weights (\Fref{F23}),
generating relations (\Fref{F27}, \Fref{F29}), a well-founded minimum
(\Fref{F42}), connectedness of \([0,1]\) (\Fref{F44}), and finite
counterexamples in the closure--reality no-go theorems.  Several laws,
such as Local--Global Obstruction (\Fref{F4}), also give a
classification into mutually exclusive cases.

\paragraph{Inherited apparatus.}
The paper uses three earlier papers of the series as background, in
the compact form recorded in \Cref{sec:apparatus}.  Foundations~I
supplies the language of closure formation, the completion of a
structure under declared rules \citep{Tsiokos2026EmergenceCalculus};
Foundations~II supplies admissibility and quotient packaging, the
standard act of retaining only the equivalence classes of a structure
\citep{Tsiokos2026FoundationsII}; and Foundations~III supplies a finite
audited interaction judgment, which records when a proposed interaction
is legitimately visible to the apparatus
\citep{Tsiokos2026FoundationsIII}.  Two rules of those papers enter
proofs here as explicit premises: the no-overread rule (NO) in
\Fref{F11} and \Fref{F12}, and the residual-coverage rule (RC) in
\Fref{F9}.

\paragraph{Organization.}
The laws are cited by F-codes, which are fixed identifiers and are not
in document order; the map of results in \Cref{subsec:reading-guide}
gives the printed number and page of each.  Six chapters group the laws
by the kind of question they answer.  The Access chapter
(\Cref{sec:access}) studies what an instrument can distinguish; its
first theorem says that an observation signature determines a canonical
access quotient.  The Stability chapter (\Cref{sec:stability}) studies
when declared future behavior is already controlled by present data and
when residuals must be accounted for.  The Transport chapter
(\Cref{sec:transport}) is the easiest entry point, since its laws are
closest to standard descent, gluing and quotient comparison.  The
Symmetry chapter (\Cref{sec:symmetry}) treats obstructions coming from
involutions, branch selection and anomalies.  The Status, Records and
Coherence chapter (\Cref{sec:status-records-coherence}) collects laws
in which records, facts, persistence, explanation and counterfactuals
are part of the mathematical object.  The Self-Reference and Limits
chapter (\Cref{sec:self-reference-limits}) records horizon,
information, singular-limit, locality, topology and completeness
boundaries.  A final chapter (\Cref{sec:anchor:closure-reality-boundary})
separates formal closure from bare physical realisability.
\Cref{sec:substrate-instances} maps the laws to the papers of the series
that instantiate them, \Cref{sec:sketches} lists candidates that are not
yet laws, and \Cref{sec:nonclaims} collects the scope limits of the
catalog.

\paragraph{Status of the results.}
Every law is proved here, with two qualifications.  The Hiddenness
Normal Form (\Fref{F13a}) is a conditional theorem: it assumes, as an
explicit hypothesis, surjectivity of the current quotient and four of
its five conclusions, and proves the fifth together with the
equivalence of predictive surplus with failure of predictive collapse.
The hypothesis is motivated by the saturated-SAT construction of
\citet{Tsiokos2026PvNP} but is not verified for it.  The
Closure--Reality Boundary resolution (\Fref{F53}) imports one
source-accounting equivalence from the paper on audited operational
realisability \citep{Tsiokos2026AOR}; its two no-go theorems are proved
here.  The substrate papers cited here are papers of the series by the same
author.  We say that a law is \emph{anchored} to a substrate paper when
it generalizes a theorem of that paper (\Fref{F6}, \Fref{F14}) or
imports one, explicitly identified (clause (b) of \Fref{F53}), and \emph{motivated} by one when that paper suggests its
hypotheses without establishing them (\Fref{F13a}, \Fref{F13b});
statements about what the substrate papers prove are those papers' own
claims.  \Cref{tab:substrate-instances} lists these connections.  Applications to particular
mathematical, physical or engineered systems appear after each law as
\emph{layer instantiations}, each marked as a special case or as an
analogy (\Cref{subsec:reading-guide}); no proof depends on them.  The
catalog is not exhaustive: \Cref{sec:sketches} records candidates that
remain sketches.  The results are formalized in Lean~4;
\Cref{app:formalization} describes the correspondence.

\paragraph{Reading path.}
We suggest reading this introduction and \Cref{sec:apparatus}, then
\Cref{sec:transport}, which depends only on \Cref{sec:apparatus} and
can be read before \Cref{sec:access,sec:stability}; the remaining
chapters can be read in any order.  Readers familiar with the earlier
foundations papers may skim \Cref{sec:apparatus} for notation and use
the summary table at the end of each chapter.  Readers coming from a
substrate paper should start from the corresponding row of
\Cref{tab:substrate-instances}.
            % sec:intro
\section{Apparatus and conventions}\label{sec:apparatus}

This section records the apparatus inherited from the earlier
foundations papers in ordinary mathematical language, fixes the
quotient notation used throughout, and states the template in which the
laws are presented.  The inherited results are used as background and
are not reproved.  On a first reading,
\Cref{def:quotient-packaging,def:split-pair,def:access-quotient}, the
descent criterion and the pairing refinement after
\Cref{def:split-pair}, the paragraph \emph{Two inherited rules}, and the
reading guide (\Cref{subsec:reading-guide}) suffice; the rest of the
section is reference material, to be consulted when a theorem uses it.

\FoundI\ supplies closure formation
\citep{Tsiokos2026EmergenceCalculus}; \FoundII\ supplies admissibility
and quotient packaging \citep{Tsiokos2026FoundationsII}; and
\FoundIII\ supplies the finite audited interaction judgment
\citep{Tsiokos2026FoundationsIII}.  The following definitions are the
forms of that apparatus used in this paper.

\begin{definition}[Closure formation]\label{def:closure-formation}
A \emph{closure formation} is the passage from initial data to the
object obtained by closing that data under a declared family of closure
rules.  Precisely, let \(X\) be a set and \(\mathcal R\) a set of rules
\((P,x)\) with \(P\subseteq X\) finite and \(x\in X\); a set \(A\subseteq X\)
is \(\mathcal R\)-closed when \(P\subseteq A\) implies \(x\in A\) for every
\((P,x)\in\mathcal R\). The formed-closure operation \(\closureForm\) sends
\(A\subseteq X\) to the intersection of the \(\mathcal R\)-closed supersets
of \(A\); it is extensive, monotone and idempotent. A \emph{formed closure} is
the completed object relative to the declared rules, that is, a set of the
form \(\closureForm(A)\), equivalently an \(\mathcal R\)-closed set.  A
\emph{closure-of-record} is a formed closure together with an
audit record of the closure data used to form it.
\end{definition}

Completing a substructure under an operation, taking the closure
of a relation, or closing a set of rules under composition are familiar
examples.  The \FoundI\ vocabulary does not introduce a new species of
closure; it records which rules have been declared and which completed
object is being used.

\begin{definition}[Current and predictive quotients]\label{def:quotient-packaging}
Quotient packaging is the standard quotient construction: an
equivalence relation is imposed on a carrier and only equivalence
classes are retained.  The quotient map is denoted \(\quotMap\) when no
more specific symbol is needed.  A quotient of a
carrier \(H\) is a map \(q:H\to Q\), read through its kernel \(q(h)=q(h')\); it
is a quotient map when surjective. The \emph{current quotient}
\(\currQ:H\to Q\) (the quotient presently visible to the apparatus) and the
\emph{predictive quotient} \(\predM:H\to M\) (classes predicted but not
necessarily visible at the current level) are two declared quotients of \(H\);
\(\predM\) \emph{refines} \(\currQ\) when \(\currQ=f\circ\predM\) for some \(f\),
and then each \(\currQ\)-fiber is a union of \(\predM\)-fibers.
\end{definition}

In theorem statements, \(\currQ\) and \(\predM\) denote maps; surjectivity
and refinement are assumed only when stated.  The gap between \(\currQ\)
and \(\predM\) is the basic way the paper records a distinction that is
predicted but not currently observable.  It appears in \Fref{F13a},
\Fref{F13b} and \Fref{F3}; \Fref{F7} makes the same comparison with a
declared family of future probes in place of \(\predM\).

\begin{definition}[BirdInt judgment]\label{def:birdint}
A \emph{\(\BirdInt\) judgment} is a judgment of the nineteen-component finite
audited interaction calculus \(\mathbf{BirdInt}^{\mathrm{aud}}_{\mathrm{fin}}\)
of \FoundIII\ (its Section~3.1).  It records that an interaction is
legitimately observable by declared finite audit data.  No hypothesis or
proof of this paper uses it.  For every law of the catalog, the data
such a judgment would read are the maps, predicates and pairs listed in
the law's setup; when carriers and codomains are finite, these are
finite records.
\end{definition}

\paragraph{Working vocabulary.} A few words recur in the statements of the
laws.  Most of them are interpretive: they say how a law is meant to be
read, and they enter a proof only where a statement displays them as a
hypothesis.
\begin{itemize}[topsep=2pt,itemsep=1pt]
  \item \emph{Carrier}: the underlying set on which a law's data live;
    a \emph{history space} or \emph{history set} is a carrier whose
    elements are read as histories, with no further structure assumed.
  \item \emph{Layer}: a carrier seen through a quotient or a readout, the
    level of description at which a law is read.
  \item \emph{Layer-agnostic}: stated in terms of declared structure
    and hypotheses, so that the law applies wherever those hypotheses hold.
  \item \emph{Substrate paper}: a paper supplying a result or construction
    on which a catalog law draws.
  \item \emph{Declared}: named as part of a law's data. A declared object
    has no content beyond what the law states about it; a declared
    predicate is assumed to hold only where a statement says so, and the
    paragraphs headed \emph{Reading the statement} say what each one is
    meant to express.
  \item \emph{Formed}, \emph{audited}: completed under declared rules,
    respectively carrying a record of the data used, as for a formed
    closure and a closure-of-record in \Cref{def:closure-formation}.
  \item \emph{Package}: a formed closure together with the quotients,
    readouts and records a law uses; a stability package, for instance,
    also carries an audit ledger and a membrane.  The \emph{apparatus}
    is the declared data of a package.
  \item \emph{Ledger}: a declared record attached to a carrier. An audit
    ledger lists audit entries; a fiber ledger assigns weights to the
    classes of a quotient.
  \item \emph{Status}: one of the declared outcome labels of a law, such
    as visible, hidden with a record, outside scope, blocked, priced or
    accepted; a residual or a claim is \emph{statused} when it carries one.
    \emph{Priced} means accepted within a declared budget rather than
    removed; an \emph{unpriced} residual is one that no budget accounts
    for.
  \item \emph{Closure carrier}, \emph{closure space}: the set on which a
    formed closure lives, respectively a set of closures or closure
    states; a theorem stated on a closure space uses only that set and
    the maps and predicates it displays.
  \item \emph{Membrane}: the declared boundary of a formed closure,
    across which a distinction of the layer may dissolve only with a
    recorded cost.
  \item \emph{Gate}: a declared map on a carrier whose effect on a layer is
    under test; it acts on the layer when it descends to the quotient.
  \item \emph{Current-visible}: constant on the fibers of the quotient the
    layer currently uses (the current quotient \(\currQ\), or the access
    quotient \(\accessQ\) in \Fref{F12}).
  \item \emph{Prop, Bool}: a map \(P:X\to\mathrm{Prop}\) is a predicate on
    \(X\), a subset written as a truth-valued map;
    \(\mathrm{Bool}=\{0,1\}\); in \Fref{F23} an event is a predicate on \(H\).
\end{itemize}
The words \emph{formed}, \emph{package}, \emph{apparatus} and
\emph{membrane} name the setting in which a law is read.  Where a
setting condition does mathematical work in a proof, it appears as a
displayed hypothesis (such as the formedness predicate \(F\) of
\Fref{F44}) or as one of the two inherited rules below; elsewhere these
words add no premise.

\begin{figure}[t]
\centering
\small
\setlength{\tabcolsep}{4pt}
\renewcommand{\arraystretch}{1.12}
\begin{tabularx}{\linewidth}{@{}X X X@{}}
\fbox{\begin{minipage}[t]{0.88\linewidth}
\textbf{Access}\\
Axis: access\\
\Fref{F10}--F12, \Fref{F13a}\(^{c}\), \Fref{F13b}, \Fref{F24}
\end{minipage}}
&
\fbox{\begin{minipage}[t]{0.88\linewidth}
\textbf{Stability}\\
Axis: closure/stability\\
\Fref{F6}\(^{a}\), \Fref{F7}--F9
\end{minipage}}
&
\fbox{\begin{minipage}[t]{0.88\linewidth}
\textbf{Transport}\\
Axis: transport\\
\Fref{F2}--F5\\
source of the split-pair repairs
\end{minipage}}
\\[0.8em]
\fbox{\begin{minipage}[t]{0.88\linewidth}
\textbf{Symmetry}\\
Axis: symmetry\\
\Fref{F14}\(^{a}\), \Fref{F15a}, \Fref{F15b}, \Fref{F40}--F41
\end{minipage}}
&
\fbox{\begin{minipage}[t]{0.88\linewidth}
\textbf{Status / Records / Coherence}\\
Axes: status, records, interfaces\\
\Fref{F17}--F23, \Fref{F25}--F31
\end{minipage}}
&
\fbox{\begin{minipage}[t]{0.88\linewidth}
\textbf{Self-reference / Limits}\\
Axis: self-reference\\
\Fref{F33}--F39, \Fref{F42}--F52
\end{minipage}}
\end{tabularx}

\vspace{0.75em}

\fbox{\begin{minipage}{0.94\linewidth}
\textbf{Dependencies between laws.}
The descent criterion of \Fref{F2} is used in the proof of \Fref{F5},
and its source-repair clause in the proofs of \Fref{F3} and \Fref{F15a};
the universal property of \Fref{F10} is used in the proofs of
\Fref{F11} and \Fref{F12}; the proof of \Fref{F12} applies (NO) as in
\Fref{F11}, and the pairing clause of \Fref{F12} reappears, with \(d=s\),
in \Fref{F24}; \Fref{F15a} is used in \Fref{F15b}; the proof of
\Fref{F13a} spells out \Fref{F7} for a single probe.
\end{minipage}}

\vspace{0.75em}

\fbox{\begin{minipage}{0.58\linewidth}
\textbf{Outside the axes:}
\Fref{F53}\(^{a}\), Closure--Reality Boundary, with the two
closure-equals-reality no-go theorems.
\end{minipage}}
\qquad
\begin{minipage}{0.31\linewidth}
\textbf{Legend.}\\
\(^{c}\) conditional on a source hypothesis\\
\(^{a}\) anchored to a substrate paper
\end{minipage}

\caption{The six chapters of the catalog and the eight axes they
group, with the Closure--Reality Boundary outside the axes.  The middle
box lists only the dependencies between laws used in proofs.}
\label{fig:axis-topology}
\end{figure}


\begin{definition}[Split-pair obstruction]\label{def:split-pair}
Let \(q:X\to Q\) be a source quotient, \(F:X\to Y\) a map, and
\(r:Y\to R\) a target identification.  The \emph{split-pair
obstruction} is
\[
  \splitObs(q,F,r)
  =
  \{(x,x')\in X\times X
    \mid q(x)=q(x')\ \text{and}\ rF(x)\ne rF(x')\}.
\]
It is the set of pairs identified by the source quotient but separated
after transport to the target quotient.
\end{definition}

\paragraph{The descent criterion.} For surjective \(q\), the map \(rF\)
descends through \(q\), that is,
\[
  \overline F\circ q=r\circ F
\]
for some \(\overline F:Q\to R\), precisely when \(\splitObs(q,F,r)=\emptyset\),
and then \(\overline F\) is unique.  Indeed, a factorization makes \(rF\)
constant on the fibers of \(q\); conversely, if \(rF\) is constant on
fibers, \(\overline F(q(x)):=rF(x)\) is well defined, and surjectivity
of \(q\) defines it on all of \(Q\).  Most proofs in the catalog apply
this argument to particular maps; the paragraph \emph{Instances of the
descent criterion} at the end of this section lists them.  When the
obstruction is nonempty, the source-side repair refines \(q\), while the
target-side repair coarsens \(R\) by the smallest equivalence relation
generated by the identifications \(rF(x)\sim rF(x')\) whenever
\(q(x)=q(x')\).  Surjectivity is assumed only where a statement says
so.

\paragraph{The pairing refinement.} For maps \(q:X\to Q\) and \(s:X\to S\),
the pairing \(\langle q,s\rangle:X\to Q\times S\),
\(x\mapsto(q(x),s(x))\), satisfies \(q=\operatorname{pr}_Q\circ\langle
q,s\rangle\) and \(s=\operatorname{pr}_S\circ\langle q,s\rangle\); if
\(q=a\circ p\) and \(s=b\circ p\) for a map \(p\), then \(\langle
q,s\rangle=\langle a,b\rangle\circ p\); and \(\langle q,s\rangle\) separates two
points of a \(q\)-fiber exactly when \(s\) does. So \(\langle q,s\rangle\) is
the coarsest map through which both \(q\) and \(s\) factor, and it strictly
refines \(q\) exactly when \(s\) splits a \(q\)-fiber.  A pairing is written
with its product codomain; it is a quotient map after corestriction to
its image, which \Cref{thm:F13b,thm:F24} state explicitly where it
matters.  The repairs and extensions of \Fref{F2}, \Fref{F3}, \Fref{F12},
\Fref{F13a}, \Fref{F13b}, \Fref{F15a}, \Fref{F17}, \Fref{F19}, \Fref{F20},
\Fref{F24}, \Fref{F36} and \Fref{F45} are built from this pairing.

\begin{definition}[Observation signature; access quotient]\label{def:access-quotient}
An \emph{observation signature} is a declared map
\(\obsSig:H\to \Sigma\) from a carrier \(H\) to observation symbols.
It determines the access kernel
\[
  h\sim_I h'
  \quad\Longleftrightarrow\quad
  \obsSig(h)=\obsSig(h').
\]
The \emph{access quotient} \(\accessQ\) is the quotient of \(H\) by
this kernel.  It is the quotient carrying exactly the distinctions made
by the declared observation signature.
\end{definition}

Any observable constant on the access kernel factors uniquely through
\(\accessQ\).  The Access chapter turns this universal property into
the Quotientality theorem; here it fixes the notation used by later
laws.

\begin{figure}[t]
\centering
\fbox{%
\begin{minipage}{0.88\linewidth}
\centering
\(\SetupBlock\)
\(\longrightarrow\)
\(\NormalForm{theorem}\)
\(\longrightarrow\)
\(\ProofSpine\)
\(\longrightarrow\)
\(\CaseEnum\)
\(\longrightarrow\)
\(\Nonclaims\)
\end{minipage}}
\caption{The order in which each catalog law is presented. The setup is the
opening paragraph, the \emph{Reading the statement} paragraph, or the
theorem's hypotheses. The proof carries the proof spine and any case
enumeration. A labelled nonclaims paragraph follows in the chapters
listed in \Cref{subsec:reading-guide}. Most laws are then followed by a
list headed \emph{Layer instantiations}, and some by a paragraph on what
a source paper supplies; neither is part of the template, and no proof
uses them.}
\label{fig:normal-form-template}
\end{figure}

\begin{definition}[Normal-form template]\label{def:normal-form}
Each catalog law is presented in a template of up to five parts:
\emph{setup}, \emph{theorem}, \emph{proof spine}, \emph{case
enumeration} and \emph{nonclaims}.  The setup names the data on which
the law operates.  The theorem states the layer-agnostic claim.  The
proof spine gives the structural argument.  The case enumeration lists
the mutually exclusive outcomes when the statement branches.  The
nonclaims field states what the law does not assert.  In the catalog
the setup precedes the theorem, the proof carries the proof spine and,
when the statement branches, the case enumeration, and the nonclaims
appear as a labelled paragraph after the proof in the chapters that
record them.  \Cref{fig:normal-form-template} gives the schematic form
and \Cref{fig:axis-topology} the chapters.
\end{definition}

The template is editorial rather than categorical.  It does not assert
that the laws are equivalent or that the axes are canonical; it puts
hypotheses, proof shape, case split and scope in the same order for
every law.

\paragraph{The conditional law.} One law, \Fref{F13a}, is conditional.
Besides its data, it assumes a source hypothesis \((\mathsf S)\):
surjectivity of the current quotient and the first four of its five
conclusions.  It proves the fifth conclusion and the equivalence of
predictive surplus with failure of predictive collapse.  The hypothesis
is motivated by the saturated-SAT construction of
\citet{Tsiokos2026PvNP}, for which it is not verified here.

\paragraph{Two inherited rules.} The rules (NO) and (RC) of earlier papers
enter three proofs of this catalog as premises: (NO) those of \Fref{F11} and
\Fref{F12}, (RC) that of \Fref{F9}.
\begin{itemize}[topsep=2pt,itemsep=1pt]
  \item \emph{No overread} (\FoundIII; used in \Fref{F11} and \Fref{F12}). A claim
    accepted as exact at the current layer may not depend on content
    that the instrument has suppressed; an explicit bridge supplying such content is declared apparatus data that enlarges the observation signature and hence \(\accessQ\), so in every case the claim may use only the distinctions that the current
    access makes. For a claim \(F\) and the access quotient \(\accessQ\):
    \[
      \text{(NO)}\qquad
      F \text{ accepted as exact}
      \ \Longrightarrow\
      \bigl(\accessQ(h)=\accessQ(h')\Rightarrow F(h)=F(h')\bigr).
    \]
    Here \emph{accepted as exact} is a declared predicate \(\mathrm{Acc}\) on
    claims \(H\to W\), for every codomain \(W\) (in \Fref{F12} it is applied with \(W=D\)), supplied by the audit discipline of \FoundIII\ and not
    analysed here; (NO) is the hypothesis that every \(F\) with
    \(\mathrm{Acc}(F)\) is constant on the fibers of \(\accessQ\).
  \item \emph{Residual coverage} (\FoundII; used in \Fref{F9}). Every
    native residual of a formed package has a declared status record.
    For the residual type \(R\) and the residual records
    \(\rho\in\mathrm{Rec}\) of the package, each naming its residual
    \(\operatorname{loc}(\rho)\) and carrying a typed status
    \(\operatorname{st}(\rho)\in\mathcal S_{\mathrm{typed}}\), the set of status values the package may assign, each tagged by the judgment family that assigns it (see the reading paragraph before \Cref{thm:F9}):
    \[
      \text{(RC)}\qquad
      \forall r\in R,\ \exists \rho\in\mathrm{Rec},\
      \operatorname{loc}(\rho)=r.
    \]
\end{itemize}

\paragraph{Notation.} A predicate whose name ends in
\emph{ObstructionEmpty} or \emph{DefectEmpty}, such as
\(\operatorname{SplitObstructionEmpty}\), says that the named set of
pairs or points is empty.  Split-pair obstruction sets are written
\(\splitObs\) in some laws and \(\mathcal O\) in others; both denote the
pairs that a quotient identifies and a readout or transport separates.
Letters are reused across chapters (for instance \(M\) is the target of
\(\predM\) and also a moduli set, \(F\) is usually a transport map, and
\(I\) indexes the access notation \(\obsSig,\accessQ,Q_I\) and also
denotes a hiddenness interface); each theorem fixes its letters in its
statement.  Every displayed line of a theorem is a definition (marked
by \(:=\), \(:\iff\) or ``when''), a hypothesis or a conclusion,
according to the surrounding sentence.  Every conclusion is proved,
except the accounting equivalence of clause (b) of \Cref{prop:F53},
which is imported from \citet{Tsiokos2026AOR}; some conclusions
follow by unfolding the definitions in the same statement.

\subsection{Reading guide}\label{subsec:reading-guide}

\paragraph{Chapters and axes.}
The catalog is organized along eight axes---access, closure or
stability, transport, status, records, symmetry, interfaces and
self-reference---which six chapters group as in the table below.  The
Status, Records, Coherence chapter carries the status, records and
interfaces axes, and the Closure--Reality Boundary lies outside the
axes.  Axes label chapters, not laws.  The headings in the summary
tables of the last two chapters (for instance Status/Records/Coherence
and Limits/Emergence/Invariants) are reading groups only.  The result
codes are fixed identifiers and are not consecutive: F1, F16 and F32 do
not occur, F13 and F15 each name two laws (a and b), and the two
closure--reality no-go theorems belong to the F53 row.  In the Access,
Stability, Symmetry and Status, Records, Coherence chapters, each law
carries a labelled nonclaims paragraph; in the Transport and
Self-Reference and Limits chapters and in the closure--reality chapter,
the scope is set by the hypotheses and by the catalog-wide nonclaims of
\Cref{sec:nonclaims}.

\begin{center}
\small
\begin{tabularx}{\linewidth}{@{}>{\raggedright\arraybackslash}p{0.33\linewidth}>{\raggedright\arraybackslash}p{0.24\linewidth}>{\raggedright\arraybackslash}X@{}}
\toprule
Chapter & Axes & Laws \\
\midrule
Access (\Cref{sec:access}) & access & \Fref{F10}--\Fref{F12}, \Fref{F13a}, \Fref{F13b}, \Fref{F24} \\
Stability (\Cref{sec:stability}) & closure or stability & \Fref{F6}--\Fref{F9} \\
Transport (\Cref{sec:transport}) & transport & \Fref{F2}--\Fref{F5} \\
Symmetry (\Cref{sec:symmetry}) & symmetry & \Fref{F14}, \Fref{F15a}, \Fref{F15b}, \Fref{F40}, \Fref{F41} \\
Status, Records, Coherence (\Cref{sec:status-records-coherence}) & status, records, interfaces & \Fref{F17}--\Fref{F23}, \Fref{F25}--\Fref{F31} \\
Self-Reference and Limits (\Cref{sec:self-reference-limits}) & self-reference & \Fref{F33}--\Fref{F39}, \Fref{F42}--\Fref{F52} \\
Closure--Reality Boundary (\Cref{sec:anchor:closure-reality-boundary}) & outside the axes & \Fref{F53} and the two closure-equals-reality (CER) no-go theorems \\
\bottomrule
\end{tabularx}
\end{center}

\paragraph{Reading path.} Sections~\ref{sec:intro} and~\ref{sec:apparatus},
then Transport (\Cref{sec:transport}), whose source-repair clause is used
in the proof of \Fref{F15a}; then the remaining chapters in any order.
\Fref{F13a} can be left for last.  A paragraph headed \emph{Reading the
statement} explains the predicates of the result that follows it; on a
first reading it may be read after that result.  No statement or proof
uses \Cref{app:further-instantiations} (further analogies) or
\Cref{app:version-history} (version history) or
\Cref{app:formalization} (formalization).

\paragraph{Layer instantiations.}
Most results are followed by a list headed \emph{Layer
instantiations}, and further analogies for most laws are collected in
\Cref{app:further-instantiations}.  Each item places a construction
from another field in the slots of the result: its carrier, quotients,
readouts and obstruction.  Each item carries one of two tags.  A
\emph{special case} fills the slots with well-defined sets and maps that
satisfy the hypotheses, so the result applies to it as stated;
statements in a special-case item beyond the slot assignment are cited
facts about the example, not consequences of the law.  An
\emph{analogy} matches the construction to the slots informally: its
data are institutional, empirical or approximate; or the construction
is exact mathematics that does not satisfy the law's hypotheses or fill
all of its slots (for instance the Schur-complement criterion beside
\Cref{thm:F6}); or the result needs a source hypothesis or a formed
package that the item does not supply.  No proof depends on these lists,
and a first reading can skip them.

\paragraph{Acronyms.}
AOR: audited operational realisability.
CER: closure equals reality, the unrestricted equivalence of Six Birds
closure and bare physical realisability, which the two no-go theorems
of the Closure--Reality Boundary chapter show to fail in each direction
under their obstruction hypotheses.
CSL: candidate structural law, in the sense of the paper
\emph{Hiddenness as a Structural Law of Emergence}
\citep{Tsiokos2026CSL}.
SDTC: self-dual trace confinement.
SSB: spontaneous symmetry breaking.

\paragraph{Map of results.}\addcontentsline{toc}{subsection}{Map of results}
The rows of the catalog carry F-codes, which the text uses as names.
The table lists each F-code with its printed number, the section that
contains it, its name, and its page.

\begingroup\small
\setlength{\tabcolsep}{3pt}
\begin{longtable}{@{}>{\raggedright\arraybackslash}p{0.12\linewidth}>{\raggedright\arraybackslash}p{0.17\linewidth}>{\raggedright\arraybackslash}p{0.08\linewidth}>{\raggedright\arraybackslash}p{0.47\linewidth}>{\raggedright\arraybackslash}p{0.07\linewidth}@{}}
\toprule
F-code & Result & Section & Name & Page \\
\midrule
\endfirsthead
\toprule
F-code & Result & Section & Name & Page \\
\midrule
\endhead
\bottomrule
\endfoot
F10 & Theorem~\ref{thm:F10} & \ref{sec:access:quotientality} & Quotientality & \pageref{thm:F10} \\
F11 & Theorem~\ref{thm:F11} & \ref{sec:access:adequacy} & Adequacy / No-Overread & \pageref{thm:F11} \\
F12 & Theorem~\ref{thm:F12} & \ref{sec:access:no-free-distinction} & No-Free-Distinction & \pageref{thm:F12} \\
F13a & Theorem~\ref{thm:F13a} & \ref{sec:access:hiddenness} & Hiddenness Normal Form & \pageref{thm:F13a} \\
F13b & Theorem~\ref{thm:F13b} & \ref{sec:access:csl-formed-closure} & CSL Formed-Closure Law & \pageref{thm:F13b} \\
F24 & Theorem~\ref{thm:F24} & \ref{sec:access:layer-multiplicity} & Layer Multiplicity & \pageref{thm:F24} \\
F6 & Theorem~\ref{thm:F6} & \ref{sec:stability:no-needles} & No-Needles Stability & \pageref{thm:F6} \\
F7 & Theorem~\ref{thm:F7} & \ref{sec:stability:sufficiency-closure} & Sufficiency Closure & \pageref{thm:F7} \\
F8 & Lemma~\ref{lem:F8} & \ref{sec:stability:strict-extension} & Strict Extension & \pageref{lem:F8} \\
F9 & Theorem~\ref{thm:F9} & \ref{sec:stability:no-unstatused-residual} & No Unstatused Residual & \pageref{thm:F9} \\
F2 & Theorem~\ref{thm:F2} & \ref{sec:transport:descent-repair} & Descent--Repair Normal Form & \pageref{thm:F2} \\
F3 & Theorem~\ref{thm:F3} & \ref{sec:transport:holonomy-memory} & Holonomy-Memory Repair Normal Form & \pageref{thm:F3} \\
F4 & Theorem~\ref{thm:F4} & \ref{sec:transport:local-global} & Local-Global Obstruction Normal Form & \pageref{thm:F4} \\
F5 & Theorem~\ref{thm:F5} & \ref{sec:transport:recombination} & Recombination Comparison & \pageref{thm:F5} \\
F14 & Theorem~\ref{thm:F14} & \ref{sec:symmetry:duality-fixity} & Duality Fixity / Self-Dual Trace Confinement & \pageref{thm:F14} \\
F15a & Theorem~\ref{thm:F15a} & \ref{sec:symmetry:selection-obstruction} & Selection Obstruction / Repair Normal Form & \pageref{thm:F15a} \\
F15b & Theorem~\ref{thm:F15b} & \ref{sec:symmetry:spontaneous-ssb} & Spontaneous Symmetry Breaking / Branch-Selection Formation & \pageref{thm:F15b} \\
F40 & Theorem~\ref{thm:F40} & \ref{sec:symmetry:anomaly} & Anomaly / Symmetry Obstruction & \pageref{thm:F40} \\
F41 & Theorem~\ref{thm:F41} & \ref{sec:symmetry:duality-equivalence} & Duality Equivalence & \pageref{thm:F41} \\
F17 & Theorem~\ref{thm:F17} & \ref{sec:status-records-coherence:objectivity} & Objectivity as Common Quotient & \pageref{thm:F17} \\
F18 & Theorem~\ref{thm:F18} & \ref{sec:status-records-coherence:lawhood-descent} & Lawhood Descent & \pageref{thm:F18} \\
F19 & Theorem~\ref{thm:F19} & \ref{sec:status-records-coherence:object-persistence} & Object Persistence & \pageref{thm:F19} \\
F20 & Theorem~\ref{thm:F20} & \ref{sec:status-records-coherence:fact-record} & Fact / Record Formation & \pageref{thm:F20} \\
F21 & Theorem~\ref{thm:F21} & \ref{sec:status-records-coherence:reflexive-nonclosure} & Reflexive Nonclosure & \pageref{thm:F21} \\
F22 & Theorem~\ref{thm:F22} & \ref{sec:status-records-coherence:interface-mediation} & Interface Mediation & \pageref{thm:F22} \\
F23 & Theorem~\ref{thm:F23} & \ref{sec:status-records-coherence:probability} & Probability as Stable Unresolved Fiber & \pageref{thm:F23} \\
F25 & Theorem~\ref{thm:F25} & \ref{sec:status-records-coherence:causality} & Causality as Stable Intervention Transport & \pageref{thm:F25} \\
F26 & Theorem~\ref{thm:F26} & \ref{sec:status-records-coherence:moduli} & Moduli / Contingency & \pageref{thm:F26} \\
F27 & Theorem~\ref{thm:F27} & \ref{sec:status-records-coherence:conservation} & Conservation as Orbit Descent & \pageref{thm:F27} \\
F28 & Theorem~\ref{thm:F28} & \ref{sec:status-records-coherence:composability} & Composability & \pageref{thm:F28} \\
F29 & Theorem~\ref{thm:F29} & \ref{sec:status-records-coherence:presentation-invariance} & Presentation Invariance & \pageref{thm:F29} \\
F30 & Theorem~\ref{thm:F30} & \ref{sec:status-records-coherence:explanation} & Explanation as Compression with Descent & \pageref{thm:F30} \\
F31 & Theorem~\ref{thm:F31} & \ref{sec:status-records-coherence:counterfactual} & Counterfactual Legitimacy & \pageref{thm:F31} \\
F33 & Theorem~\ref{thm:F33} & \ref{sec:self-reference-limits:horizon-sufficiency} & Horizon Sufficiency / Holographic Compression & \pageref{thm:F33} \\
F34 & Theorem~\ref{thm:F34} & \ref{sec:self-reference-limits:information-loss} & Information Loss Normal Form & \pageref{thm:F34} \\
F35 & Theorem~\ref{thm:F35} & \ref{sec:self-reference-limits:scale-descent} & Scale Descent / Renormalization & \pageref{thm:F35} \\
F36 & Theorem~\ref{thm:F36} & \ref{sec:self-reference-limits:singular-limit} & Singular Limit Completion & \pageref{thm:F36} \\
F37 & Theorem~\ref{thm:F37} & \ref{sec:self-reference-limits:complementarity} & Complementarity as Non-Joint Access & \pageref{thm:F37} \\
F38 & Theorem~\ref{thm:F38} & \ref{sec:self-reference-limits:arrow} & Arrow from Record Monotonicity & \pageref{thm:F38} \\
F39 & Theorem~\ref{thm:F39} & \ref{sec:self-reference-limits:entropy} & Entropy as Fiber Volume & \pageref{thm:F39} \\
F42 & Theorem~\ref{thm:F42} & \ref{sec:self-reference-limits:irreducibility} & Irreducibility / No Short Closed Description & \pageref{thm:F42} \\
F43 & Theorem~\ref{thm:F43} & \ref{sec:self-reference-limits:continuum-emergence} & Continuum Emergence & \pageref{thm:F43} \\
F44 & Theorem~\ref{thm:F44} & \ref{sec:self-reference-limits:criticality} & Criticality / Phase Boundary & \pageref{thm:F44} \\
F45 & Theorem~\ref{thm:F45} & \ref{sec:self-reference-limits:locality} & Locality / Domain of Dependence & \pageref{thm:F45} \\
F46 & Theorem~\ref{thm:F46} & \ref{sec:self-reference-limits:law-state-split} & Law--State Split & \pageref{thm:F46} \\
F47 & Theorem~\ref{thm:F47} & \ref{sec:self-reference-limits:fine-tuning} & Fine-Tuning as Small Selector Region & \pageref{thm:F47} \\
F48 & Theorem~\ref{thm:F48} & \ref{sec:self-reference-limits:topology} & Topology as Global Gluing Invariant & \pageref{thm:F48} \\
F49 & Theorem~\ref{thm:F49} & \ref{sec:self-reference-limits:common-source} & Common-Source / Nonlocal Correlation Normal Form & \pageref{thm:F49} \\
F50 & Theorem~\ref{thm:F50} & \ref{sec:self-reference-limits:vacuum} & Vacuum / Background Selection & \pageref{thm:F50} \\
F51 & Theorem~\ref{thm:F51} & \ref{sec:self-reference-limits:unification} & Unification as Common Refinement & \pageref{thm:F51} \\
F52 & Theorem~\ref{thm:F52} & \ref{sec:self-reference-limits:closure-completeness} & Closure Completeness Boundary & \pageref{thm:F52} \\
CER (direct) & Theorem~\ref{thm:CER-direct-no-go} & \ref{sec:anchor:closure-reality-boundary} & Direct CER no-go & \pageref{thm:CER-direct-no-go} \\
CER (converse) & Theorem~\ref{thm:CER-converse-no-go} & \ref{sec:anchor:closure-reality-boundary} & Converse CER no-go & \pageref{thm:CER-converse-no-go} \\
F53 & Proposition~\ref{prop:F53} & \ref{sec:anchor:closure-reality-boundary} & Closure--Reality Boundary resolution & \pageref{prop:F53} \\
\end{longtable}
\endgroup

\paragraph{Instances of the descent criterion.} The following laws contain the descent criterion, applied to the source map and readout indicated in parentheses (with surjectivity, or another hypothesis supplying the descended map, as stated in each law) (for \Fref{F17}, in its third displayed line); it is also the descent clause of \Cref{thm:F2} with \(F=\mathrm{id}\): \Fref{F10} (\(\accessQ\), \(o\)), \Fref{F11} (\(\accessQ\), \(F\)), \Fref{F12} (\(\accessQ\), \(d\)), \Fref{F13a} (\(\currQ_I\), \(\predM_I\)), \Fref{F13b} (\(\langle\currQ,d'\rangle\), \(\predM\)), \Fref{F24} (\(q\), \(s\)), \Fref{F7} (\(q\), \(\Phi\)), \Fref{F8} (\(q_{\mathrm{old}}\), \(q_{\mathrm{new}}\)), \Fref{F5} (\(K\), \(R\)), \Fref{F15a} (\(\pi\), \(s\)), \Fref{F40} (\(q\), \(q\circ\alpha(g,\cdot)\) for each \(g\)), \Fref{F17} (\(p\), \(r\)), \Fref{F18} (\(\accessQ\), \(P\)), \Fref{F19} (\(q_i\), \(q_j\circ\gamma\)), \Fref{F20} (\(s\), \(\varepsilon\)), \Fref{F23} (\(q\), \(r\)), \Fref{F27} (\(q_{\mathrm{orb}}\), \(C\)), \Fref{F29} (\(\pi\), \(T\)), \Fref{F30} (\(b\), \(e\)), \Fref{F33} (\(\pi\), \(t\)), \Fref{F34} (\(\rho\circ\tau\), \(\sigma\)), \Fref{F35} (\(q_n\), \(t\)), \Fref{F36} (\(\ell\), \(u\)), \Fref{F38} (\(r_s\), \(r_t\circ\kappa\)), \Fref{F39} (\(q\), \(t\)), \Fref{F43} (\(\ell\), \(r_b\circ p_b\)), \Fref{F44} (\(c\), \(r\)), \Fref{F45} (\(\ell\), \(t\)), \Fref{F46} (\(\mu\), \(r\)), \Fref{F48} (\(\gamma\), \(r\)), \Fref{F49} (\((\sigma,q_A)\), \(r_A\), and the interface and route quotients, \(r\)) and \Fref{F50} (\(\mu\), \(b_0\)); \Fref{F22}, \Fref{F25}, \Fref{F28} and \Fref{F31} apply it to a product source such as \(q\times\beta\).
        % sec:apparatus
\section{Access}\label{sec:access}

The Access chapter is about what an instrument can distinguish.  An
observation signature \(\obsSig\) (\Cref{def:access-quotient}) is a
declared map from states to observation symbols; its access quotient
\(\accessQ\) identifies exactly those states that no declared
observation separates.  The laws of this chapter say how such
signatures determine quotients, when those quotients are adequate for
accepted claims, when a new distinction requires new apparatus data
rather than a free relabeling, and when structure remains hidden behind
the current quotient \(\currQ\) although a predictive quotient
\(\predM\) separates it.  In the running example of \Cref{sec:intro},
\(q\) is the access quotient of the signature that reads each state's
class \(a\), \(b\) or \(c\); the states \(1\) and \(2\) are
indistinguishable to it, which is why \(F\) cannot be read off \(q\).

Quotientality (\Fref{F10}) says that an observation signature
determines a canonical access quotient, through which every observable
constant on the signature factors.  Adequacy / No-Overread (\Fref{F11})
says that accepted claims must factor through the available access.
No-Free-Distinction (\Fref{F12}) says that a readout splits an access
fiber exactly when it does not factor through the access quotient, that
such a readout cannot be accepted as exact under the no-overread rule,
and that the pairing \(\langle\accessQ,d\rangle\) is the coarsest
refinement carrying it.  Layer Multiplicity (\Fref{F24}) closes the
chapter: a role readout that the access quotient does not determine is
carried by a strict refinement of that quotient, an additional layer.

Between these, the Hiddenness Normal Form (\Fref{F13a}) and the CSL
Formed-Closure Law (\Fref{F13b}) compare current and predictive
observation.  Hiddenness is the situation in which states agree under
\(\currQ\) but differ under \(\predM\).  \Fref{F13a} is conditional on a
source hypothesis motivated by \citet{Tsiokos2026PvNP}; \Fref{F13b} is
a direct theorem under an explicit determining hypothesis, motivated by
\citet{Tsiokos2026CSL}.  No result of this paper assumes the
P~\(\neq\)~NP claim of the former.  On a first reading these two laws
can be skipped; the rest of the chapter does not use them.

\subsection{Quotientality (Layer-Sameness via Quotient) [\texorpdfstring{\Fref{F10}}{F10}]}\label{sec:access:quotientality}

\Cref{def:obs-signature} restates \Cref{def:access-quotient} for a closure
carrier, names the quotient set \(Q_I\), and adds the calibration-family case.

\begin{definition}[Observation signature]\label{def:obs-signature}
Let \(H\) be a closure carrier.  An observation signature on \(H\)
is a declared map
\[
  \obsSig:H\to\Sigma
\]
from states to observation symbols.  Its access kernel is the
equivalence relation
\[
  h\sim_I h'
  \quad\Longleftrightarrow\quad
  \obsSig(h)=\obsSig(h'),
\]
and its access quotient is the quotient map
\[
  \accessQ:H\to Q_I:=H/{\sim_I}.
\]
For a family \((\sigma_j)_{j\in J}\) of observation signatures (a calibration
family), the joint access kernel is \(h\sim h'\iff\sigma_j(h)=\sigma_j(h')\) for
every \(j\), the kernel of the bundled signature \(h\mapsto(\sigma_j(h))_j\), and
the same construction is applied to it.
\end{definition}

\begin{theorem}[Quotientality]\label{thm:F10}
Let \(\obsSig:H\to\Sigma\) be an observation signature with access
quotient \(\accessQ:H\to Q_I\).  Every observable \(o:H\to R\) that
is constant on access classes factors uniquely through
\(\accessQ\):
\[
  o=\overline o\circ \accessQ
\]
for a unique \(\overline o:Q_I\to R\).  Moreover, \(Q_I\) is the
coarsest quotient carrying the declared observations: if
\(\pi:H\to Q'\) is any surjective map on whose fibers \(\obsSig\) is constant,
then \(\accessQ\) factors through \(\pi\).
\end{theorem}
\begin{proof}
Define \(\overline o([h]_I)=o(h)\).  This is well defined because
\([h]_I=[h']_I\) means \(h\sim_I h'\), and the hypothesis on \(o\)
gives \(o(h)=o(h')\).  The equation
\(o=\overline o\circ\accessQ\) follows from the definition on
representatives.  If \(\overline o'\) is another such map, then for
every access class \([h]_I\) we have
\[
  \overline o'([h]_I)
  =
  \overline o'(\accessQ(h))
  =
  o(h)
  =
  \overline o([h]_I),
\]
so \(\overline o'=\overline o\).

For coarsestness, suppose \(\pi:H\to Q'\) is a surjective map on whose fibers
\(\obsSig\) is constant: \(\pi(h)=\pi(h')\) implies
\(\obsSig(h)=\obsSig(h')\).  Since \(\pi\) is surjective, define
\(u:Q'\to Q_I\) by \(u(\pi(h))=\accessQ(h)\).  The same invisibility
hypothesis makes this assignment independent of the representative.
Hence \(\accessQ=u\circ\pi\).  Thus any quotient already carrying the
declared observations refines the canonical access quotient.
\end{proof}

\paragraph{Nonclaims.} The
theorem does not produce the coarsest quotient possible for all
observers; it produces the coarsest quotient carrying the given
signature.  It also does not create new observable distinctions and
does not instantiate any particular carrier theory.

\subsection*{Layer instantiations}\label{layer-inst:F10}

\begingroup\small
\begin{description}
\item[\emph{Myhill--Nerode minimal {DFA} (formal language theory).}] \textit{Special case.}
For a regular language \(L\subseteq\Sigma^*\), take \(H=\Sigma^*\) and
\(\obsSig(w):=w^{-1}L=\{z:wz\in L\}\), valued in the power set of \(\Sigma^*\) (here \(\Sigma\) is the alphabet, not the symbol set of \Cref{def:obs-signature}). Its access kernel is the Nerode
relation \(\equiv_L\), defined by
\(w\equiv_L w'\Leftrightarrow\forall z.\,(wz\in L\Leftrightarrow w'z\in L)\),
so \(Q_I=\Sigma^*/\!\!\equiv_L\) is the state set of the minimal DFA
recognising \(L\). \Fref{F10} then says that every observable constant on
\(\equiv_L\) factors uniquely through \(Q_I\), and that every surjection with
\(\obsSig\)-constant fibers refines \(\equiv_L\). That \(\equiv_L\) is also the
coarsest right-congruence making \(\mathbf 1_L\) class-constant is the
Myhill--Nerode theorem, not part of \Fref{F10}
\citep{Nerode1958LinearAutomatonTransformations}.

\item[\emph{Kalman observability canonical decomposition (control
engineering).}] \textit{Special case.}
Take \(H=\mathbb R^n\) and \(\obsSig(x):=(CA^kx)_{k\ge0}\); its access kernel
is \(x-x'\in N\). For a linear time-invariant system \(\dot x=Ax+Bu,\ y=Cx\), the
unobservable subspace
\(N=\bigcap_{k\ge 0}\ker(CA^k)\) is the kernel of the observability
matrix.  Two state vectors agree under \(\mathbb R^n/N\) iff they
produce identical output trajectories for every admissible input,
and the Kalman canonical form factors every input--output
observable uniquely through \(\mathbb R^n/N\); every surjection of \(\mathbb R^n\) on whose fibres the output trajectory is constant refines \(\mathbb R^n/N\) (the quotient map \(\mathbb R^n\to\mathbb R^n/N\) factors through it), which is the coarsest-quotient clause of \Fref{F10}
\citep{Kalman1963MathematicalDescription}.

\item[\emph{Bernstein 3{NF} synthesis (database engineering).}] \textit{Analogy.}
Given a relation schema with attributes \(U\) and a
functional-dependency set \(F\), the 3NF decomposition produced by Bernstein
synthesis from a chosen minimal cover of \(F\) is a
dependency-preserving 3NF decomposition, which depends on the chosen
cover: every FD-checker for the cover factors through the projection
onto a single component schema, and coalescing components keeps
every declared FD checkable but can introduce a 3NF violation
\citep{Bernstein1976Synthesizing3NF}.

\end{description}
\endgroup


\subsection{Adequacy / No-Overread [\texorpdfstring{\Fref{F11}}{F11}]}\label{sec:access:adequacy}

Fix the same closure carrier \(H\) and declared observable family as
in the previous subsection, and let \(\accessQ:H\to Q_I\) be the
access quotient constructed from that family.  A claim is \emph{accepted as an exact current claim} when it passes
the finite audited status discipline of \FoundIII\ as an exact claim
about the current layer \citep{Tsiokos2026FoundationsIII}.  By the
no-overread rule (NO) of \Cref{sec:apparatus}, it may then use only the
distinctions available to the current access; this is the premise the
proof uses.  The theorem records what the discipline guarantees: an
accepted exact claim is constant on the fibers of \(\accessQ\), and
hence factors through it.

\begin{theorem}[Adequacy / No-Overread]\label{thm:F11}
Let \(F:H\to W\) be a map (in the intended reading, a current-layer claim over a
formed access package), and let
\(\mathcal O(\accessQ,F):=\{(h,h'):\accessQ(h)=\accessQ(h'),\ F(h)\ne F(h')\}\)
be its overread obstruction.
\begin{enumerate}[label=(\alph*)]
\item For every \(F\),
\[
  F\ \text{constant on the fibers of}\ \accessQ
  \iff \mathcal O(\accessQ,F)=\emptyset
  \iff F=\overline F\circ\accessQ\ \text{for a unique}\ \overline F:Q_I\to W .
\]
\item Under the no-overread rule (NO) of \Cref{sec:apparatus}, if \(F\) is
  accepted as an exact current claim, then it is constant on the fibers of
  \(\accessQ\), and there is a unique map \(\overline F:Q_I\to W\) such that
  \(F=\overline F\circ\accessQ\).
\item Hence the no-overread rule (NO) holds if and only if every accepted
  exact current claim factors uniquely through
  \(\accessQ\).
\end{enumerate}
The first conclusion of (b) is (NO) applied to \(F\); the uniqueness uses the
universal property of \(\accessQ\) (\Cref{thm:F10}).
\end{theorem}
\begin{proof}
For (b), (NO) of \Cref{sec:apparatus} applied to the accepted claim \(F\) states exactly that \(F(h)=F(h')\) whenever \(\accessQ(h)=\accessQ(h')\).

By Quotientality, every map constant on access classes factors
uniquely through \(\accessQ\).  Applying \Cref{thm:F10} to \(F\)
gives a unique \(\overline F:Q_I\to W\) with
\(F=\overline F\circ\accessQ\).  Thus an accepted current claim has
no overread defect. For the equivalences, a pair in \(\mathcal O(\accessQ,F)\)
is exactly a failure of constancy on a fiber, and Quotientality, with
surjectivity of \(\accessQ\), gives the unique factorization of a map constant
on fibers; conversely a factored map is constant on fibers. Applying these
equivalences to each accepted exact current claim gives the final
statement.
\end{proof}

\paragraph{Nonclaims.} The theorem does not forbid recording distinctions beyond
current access; such distinctions may belong to the predictive
quotient \(\predM\), to memory, or to a later access extension.  It
also does not determine the audit data, bridge strengths, thresholds,
or residual budgets required for acceptance.

\subsection*{Layer instantiations}\label{layer-inst:F11}

\begingroup\small
\begin{description}
\item[\emph{Rao--Blackwell theorem (statistical decision theory).}] \textit{Analogy.}
For a sufficient statistic \(T:X\to\mathcal T\) on a parametric
family \(\{P_\theta\}\), the Rao--Blackwell construction
\(\delta_{\mathrm{RB}}(x)=\mathbb E[\delta(X)\mid T(X)=T(x)]\) takes any
decision rule \(\delta\) and produces a \(\sigma(T)\)-measurable
rule whose risk under a convex loss is at most that of \(\delta\)
at every \(\theta\), and strictly smaller at some \(\theta\) under a
strictly convex loss unless \(\delta\) is already a function of \(T\)
almost surely.  Hence under a strictly convex loss any admissible
decision rule with finite risk factors as
\(\delta=\bar\delta\circ T\) a.s.; a rule that uses within-fibre
detail of \(X\) beyond \(T(X)\) is the overread failure
\citep{Blackwell1947ConditionalExpectation}.

\end{description}
\endgroup


\subsection{No-Free-Distinction [\texorpdfstring{\Fref{F12}}{F12}]}\label{sec:access:no-free-distinction}

Keep the formed access quotient \(\accessQ:H\to Q_I\) from
Quotientality, together with the no-overread discipline from
Adequacy.  A proposed new current distinction is a readout
\(d:H\to D\).  Its split-pair obstruction at fixed access is
\[
  \splitObs_{\mathrm{dist}}(\accessQ,d)
  =
  \{(h,h')\mid \accessQ(h)=\accessQ(h')\ \text{and}\ d(h)\ne d(h')\}.
\]

Call \(d\) \emph{current-visible} when \(\accessQ(h)=\accessQ(h')\) implies
\(d(h)=d(h')\). The \emph{financed refinement} of \(d\), the refinement that adjoins \(d\)
as declared apparatus data, is
\(\langle\accessQ,d\rangle:H\to Q_I\times D\).

\begin{theorem}[No-Free-Distinction]\label{thm:F12}
Let \(d:H\to D\) be a readout (the access quotient \(\accessQ\) is surjective by construction). Then
\(d\) is current-visible if and only if
\[
  \splitObs_{\mathrm{dist}}(\accessQ,d)=\varnothing,
  \qquad\text{equivalently}\qquad
  d=\overline d\circ\accessQ
\]
for some \(\overline d:Q_I\to D\); in particular a readout obtained by
applying a map to the access quotient cannot split an access fiber. The
financed refinement \(\langle\accessQ,d\rangle\) retains \(\accessQ\) and
carries \(d\) (both factor through it), it strictly refines \(\accessQ\) when
the obstruction is nonempty (it separates two histories that \(\accessQ\)
identifies, so it does not factor through \(\accessQ\)), and every map \(p\) through which both
\(\accessQ\) and \(d\) factor also has \(\langle\accessQ,d\rangle\) factoring
through it. If the obstruction is nonempty and the no-overread rule (NO) of
\Cref{sec:apparatus} holds, \(d\) is not accepted as an exact current claim.
\end{theorem}
\begin{proof}
By definition, \(d\) is current-visible exactly when no pair identified by
\(\accessQ\) is separated by \(d\), that is, when the displayed obstruction is
empty. By \Cref{thm:F10}, this is equivalent to the existence of a
factorization \(d=\overline d\circ\accessQ\). A readout \(g\circ\accessQ\)
depends only on the access class, so it is current-visible.

The projections of \(Q_I\times D\) give
\(\accessQ=\operatorname{pr}_1\circ\langle\accessQ,d\rangle\) and
\(d=\operatorname{pr}_2\circ\langle\accessQ,d\rangle\). If
\((h,h')\) lies in the obstruction, then \(\accessQ(h)=\accessQ(h')\) while
\(\langle\accessQ,d\rangle(h)\ne\langle\accessQ,d\rangle(h')\), so the
refinement is strict. If \(\accessQ=a\circ p\) and \(d=b\circ p\), then
\(\langle\accessQ,d\rangle=\langle a,b\rangle\circ p\).

Finally, if the obstruction is nonempty, \(d\) separates two histories that
the current access quotient identifies, so \(d\) is not constant on the fibers
of \(\accessQ\). By (NO), as in \Cref{thm:F11}, an exact current claim is
constant on those fibers; hence \(d\) is not accepted as exact.
\end{proof}

A distinction that the access quotient does not support can still be used
with declared apparatus data: memory, a bridge, a budgeted residual, a
scope restriction, or stronger calibration. Among quotient refinements
through which both \(\accessQ\) and \(d\) factor,
\(\langle\accessQ,d\rangle\) is coarsest, by the factorization proved above.

\paragraph{Nonclaims.} The theorem is not anti-refinement.  Calibration families may enrich
the apparatus, and strict extensions may add real distinctions.  The
claim is only that the enrichment must be declared and statused, not
smuggled in as a free consequence of the old access quotient.  Nor
does this theorem classify every possible source of new data; later
axes treat memory, bridges, residual budgets, scope changes, and
predictive refinements in their own normal forms.

\subsection*{Layer instantiations}\label{layer-inst:F12}

\begingroup\small
\begin{description}
\item[\emph{Landauer's principle (thermodynamics of computation).}] \textit{Analogy.}
Resetting a one-bit register whose prior content is an unknown,
unbiased bit to a fixed value, by a process that exchanges heat with a
thermal reservoir at temperature \(T\) and starts uncorrelated with it,
dissipates at least \(k_B T\ln 2\) of heat into the reservoir; in
general the bound is \(k_BT\) times the decrease of the register's
entropy, measured in nats, so the lower bound is smaller for a biased prior bit
\citep{Landauer1961IrreversibilityHeatGeneration,ReebWolf2014ImprovedLandauer}.
The analogy is that a change in the distinctions a register carries is
not free: it is paid for by declared apparatus, here heat exported to
the reservoir.

\end{description}
\endgroup



\subsection{Hiddenness Normal Form (conditional) [\texorpdfstring{\Fref{F13a}}{F13a}]}\label{sec:access:hiddenness}

\begin{definition}[Hiddenness interface]\label{def:hiddenness-exposure}
A \emph{hiddenness interface} \(I\) consists of a closure carrier
\(\mathcal H_I\), a current quotient
\(\currQ_I:\mathcal H_I\to Q_I^{\mathrm{cur}}\), a predictive quotient
\(\predM_I:\mathcal H_I\to M_I\) refining \(\currQ_I\), an upstream
readout \(u_I:\mathcal H_I\to U_I\) of the determining state, and three
declared propositions
\[
  \operatorname{determiningStateTyped}(I),\quad
  \operatorname{minimalRelevantUpstream}(I),\quad
  \operatorname{lawfulCurrentExposure}(I).
\]
From these data we define \(\operatorname{QuotientExposed}(I)\),
\(\operatorname{PredictiveCollapse}(I)\), \(\operatorname{PredictiveSurplus}(I)\)
and \(\operatorname{HiddenFromCurrent}(I)\) below.
\Fref{F13b} defines a
distinct pair-level predicate using a readout \(d\) and proves an equivalence
under an explicit determining hypothesis. \(\operatorname{QuotientExposed}\),
\(\operatorname{PredictiveCollapse}\), \(\operatorname{PredictiveSurplus}\) and
\(\operatorname{HiddenFromCurrent}\) are defined from the maps and the declared
predicates:
\begin{align*}
  \operatorname{QuotientExposed}(I)
  &:\Longleftrightarrow \exists f,\ u_I=f\circ\currQ_I,\\
  \operatorname{PredictiveCollapse}(I)
  &:\Longleftrightarrow \exists g,\ \predM_I=g\circ\currQ_I,\\
  \operatorname{PredictiveSurplus}(I)
  &:\Longleftrightarrow \exists h,h',\ \currQ_I(h)=\currQ_I(h')
    \wedge \predM_I(h)\ne\predM_I(h'),\\
  \operatorname{HiddenFromCurrent}(I)
  &:\Longleftrightarrow \neg\operatorname{lawfulCurrentExposure}(I).
\end{align*}
The hiddenness normal form is the
five-conjunct property
\begin{align*}
  \operatorname{HiddennessNormalForm}(I)
  :={}&
  \operatorname{determiningStateTyped}(I)
  \wedge
  \operatorname{minimalRelevantUpstream}(I)\\
  &{}\wedge
  \bigl(\operatorname{lawfulCurrentExposure}(I)
    \leftrightarrow
    \operatorname{QuotientExposed}(I)\bigr)\\
  &{}\wedge
  \bigl(\operatorname{PredictiveCollapse}(I)
    \leftrightarrow
    \operatorname{lawfulCurrentExposure}(I)\bigr)\\
  &{}\wedge
  \bigl(\operatorname{PredictiveSurplus}(I)
    \leftrightarrow
    \operatorname{HiddenFromCurrent}(I)\bigr).
\end{align*}
\end{definition}

\paragraph{Reading the statement.}
The current quotient \(\currQ_I\) records what is observable now, and
the predictive quotient \(\predM_I\) records what matters for
prediction.  \(\operatorname{PredictiveCollapse}(I)\) says that
\(\predM_I\) factors through \(\currQ_I\), so that prediction needs
nothing beyond the current class; \(\operatorname{PredictiveSurplus}(I)\)
says that two histories with the same current class have different
predictive classes; and \(\operatorname{QuotientExposed}(I)\) says that
the whole upstream readout factors through \(\currQ_I\).  The predicates
\(\operatorname{determiningStateTyped}(I)\) (the determining state has
been given a type), \(\operatorname{minimalRelevantUpstream}(I)\) (the
relevant upstream data are minimal) and
\(\operatorname{lawfulCurrentExposure}(I)\) (the current quotient is a
lawful exposure of the determining state) are declared data of the
interface, not reduced to the quotients.

The theorem has two parts.  Part (i) assumes the source hypothesis
\((\mathsf S)\): \(\currQ_I\) is surjective and the first four conjuncts
of the normal form hold.  It proves the fifth conjunct and the chain of
equivalences linking exposure, collapse and the absence of surplus.  In
the intended application, \(I\) would be an interface extracted from
the saturated-SAT construction of \citet{Tsiokos2026PvNP}, and
\((\mathsf S)\) a property of that construction; that paper constructs
no such interface, and \((\mathsf S)\) is not verified for it here.
Hypothesis \((\mathsf S)\) is satisfiable, also together with a
predictive surplus: take two histories, a one-point current quotient,
the identity as predictive quotient and as upstream readout,
\(\operatorname{lawfulCurrentExposure}(I)\) false and the other two
declared predicates true.  Part (ii) assumes nothing beyond the
interface: the equivalence of surplus with non-collapse holds exactly
when \(\currQ_I\) is surjective or \(M_I\) is nonempty.  The refinement
of \(\currQ_I\) by \(\predM_I\) required in
\Cref{def:hiddenness-exposure} is not used by the theorem; it serves
only the remark, in the nonclaims below, that
\(\langle\currQ_I,\predM_I\rangle\) has the fibers of \(\predM_I\).

\begin{theorem}[Hiddenness Normal Form]\label{thm:F13a}
Let \(I\) be a hiddenness interface.
\begin{enumerate}[label=(\roman*)]
\item Assume the source hypothesis
  \((\mathsf S)\): \(\currQ_I\) is surjective, and
  \begin{align*}
    &\operatorname{determiningStateTyped}(I),\\
    &\operatorname{minimalRelevantUpstream}(I),\\
    &\operatorname{lawfulCurrentExposure}(I)
      \iff \operatorname{QuotientExposed}(I),\\
    &\operatorname{PredictiveCollapse}(I)
      \iff \operatorname{lawfulCurrentExposure}(I).
  \end{align*}
  Then
  \[
    \operatorname{PredictiveSurplus}(I)
      \iff \operatorname{HiddenFromCurrent}(I),
  \]
  so \(I\) satisfies \(\operatorname{HiddennessNormalForm}(I)\), and
  \[
    \operatorname{QuotientExposed}(I)\iff
    \operatorname{PredictiveCollapse}(I)\iff
    \neg\operatorname{PredictiveSurplus}(I):
  \]
  the upstream readout descends through the current quotient exactly
  when prediction does.
\item Without \((\mathsf S)\), \(I\) satisfies
  \(\operatorname{PredictiveSurplus}(I)\iff\neg\operatorname{PredictiveCollapse}(I)\)
  if and only if \(\currQ_I\) is surjective or \(M_I\) is nonempty, in
  particular whenever \(\currQ_I\) is surjective; and for every \(I\)
  satisfying this equivalence, the fourth conjunct of the normal form
  holds if and only if the fifth does.
\end{enumerate}
\end{theorem}
\begin{proof}
We first show that a surjective \(\currQ_I\) gives
\(\operatorname{PredictiveSurplus}(I)\iff\neg\operatorname{PredictiveCollapse}(I)\);
this is \Cref{thm:F7} for the single probe \(\predM_I\), which we spell
out.  A surplus pair \(h,h'\) rules out \(\predM_I=g\circ\currQ_I\),
since it would give
\(\predM_I(h)=g(\currQ_I h)=g(\currQ_I h')=\predM_I(h')\).  Conversely,
if there is no surplus pair, \(\predM_I\) is constant on the fibers of
\(\currQ_I\), and since \(\currQ_I\) is surjective,
\(g(\currQ_I h):=\predM_I(h)\) defines \(g\) with
\(\predM_I=g\circ\currQ_I\).

(i) By the fourth conjunct, \(\neg\operatorname{PredictiveCollapse}(I)\)
is equivalent to
\(\neg\operatorname{lawfulCurrentExposure}(I)=\operatorname{HiddenFromCurrent}(I)\),
which with the equivalence just proved gives the fifth conjunct.  The
final chain combines the third and fourth conjuncts with the same
equivalence.

(ii) If \(\currQ_I\) is surjective, the equivalence was proved above.
If \(M_I\ne\emptyset\), the same argument applies, with the
fiber-defined map extended outside the image of \(\currQ_I\) by any
chosen value.  If \(M_I=\emptyset\), then \(\mathcal H_I=\emptyset\),
there is no surplus, and collapse holds exactly when the current
codomain is empty, that is, when \(\currQ_I\) is surjective.  Given the
equivalence, the fourth conjunct is equivalent to
\(\neg\operatorname{PredictiveSurplus}(I)\iff\operatorname{lawfulCurrentExposure}(I)\),
which is the fifth conjunct negated on both sides.
\end{proof}

\paragraph{Nonclaims.} The theorem does not prove \((\mathsf S)\) for
any interface derived from the saturated-SAT construction of
\citet{Tsiokos2026PvNP}, does not re-prove the closure result of that
paper, and asserts nothing about a complexity class.  Its unconditional
content is part (ii).  The upstream factorization here concerns the
whole declared readout; it is distinct from the exposure predicate of
the Hiddenness paper \citep{Tsiokos2026CSL}, which asks for a single
nonconstant visible function of state.  Since \(\predM_I\) refines
\(\currQ_I\), the pairing \(\langle\currQ_I,\predM_I\rangle\), the
coarsest refinement of \(\currQ_I\) through which \(\predM_I\) descends
(\Cref{sec:apparatus}), has the same fibers as \(\predM_I\).  The
Kochen--Specker theorem \citep{KochenSpecker1967Hidden}, a
non-contextual hidden-variable obstruction in quantum foundations, is
a substrate-specific result, not a special case of this law.

\subsection*{Layer instantiations}\label{layer-inst:F13a}

\begingroup\small
\begin{description}
\item[\emph{Hidden Markov model identifiability (statistical
latent-variable analysis).}] \textit{Analogy.}
In a hidden Markov model, two observation histories ending in the
same observed symbol can predict different distributions of future
observations, because they lead to different posterior distributions
over the hidden state; such a pair is a predictive surplus over the
current observation, and the hidden state is the determining state
upstream of it.  For a stationary hidden Markov model with \(r\) hidden
states and \(\kappa\ge2\) observation symbols, Allman, Matias and
Rhodes show that the parameters are generically identifiable, up to
relabeling of the hidden states, from the joint distribution of
\(2k+1\) consecutive observations whenever
\(\binom{k+\kappa-1}{\kappa-1}\ge r\); \emph{generically}
excludes a set of parameters of measure zero.  Identifiability concerns
the model, not the success of an estimation procedure such as
Baum--Welch EM \citep{AllmanMatiasRhodes2009Identifiability}.

\end{description}
\endgroup


\subsection{CSL Formed-Closure Law [\texorpdfstring{\Fref{F13b}}{F13b}]}\label{sec:access:csl-formed-closure}

\begin{definition}[Diagnosis source data]\label{def:diagnosis-source}
Diagnosis source data on a carrier \(H\) is a triple \((\currQ,\predM,d)\) of
a current quotient \(\currQ:H\to Q\), a predictive quotient
\(\predM:H\to M\) and a readout \(d:H\to S\) that is \emph{determining}:
\(\currQ(h)=\currQ(h')\wedge d(h)=d(h')\Rightarrow\predM(h)=\predM(h')\). In
the CSL case, \PaperCSL\ \citep{Tsiokos2026CSL} identifies an internal-state
readout; the determining property is an additional hypothesis that must be
verified for a proposed diagnosis source.
\end{definition}

\Cref{thm:F13b} applies to exactly such data.

\paragraph{Reading the statement.} The theorem has three parts: the surplus equivalence for the given determining \(d\); for an arbitrary \(d'\), a factorization test for being determining and a two-case split of predictive surplus; and the statement that \(\langle\currQ,\predM\rangle\) is the coarsest pairing \(\langle\currQ,d'\rangle\) with \(d'\) determining.
The closure carries a declared determining-state readout \(d\), the hidden
state that the diagnosis of \PaperCSL\ identifies.
\(\operatorname{PredictiveSurplus}(\mathcal C)\) says that two histories with
the same current class have different predictive classes;
\(\operatorname{HiddenStateSurplus}(\mathcal C)\) says that two histories with
the same current class and different determining states have different
predictive classes, so that the predictive split is carried by the hidden
determining state. The hypothesis that \(d\) is determining says that the
current class and the determining state together fix the predictive class.
This pairwise property is assumed here; it is not established by the cited CSL
diagnosis.

\begin{theorem}[CSL Formed-Closure Law]\label{thm:F13b}
Let \(\mathcal C=(H,\currQ,\predM,d)\) consist of a carrier \(H\) with
current quotient \(\currQ:H\to Q\), predictive quotient \(\predM:H\to M\) and a
declared determining-state readout \(d:H\to S\) that is determining:
\[
  \currQ(h)=\currQ(h')\ \wedge\ d(h)=d(h')
  \ \Longrightarrow\ \predM(h)=\predM(h').
\]
Define
\begin{align*}
  \operatorname{PredictiveSurplus}(\mathcal C)
  &:\Longleftrightarrow
  \exists h,h'\in H:\ \currQ(h)=\currQ(h')\ \wedge\ \predM(h)\ne\predM(h'),\\
  \operatorname{HiddenStateSurplus}(\mathcal C)
  &:\Longleftrightarrow
  \exists h,h'\in H:\ \currQ(h)=\currQ(h')\ \wedge\ d(h)\ne d(h')\\
  &\qquad{}\wedge\ \predM(h)\ne\predM(h').
\end{align*}
Then
\[
  \operatorname{PredictiveSurplus}(\mathcal C)
  \iff
  \operatorname{HiddenStateSurplus}(\mathcal C),
\]
and more precisely every pair \(h,h'\) with \(\currQ(h)=\currQ(h')\) and
\(\predM(h)\ne\predM(h')\) has \(d(h)\ne d(h')\): predictive surplus across
the \(\predM/\currQ\) gap lies only on pairs that the hidden determining state
separates; equivalently, the surplus pairs are the current-fiber pairs that both \(d\) and
\(\predM\) separate. Moreover, for an arbitrary readout \(d':H\to S'\), not assumed determining, \(d'\) is determining if and only if
\(\predM=g\circ\langle\currQ,d'\rangle\) for some map \(g\) from the image of
\(\langle\currQ,d'\rangle\) to \(M\) (equivalently, when \(M\) is nonempty, for some
\(g:Q\times S'\to M\)), and
\(\operatorname{PredictiveSurplus}\) holds if and only if some pair of one current
fiber is separated by both \(d'\) and \(\predM\), or some pair of one current
fiber agrees under \(d'\) and is separated by \(\predM\) (this second disjunct
says exactly that \(d'\) is not determining). In particular
\(d'=\predM\) is determining, and \(d'\) is determining exactly when
\(\langle\currQ,\predM\rangle\) is constant on the fibers of
\(\langle\currQ,d'\rangle\); hence, after corestricting the pairings to their images,
\(\langle\currQ,\predM\rangle\) is the coarsest pairing \(\langle\currQ,d'\rangle\) with
\(d'\) determining. When \(M\) is nonempty, the canonical pairing also factors
through the full-codomain pairing \(\langle\currQ,d'\rangle:H\to Q\times S'\) for
every determining \(d'\).
\end{theorem}
\begin{proof}
If \(\currQ(h)=\currQ(h')\) and \(\predM(h)\ne\predM(h')\), then
\(d(h)\ne d(h')\), since otherwise the determining hypothesis would give
\(\predM(h)=\predM(h')\). So every surplus pair is a hidden pair, which gives
\(\operatorname{PredictiveSurplus}\Rightarrow\operatorname{HiddenStateSurplus}\);
the converse holds because a hidden pair is in particular a surplus pair.
For an arbitrary \(d'\), if \(d'\) is determining, define \(g(a,b)\) as
\(\predM(h)\) for any \(h\) with \(\langle\currQ,d'\rangle(h)=(a,b)\), which is
well defined by the determining property; when \(M\) is nonempty, \(g\) extends
to \(Q\times S'\) by a fixed value off the image. Conversely such a \(g\) makes
\(d'\) determining. The surplus disjunction splits a surplus pair by
whether \(d'\) separates it.
For the last clause, \(\predM\) is determining trivially, and since the \(\currQ\) coordinates of two histories with equal \(\langle\currQ,d'\rangle\) agree, determination of \(d'\) is constancy of \(\langle\currQ,\predM\rangle\) on the fibers of \(\langle\currQ,d'\rangle\).
\end{proof}

\paragraph{What the source supplies.} \PaperCSL\ \citep{Tsiokos2026CSL} defines an internal-state readout and studies its exposure to current observables and its behavior under trace-fingerprint loops. It does not prove that \((\currQ,d)\) determines \(\predM\) for every pair of histories. That determination is an explicit hypothesis of \Cref{thm:F13b}; applying the theorem to a CSL carrier requires a separate proof of it.


\paragraph{Nonclaims.} This subsection does not instantiate the CSL diagnosis source and does
not assert a measurement-theoretic contextuality theorem at the
\FIV\ layer.

\subsection*{Layer instantiations}\label{layer-inst:F13b}

\begingroup\small
\begin{description}
\item[\emph{Bayesian clinical diagnostic testing (medicine).}] \textit{Analogy.}
For a clinical assay with declared sensitivity \(\mathrm{Se}\),
specificity \(\mathrm{Sp}\) and cohort prevalence, Bayes' rule gives
\(P(\text{disease}\mid r)\) for each possible result \(r\). Take the current
class to record the cohort and assay parameters, and the determining readout
to record \(r\); together they determine the posterior. If two possible
results give different posteriors within one current class, they witness
predictive and hidden-state surplus. Disease status alone need not determine
the posterior, and this analogy does not infer disease status from that
surplus \citep{Sackett2000EvidenceBasedMedicine}.

\end{description}
\endgroup


\subsection{Layer Multiplicity [\texorpdfstring{\Fref{F24}}{F24}]}\label{sec:access:layer-multiplicity}

Let \(H\) be a closure carrier with a formed current access quotient
\(q:H\to Q\).  A calibration or role signature is another declared
map \(s:H\to S\): it may record a future probe, an interaction role,
a lawhood signature, a branch readout, or a calibrated instrument
response.  The signature is already carried by the old layer exactly
when it descends through \(q\), equivalently when it is constant on
the \(q\)-fibers.

\paragraph{Reading the statement.}
\(\operatorname{RoleSplit}(q,s)\) says that the old layer merges two
histories whose roles differ, that is, some \(q\)-fiber contains two
values of \(s\). \(\operatorname{StrictLayerExtension}(q,q_s)\) says
that \(q_s\) retains \(q\), that is \(q=\mathrm{pr}_Q\circ q_s\), and that some
\(q\)-fiber is split by \(q_s\): there are \(h,h'\) with \(q(h)=q(h')\) and
\(q_s(h)\ne q_s(h')\). The theorem shows that a role split
is always resolved by this extension. The layer-formation discipline also
records eight other ways a split may be handled: by a memory layer, by a
hidden upstream role, through a bridge, within a budget, within a declared
scope, by coarsening, by placing the role outside the declared scope, or as
blocked nonclosure. These are declared statuses, not determined by \(q\) and
\(s\), and the theorem makes no claim about them. In ordinary notation: \(s\) factors
through \(q\) exactly when \(\mathcal O_s=\emptyset\); some \(q\)-fiber
carries two values of \(s\) exactly when \(\mathcal O_s\neq\emptyset\),
which happens exactly when \(q_s\) strictly refines \(q\); and a split
is then resolved at least by the strict extension \(q_s\). With \(d=s\), the
descent criterion, the strict refinement under a split and the coarsest-map
property are those of \Cref{thm:F12}; the converse (a strict extension
implies a role split) and the corestriction are added here.

\begin{theorem}[Layer Multiplicity]\label{thm:F24}
Let \(H\) be a history space, let \(q:H\to Q\) be a surjective access
quotient, and let \(s:H\to S\) be a declared role readout.  Define
the role obstruction and joint role-layer quotient by
\[
  \mathcal O_s
  :=
  \{(h,h')\in H^2:q(h)=q(h')\ \text{and}\ s(h)\ne s(h')\},
  \qquad
  q_s(h):=(q(h),s(h)).
\]
Then
\begin{align*}
  s\ \text{descends through}\ q
  &\iff \mathcal O_s=\emptyset,\\
  \operatorname{RoleSplit}(q,s)
  &\iff \mathcal O_s\ne\emptyset\\
  &\iff \operatorname{StrictLayerExtension}(q,q_s).
\end{align*}
In particular,
\[
  \operatorname{RoleSplit}(q,s)\Longrightarrow
  \operatorname{StrictLayerExtension}(q,q_s):
\]
a role split is resolved by the strict layer extension \(q_s\). Moreover, \(q_s:H\to Q\times S\) is coarsest in the factorization
preorder among maps through which both \(q\) and \(s\) factor: if
\(q=a\circ p\) and \(s=b\circ p\) for a map \(p\), then
\(q_s=\langle a,b\rangle\circ p\). Its corestriction
\(H\to\operatorname{im}(q_s)\) is the corresponding coarsest surjective
quotient among surjective maps \(p\) through which \(q\) and \(s\) factor.
\end{theorem}
\begin{proof}
Assume \(q:H\to Q\) is surjective and \(s:H\to S\) is declared.  If
\(s\) descends through \(q\), then \(s=\bar s\circ q\) for some
\(\bar s:Q\to S\), so
\[
  q(h)=q(h')\Longrightarrow s(h)=s(h')
\]
and \(\mathcal O_s=\emptyset\).  Conversely, if
\(\mathcal O_s=\emptyset\), define
\[
  \bar s(q(h)):=s(h).
\]
The empty obstruction makes \(\bar s\) independent of the
representative \(h\), and surjectivity of \(q\) supplies all quotient
classes.  Thus \(s\) descends through \(q\).

If \(\mathcal O_s\ne\emptyset\), choose
\((h,h')\in\mathcal O_s\).  Then \(q(h)=q(h')\) but
\[
  q_s(h)=(q(h),s(h))\ne(q(h'),s(h'))=q_s(h'),
\]
so \(q_s\) strictly refines \(q\).  Conversely, a strict refinement of
\(q\) by \(q_s\) means that some \(q\)-fiber is split by the second
coordinate \(s\), giving an element of \(\mathcal O_s\).

Finally, if \(\operatorname{RoleSplit}(q,s)\) holds, the previous
paragraph shows that \(q_s\) strictly refines \(q\), which is
\(\operatorname{StrictLayerExtension}(q,q_s)\).
For the last clause, \(q_s(h)=(q(h),s(h))=(a(p(h)),b(p(h)))=\langle a,b\rangle(p(h))\).
Corestricting \(q_s\) to its image makes it surjective. If \(p\) is
surjective, every \(t\) is \(p(h)\) for some \(h\), so
\((a(t),b(t))=q_s(h)\) lies in that image and gives the asserted
quotient factorization.
\end{proof}

\paragraph{Nonclaims.} The theorem proves that a role split forces
the strict extension \(q_s\); adopting \(q_s\) as a new layer still
requires admissibility, closure formation, and status.
Layer Multiplicity does not classify the number of layers in a
carrier, does not prescribe how layers must be presented, and does
not assert that multiplicity is finite in every setting.

\subsection*{Layer instantiations}\label{layer-inst:F24}

\begingroup\small
\begin{description}
\item[\emph{Operating-system privilege protection (computer
security).}] \textit{Analogy.}
In a multi-user operating system, the privilege quotient classifies
operations by the trust level of the executing principal.  When a
process needs an operation outside its declared privilege, the
layer-formation discipline names a resolution: a user-to-kernel
system-call trap, a capability-bridge invocation, an
inter-process-communication transfer, a sandbox or namespace
restriction, or
outright denial --- corresponding as follows to the declared statuses memory layer (the trap's
saved context), bridge (capability invocation and IPC transfer), declared scope
(sandbox) and blocked nonclosure (denial) listed in the reading paragraph of
\Fref{F24}, about which \Fref{F24} makes no claim
\citep{SaltzerSchroeder1975ProtectionInformation}.

\end{description}
\endgroup


\subsection*{Access at a glance}\label{sec:access:summary}\addcontentsline{toc}{subsection}{Access at a glance}
\begin{table}[H]
  \centering
  \footnotesize
  \begin{tabularx}{\textwidth}{@{}
      >{\raggedright\arraybackslash}p{0.24\textwidth}
      >{\raggedright\arraybackslash}X
      >{\raggedright\arraybackslash}p{0.12\textwidth}@{}}
    \toprule
    Name & One-line claim & Substrate cite \\
    \midrule
    Quotientality theorem (\Fref{F10}) &
    Every observable constant on access classes factors uniquely through the access quotient, and the access quotient factors through every surjection on whose fibers the signature is constant (it is the coarsest such quotient).  &
    \textemdash \\
    Adequacy / No-Overread theorem (\Fref{F11}) &
    Under the no-overread rule, accepted exact claims factor uniquely through the access quotient.  &
    \textemdash \\
    No-Free-Distinction theorem (\Fref{F12}) &
    A readout is current-visible iff it factors through the access
    quotient; the pairing with it is always the coarsest map through which
    both factor, strict when the readout splits a fiber, and then, under
    (NO), the readout is not accepted as exact.  &
    \textemdash \\
    Hiddenness Normal Form (\Fref{F13a}), conditional &
    Under the source hypothesis \((\mathsf S)\) (surjectivity and four of
    the five conjuncts), the fifth conjunct holds, and exposure, collapse
    and absence of surplus are equivalent; without \((\mathsf S)\),
    predictive surplus holds exactly when predictive collapse fails, for a
    surjective current quotient. &
    \PaperPvNP \\
    CSL Formed-Closure (\Fref{F13b}) &
    With a determining-state readout, every predictive-surplus pair is separated by the determining state.  &
    \PaperCSL \\
    Layer-Multiplicity Theorem (\Fref{F24}) &
    A role readout splits an access quotient exactly when its role
    obstruction is nonempty, and a role split implies the strict
    layer extension by the joint role-layer quotient.  &
    \textemdash \\
    \bottomrule
  \end{tabularx}
  \caption{Access at a glance: the six laws of \Cref{sec:access}.}
  \label{tab:axis-access-summary}
\end{table}

\FloatBarrier
     % sec:access
\section{Stability}\label{sec:stability}

The Stability axis asks when declared future behavior is already
controlled by present data, and when a residual must be recorded
instead of ignored.  In the language of \Cref{sec:apparatus}, the
question is whether the current quotient carries the probes the
package has declared, or whether closure-stability fails under a
structural perturbation.  A perturbation may be a new future probe,
a layer-dissolving move at a membrane, a strict extension of a
quotient, or an unpriced residual.  The axis records the normal
forms that decide which of these outcomes is legitimate.
In the running example of \Cref{sec:intro}, the value of \(F\) is a
declared future probe that the current quotient \(q\) does not carry:
\(q\) identifies \(1\) and \(2\) while \(F\) separates them, so the
split pair must be resolved or recorded as a residual.

The chapter has four laws.  No-Needles Stability is the first theorem of the chapter: when the
currency, residual and endpoint budgets hold together with the stated block
identities, every transported dissolving direction is bounded by its budget,
so no needle exists.  Sufficiency Closure then gives the
factorization criterion for declared futures: the current quotient is
sufficient exactly when every declared future probe descends through
it.  Strict Extension records when genuine layer growth has occurred,
namely when the old quotient is retained and a new quotient splits at
least one old fiber.  No Unstatused Residual closes the axis: under the
residual-coverage rule and agreement of statuses at a common locus, every
native residual carries a well-defined typed status, and neither hypothesis
can be dropped.

No-Needles Stability is the layer-agnostic form of the layer-dissolving
membrane theorem of \citet{Tsiokos2026NeedleKiller}, which
\citet{Tsiokos2026NavierStokes} applies to the Navier--Stokes
equations.  Here a membrane means the
layer-dissolving boundary of a closure: the place where a distinction
between ``inside'' and ``outside'' the layer can cease to be
meaningful.  Strict Extension also has descriptive substrate context,
but it is a direct law here: the Schur-calculus intuition comes from
\PaperAR\ \citep{Tsiokos2026XiCompanion}, where residuals compose
along refinement chains.  The Stability axis can be read
independently of the Access axis; its only prerequisite is the
notation of \Cref{sec:apparatus}.

\subsection{No-Needles Stability [\texorpdfstring{\Fref{F6}}{F6}]}\label{sec:stability:no-needles}

The setup is a formed stability package: a closure carrier equipped
with the inherited closure apparatus, quotient discipline, audit
ledger, and a layer-dissolving membrane.  The membrane is the
boundary across which a native layer distinction is allowed to
dissolve only when the package has recorded the appropriate
apparatus footprint.  A \emph{needle} is a transported dissolving direction whose effect exceeds
its declared budget. In the matrix data of the theorem below, it is a vector
\(v\in\mathbb R^d\) with
\(v^{\!\top}S\,K_{DD}\,S^{\!\top}v>v^{\!\top}\Theta_D v\) (the block \(K_{DD}\), the selector \(S\) and the bound \(\Theta_D\in\mathbb R^{d\times d}\) are introduced in the paragraphs \emph{Notation} and \emph{Bounds and order} below).

\paragraph{Notation.} The theorem itself concerns four blocks
\(K_{LL},K_{LD},K_{DL},K_{DD}\) and a matrix \(K_{LL}^{\dagger}\) subject to the
identities displayed in it; the next paragraph describes their
intended origin. The data of the theorem are the matrices displayed in it.
\paragraph{Intended origin (not used by the theorem).} In the intended origin, which the theorem does not use, the currency
block \(C\) is an \(e\times e\) matrix on the exterior carrier, and \(L\)
and \(D\) are \(y\times e\) and \(z\times e\) blocks for the layer-bridge
and dissolving directions. Let \(C^{+}\) be any symmetric \(e\times e\) matrix (for instance the
Moore--Penrose pseudoinverse of a symmetric \(C\); the theorem uses \(C\) only through \(C^{+}\), and uses the blocks only
through the stated identities); the currency pieces are \(K_{LL}=L\,C^{+}L^{\!\top}\),
\(K_{LD}=L\,C^{+}D^{\!\top}\), \(K_{DL}=D\,C^{+}L^{\!\top}\) and
\(K_{DD}=D\,C^{+}D^{\!\top}\), and symmetry of \(C^{+}\) gives
\(K_{DL}=K_{LD}^{\!\top}\). Let \(K_{LL}^{\dagger}\) be any symmetric \(y\times y\)
matrix with \(K_{LL}K_{LL}^{\dagger}K_{LD}=K_{LD}\) (for instance the
Moore--Penrose pseudoinverse of \(K_{LL}\) when \(C^{+}\succeq0\)). The declared native explanation
\(\widehat E=K_{DL}K_{LL}^{\dagger}\) predicts the dissolving block from the
layer-bridge block, and the adequacy residual
\(\Xi=K_{DD}-K_{DL}K_{LL}^{\dagger}K_{LD}\), a generalized Schur complement,
is the algebraic remainder of that explanation. When \(C^{+}\succeq0\), the
block matrix \(\bigl[\begin{smallmatrix}K_{LL}&K_{LD}\\K_{DL}&K_{DD}\end{smallmatrix}\bigr]\)
is positive semidefinite; then \(\Xi\succeq0\), and for every \(z\times y\)
matrix \(E\),
\(K_{DD}-EK_{LD}-K_{DL}E^{\!\top}+EK_{LL}E^{\!\top}=\Xi+(E-\widehat E)K_{LL}(E-\widehat E)^{\!\top}\succeq\Xi\),
so \(\widehat E\) is the Loewner-optimal explanation. For arbitrary
symmetric \(C^{+}\), neither positivity nor optimality is asserted; the endpoint proof uses the block identities and
budgets alone.
\paragraph{Bounds and order.} \(\Theta_Y\in\mathbb R^{y\times y}\), \(\Omega_Z\in\mathbb R^{z\times z}\) and
\(\Theta_D\in\mathbb R^{d\times d}\) are declared bounds; no construction from the blocks \(L\) or \(D\) is required,
\(\preceq\) is the Loewner preorder of quadratic forms
(\(A\preceq B\) when \(v^{\!\top}(B-A)v\ge0\) for every vector \(v\); on
symmetric matrices this is the Loewner order, and in general it compares
symmetric parts), and \(S\) is any \(d\times z\) matrix, for instance
one selecting \(d\) of the dissolving directions.
In \citet{Tsiokos2026NeedleKiller} the audit form \(C\) is a positive
semidefinite operator on a finite-dimensional real or complex space, and
\(C^{+}\) and \(K_{LL}^{\dagger}\) are Moore--Penrose pseudoinverses.  The
proof below uses only the displayed matrix identities and budgets; the
package and its membrane name the setting in which they are declared.

\begin{theorem}[No-Needles Stability]\label{thm:F6}
Suppose given the following data (in the intended reading, declared by a
formed stability package \(\mathcal P\), whose other properties, including its
membrane, are not used):
\begin{itemize}
  \item block matrices \(K_{LL}\in\mathbb R^{y\times y}\),
  \(K_{LD}\in\mathbb R^{y\times z}\), \(K_{DL}\in\mathbb R^{z\times y}\) and
  \(K_{DD}\in\mathbb R^{z\times z}\) (in the intended reading, the currency pieces
  formed from a symmetric \(C^{+}\) and the blocks \(L\), \(D\) as in the paragraph \emph{Intended origin}) and a matrix \(K_{LL}^{\dagger}\in\mathbb R^{y\times y}\);
  \item the symmetry conditions \(K_{DL}=K_{LD}^{\!\top}\) (automatic when the
  pieces are formed from a symmetric \(C^{+}\)) and
  \(K_{LL}^{\dagger}=(K_{LL}^{\dagger})^{\!\top}\), together with the
  projection identity
  \[
    K_{LL}\,K_{LL}^{\dagger}\,K_{LD}=K_{LD};
  \]
  \item currency and residual budgets
  \[
    K_{LL}\preceq \Theta_Y,
    \qquad
    \Xi\preceq \Omega_Z,
  \]
  where \(\Xi:=K_{DD}-K_{DL}K_{LL}^{\dagger}K_{LD}\) is the adequacy residual;
  \item a transport selector \(S\) of size \(d\times z\) satisfying
  the endpoint budget
  \[
    S\bigl(\widehat E\,\Theta_Y\,\widehat E^{\!\top}
      +\Omega_Z\bigr)S^{\!\top}\preceq\Theta_D,
  \]
  where \(\widehat E:=K_{DL}K_{LL}^{\dagger}\) is the declared native
  explanation.
\end{itemize}
Then the transported dissolving block satisfies the Loewner endpoint
\[
  S\,K_{DD}\,S^{\!\top}\preceq\Theta_D .
\]
Equivalently, no needle exists: every transported dissolving direction is
bounded by its declared budget.
\end{theorem}
\begin{proof}
Assume the symmetry, projection, currency-budget, residual-budget,
and transport-budget hypotheses displayed in the theorem.  The
projection identity
\[
  K_{LL}\,K_{LL}^{\dagger}\,K_{LD}=K_{LD}
\]
gives the layer-dissolving decomposition of the substrate master theorem.
Indeed, with \(\widehat E:=K_{DL}\,K_{LL}^{\dagger}\), the symmetry conditions
give \(\widehat E^{\!\top}=K_{LL}^{\dagger}K_{LD}\), so
\[
  \widehat E\,K_{LL}\,\widehat E^{\!\top}
  =K_{DL}K_{LL}^{\dagger}\bigl(K_{LL}K_{LL}^{\dagger}K_{LD}\bigr)
  =K_{DL}K_{LL}^{\dagger}K_{LD},
\]
and by the definition of \(\Xi\),
\[
  K_{DD}
  =
  \widehat E\,K_{LL}\,\widehat E^{\!\top}
  +\Xi.
\]
Here \(\widehat E\) is the declared native explanation and \(\Xi\) is
the adequacy residual.  Applying Loewner monotonicity under the
declared bounds \(K_{LL}\preceq\Theta_Y\) and
\(\Xi\preceq\Omega_Z\) yields
\[
  K_{DD}
  \preceq
  \widehat E\,\Theta_Y\,\widehat E^{\!\top}+\Omega_Z .
\]
Here the first step also uses that congruence preserves the order: for any
matrix \(A\) and any \(v\),
\(v^{\!\top}A(Y-X)A^{\!\top}v=(A^{\!\top}v)^{\!\top}(Y-X)(A^{\!\top}v)\), so
\(X\preceq Y\) implies \(AXA^{\!\top}\preceq AYA^{\!\top}\), applied with
\(A=\widehat E\). Conjugation by \(S\) preserves the Loewner order in the same
way, so
\[
  S\,K_{DD}\,S^{\!\top}
  \preceq
  S\bigl(\widehat E\,\Theta_Y\,\widehat E^{\!\top}
    +\Omega_Z\bigr)S^{\!\top}.
\]
The declared transport-side endpoint budget gives
\[
  S\bigl(\widehat E\,\Theta_Y\,\widehat E^{\!\top}
    +\Omega_Z\bigr)S^{\!\top}
  \preceq \Theta_D .
\]
Transitivity of \(\preceq\) therefore gives the endpoint
\[
  S\,K_{DD}\,S^{\!\top}\preceq\Theta_D .
\]

\end{proof}

The proof uses only the arithmetic of an ordered field and the order
on quadratic forms, so the theorem holds verbatim with \(\mathbb R\)
replaced by any ordered field, for instance \(\mathbb Q\).  The complex
instances of \citet{Tsiokos2026NeedleKiller} need the analogous
statement with conjugate transposes and Hermitian forms, which is not
stated here.

\paragraph{What the sources supply.} The layer-dissolving membrane
theorem of \citet{Tsiokos2026NeedleKiller} applies the same argument at
every accepted stage of a refinement, with positive semidefinite audit
forms and Moore--Penrose pseudoinverses, and identifies the endpoint
with a bound on delivered capacity; \citet{Tsiokos2026NavierStokes}
applies it to shell readouts of the Navier--Stokes equations.
\Cref{thm:F6} isolates the matrix inequality common to these
instances.

\paragraph{Nonclaims.} The theorem does not re-prove the Needle Killer or Navier--Stokes
substrate instances, does not classify every possible
layer-dissolving move, and does not assert that the substrate carrier
is unique.  It records the common stability normal form: under the
declared budgets, the transported dissolving block is dominated in the Loewner
order by its endpoint budget.

\subsection*{Layer instantiations}\label{layer-inst:F6}

\begingroup\small
\begin{description}
\item[\emph{Schur-complement positivity criterion (linear algebra).}] \textit{Analogy.}
A symmetric block matrix
\(\bigl[\begin{smallmatrix}A&B\\B^{\!\top}&C\end{smallmatrix}\bigr]\)
is positive semidefinite iff \(A\succeq0\), \((I-AA^{+})B=0\), and
\(C-B^{\!\top}A^{+}B\succeq0\).  The criterion is not a case of \Cref{thm:F6}; it is the positivity fact stated in the paragraph \emph{Intended origin}, with \(A,B,C,A^{+}\) in the roles of \(K_{LL},K_{LD},K_{DD},K_{LL}^{\dagger}\): \((I-AA^{+})B=0\) is the projection identity \(K_{LL}K_{LL}^{\dagger}K_{LD}=K_{LD}\), and the Schur complement \(C-B^{\!\top}A^{+}B\) is the adequacy residual \(\Xi\), whose bound \(\Xi\preceq\Omega_Z\) is a hypothesis of the theorem; the transported endpoint is \(SK_{DD}S^{\!\top}\) \citep{Zhang2005SchurApplications}.

\end{description}
\endgroup


\subsection{Sufficiency Closure [\texorpdfstring{\Fref{F7}}{F7}]}\label{sec:stability:sufficiency-closure}

Fix a closure package with history carrier \(H\), current quotient
map \(q:H\to Q\), and a declared family \(\mathcal F\) of future
probes.  The family is the apparatus's announcement of which future
observables will be tested at this interface; it is not a claim of
sufficiency for every possible future question.

Write \(\Phi:H\to\prod_{f\in\mathcal F}F_f\), \(\Phi(h)=(f(h))_{f\in\mathcal F}\),
for the bundled future signature, and call \(q\) \emph{sufficient} for
\(\mathcal F\) when \(q(h)=q(h')\) implies \(f(h)=f(h')\) for every
\(f\in\mathcal F\).

\begin{theorem}[Sufficiency Closure]\label{thm:F7}
Assume \(q\) is surjective. The current quotient \(q\) is sufficient
for the declared future family \(\mathcal F\) if and only if every declared
future probe \(f\in\mathcal F\) descends through \(q\), that is, for every
\(f:H\to F_f\) in \(\mathcal F\) there is a map
\(\overline f:Q\to F_f\) such that
\[
  f=\overline f\circ q ,
\]
equivalently \(\Phi=\overline\Phi\circ q\) for some
\(\overline\Phi:Q\to\prod_{f\in\mathcal F}F_f\), and if and only if no pair \((h,h')\) has \(q(h)=q(h')\) and
\(\Phi(h)\ne\Phi(h')\). Moreover exactly one of the following four statuses
holds: \emph{exact closure} (\(q\) is sufficient and \(\Phi\) determines \(q\): \(\Phi(h)=\Phi(h')\Rightarrow q(h)=q(h')\) for all \(h,h'\in H\)); \emph{redundant closure} (\(q\) sufficient, \(\Phi\) does not determine \(q\)); \emph{predictive non-closure} (\(\Phi\) determines \(q\), \(q\) not sufficient); and \emph{incomparable} (neither).
\end{theorem}
\begin{proof}
Assume first that every declared future probe descends through
\(q\).  Then each future value \(f(h)\) is determined by the current
class \(q(h)\), so two histories with the same current class have the
same value under every declared probe; this is sufficiency.

Conversely, suppose \(q\) is sufficient for the declared future family.
For a probe \(f:H\to F_f\) and a class \(x\in Q\), choose \(h\) with
\(q(h)=x\), which exists because \(q\) is surjective, and set
\(\overline f(x)=f(h)\). Sufficiency makes this independent of the chosen
\(h\): if \(q(h)=q(h')\), then \(f(h)=f(h')\). Therefore \(\overline f\) is
well-defined and satisfies \(f=\overline f\circ q\). The bundled form follows
by taking \(\overline\Phi=(\overline f)_{f\in\mathcal F}\); conversely
\(\overline f=\mathrm{pr}_f\circ\overline\Phi\).

Sufficiency says that no pair in a common current fiber is separated by a
probe, which is the statement that no pair has \(q(h)=q(h')\) and
\(\Phi(h)\ne\Phi(h')\). Finally, the four statuses are the four
combinations of two Boolean conditions, sufficiency and the implication
\(\Phi(h)=\Phi(h')\Rightarrow q(h)=q(h')\); every configuration satisfies
exactly one combination.
\end{proof}

\paragraph{Nonclaims.} The theorem does not enumerate the
declared future probes; that choice belongs to the apparatus.  It
also does not assert sufficiency for undeclared probes, which lie
outside the current audit scope.

\subsection*{Layer instantiations}\label{layer-inst:F7}

\begingroup\small
\begin{description}
\item[\emph{Information bottleneck (information theory / machine
learning).}] \textit{Special case.}
For input \(X\) and a declared prediction target \(Y\), the
Tishby--Pereira--Bialek information bottleneck trades the compression cost \(I(X;T)\) against the relevance
\(I(T;Y)\) over stochastic encoders; its limiting case, with a
deterministic encoder \(T=q(X)\), \(q\) corestricted to its image as \Fref{F7} requires, minimises \(I(X;T)\) under the
sufficiency constraint \(I(T;Y)=I(X;Y)\).  For a finite input alphabet
with \(P_X(x)>0\) for every \(x\), taken as \(H\), the constraint holds exactly
when \(P(Y\mid X=x)=P(Y\mid T=q(x))\) for every \(x\), which is the \Fref{F7}
factorization for \(\mathcal F=\{x\mapsto P(Y\mid X=x)\}\).  Over-compressing the
encoder breaks the factorization: the \Fref{F7} obstruction is nonempty
exactly when the gap \(I(X;Y)-I(T;Y)\) is positive
\citep{TishbyPereiraBialek1999InformationBottleneck}.

\item[\emph{Predictive state representations (reinforcement
learning).}] \textit{Special case.}
A PSR summarizes the history of a partially observable dynamical
system by the vector \(p(h)\), corestricted to its image as \Fref{F7} requires, of probabilities of a declared family of
core future tests.  The Singh--James--Rudary sufficiency theorem
states that every declared future test factors linearly through
\(p(h)\): \(\Pr(t\mid h)=m_t\cdot p(h)\).  The summary is sufficient
for the declared test family while remaining hidden-state-free, and a
summary built from fewer tests than the system's linear dimension does not make
every test a linear function of it
\citep{SinghJamesRudary2004PredictiveStateRepresentations}.

\item[\emph{Query-complete materialized views (databases).}] \textit{Special case.}
For a base schema \(R\) and a declared workload \(\mathcal Q\) of
queries, a set of materialized views \(V\) \emph{determines} \(Q\in\mathcal Q\)
when \(V(I)=V(I')\) implies \(Q(I)=Q(I')\) for all instances \(I,I'\). By
\Fref{F7}, with the view map corestricted to its image, this holds iff
\(Q=Q'\circ V\) for some map \(Q'\) on view extents; the view materialization
plays the summary and \(Q'\) the lower map. A rewriting of \(Q'\) in a query
language is a stronger requirement: conjunctive views can determine a
conjunctive query that has no conjunctive rewriting \citep{NashSegoufinVianu2010Views}.  The
sufficiency is workload-relative: a query needing an attribute
projected away by \(V\) does not factor through the views
\citep{Halevy2001AnsweringQueriesUsingViewsSurvey}.

\item[\emph{Quantum sufficient subalgebras (quantum information).}] \textit{Analogy.}
For a parametric family of density operators
\(\{\rho_\theta:\theta\in\Theta\}\) and a CPTP quantum channel \(T\),
Petz's theorem says \(T\) is sufficient for the family iff there is a
recovery channel \(R\) (the Petz recovery map) with
\(\rho_\theta=R(T(\rho_\theta))\) for every \(\theta\).  Every declared
measurement statistic \(\operatorname{Tr}(\rho_\theta M_i)\) then
factors through \(T\), and the obstruction is strict contraction of
relative entropy by \(T\), which prevents recovery
\citep{Petz1986SufficientSubalgebras}.

\item[\emph{Noninterference for low observers.}] \textit{Special case.}
For a deterministic program \(P:H\to H\) on states with high-security (secret) and low-security (public) components, take the current quotient to be the low projection \(q:H\to H{\restriction}L\) of the initial state and the declared probes to be \(h\mapsto o(P(h))\) for the declared low observers \(o\) of the final state; \(P\) satisfies noninterference iff every such probe factors through \(q\).  Side-channel attacks --- timing, cache,
fault --- appear as family-extension bridges: the declared
low-observer family is enlarged to include an observer for which the
low-projection is no longer sufficient
\citep{GoguenMeseguer1982SecurityPolicies,SabelfeldMyers2003LanguageBasedInformationFlow}.
\end{description}
\endgroup

\subsection{Strict Extension [\texorpdfstring{\Fref{F8}}{F8}]}\label{sec:stability:strict-extension}

Let \(H\) be the carrier under discussion, let
\[
  q_{\mathrm{old}}:H\to Q_{\mathrm{old}}
\]
be the old quotient, and let
\[
  q_{\mathrm{new}}:H\to Q_{\mathrm{new}}
\]
be a candidate new quotient.  A new presentation is not
automatically a new layer: it may be a relabeling, a coarsening, or
an incomparable rewrite.  The Schur-calculus intuition from
\PaperAR\ \citep{Tsiokos2026XiCompanion} is that residuals compose
along refinement chains, so genuine layer growth preserves the old
quotient while accounting for the new split.

An extension is \emph{strict} when it retains the old quotient, so that the
old classes can be read off the new ones, and adds genuine growth, so that
it is not a relabeling of the old quotient. Precisely, the extension from
\(q_{\mathrm{old}}\) to \(q_{\mathrm{new}}\) \emph{retains} the old quotient
when \(q_{\mathrm{old}}=\rho\circ q_{\mathrm{new}}\) for some map \(\rho\), the
new quotient is \emph{definable from the old} when
\(q_{\mathrm{new}}=\lambda\circ q_{\mathrm{old}}\) for some map \(\lambda\), and
the extension is strict when it retains the old quotient and the new one is
not definable from the old. The lemma turns the second condition into a
split fiber and places strictness among the four possible comparisons.

\begin{lemma}[Strict Extension]\label{lem:F8}
Let \(q_{\mathrm{old}}\) be surjective. The extension from \(q_{\mathrm{old}}\) to \(q_{\mathrm{new}}\) is
strict if and only if there is a retention map
\[
  \rho:Q_{\mathrm{new}}\to Q_{\mathrm{old}},
  \qquad
  q_{\mathrm{old}}=\rho\circ q_{\mathrm{new}},
\]
and there are histories \(h,h'\in H\) such that
\[
  q_{\mathrm{old}}(h)=q_{\mathrm{old}}(h')
  \quad\text{but}\quad
  q_{\mathrm{new}}(h)\ne q_{\mathrm{new}}(h').
\]
Moreover, for any \(q_{\mathrm{old}}\) and \(q_{\mathrm{new}}\), surjective or not, exactly one of
four comparisons holds: \emph{equivalent} (retains and definable),
\emph{strict} (retains and not definable), \emph{coarsening} (definable and
not retains) or \emph{incomparable} (neither).
\end{lemma}
\begin{proof}
Assume the extension is strict.  By the retention part of strictness,
there is a map \(\rho:Q_{\mathrm{new}}\to Q_{\mathrm{old}}\) with
\(q_{\mathrm{old}}=\rho\circ q_{\mathrm{new}}\).  If no old fiber were
split, then \(q_{\mathrm{new}}\) would be constant on each fiber of
\(q_{\mathrm{old}}\); since \(q_{\mathrm{old}}\) is surjective, every
\(y\in Q_{\mathrm{old}}\) is \(q_{\mathrm{old}}(h)\) for some \(h\), and
\(\lambda(y):=q_{\mathrm{new}}(h)\) would be a well-defined map with
\(q_{\mathrm{new}}=\lambda\circ q_{\mathrm{old}}\), contradicting
non-definability.  So there are \(h,h'\in H\) with
\(q_{\mathrm{old}}(h)=q_{\mathrm{old}}(h')\) but
\(q_{\mathrm{new}}(h)\ne q_{\mathrm{new}}(h')\).

Conversely, suppose such a retention map and split old fiber are
given.  The equation \(q_{\mathrm{old}}=\rho\circ q_{\mathrm{new}}\)
says that every new class lies inside an old class, so the candidate
retains the old quotient rather than replacing it.  The displayed
pair \(h,h'\) shows that at least one old class has been split by
the new quotient, so no \(\lambda\) satisfies
\(q_{\mathrm{new}}=\lambda\circ q_{\mathrm{old}}\): it would give
\(q_{\mathrm{new}}(h)=\lambda(q_{\mathrm{old}}(h))=\lambda(q_{\mathrm{old}}(h'))=q_{\mathrm{new}}(h')\).
Thus the extension is strict.  The four comparisons are the four
combinations of the two conditions, so exactly one of them holds.
\end{proof}

\paragraph{Nonclaims.} In the strict case the lemma supplies a
retention map \(\rho\) and a split pair \((h,h')\), which a finite audit
can check.  It does not re-derive the Schur calculus,
and it does not classify every kind of layer growth.  It only gives
the quotient test for genuine retentive growth.

\subsection*{Layer instantiations}\label{layer-inst:F8}

\begingroup\small
\begin{description}
\item[\emph{Ring species (evolutionary biology).}] \textit{Analogy.}
For a ring-distributed population complex (the greenish-warbler ring around the Tibetan Plateau, the
\emph{Ensatina} salamander complex around California's Central
Valley), the local interbreeding equivalence
\(q_{\mathrm{old}}\) identifies adjacent populations that interbreed
in the wild --- yielding a single ring-wide class.  Reproductive compatibility is not
transitive around the ring, so it is not itself an equivalence; the
reproductive partition \(q_{\mathrm{new}}\) refines
\(q_{\mathrm{old}}\) by assigning the two reproductively isolated
terminal forms to distinct classes at a declared boundary along the
chain.  The
retention map is the inclusion of new classes into the single old
class, and the witness split is the closure pair, reproductively
isolated yet connected through the chain of locally interbreeding
intermediates \citep{IrwinBenschPrice2001SpeciationInARing}.

\item[\emph{Refinement mappings in formal verification (computer
science).}] \textit{Analogy.}
For a high-level specification \(S_2\) and a candidate implementation
\(S_1\), the Abadi--Lamport refinement-mapping theorem states that, when
\(S_1\) is machine-closed and \(S_2\) has finite invisible
nondeterminism and is internally continuous, \(S_1\) implements
\(S_2\) if and only if, possibly after adjoining history and prophecy
variables, there exists a state-level map
\(f:\mathrm{States}(S_1)\to\mathrm{States}(S_2)\) sending each \(S_1\)
step to a legal \(S_2\) step or a stuttering step and preserving
liveness.  The map \(f\) is the \Fref{F8} retention map;
the witness split is two implementation behaviors with identical
\(S_2\)-observable trace but distinct internal scheduling or
auxiliary-variable values
\citep{AbadiLamport1991RefinementMappings}.

\end{description}
\endgroup


\subsection{No Unstatused Residual [\texorpdfstring{\Fref{F9}}{F9}]}\label{sec:stability:no-unstatused-residual}

A formed stability package has a native residual registry: the
defects, gaps, budget failures, split pairs, critical pairs, or
endpoint failures that its own claims make relevant.  Each native
residual is not merely a number or a warning flag.  It comes with a
record saying which judgment family is responsible for it and which
status classifier applies.

\paragraph{Reading the statement.}
\(\mathcal S_{\mathrm{typed}}\) is the set of typed statuses the package
may assign: a status value, such as visible, hidden with a record, outside scope,
blocked, priced or accepted, tagged by the judgment family that assigns it.
The package keeps a set \(\mathrm{Rec}\) of residual records; each record
\(\rho\) names a residual \(\operatorname{loc}(\rho)\in R\), its locus, and
carries a typed status \(\operatorname{st}(\rho)\in\mathcal S_{\mathrm{typed}}\).
A residual \(r\) has a typed status when it is the locus of some record, and
it is \emph{unstatused}, \(\operatorname{UnstatusedResidual}_{\mathcal P}(r)\),
when it is the locus of no record.

\begin{theorem}[No Unstatused Residual]\label{thm:F9}
Let \(R\) be a residual type and let \(\mathcal P\) be a stability package over
\(R\) (in the intended reading, a formed one) with residual records \(\mathrm{Rec}\), locus map
\(\operatorname{loc}:\mathrm{Rec}\to R\) and status map
\(\operatorname{st}:\mathrm{Rec}\to\mathcal S_{\mathrm{typed}}\). Assume
the residual-coverage rule (RC) of \Cref{sec:apparatus}, so every residual is
the locus of a record, and assume records with the same locus carry the same
status: \(\operatorname{loc}(\rho)=\operatorname{loc}(\rho')\) implies
\(\operatorname{st}(\rho)=\operatorname{st}(\rho')\). (This holds in particular
when records are unique at their locus.) Then, for every \(r\in R\):
\begin{enumerate}[label=(\roman*)]
\item \(r\) has a typed status;
\item all records with locus \(r\) carry the same typed status, so
  \(\operatorname{st}_{\mathcal P}(r):=\operatorname{st}(\rho)\) for any
  record \(\rho\) with locus \(r\) defines a map
  \(\operatorname{st}_{\mathcal P}:R\to\mathcal S_{\mathrm{typed}}\);
\item \(\neg\operatorname{UnstatusedResidual}_{\mathcal P}(r)\).
\end{enumerate}
Clauses (i) and (iii) restate (RC) at \(r\); clause (ii) uses the agreement
of statuses at a common locus. Both hypotheses are needed: under (RC), a map
\(\operatorname{st}_{\mathcal P}:R\to\mathcal S_{\mathrm{typed}}\) with
\(\operatorname{st}_{\mathcal P}(\operatorname{loc}\rho)=\operatorname{st}(\rho)\)
for every record exists if and only if records with a common locus carry the
same status; and (RC) is equivalent to the absence of unstatused residuals.
\end{theorem}
\begin{proof}
Let \(r\in R\). By (RC) there is a record \(\rho\) with
\(\operatorname{loc}(\rho)=r\), and \(\operatorname{st}(\rho)\) is a typed
status of \(r\); this is (i). If \(\rho'\) is another record with
\(\operatorname{loc}(\rho')=r\), the agreement hypothesis gives
\(\operatorname{st}(\rho')=\operatorname{st}(\rho)\); this
is (ii), and \(\operatorname{st}_{\mathcal P}(r)\) does not depend on the
record chosen. Finally, \(\operatorname{UnstatusedResidual}_{\mathcal P}(r)\)
says that \(r\) is the locus of no record, which contradicts (RC); this is
(iii). For the converses, if a map \(\operatorname{st}_{\mathcal P}\) with
\(\operatorname{st}_{\mathcal P}(\operatorname{loc}\rho)=\operatorname{st}(\rho)\)
exists and \(\operatorname{loc}(\rho)=\operatorname{loc}(\rho')\), then
\(\operatorname{st}(\rho)=\operatorname{st}_{\mathcal P}(\operatorname{loc}\rho)
=\operatorname{st}_{\mathcal P}(\operatorname{loc}\rho')=\operatorname{st}(\rho')\);
and (RC) says exactly that every residual is the locus of some record, that is,
that no residual is unstatused.
\end{proof}

\paragraph{Nonclaims.} The theorem does not
enumerate or redefine the inherited status families; that belongs to
the admissibility apparatus.  Formedness is not used: the conclusions
hold for any locus and status maps satisfying (RC) and status agreement.

\subsection*{Layer instantiations}\label{layer-inst:F9}

\begingroup\small
\begin{description}
\item[\emph{{GHS} chemical hazard classification (regulatory
registry).}] \textit{Analogy.}
GHS classification records the applicable hazard category or subcategory
using class-specific criteria, tiered assessment and weight of evidence.
Conflicting data are resolved by the applicable assessment rules. For a declared
registry, take \(R\) to be the hazard endpoints under review, \(\mathrm{Rec}\)
its endpoint-status records, \(\operatorname{loc}(\rho)\) the endpoint assessed and
\(\operatorname{st}(\rho)\) its recorded assessment outcome, including a hazard
category or an explicit nonclassification or insufficient-data status.
\Fref{F9} applies when the registry records every endpoint, satisfying (RC),
and all records at a common endpoint agree. Coverage and agreement are
conditions on this registry. EU CLP and
OSHA HazCom implement GHS-based classification requirements
\citep{UN2025GHSRev11}.

\end{description}
\endgroup


\subsection*{Stability at a glance}\label{sec:stability:summary}\addcontentsline{toc}{subsection}{Stability at a glance}
\begin{table}[H]
  \centering
  \footnotesize
  \begin{tabularx}{\textwidth}{@{}
      >{\raggedright\arraybackslash}p{0.24\textwidth}
      >{\raggedright\arraybackslash}X
      >{\raggedright\arraybackslash}p{0.12\textwidth}@{}}
    \toprule
    Name & One-line claim & Substrate cite \\
    \midrule
    No-Needles Theorem (\Fref{F6}) &
    Under the block symmetry conditions, the projection identity
    \(K_{LL}K_{LL}^{\dagger}K_{LD}=K_{LD}\) and its three budgets, no transported
    dissolving direction exceeds its declared budget (no needle).  &
    \PaperNeedleKiller, \PaperNS \\
    Sufficiency Closure Theorem (\Fref{F7}) &
    For surjective \(q\), \(q\) is sufficient for declared futures iff each future
    probe descends through \(q\).  &
    \textemdash \\
    Strict Extension Lemma (\Fref{F8}) &
    For surjective \(q_{\mathrm{old}}\), an extension is strict iff it
    retains the old quotient and splits an old fiber.  &
    \PaperAR \\
    No-Unstatused-Residual Theorem (\Fref{F9}) &
    Under residual coverage and agreement of statuses at a common locus,
    the typed status is a well-defined map on residuals and no residual is
    unstatused; under coverage, status agreement is equivalent to a
    well-defined status map, and coverage is equivalent to no unstatused
    residual.  &
    \textemdash \\
    \bottomrule
  \end{tabularx}
  \caption{Stability at a glance: the four laws of \Cref{sec:stability}.}
  \label{tab:axis-stability-summary}
\end{table}

\FloatBarrier
           % sec:stability
\section{Transport}\label{sec:transport}

The Transport axis is about moving structure across quotients.  Its
basic question is the familiar one: when a map, comparison, or family
of local data is defined before quotienting, does it still make sense
after the quotient has forgotten some distinctions?  This is the
easiest entry point into the catalog because the laws are closest to
standard descent, gluing, and quotient comparison.  The closure
apparatus supplies the bookkeeping, but the mathematical shape is the
ordinary universal-property shape: a construction is legitimate only
when it respects the equivalence relation through which it is being
read.  The running example of \Cref{sec:intro} is the basic instance:
\(F\) does not descend through \(q\) because of the split pair
\((1,2)\), and its two repairs are those of the Descent--Repair Normal
Form (\Fref{F2}).

The chapter has four direct laws.  Descent--Repair Normal Form is
the first theorem of the chapter: it makes formal the scenario from the
introduction, where a
map fails to descend through a quotient and the obstruction is exactly
the split-pair set \(\splitObs\).  Holonomy-Memory Repair Normal Form
then treats the case where two routes look identical at the current
quotient \(\currQ\) but leave different residue for a predictive
quotient \(\predM\).  Local--Global Obstruction Normal Form records
the gluing situation: a family that globalizes is compatible, every local
family has exactly one of four statuses, and every compatible family glues
uniquely exactly when the restriction map is injective with image the
compatible families (the sheaf condition).  Recombination
Comparison treats branch-merge transport, where branchwise equality
need not determine merged behavior unless the recombination quotient
factors through the branchwise quotient.

All four laws are direct quotient-theoretic theorems, and the chapter
depends only on \Cref{sec:apparatus}.  It is also a toolbox for the
rest of the paper: the split-pair repairs of the Descent--Repair Normal
Form reappear in the recombination law and later in the Symmetry and
Self-Reference and Limits chapters.  Readers new to the framework
should read this chapter before choosing another.

\subsection{Descent--Repair Normal Form [\texorpdfstring{\Fref{F2}}{F2}]}\label{sec:transport:descent-repair}

The first transport law is the quotient descent criterion already
introduced in \Cref{sec:intro}.  We use the split-pair obstruction of
\Cref{def:split-pair}.

\begin{definition}[Canonical target-repair quotient]\label{def:target-repair}
Let \(q:X\to Q\), \(F:X\to Y\), and \(r:Y\to R\).  The canonical
target-repair relation \(\equiv_F\) on \(R\) is the smallest
equivalence relation generated by the identifications
\[
  r(Fx)\equiv_F r(Fx')
  \qquad\text{whenever}\qquad
  q(x)=q(x').
\]
The \emph{canonical target-repair quotient} is
\[
  R_F:=R/{\equiv_F},
\]
with quotient map \(c_F:R\to R_F\).
\end{definition}

\begin{theorem}[Descent-Repair Normal Form]\label{thm:F2}
Let \(q:X\to Q\) be a surjective quotient map, let \(F:X\to Y\) be a map, and
let \(r:Y\to R\) be a target readout.  Then \(r\circ F\) descends
through \(q\) if and only if
\[
  \splitObs(q,F,r)=\emptyset .
\]
The two canonical repairs are the following. The source repair
\(\langle q,r\circ F\rangle:X\to Q\times R\) is the coarsest refinement of
\(q\) through which \(r\circ F\) descends: it refines \(q\), \(r\circ F\)
descends through it, and it factors through every refinement of \(q\)
through which \(r\circ F\) descends: if \(q=a\circ q'\) and
\(r\circ F=b\circ q'\) for maps \(q'\), \(a\) and \(b\), then
\(\langle q,r\circ F\rangle=\langle a,b\rangle\circ q'\). The target repair is the quotient
\(c_F:R\to R_F\) of \Cref{def:target-repair}: a map \(c:R\to C\) makes
\(c\circ r\circ F\) descend through \(q\) exactly when \(c\) factors through
\(c_F\), so \(R_F\) is the finest coarsening of \(R\) that repairs descent.
\end{theorem}

\begin{proof}
Suppose first that \(r\circ F\) descends through \(q\).  Then there is
a map \(\overline F:Q\to R\) such that
\(\overline F\circ q=r\circ F\).  If \(q(x)=q(x')\), then applying
\(\overline F\) gives
\[
  r(Fx)=\overline F(qx)=\overline F(qx')=r(Fx').
\]
Thus no pair \((x,x')\) with \(q(x)=q(x')\) can be separated by
\(rF\), so \(\splitObs(q,F,r)=\emptyset\).

Conversely, assume \(\splitObs(q,F,r)=\emptyset\).  Then whenever
\(q(x)=q(x')\), the pair \((x,x')\) is not a split pair, so
\(r(Fx)=r(Fx')\).  Therefore \(r\circ F\) is constant on the fibers
of \(q\).  By the quotient universal property, there is a unique map
\(\overline F:Q\to R\) satisfying
\(\overline F\circ q=r\circ F\).

For the source repair, the first projection gives
\(q=\operatorname{pr}_Q\circ\langle q,r\circ F\rangle\), so
\(\langle q,r\circ F\rangle\) refines \(q\), and the second projection gives
\(r\circ F=\operatorname{pr}_R\circ\langle q,r\circ F\rangle\), so
\(r\circ F\) descends through it. If \(q'\) refines \(q\), say
\(q=a\circ q'\), and \(r\circ F=b\circ q'\), then
\(\langle q,r\circ F\rangle=\langle a,b\rangle\circ q'\); hence
\(\langle q,r\circ F\rangle\) factors through every refinement of \(q\)
through which \(r\circ F\) descends, and is the coarsest one.  For the
target repair, quotient \(R\) by the generated
relation \(\equiv_F\).  If \(q(x)=q(x')\), then
\(r(Fx)\equiv_F r(Fx')\), hence
\[
  c_F(r(Fx))=c_F(r(Fx')).
\]
Thus \(c_F\circ r\circ F\) is constant on \(q\)-fibers and descends
through \(q\).

It remains to check minimality.  Let \(c:R\to C\) be any target
coarsening such that \(c\circ r\circ F\) descends through \(q\).  Then
for every generator \(r(Fx)\equiv_F r(Fx')\), the equality
\(q(x)=q(x')\) implies
\[
  c(r(Fx))=c(r(Fx')).
\]
Since equality under \(c\) is an equivalence relation on \(R\), it
contains the equivalence relation generated by those identifications.
Therefore \(c\) is constant on \(\equiv_F\)-classes and factors
uniquely through \(c_F:R\to R_F\).  Conversely, if \(c=d\circ c_F\), then
\(c\circ r\circ F=d\circ(c_F\circ r\circ F)\) descends through \(q\),
because \(c_F\circ r\circ F\) does.  Hence \(R_F\) is the finest
coarsening that repairs descent.
\end{proof}

\subsection*{Layer instantiations}\label{layer-inst:F2}

\begingroup\small
\begin{description}
\item[\emph{Coequalizer in a category (mathematics; target
coarsening).}] \textit{Special case.}
In the category of sets, the coequalizer of a parallel pair
\(u, v : A \rightrightarrows B\) is the universal arrow
\(c : B \to C\) with \(c \circ u = c \circ v\), the quotient of \(B\)
by the equivalence generated by \(u(a)\sim v(a)\); in a general
category a coequalizer need not be such a quotient. Given a candidate target readout \(r : B \to R\), the
canonical minimal target coarsening that restores descent through
the \(\langle u, v\rangle\)-generated equivalence on \(B\) is
\(R \to \operatorname{coeq}(r \circ u, r \circ v)\): the
coequalizer of the parallel pair pushed forward into \(R\).
\citep{MacLane1998Categories}.

\item[\emph{Gauge-invariant observables (physics; source
refinement).}] \textit{Special case.}
In a gauge theory the gauge group \(G\) is given a priori by the
theory's symmetry structure, and the orbit quotient
\(q : X \to X/G\) on field configurations is the discipline-fixed
source quotient. A candidate observable \(r : X \to \mathbb{R}\)
descends through \(q\) iff it is gauge-invariant. \Fref{F2}'s canonical
source refinement is the joint quotient
\(\langle q, r\rangle : X \to (X/G) \times \mathbb{R}\); a gauge fixing, a section \(s\) of \(q\), gives instead the descended readout \(r\circ s\circ q\), which agrees with \(r\) on the chosen representatives. \Fref{F2}'s target
coarsening of \(r\) is the quotient of \(\mathbb{R}\) by the
equivalence generated by \(r(x)\sim r(g\cdot x)\), and the
gauge-invariant observables are exactly those for which this
quotient map is a bijection
\citep{PeskinSchroeder1995QFT}.

\item[\emph{Hash collisions in dictionary lookup (computer
science; source refinement).}] \textit{Special case.}
A hash function \(q : K \to B\), corestricted to its set \(q(K)\) of attained
buckets, is a surjective source quotient on a key universe. A downstream lookup operation
\(r : K \to R\) descends through \(q\) iff no two collision-merged
keys carry distinct intended values. When they do, \Fref{F2}'s canonical
source refinement is the joint quotient
\(\langle q, r\rangle : K \to B \times R\); FKS two-level perfect
hashing refines further: its secondary hash is injective on the
stored keys of each bucket, so on the stored key set it separates
every collision and \(\langle q, r\rangle\) factors through it
\citep{CormenLeisersonRivestStein2009CLRS}.

\item[\emph{Lehmann--Scheff\'e minimal sufficient statistic
(statistics; source refinement).}] \textit{Special case.}
On a countable sample space \(X\), let \(L(\theta;x)\) be probability
masses, and assume some \(\theta_0\) has \(L(\theta_0;x)>0\) for every \(x\) (as when all members share one support; no such \(\theta_0\) exists for the uniform laws on \(\{1,\dots,\theta\}\), \(\theta\in\mathbb N\)); fix such a \(\theta_0\).
Regard a candidate statistic \(T\) as a surjection onto \(T(X)\), and ask
whether \(\lambda(x):=\bigl(\theta\mapsto L(\theta;x)/L(\theta_0;x)\bigr)\)
descends through \(T\). When \(T\) is not sufficient, the \Fref{F2} split pairs are sample
pairs with the same \(T\) but distinct values of \(\lambda\). \Fref{F2}'s
canonical source refinement is the joint quotient \(\langle T,\lambda\rangle\);
\(\lambda\) has the fibers of the minimal sufficient statistic \(T^{*}\), so
when \(T\) is coarser than \(T^{*}\) this joint quotient has the fibers of
\(T^{*}\) (Lehmann--Scheff\'e)
\citep{LehmannScheffe1950Completeness}.

\item[\emph{Genetic code, codon-to-amino-acid quotient (biology;
source refinement).}] \textit{Analogy.}
The standard genetic code is a discipline-given quotient
\(q : \mathrm{Codons} \to \mathrm{AminoAcids}\) (61 sense codons
to 20 amino acids). A codon-level cellular readout such as
translation efficiency \(\tau : \mathrm{Codons} \to \mathbb{R}\)
does not in general descend through \(q\): synonymous codons can
differ in translation rate due to tRNA-pool composition (the
\emph{codon usage bias} literature). \Fref{F2}'s canonical source
refinement is the joint quotient
\(\langle q, \tau\rangle : \mathrm{Codons} \to \mathrm{AminoAcids}
\times \mathbb{R}\); codon usage tables and codon-adaptation
indices are operational approximations \citep{Crick1968GeneticCode}.
\end{description}
\endgroup

\subsection{Holonomy-Memory Repair Normal Form [\texorpdfstring{\Fref{F3}}{F3}]}\label{sec:transport:holonomy-memory}

The second transport law concerns route dependence: two routes can be
indistinguishable at the current quotient while still leaving a
predictive difference.

\begin{definition}[Route residue]\label{def:route-residue}
Let \(H\) be a carrier, let \(\gamma,\eta:H\to Z\) be two declared
routes, let \(\currQ:Z\to Q\) be the current quotient, and let
\(\predM:Z\to M\) be a predictive quotient (refinement of \(\currQ\) is not
assumed).  The \emph{route residue} of \(\gamma\) and \(\eta\) at
\(h\in H\) is present when
\[
  \currQ(\gamma h)=\currQ(\eta h)
  \quad\text{but}\quad
  \predM(\gamma h)\ne\predM(\eta h).
\]
\end{definition}

Say that \emph{memory is forced} at \(h\in H\) when no map
\(g:Q\to M\) satisfies \(g(\currQ(z))=\predM(z)\) for both route outputs
\(z\in\{\gamma h,\eta h\}\): the current class alone does not determine the
predictive class of the two outputs.

\begin{theorem}[Holonomy-Memory Repair Normal Form]\label{thm:F3}
For two declared routes \(\gamma,\eta:H\to Z\), current quotient
\(\currQ:Z\to Q\), and predictive quotient \(\predM:Z\to M\), memory is forced
at \(h\) exactly when there is route residue at \(h\). The memory repair
\(\langle\currQ,\predM\rangle:Z\to Q\times M\) refines \(\currQ\),
\(\predM\) descends through it, and it is the coarsest such refinement:
for every refinement \(q'\) of \(\currQ\) through which \(\predM\) descends,
equality under \(q'\) implies equality under \(\langle\currQ,\predM\rangle\).
In particular, route outputs with route residue receive different
repaired classes.
\end{theorem}

\begin{proof}
If \(\currQ(\gamma h)\ne\currQ(\eta h)\), a map \(g\) taking the two current
classes to \(\predM(\gamma h)\) and \(\predM(\eta h)\) exists, so memory is
not forced; there is no route residue either. If
\(\currQ(\gamma h)=\currQ(\eta h)\), any such \(g\) gives
\(\predM(\gamma h)=g(\currQ(\gamma h))=g(\currQ(\eta h))=\predM(\eta h)\), so a
suitable \(g\) exists exactly when \(\predM(\gamma h)=\predM(\eta h)\). Hence
memory is forced at \(h\) exactly when \(\currQ(\gamma h)=\currQ(\eta h)\) and
\(\predM(\gamma h)\ne\predM(\eta h)\), which is route residue at \(h\).

Write \(m=\langle\currQ,\predM\rangle\). Then \(\currQ=\operatorname{pr}_Q\circ m\)
and \(\predM=\operatorname{pr}_M\circ m\). If \(\currQ=a\circ q'\) and
\(\predM=b\circ q'\), then \(m=\langle a,b\rangle\circ q'\); thus equality under
\(q'\) implies equality under \(m\). Route residue separates the
\(\predM\)-values of the outputs, hence their \(m\)-values.
The memory repair is the source repair of \Cref{thm:F2} with \(q=\currQ\), \(F=\mathrm{id}_Z\) and \(r=\predM\); that argument does not use surjectivity of \(q\).
\end{proof}

\subsection*{Layer instantiations}\label{layer-inst:F3}

\begingroup\small
\begin{description}
\item[\emph{Holonomy of a connection on a fiber bundle
(mathematics).}] \textit{Special case.}
Given a principal \(G\)-bundle \(E\to M\) with a principal connection \(\nabla\), choose a point of \(E_x\) to identify holonomy with a subgroup of \(G\), and take two
piecewise-smooth paths \(\gamma, \eta\) in \(M\) with common
endpoints, assume parallel transport along both paths is defined on every
point of the initial fiber; take \(H=E_x\), the fiber over the common initial point \(x\); the
routes \(P_\gamma,P_\eta:E_x\to Z=E\) (parallel transport); the current
quotient the projection \(\pi:E\to M\) (here the base \(M\) plays the role of \(Q\) in \Cref{thm:F3}, and the predictive codomain is \(E\)); and the predictive quotient
\(\mathrm{id}_E\). Both routes land over the common endpoint, and route
residue at \(e\) is \(P_\gamma e\ne P_\eta e\). Route residue is detected by
the holonomy
\(P_\eta^{-1} \circ P_\gamma \in \operatorname{Hol}(\nabla, x)\),
defined on based loops at \(x\); for a flat connection it depends
only on homotopy classes and gives the holonomy representation
\(\pi_1(M, x) \to G\), with \(G\) the structure group
\citep{KobayashiNomizu1963Foundations}.

\item[\emph{Berry phase in adiabatic quantum mechanics (physics).}] \textit{Special case.}
For a smooth family \(H(\lambda)\) with a nondegenerate eigenvalue
\(E_n(\lambda)\), let \(Z\) be the eigenline bundle. The routes send a unit
eigenvector \(\psi_0\) at \(\lambda_0\) to its parallel transports along two
loops \(\gamma,\eta\) based at \(\lambda_0\), for the Berry connection
\(\mathbf A = i\langle \psi \mid \nabla \psi\rangle\); take \(H=\{\psi_0\}\). The current quotient is the base point, and the predictive readout
is the transported vector. The outputs are \(e^{i\gamma_B(\gamma)}\psi_0\) and
\(e^{i\gamma_B(\eta)}\psi_0\), so route residue is present exactly when the
Berry phases differ modulo \(2\pi\). Adiabatic evolution realizes this
transport up to the dynamical phase \citep{Berry1984QuantalPhase}.

\item[\emph{Event sourcing in distributed systems (computer
science).}] \textit{Analogy.}
In an event-sourced system the authoritative state is an
append-only log of domain events; the current materialized snapshot
is a fold over that log.  Two distinct event histories can produce
identical snapshots while remaining distinguishable under
retroactive policy replay, audit, or future business rules — the
enriched \((\text{snapshot}, \text{retro\,results})\) readout
differs even when the snapshot agrees.  The event log itself is the
canonical memory witness retained precisely because the snapshot
is lossy \citep{Kleppmann2017DDIA}.

\item[\emph{Epigenetic memory of cell fate (biology).}] \textit{Analogy.}
Two somatic cells derived from a common zygote share an essentially
identical genome (the current quotient) but exhibit divergent
transcriptional responses to the same regulatory cues.  The
enriched
\((\text{genome}, \text{chromatin}, \text{predicted response})\)
readout differs because chromatin marks (DNA methylation, histone
modifications, bivalent domains) encode the developmental lineage.
The chromatin state, propagated through mitosis, is the canonical
memory witness \citep{Reik2007Epigenetic}.

\item[\emph{Preisach hysteresis and return-point memory
(engineering).}] \textit{Special case.}
The current magnetization \(M(t)\) of a hysteretic material is
reachable by many distinct input-field histories \(H(t)\); two
histories arriving at the same \((H, M)\) state predict divergent
responses to the next field reversal. Take the history carrier a one-point
set, \(Z\) the set of input histories, the routes choosing two histories, the
current quotient the present pair \((H,M)\) (here \(H\) and \(M\) are the applied field and the magnetization, not the carrier and the predictive codomain of \Cref{thm:F3}), and the predictive quotient the
output response to every future input continuation.  The Preisach model's
canonical memory witness is the \emph{return-point staircase}: the
stack of dominant past extrema of the input history that the
Preisach plane records.  Mayergoyz's wiping-out property
establishes this staircase as the residue that determines future
responses, and for a Preisach density positive on the Preisach
half-plane it is the minimal residue distinguishing
otherwise-equivalent presents \citep{Mayergoyz2003Hysteresis}.
\end{description}
\endgroup

\subsection{Local--Global Obstruction Normal Form [\texorpdfstring{\Fref{F4}}{F4}]}\label{sec:transport:local-global}

The third transport law is the local-to-global image criterion.  It
separates overlap tests from genuine global existence.

\begin{definition}[Local-global compatibility]\label{def:local-global-compat}
Let \(I\) index a declared patch system (finite in the intended reading), with local object sets
\((Q_i)_{i\in I}\).  Let \(\mathcal C\subseteq \prod_i Q_i\) be the
subset of local families satisfying all declared overlap comparisons.
Let \(Q\) be the global object set and
\[
  R:Q\to \prod_i Q_i
\]
the declared restriction map.  A local family \(a\in\prod_i Q_i\) is
\emph{compatible} if \(a\in\mathcal C\), and it \emph{globalizes} if
\(a\in\operatorname{im}(R)\).
\end{definition}

\paragraph{Reading the statement.}
A local family \(a\) is \emph{overlap-inconsistent} when
\(\operatorname{Comp}(a)\) fails; \emph{compatible non-global} when it
is compatible but no global state restricts to it; \emph{globally
ambiguous} when at least two global states restrict to it; and
\emph{globally exact} when exactly one does.
\(\operatorname{ExactlyOneStatus}(a)\) says that exactly one of these
four statuses holds for \(a\). In the theorem, \(\mathcal G\) plays the role
of \(Q\), \(\mathcal F\) that of \(\prod_i Q_i\), and
\(\operatorname{Comp}(a)\) that of membership \(a\in\mathcal C\) in
\Cref{def:local-global-compat}.

\begin{theorem}[Local-Global Obstruction Normal Form]\label{thm:F4}
Let \(\mathcal G\) be a set of global states, let \(\mathcal F\) be a
set of local-family data, let
\(\operatorname{Comp}:\mathcal F\to\mathrm{Prop}\) be the
overlap-compatibility predicate, and let
\[
  R:\mathcal G\longrightarrow\mathcal F
\]
be a declared restriction map such that
\(\operatorname{Comp}(R(g))\) for every \(g\in\mathcal G\).  For
\(a\in\mathcal F\), write
\[
  \operatorname{Globalizes}_R(a)
  \quad:\Longleftrightarrow\quad
  \exists g\in\mathcal G,\ R(g)=a,
\]
and define the four statuses
\begin{align*}
  \operatorname{OverlapInconsistent}(a)&:\iff\neg\operatorname{Comp}(a),\\
  \operatorname{CompatibleNonGlobal}(a)&:\iff\operatorname{Comp}(a)\wedge a\notin R(\mathcal G),\\
  \operatorname{GloballyAmbiguous}(a)&:\iff\lvert R^{-1}(a)\rvert\ge 2,\\
  \operatorname{GloballyExact}(a)&:\iff\lvert R^{-1}(a)\rvert=1.
\end{align*}
Then \(\operatorname{Globalizes}_R(a)\) implies \(\operatorname{Comp}(a)\), so
\(a\) globalizes exactly when it is compatible and lies in \(R(\mathcal G)\),
and \(\operatorname{ExactlyOneStatus}(a)\) holds: exactly one of the
four statuses holds for \(a\). Moreover, every \(a\) with \(\operatorname{Comp}(a)\)
is globally exact if and only if \(R\) is injective and
\(R(\mathcal G)=\{a:\operatorname{Comp}(a)\}\); in the setting of
\Cref{def:local-global-compat}, with \(\operatorname{Comp}(a)\) the agreement of
\(a\) on all declared overlaps and \(R\) the restriction map of a presheaf of
sets, this is the sheaf (equalizer) condition for the cover, and
\(\operatorname{CompatibleNonGlobal}\) and \(\operatorname{GloballyAmbiguous}\)
are the two ways it fails at \(a\).
\end{theorem}

\begin{proof}
Assume the compatibility predicate \(\operatorname{Comp}\), the
restriction map \(R\), and the hypothesis
\(\operatorname{Comp}(R(g))\) for every \(g\in\mathcal G\).  Fix
\(a\in\mathcal F\).  If \(a\) globalizes, then \(a=R(g)\) for some
\(g\in\mathcal G\), and the hypothesis gives \(\operatorname{Comp}(a)\).
Since \(\operatorname{Globalizes}_R(a)\) means \(a\in R(\mathcal G)\), \(a\)
globalizes exactly when it is compatible and lies in \(R(\mathcal G)\). This
proves the local-to-global image criterion.

The status enumeration is a case split on compatibility, image
membership, and the fiber \(R^{-1}(a)\).  If
\(\neg\operatorname{Comp}(a)\), then \(a\) is overlap-inconsistent.
If \(\operatorname{Comp}(a)\) and \(a\notin R(\mathcal G)\), then
\[
  R^{-1}(a)=\emptyset,
\]
so \(a\) is compatible non-global.  If
\(\operatorname{Comp}(a)\), \(a\in R(\mathcal G)\), and
\(|R^{-1}(a)|>1\), then \(a\) is globally ambiguous.  If
\(\operatorname{Comp}(a)\), \(a\in R(\mathcal G)\), and
\(|R^{-1}(a)|=1\), then \(a\) is globally exact.  These four cases
are exhaustive. They are mutually exclusive: CompatibleNonGlobal(\(a\)) includes \(\operatorname{Comp}(a)\), and GloballyAmbiguous(\(a\)) and GloballyExact(\(a\)) give \(R^{-1}(a)\ne\emptyset\), hence \(\operatorname{Comp}(a)\) by the hypothesis \(\operatorname{Comp}(R(g))\); so none of the last three holds together with \(\operatorname{OverlapInconsistent}(a)\). The second needs
\(R^{-1}(a)=\emptyset\), and the last two differ in \(|R^{-1}(a)|\).
For the last clause, if every compatible \(a\) is globally exact, then every
compatible \(a\) lies in \(R(\mathcal G)\), which is contained in the compatible
set by the first part, and \(R(g)=R(g')=a\) makes \(a\) compatible, hence
exact, so \(g=g'\). Conversely, if \(R\) is injective with image the compatible
set, every compatible \(a\) has exactly one preimage.
\end{proof}

\subsection*{Layer instantiations}\label{layer-inst:F4}

\begingroup\small
\begin{description}
\item[\emph{\v{C}ech cohomology and sheaf gluing (mathematics).}] \textit{Special case.}
For an abelian presheaf \(\mathcal F\) on a space \(X\) with finite
open cover \(\{U_i\}\), restriction sends a global section to a family
of local sections that agree pairwise on overlaps, the 0-cocycle
condition. In the slots of \Cref{thm:F4}, the global states are \(\mathcal F(X)\), the local-family data are \(\prod_i\mathcal F(U_i)\), \(\operatorname{Comp}\) is agreement on every \(U_i\cap U_j\), and \(R\) is restriction; in this item \(\mathcal F,\mathcal G,\mathcal K\) denote sheaves, not the sets of the theorem.  The four \Fref{F4} statuses correspond to families that
disagree on some overlap, compatible families that glue to no global
section, compatible families that glue to several, and families that
glue uniquely; the sheaf axioms exclude the middle two, so for a
sheaf every compatible family glues uniquely.  For sheaves the
obstruction moves one level up: given a surjection of sheaves
\(\mathcal F\to\mathcal G\) with kernel \(\mathcal K\), a global
section of \(\mathcal G\) has local lifts on a fine enough cover, and
their differences on overlaps form a \v{C}ech 1-cocycle in
\(\mathcal K\) whose class in \(\check H^1\) obstructs a compatible
choice of lifts (canonical example: the section \(z\) of
\(\mathcal O^*\) on \(\mathbb C^*\) has local logarithms but no global
one)
\citep{Hartshorne1977AlgebraicGeometry}.

\item[\emph{Topological defects in ordered media (physics).}] \textit{Analogy.}
In a system with order-parameter space \(M = G/H\), patchwise
order-parameter sections must respect the group law of \(\pi_n(M)\)
on defect linking charges across patch overlaps.  Forbidden total
charges (compatible local but no global texture), non-abelian
disclination entanglements (multiple distinct globals with
identical local tabulation), and defect-free uniform textures
illustrate the four-status partition in condensed-matter physics
\citep{Mermin1979TopologicalDefects}.

\end{description}
\endgroup


\subsection{Recombination Comparison [\texorpdfstring{\Fref{F5}}{F5}]}\label{sec:transport:recombination}

The final transport law compares branchwise equality with equality
after a declared branch merge.

\begin{definition}[Bridge composition]\label{def:bridge-composition}
Let \(B\) be the carrier of branch data, let \(K:B\to Q_K\) be the
branchwise quotient, and let \(R:B\to Q_R\) be the recombination
quotient.  The \emph{bridge composition} is the use of \(K\) as the
source quotient and \(R\) as the target readout in the split-pair
criterion of \Cref{def:split-pair}.  Its obstruction is
\[
  \{(b,b')\in B\times B:
      K(b)=K(b')\ \text{and}\ R(b)\ne R(b')\}.
\]
\end{definition}

\begin{theorem}[Recombination Comparison]\label{thm:F5}
Let \(K:B\to Q_K\) be a surjective branchwise quotient and
\(R:B\to Q_R\) a surjective recombination quotient on the same branch
carrier.
Then branchwise equality determines recombination equality if and only
if \(R\) factors through \(K\), that is, if and only if there exists
\(\rho:Q_K\to Q_R\) such that
\[
  R=\rho\circ K .
\]
Equivalently, the branch-merge split-pair obstruction of
\Cref{def:bridge-composition} is empty. Symmetrically, \(K\) factors through
\(R\) if and only if the reverse obstruction
\(\{(b,b'):R(b)=R(b'),\ K(b)\ne K(b')\}\) is empty. Hence exactly one of four
statuses holds: \emph{collapse} (both obstructions empty, so \(K\) and \(R\)
have the same fibers); \emph{recombination-sensitive refinement} (only the
branch-merge obstruction nonempty); \emph{washout} (only the reverse
obstruction nonempty); and \emph{incomparable} (both nonempty).
\end{theorem}

\begin{proof}
If \(R=\rho\circ K\), then any two branch records with
\(K(b)=K(b')\) satisfy
\[
  R(b)=\rho(K(b))=\rho(K(b'))=R(b').
\]
Thus the bridge-composition obstruction is empty: branchwise equality
determines equality after recombination.

Conversely, assume the obstruction is empty.  Then \(R\) is constant
on the fibers of \(K\): whenever \(K(b)=K(b')\), emptiness of the
obstruction gives \(R(b)=R(b')\).  By the quotient universal property,
there is a unique map \(\rho:Q_K\to Q_R\) with \(R=\rho\circ K\).
This proves the factorization criterion. The reverse clause is the same
argument with the roles of \(K\) and \(R\) exchanged, using that \(R\) is
surjective, and the four statuses are the four combinations of the two
emptiness conditions.

This is exactly \Cref{thm:F2} applied with \(K\) as source quotient,
the identity branch transport as the pre-quotient map, and \(R\) as
target readout.  A recombination failure is therefore not a one-sided
strict-refinement slogan; it is a two-sided quotient comparison
witnessed by a split pair at the merged level.
\end{proof}

\subsection*{Layer instantiations}\label{layer-inst:F5}

\begingroup\small
\begin{description}
\item[\emph{Meiotic recombination and genetic linkage (biology).}] \textit{Special case.}
In meiosis the parental homolog from which a gamete's chromosome
descends at a reference locus (branchwise quotient \(K\), corestricted to the homologs that occur) determines the
gamete genotype at a linked locus (recombination quotient \(R\), corestricted to the genotypes that occur), for a parent
heterozygous at both loci, iff all gametes that inherit the same parental
homolog at the reference locus have the same parity of crossovers between the
two loci (in particular, when no gamete carries an odd number).  The bridge-composition
obstruction is the \emph{recombinant gamete}: two gametes whose
chromosomes descend from the same parental homolog at the reference
locus but carry distinct genotypes at the linked locus
\citep{Morgan1911Coupling}.

\item[\emph{Church--Rosser confluence (computer science / logic).}] \textit{Special case.}
For a terminating term-rewriting system, take as branch carrier the pairs
\((s,\rho)\) of a term and a maximal reduction sequence from it, with
\(K(s,\rho)=s\) and \(R(s,\rho)\) the normal form at which \(\rho\) ends, valued in the set of normal forms (each normal form \(n\) is \(R(n,\varnothing)\), so \(R\) is surjective); the source term being reduced
(branchwise quotient \(K\)) determines the normal form
(recombination target \(R\)) independently of the reduction
strategy iff the system is confluent; without termination, take as branch
carrier the pairs \((s,\rho)\) with \(\rho\) a finite reduction sequence from \(s\)
ending in a normal form, and corestrict \(K\) to the terms that have a normal
form; the condition is then the unique-normal-form property, which is weaker
than confluence.  Bridge obstructions are
\emph{non-confluence witnesses}: two reductions from the same
source term producing distinct normal forms
\citep{ChurchRosser1936Conversion}.

\item[\emph{Path integral with gauge equivalence (physics).}] \textit{Special case.}
For a quantum-mechanical path integral with gauge-orbit
quotient \(K\) on field-configuration space, the exponentiated
action \(R = \exp(iS[\text{path}])\), corestricted to its image in \(U(1)\), descends through \(K\) iff
the classical action is gauge-invariant modulo \(2\pi\) (with
\(\hbar=1\)), as for a Chern--Simons action at integer level under
large gauge transformations.  Bridge obstructions are gauge-related
paths whose action values differ by a non-multiple of \(2\pi\), signalling a
\emph{gauge-invariance failure} of the action
\citep{FeynmanHibbs1965PathIntegrals}.

\item[\emph{Tversky--Kahneman framing effect (cognitive psychology
/ behavioural economics).}] \textit{Analogy.}
For a choice problem presented to separate groups of participants
under various framings, the objective lottery underlying the problem
(branchwise quotient \(K\), with frame forgotten) and the group's
modal choice (recombination target \(R\)) coincide --- i.e., \(R\)
descends through \(K\) --- iff choices are frame-invariant.  The Tversky--Kahneman Asian disease problem
demonstrates that participants systematically violate
frame-invariance: presenting the same objective lottery to one group in a gain frame
and to another in a loss frame yields different modal choices, producing the bridge-composition obstruction --- a
presentation pair sharing the underlying lottery but differing in
choice
\citep{TverskyKahneman1981Framing}.
\end{description}
\endgroup

\subsection*{Transport at a glance}\label{sec:transport:summary}\addcontentsline{toc}{subsection}{Transport at a glance}
\begin{table}[H]
  \centering
  \footnotesize
  \begin{tabularx}{\textwidth}{@{}
      >{\raggedright\arraybackslash}p{0.24\textwidth}
      >{\raggedright\arraybackslash}X
      >{\raggedright\arraybackslash}p{0.12\textwidth}@{}}
    \toprule
    Name & One-line claim & Substrate cite \\
    \midrule
    Descent--Repair Theorem (\Fref{F2}) &
    For surjective \(q\), \(\splitObs\) is empty iff descent holds; canonical target
    repair coarsens \(R\) by the generated equivalence.  &
    \textemdash \\
    Holonomy-Memory Theorem (\Fref{F3}) &
    Memory is forced at \(h\) exactly when there is route residue at
    \(h\); \(\langle\currQ,\predM\rangle\) is the coarsest refinement of
    \(\currQ\) through which \(\predM\) descends.  &
    \textemdash \\
    Local--Global Obstruction Theorem (\Fref{F4}) &
    When every restricted global state \(R(g)\) is compatible, globalizing
    data is compatible, and every local family has exactly
    one of four statuses: overlap-inconsistent, compatible but not
    global, globally ambiguous or globally exact.  &
    \textemdash \\
    Recombination Comparison (\Fref{F5}) &
    For surjective branchwise and recombination quotients, branchwise equality determines merged equality iff the
    recombination quotient factors through the branchwise quotient.  &
    \textemdash \\
    \bottomrule
  \end{tabularx}
  \caption{Transport at a glance: the four laws of \Cref{sec:transport}.}
  \label{tab:axis-transport-summary}
\end{table}

\FloatBarrier
   % sec:transport
\section{Symmetry}\label{sec:symmetry}

The Symmetry axis collects obstructions caused by involutions,
branch selection, and anomaly-style failures.  In this chapter a
symmetry is not merely a convenient change of notation.  It is part
of the apparatus that controls which branches, traces, and quotient
classes may be distinguished without adding new breaking data.  The
basic question is therefore familiar: when a construction is
invariant under the declared symmetry, what can it select, transport,
or record?  In the running example of \Cref{sec:intro}, the involution
exchanging \(1\) and \(2\) fixes every \(q\)-class; its orbit \(\{1,2\}\) has no fixed point, so by \Fref{F15a} no equivariant section of the orbit map exists (every equivariant map from the orbit set to \(X\) takes values in the fixed points \(3,4\)); choosing \(1\) or \(2\) needs breaking data, such as the readout \(F\), which separates them.

The five laws proceed from fixed data to failure modes and then to
equivalence.  Duality Fixity, or Self-Dual Trace Confinement, is the
layer-agnostic form of the confinement theorem of
\citet{Tsiokos2026SDTC}, which \citet{Tsiokos2026RH} applies to the
zeros of the Riemann zeta function.  For a readout that is strongly
measurable and square-integrable when the visible weight is a measure,
a positive anti-invariant ledger dominated by bounds whose traces have
infimum at most \(0\) is zero, and a separating readout then confines
visible mass to the fixed locus.  Such bounds exist exactly when
\(\psi^-=0\) \(\mu\)-almost everywhere (last clause of \Cref{thm:F14}),
so all content lies in constructing them; the source papers take them
as input, and no result of this paper assumes the Riemann hypothesis.  Selection Obstruction / Repair
Normal Form is a direct theorem: exact symmetry cannot select a
non-fixed branch without declared breaking data.  Spontaneous
Symmetry Breaking / Branch-Selection Formation characterizes when a
branch-selection record exists: exactly when some ground state is not invariant
under the symmetry.  Anomaly / Symmetry Obstruction and Duality
Equivalence are direct theorems about failed symmetry descent and
role-preserving equivalence of formed packages.

The direct selection-obstruction theorem reuses the
Descent--Repair split-pair discipline: a forbidden selection is a
split-pair obstruction at the level of symmetry orbits.  The later
anomaly and duality-equivalence subsections keep the same posture:
the apparatus must say when symmetry descends, when it fails, and
when two symmetry presentations preserve the same formed roles.

\subsection{Duality Fixity (Self-Dual Trace Confinement) [\texorpdfstring{\Fref{F14}}{F14}]}\label{sec:symmetry:duality-fixity}

\begin{definition}[Involutive ledger; anti-invariant readout]\label{def:involutive-ledger}
An involutive ledger is a carrier \(X\) equipped with an involution
\(\Jinv:X\to X\), so \(\Jinv^2=\mathrm{id}\), and a visible weight
\(\mu\ge 0\) on \(X\) (a finite weight function, or a measure).  Its fixed
locus is
\[
  \Fixinv=\{x\in X:\Jinv(x)=x\}.
\]
An anti-invariant readout is a map \(\psi^-\) from \(X\) to a Hilbert space
\(\mathcal H\), with a declared unitary involution \(U\) of \(\mathcal H\), such
that \(\psi^-(\Jinv x)=-U\psi^-(x)\) for every \(x\); when \(\mu\) is a
measure, \(\psi^-\) is strongly (Bochner) measurable, for instance Borel measurable with
\(\mathcal H\) separable, with \(\int_X\|\psi^-\|^2\,d\mu<\infty\).
A map \(\psi^-:X\to\mathcal H\), anti-invariant or not, is \emph{separating} when \(\{x\notin\Fixinv:\psi^-(x)=0\}\) is contained in a
\(\mu\)-null set; for a finite weight this means that
every \(x\) with \(\mu(x)>0\) and \(\psi^-(x)=0\) lies in \(\Fixinv\).
\end{definition}

\paragraph{Reading the statement.}
The objects are those of \citet{Tsiokos2026SDTC}, over a real Hilbert
space.  There the odd part \(\psi^-\) of an observable satisfies
\(\psi^-(\Jinv x)=-\psi^-(x)\), which is the case \(U=\mathrm{id}\) of
\Cref{def:involutive-ledger}; the response involution of that paper,
under which \(\psi^-\) is odd, is a different map and is not the \(U\)
of this definition, so the restriction on \(U\) below is not imposed by
that paper. The anti-invariant ledger
is the positive element
\(A_X=\int_X\psi^-(x)\psi^-(x)^*\,d\mu(x)\), a Bochner integral of
trace-class operators, which converges because
\(\int_X\|\psi^-\|^2\,d\mu<\infty\); for a finite weighted ledger it is the
finite sum \(A_X=\sum_{x}\mu(x)\,\psi^-(x)\psi^-(x)^*\) (Definition~4.1
there). Anti-invariance identifies \(A_X\) as the anti-invariant ledger.
The trace-squeeze argument does not use anti-invariance; the implication from confinement to \(A_X=0\), under the stated restriction on \(U\), does. No clause uses \(\Jinv^2=\mathrm{id}\), nor that \(U\) is unitary or an involution: \(\Jinv\) enters only through \(\Fixinv\), and \(U\) only through \(U\psi^-(x)=-\psi^-(x)\) at fixed points, so every clause holds for an arbitrary map \(\Jinv:X\to X\) and an arbitrary map \(U:\mathcal H\to\mathcal H\). The ledger \(A_X\) lies
in a positive cone with order \(\preceq\) and trace \(\trAX\); positive
trace-class operators on a Hilbert space, or positive semidefinite matrices
with the usual trace, are the classical model. The theorem uses two
properties of the cone and one property of \(A_X\), all of which hold in that
model:
(C1) the trace is monotone, \(A\preceq B\) implies \(\trAX A\le\trAX B\);
(C2) a positive element of trace zero is zero;
(C3) the trace identity \(\trAX A_X=\int_X\|\psi^-(x)\|^2\,d\mu(x)\)
(Lemma~4.3 there; in the classical model it is the trace of a rank-one
operator, \(\trAX(vv^*)=\|v\|^2\)).
The theorem is a reformulation, not a way of establishing
confinement: by its last clause, suitable bounds exist exactly when
\(\psi^-=0\) \(\mu\)-almost everywhere.  Its value lies in isolating the
properties (C1)--(C3) that the squeeze uses.  For a linear involution
\(U\), the condition in the theorem that no nonzero \(v\) satisfies
\(Uv=-v\) holds exactly when \(U=\mathrm{id}\).
The \(B_n\) are the bounds of a domination ladder, a sequence of elements of
the cone with \(A_X\preceq B_n\) for every \(n\); in the source paper they
have the form \(B_n=K^-_n+E_n\) (Definition~6.1 there).

\begin{theorem}[Duality Fixity / Self-Dual Trace Confinement]\label{thm:F14}
Let \(X\) carry a map \(\Jinv:X\to X\) (an involution in the intended reading; \(\Jinv^2=\mathrm{id}\) is not required) and a visible weight \(\mu\) as in \Cref{def:involutive-ledger} (in the intended reading, the ledger of a formed closure; formedness is not used), with
separating readout \(\psi^-:X\to\mathcal H\), strongly measurable and square-integrable when \(\mu\) is a measure, as in \Cref{def:involutive-ledger} but not necessarily anti-invariant, and let its
ledger \(A_X\) (the anti-invariant ledger when \(\psi^-\) is anti-invariant) lie in a positive cone with order \(\preceq\) and trace
\(\trAX\) satisfying (C1)--(C3). If
\[
  A_X \preceq B_n\ \ \text{for every } n
  \qquad\text{and}\qquad
  \inf_n\trAX(B_n)\le 0
  \quad(\text{for every }\varepsilon>0\text{ some }\trAX(B_n)<\varepsilon),
\]
then \(A_X=0\) and \(X\setminus\Fixinv\) is contained in a \(\mu\)-null set,
so \(\mu(X\setminus\Fixinv)=0\) whenever \(\Fixinv\) is measurable (for
instance, when \(\Jinv\) is measurable on a standard Borel space): visible
mass is confined to the fixed locus \(\Fixinv\). Moreover, if \(\psi^-(\Jinv x)=-U\psi^-(x)\) for a declared map \(U:\mathcal H\to\mathcal H\) (a unitary involution in the intended reading), \(\trAX(0)=0\) (true in the classical model), and no nonzero \(v\in\mathcal H\) has \(Uv=-v\) (for instance
\(U=\mathrm{id}\)), then \(A_X=0\) if and only if \(X\setminus\Fixinv\) is
contained in a \(\mu\)-null set; without this condition on \(U\) the
equivalence can fail. Conversely, if \(\psi^-=0\) \(\mu\)-almost everywhere, then \(\trAX A_X=0\) by (C3), and the constant sequence \(B_n:=A_X\) satisfies both displayed hypotheses; hence such a sequence exists if and only if \(\psi^-=0\) \(\mu\)-almost everywhere.
\end{theorem}
\begin{proof}
By (C3), \(\trAX A_X=\int_X\|\psi^-\|^2\,d\mu\ge 0\), and by (C1),
\(\trAX A_X\le\trAX B_n\) for every \(n\). Taking the infimum over \(n\) gives
\(\trAX A_X\le 0\), so \(\trAX A_X=0\). Since \(A_X\) is positive, (C2)
gives \(A_X=0\). By (C3) again, \(\int_X\|\psi^-\|^2\,d\mu=0\), and the
integrand is nonnegative, so \(\psi^-(x)=0\) for \(\mu\)-almost every
\(x\). Then \(X\setminus\Fixinv\) is contained in the union of the null set
\(\{\psi^-\neq0\}\) and the set \(\{x\notin\Fixinv:\psi^-(x)=0\}\), which is
contained in a null set by separation; so \(X\setminus\Fixinv\) is contained
in a \(\mu\)-null set, and \(\mu(X\setminus\Fixinv)=0\) when \(\Fixinv\) is
measurable. For the equivalence under the restriction on \(U\), if \(A_X=0\) then \(\trAX A_X=\trAX(0)=0\), and the argument just given, from \(\trAX A_X=0\) on, confines \(X\setminus\Fixinv\) to a null set. Conversely, let
\(U\) have no nonzero \(v\) with \(Uv=-v\), and let \(X\setminus\Fixinv\) lie in
a null set. For \(x\in\Fixinv\), anti-invariance gives
\(\psi^-(x)=\psi^-(\Jinv x)=-U\psi^-(x)\), so \(U\psi^-(x)=-\psi^-(x)\) and
\(\psi^-(x)=0\); hence \(\psi^-=0\) almost everywhere,
\(\int_X\|\psi^-\|^2\,d\mu=0\), \(\trAX A_X=0\) by (C3), and \(A_X=0\) by (C2).
Without the condition, take one point \(x\) of mass one with \(\Jinv x=x\),
\(Uv=-v\) and \(\psi^-(x)=v\neq0\): then \(X\setminus\Fixinv=\varnothing\) but
\(A_X=vv^*\neq0\). For the converse, if \(\psi^-=0\) \(\mu\)-almost everywhere, (C3) gives \(\trAX A_X=\int_X\|\psi^-\|^2\,d\mu=0\); the constant sequence \(B_n:=A_X\) satisfies \(A_X\preceq B_n\) by reflexivity of \(\preceq\), and \(\trAX(B_n)=0\).
\end{proof}

\paragraph{What the sources supply.}
The theorem assumes the domination ladder; it does not construct one.
\PaperSDTC\ \citep{Tsiokos2026SDTC} proves the same squeeze for real
Hilbert spaces, adds a quantitative form of the confinement and
variants on finite windows, and names as a hypothesis, Self-Dual Trace
Confinement, that in specified self-dual settings such domination
records are supplied by independent structure.  \PaperRH\
\citep{Tsiokos2026RH} applies the squeeze to the reflection
\(s\mapsto1-\bar s\) in the critical line.  Its conditional theorem uses
a named hypothesis on an auxiliary ledger of represented zeros, with
explicit encoding conditions, rather than the global ledger
\(A_Z=\sum_\rho m_\rho(\operatorname{Re}\rho-\tfrac12)^2\) over the
nontrivial zeros, whose vanishing is equivalent to the Riemann
hypothesis for positive integer weights \(m_\rho\).  Its concrete
version uses finite windows of actual zeros that are closed under
\(s\mapsto1-\bar s\), carry unit weights and eventually contain every
nontrivial zero.  The named hypothesis holds exactly when every
represented zero has critical coordinate; it is therefore equivalent to
the Riemann hypothesis on the bare shell, which represents exactly the
nontrivial zeros with coordinate \(\operatorname{Re}\rho\), and the
concrete window version is likewise equivalent to the Riemann
hypothesis.  That paper presents its conditional theorem as a
reformulation, not as evidence.  \Cref{thm:F14} states the
confinement law for any involutive ledger with a separating readout.


\paragraph{Nonclaims.} This subsection does not construct a domination
ladder, does not prove the RH specialization, and does not instantiate the
law to a specific arithmetic carrier.

\subsection*{Layer instantiations}\label{layer-inst:F14}

\begingroup\small
\begin{description}
\item[\emph{Symmetric/skew decomposition under transpose involution
(linear algebra).}] \textit{Analogy.}
For any \(A\in\mathrm{Mat}_n(\mathbb R)\), the transpose
\(J(A)=A^{\!\top}\) is an involution; the unique decomposition
\(A=\tfrac12(A+A^{\!\top})+\tfrac12(A-A^{\!\top})\) separates
\(J\)-symmetric and \(J\)-skew parts.  The skew part \(A-A^{\!\top}\), which \(J\) sends to its negative,
is the \Fref{F14} anti-invariant readout; its Frobenius norm
\(\lVert A-A^{\!\top}\rVert_F\) is \(J\)-invariant, and its vanishing
places \(A\) in the fixed locus \(\mathrm{Sym}_n(\mathbb R)\), on
which the Loewner order is defined
\citep{HornJohnson2013MatrixAnalysis}.

\end{description}
\endgroup


\subsection{Selection Obstruction / Repair Normal Form [\texorpdfstring{\Fref{F15a}}{F15a}]}\label{sec:symmetry:selection-obstruction}

\begin{definition}[Selection obstruction]\label{def:selection-obstruction}
Let a declared symmetry group \(G\) act on a carrier \(X\), and let
\[
  p:X\to X/G
\]
be the orbit quotient.  For a selectedness predicate
\(\chi:X\to\{0,1\}\), the selection obstruction is the split-pair set
\[
  \splitObs(p,\mathrm{id}_X,\chi) =
  \{(x,x')\in X\times X:
      p(x)=p(x')\ \text{and}\ \chi(x)\ne\chi(x')\}.
\]
\end{definition}

\paragraph{Reading the statement.}
The symbols are those of \Cref{thm:F15a} below: \(G\) acts on \(X\)
by \(\rho_X\) and on \(O\) by \(\rho_O\), \(\pi:X\to O\) is
surjective, \(s\) is the selectedness predicate, and
\(\sigma:O\to X\) is a map, in the intended reading a section picking one point in each \(\pi\)-fiber; \(\pi\circ\sigma=\mathrm{id}_O\) is assumed only in the equivariant-section clause.
\(\operatorname{EquivariantSelector}(\rho_O,\rho_X,\sigma)\) says
that this choice commutes with the group:
\(\sigma(\rho_O(g,o))=\rho_X(g,\sigma(o))\) for all \(g,o\).
\(\operatorname{TrivialOrbitAction}(\rho_O)\) says that every group
element fixes every element of \(O\): \(\rho_O(g,o)=o\).
\(\operatorname{HasSelectionSplit}(\pi,s)\) says that some
\(\pi\)-fiber contains a selected and an unselected point.
Under the orbit-quotient hypothesis of the theorem's moreover
clause, these fibers are the orbits, \(O\) plays the role of \(X/G\),
\(\pi\) that of \(p\), and \(s\) that of \(\chi\), so
\(\mathcal O_s\) is the selection obstruction of
\Cref{def:selection-obstruction}.

\begin{theorem}[Selection Obstruction / Repair Normal Form]\label{thm:F15a}
Let \(G\) act on a set \(X\) and on an orbit set \(O\) by
\(\rho_X:G\times X\to X\) and \(\rho_O:G\times O\to O\).  Let
\(\pi:X\to O\) be a surjective map (in the intended reading the orbit map; the
orbit-quotient property is assumed only in the clause that states it), let
\(s:X\to\{0,1\}\) be a selectedness predicate, let
\(\sigma:O\to X\) be a map (in the intended reading a section of \(\pi\), though \(\pi\circ\sigma=\mathrm{id}_O\) is not assumed; the equivariant-section clause below concerns maps that satisfy it).  Define
\[
  \mathcal O_s
  :=
  \{(x,x')\in X^2:\pi(x)=\pi(x')\ \text{and}\ s(x)\ne s(x')\}.
\]
Then
\begin{align*}
  &\operatorname{EquivariantSelector}(\rho_O,\rho_X,\sigma)
    \wedge \operatorname{TrivialOrbitAction}(\rho_O)\\
  &\hspace{4em}\Longrightarrow
    \forall o\in O,\ \forall g\in G,\
      \rho_X(g,\sigma(o))=\sigma(o),\\
  s\ \text{descends through}\ \pi
  &\iff
  \mathcal O_s=\emptyset
  \iff
  \exists\,\bar s:O\to\{0,1\}\ \text{with}\ \bar s\circ\pi=s,\\
  \operatorname{HasSelectionSplit}(\pi,s)
  &\iff
  \mathcal O_s\ne\emptyset.
\end{align*}
If moreover \(\pi\) is the orbit quotient,
\(\pi(x)=\pi(x')\iff x'=\rho_X(g,x)\) for some \(g\in G\), and \(\rho_O\) is
the induced action, \(\rho_O(g,\pi(x))=\pi(\rho_X(g,x))\), then
\(\operatorname{TrivialOrbitAction}(\rho_O)\) holds, every equivariant map
\(\sigma:O\to X\), in particular every equivariant section, takes values in the
common fixed-point set \(X^G\), and if \(X^G=\emptyset\) and
\(O\ne\emptyset\), no equivariant map \(O\to X\) exists; and an equivariant section exists if and only if \(G\) acts trivially on \(X\), and it is then the inverse of \(\pi\).
For every breaker readout \(b:X\to B\), the breaker extension
\(\langle\pi,b\rangle:X\to O\times B\) retains \(\pi\) and carries \(b\), strictly
refines \(\pi\) when \(b\) separates two points of one \(\pi\)-fiber, and factors
through every map \(p\) through which both \(\pi\) and \(b\) factor. Conversely to the fixed-point clause above, under the orbit-quotient and induced-action hypotheses every map \(O\to X^G\) is equivariant, so an equivariant map \(O\to X\) exists if and only if \(X^G\ne\emptyset\) or \(O=\emptyset\). Adjoining \(b\) records the
distinction carried by \(b\); taking \(b=s\) carries selectedness, but the
extension alone does not construct a selector.
\end{theorem}
\begin{proof}
Assume the actions \(\rho_X,\rho_O\), the surjective map
\(\pi:X\to O\), the selectedness predicate \(s\), and the map
\(\sigma:O\to X\).  If \(\sigma\) is equivariant and the \(G\)-action on
\(O\) is trivial, then for every \(g\in G\) and \(o\in O\), equivariance of
the selector and triviality of the action on \(O\) give
\(\rho_X(g,\sigma(o))=\sigma(\rho_O(g,o))=\sigma(o)\), so
\(\sigma(o)\) lies in the common fixed locus of the \(G\)-action on
\(X\).

For selectedness, if \(s=\bar s\circ\pi\), then
\(\pi(x)=\pi(x')\) implies \(s(x)=s(x')\), hence
\(\mathcal O_s=\emptyset\).  Conversely, if
\(\mathcal O_s=\emptyset\), define
\[
  \bar s(\pi(x)):=s(x).
\]
The empty obstruction makes the definition independent of the
representative \(x\), and surjectivity of \(\pi\) supplies
representatives for all \(o\in O\).  If \(\mathcal O_s=\emptyset\), selectedness is a quotient-level
predicate. If \(\mathcal O_s\ne\emptyset\), a split pair shows that no
\(\bar s\) exists; by the breaker clause with \(b=s\), the extension
\(\langle\pi,s\rangle\) removes every split pair, since \(s\) factors through
it.

For the orbit-quotient clause, \(\rho_O(g,\pi(x))=\pi(\rho_X(g,x))=\pi(x)\),
since \(\rho_X(g,x)\) lies in the orbit of \(x\); as \(\pi\) is surjective, the
action on \(O\) is trivial. The first display then places every value of an
equivariant map \(\sigma:O\to X\) in \(X^G\), and if \(X^G=\emptyset\) while \(O\) has an
element \(o\), the value \(\sigma(o)\) cannot exist. For an equivariant section \(\sigma\), the value \(\sigma(\pi x)\) is fixed and lies in the orbit of \(x\), since \(\pi(\sigma(\pi x))=\pi x\); a point of an orbit is fixed only if every point of that orbit is, so \(x\) is fixed and \(\sigma(\pi x)=x\). Conversely, a trivial action makes every orbit a point, so \(\pi\) is bijective and its inverse is an equivariant section. As \(\rho_O\) is trivial, every \(\sigma:O\to X^G\) satisfies \(\sigma(\rho_O(g,o))=\sigma(o)=\rho_X(g,\sigma(o))\), so the equivariant maps are exactly those with values in \(X^G\).

The breaker clauses are the source-repair part of \Cref{thm:F2} with source \(\pi\),
identity transport and target readout \(b\). For the selectedness reading, take
\(b=s\).
\end{proof}

\paragraph{Nonclaims.} The theorem does not classify which
breaking data are allowed and does not instantiate the obstruction to
a particular symmetry group.

\subsection*{Layer instantiations}\label{layer-inst:F15a}

\begingroup\small
\begin{description}
\item[\emph{Vitali sets and the axiom of choice (mathematical
logic).}] \textit{Special case.}
For the free \(\mathbb Q\)-action on \(\mathbb R\) by translation,
no equivariant section of the orbit projection
\(\mathbb R\to\mathbb R/\mathbb Q\) can exist (there are no fixed
points).  The axiom of choice supplies a non-equivariant Vitali
set of representatives, which is necessarily non-measurable: the
\Fref{F15a} declared breaker is the readout \(b(x)=[x=v(\pi(x))]\) for a choice of
representatives \(v:\mathbb R/\mathbb Q\to\mathbb R\), which the axiom of choice supplies, licensing a
non-equivariant selection where no equivariant one is available
\citep{Vitali1905ProblemaMisura}.

\item[\emph{Faddeev--Popov gauge fixing and the Gribov--Singer
obstruction (gauge theory).}] \textit{Analogy.}
Singer's theorem states that the gauge orbit space
\(\mathcal A/\mathcal G\) of Yang--Mills theory admits no
continuous global section.  Faddeev--Popov gauge fixing provides
the \Fref{F15a} declared breaker through a local gauge condition
\(F(A)=0\) together with the Faddeev--Popov determinant
\(\det(\delta F/\delta\omega)\) and ghost fields, licensing local
path-integral computations; the global Gribov--Singer obstruction
witnesses the non-existence of a clean equivariant selector
\citep{FaddeevPopov1967YangMills}.

\item[\emph{Buridan's ass and symmetric tie-breaking (decision
theory).}] \textit{Special case.}
Under exact symmetry between two equally optimal choices, no
equivariant deterministic decision rule selects one option.  Take
\(X=\{\ell,r\}\), \(G=\mathbb Z/2\) exchanging the options, \(O\) a one-point set
and \(\pi\) constant with the induced action; since \(X^G=\varnothing\) and
\(O\ne\varnothing\), \Fref{F15a} gives no equivariant section.  The
declared breaker is a tie-breaking rule --- random coin flip, a
lexicographic preference axis, or a memory of past choices --- that
supplies the non-symmetric apparatus needed for a definite
selection
\citep{Rescher1959ChoiceWithoutPreference}.
\end{description}
\endgroup


\subsection{Spontaneous Symmetry Breaking / Branch-Selection Formation [\texorpdfstring{\Fref{F15b}}{F15b}]}\label{sec:symmetry:spontaneous-ssb}

The setup is a formed closure under an exact declared symmetry, with
the orbit quotient and selection interface already in place as in the
previous subsection.  A spontaneous symmetry-breaking claim asks for
more than a set-theoretic section of the orbit map.  It asks for a
branch-selection record: a map that keeps each state in its orbit and
selects, for some ground state, a branch that the symmetry moves.

\paragraph{Reading the statement.}
A \emph{branch-selection formation record} for \((\mathcal C,G,\Omega)\) is a
triple \((F,c,g)\) with \(F:\mathcal C\to\mathcal C\), \(q\circ F=q\), \(F\)
constant on orbits (\(F(h\cdot x)=F(x)\) for all \(h\in G\), so that
\(F=\sigma\circ q\) for a section \(\sigma\) of \(q\)), \(\Omega(c)\) and
\(g\cdot F(c)\ne F(c)\). The identity map is not such a record, since it is not
constant on a nontrivial orbit: the record selects one branch in each orbit. Clause (o) says that such a record
exists exactly when some ground state is not invariant under \(G\). When
\(\Omega\) is \(G\)-invariant, as in clause (iii), this is the usual definition
of a spontaneously broken symmetry; without that invariance the asymmetry may
already be in \(\Omega\). Clause (i) locates
the selected branch, clause (ii) says that no symmetric selector
produces it, and clause (iii) says that when the ground-state predicate is
symmetric the record exhibits two distinct ground states in one orbit, the
selected branch and its image under \(g\).

\begin{theorem}[Spontaneous Symmetry Breaking / Branch-Selection Formation]\label{thm:F15b}
Let \(\mathcal C\) be a closure space with a \(G\)-action, let
\(\Omega:\mathcal C\to\mathrm{Prop}\) be the ground-state predicate,
and let \(q:\mathcal C\to\mathcal C/G\) be the orbit map onto the
branch quotient.  Then:
\begin{enumerate}[label=(\roman*)]
\item[(o)] a branch-selection formation record exists if and only if some
  ground state is not fixed by \(G\): there are \(c\in\mathcal C\) with
  \(\Omega(c)\) and \(g\in G\) with \(g\cdot c\ne c\);
\end{enumerate}
and for every branch-selection formation record \((F,c,g)\):
\begin{enumerate}[label=(\roman*)]
\item \(F\) maps each orbit into itself, and the branch \(F(c)\) lies in the
  orbit of the ground state \(c\) and is not fixed by \(G\);
\item no \(G\)-equivariant map \(\sigma:\mathcal C/G\to\mathcal C\), for the
  trivial action of \(G\) on \(\mathcal C/G\), takes the value \(F(c)\) at
  \(q(c)\);
\item if \(\Omega\) is \(G\)-invariant, that is \(\Omega(h\cdot x)\iff\Omega(x)\)
  for all \(h\in G\) and \(x\in\mathcal C\), then \(F(c)\) and \(g\cdot F(c)\)
  are distinct ground states in the orbit of \(c\).
\end{enumerate}
So the branch \(F(c)\) cannot be chosen by a symmetric selector: the record
\(F\) is breaker data in the sense of \Cref{thm:F15a}.
\end{theorem}
\begin{proof}
(o) If \(c\) is a ground state and \(g\cdot c\ne c\), choose a section
\(\sigma\) of \(q\) with \(\sigma(q(c))=c\), and let \(F=\sigma\circ q\). Then
\(q\circ F=q\), \(F\) is constant on orbits, and \(F(c)=c\), so
\(g\cdot F(c)\ne F(c)\) and \((F,c,g)\) is a formation record. Conversely, let \((F,c,g)\) be a formation record.
Since \(q(F(c))=q(c)\), there is \(h\in G\) with \(F(c)=h\cdot c\). If \(G\)
fixed \(c\), then \(F(c)=c\) and \(g\cdot F(c)=F(c)\), which is not the case;
so the ground state \(c\) is not fixed by \(G\).

(i) Since \(q\circ F=q\), each \(F(x)\) lies in the orbit of \(x\); in
particular \(F(c)\) lies in the orbit of \(c\), and \(g\cdot F(c)\ne F(c)\)
says that \(G\) does not fix it.  (ii) Let \(\sigma\) be \(G\)-equivariant for
the trivial action on \(\mathcal C/G\).  Then for every \(h\in G\) and every
orbit \(o\), \(h\cdot\sigma(o)=\sigma(h\cdot o)=\sigma(o)\), so every value of
\(\sigma\) is fixed by \(G\), as in the first clause of \Cref{thm:F15a}.  Since
\(g\cdot F(c)\ne F(c)\), \(\sigma(q(c))\ne F(c)\).  The asymmetry is therefore
carried by the selection record \(F\), not by the symmetric quotient.
(iii) By (i), \(F(c)=h\cdot c\) for some \(h\in G\), so \(\Omega(F(c))\) by
invariance, and then \(\Omega(g\cdot F(c))\) by invariance again. Both lie in
the orbit of \(c\), and they are distinct because \(g\cdot F(c)\ne F(c)\).
\end{proof}

\paragraph{Where the breaking comes from.}
The theorem does not produce a non-invariant ground state; whether one exists
is a property of the closure and its ground-state predicate. In a physical
system this is the content of a spontaneous symmetry-breaking result, and the
theorem records where the asymmetry is carried once it holds: in the
selection record, not in the symmetric quotient. \PaperCSL\
\citep{Tsiokos2026CSL} gives a parallel diagnosis in the access setting.


\paragraph{Nonclaims.} This subsection does not exhibit a non-invariant
ground state for any particular closure and does not classify spontaneous
symmetry breaking at a substrate level.  The canonical particle-physics
instances of spontaneous symmetry breaking are surveyed in Goldstone, Salam,
and Weinberg \citep{Goldstone1962BrokenSymmetries}, and the broader
emergent-layer reading of multi-level breaking appears in Anderson's
\emph{More is Different} \citep{Anderson1972MoreIsDifferent}.

\subsection*{Layer instantiations}\label{layer-inst:F15b}

\begingroup\small
\begin{description}
\item[\emph{Bose--Einstein condensation (atomic physics).}] \textit{Analogy.}
When a dilute gas of bosons is cooled below its critical
temperature, a macroscopic fraction of atoms condenses into a
single one-particle quantum state.  The condensate's overall
\(U(1)\) phase is arbitrary --- a continuous symmetry of the
underlying dynamics --- yet each realisation selects a definite
phase, breaking the symmetry.  The 1995 Anderson--Ensher--
Matthews--Wieman--Cornell experiment on rubidium-87 confirmed
Einstein's 1925 prediction; the macroscopic phase selection is the
\Fref{F15b} formation operator
\citep{AndersonEnsherMatthewsWiemanCornell1995BEC}.

\end{description}
\endgroup


\subsection{Anomaly / Symmetry Obstruction [\texorpdfstring{\Fref{F40}}{F40}]}\label{sec:symmetry:anomaly}

Let a formed closure have carrier \(H\), layer quotient
\(q:H\to Q\), declared readouts and audit records, and a would-be
symmetry, an action \(\alpha\) of a group \(G\) on the raw carrier.  The
action is a genuine layer symmetry only when it acts functorially on the quotient and
preserves the declared structure of the layer.

\paragraph{Reading the statement.}
\(\operatorname{SymmetryDescends}(q,\alpha)\) says that each group
element induces a well-defined map on the quotient: for every \(g\),
the map \(h\mapsto q(\alpha(g,h))\) factors through \(q\).
\(\operatorname{DomainDescends}(q,D)\) says that whether \(g\) applies
to \(h\) depends only on the quotient class of \(h\): for every \(g\),
the predicate \(h\mapsto D(g,h)\) factors through \(q\). The other two
conditions are defined by their displayed equations.

\begin{theorem}[Anomaly / Symmetry Obstruction]\label{thm:F40}
Let \(G\) act on a history space \(H\) by
\(\alpha:G\times H\to H\), let \(q:H\to Q\) be a surjective quotient,
let \(\bar\alpha:G\times Q\to Q\) be a candidate descended action,
let \(D:G\times H\to\{0,1\}\) be a domain predicate, and let
\(r:H\to V\) be a readout with value action
\(\rho:G\times V\to V\).  Define
\[
  \mathcal O_{\alpha}
  :=
  \{(g,h,h'):\ q(h)=q(h')\ \text{and}\
    q(\alpha(g,h))\ne q(\alpha(g,h'))\},
\]
and define \(\mathcal O_D\) analogously by replacing
\(q(\alpha(g,h))\) with \(D(g,h)\), and define
\begin{align*}
  \operatorname{GroupLawPreserved}(\bar\alpha)
  &:\Longleftrightarrow
  \forall g,g'\in G,\forall \bar h\in Q,\,
    \bar\alpha(gg',\bar h)=
    \bar\alpha(g,\bar\alpha(g',\bar h)),\\
  \operatorname{ReadoutPreserved}(r,\alpha,\rho)
  &:\Longleftrightarrow
  \forall g\in G,\forall h\in H,\,
    r(\alpha(g,h))=\rho(g,r(h)).
\end{align*}
Then
\begin{align*}
  \operatorname{SymmetryDescends}(q,\alpha)
  &\iff \mathcal O_{\alpha}=\emptyset,\\
  \operatorname{DomainDescends}(q,D)
  &\iff \mathcal O_D=\emptyset.
\end{align*}
Moreover, if \(\bar\alpha\) is a descended action,
\(\bar\alpha(g,q(h))=q(\alpha(g,h))\) for all \(g\) and \(h\), then
\(\mathcal O_\alpha=\emptyset\), \(\operatorname{GroupLawPreserved}(\bar\alpha)\)
holds, and \(\bar\alpha(e,\bar h)=\bar h\) for the identity \(e\) and every
\(\bar h\in Q\), so \(\bar\alpha\) is a group action. Let the would-be symmetry carry a declared record
\(\mathsf{broken}\in\{0,1\}\), which is \(1\) when it is recorded as
explicitly broken. When \(\mathsf{broken}=0\), the \emph{anomaly} is defined by
\[
  \operatorname{Anomaly}
  \;:\iff\;
  \mathcal O_\alpha\ne\emptyset\ \lor\ \mathcal O_D\ne\emptyset\ \lor\
  \neg\operatorname{ReadoutPreserved}(r,\alpha,\rho).
\]
If \(\mathcal O_\alpha\ne\emptyset\), action descent fails and no descended
action exists; when \(\mathsf{broken}=0\) this is an anomaly. If \(\mathcal O_\alpha=\emptyset\), the
\emph{induced action} \(\alpha^{\flat}(g,q(h)):=q(\alpha(g,h))\) is well
defined and is a descended action, hence preserves the group law. An anomaly
then occurs exactly when one of \(\operatorname{SymmetryDescends}(q,\alpha)\),
\(\operatorname{DomainDescends}(q,D)\),
\(\operatorname{GroupLawPreserved}(\alpha^{\flat})\) and
\(\operatorname{ReadoutPreserved}(r,\alpha,\rho)\) fails; since the first and
third hold, this is exactly when domain descent or readout preservation
fails. The record \(\mathsf{broken}\) affects only the anomaly designation:
the descent equivalences and descended-action conclusions hold whatever
its value, while the anomaly characterizations apply only when
\(\mathsf{broken}=0\).
\end{theorem}
\begin{proof}
Assume the action \(\alpha\), quotient \(q\), candidate descended
action \(\bar\alpha\), domain predicate \(D\), readout \(r\), and
value action \(\rho\).  The action descends through \(q\) exactly when
the formula
\[
  \alpha^{\flat}(g,q(h)):=q(\alpha(g,h))
\]
is independent of the representative \(h\).  That independence is
equivalent to \(\mathcal O_{\alpha}=\emptyset\).  The domain predicate
descends by the identical argument applied to \(D(g,h)\), giving
\(\mathcal O_D=\emptyset\).

The displayed equation defines \(\operatorname{GroupLawPreserved}\):
\[
  \bar\alpha(gg',\bar h)
  =
  \bar\alpha(g,\bar\alpha(g',\bar h))
  \quad
  \forall g,g'\in G,\ \bar h\in Q.
\]
If \(\alpha\) is a group action and \(\bar\alpha\) is the descended action,
then for \(\bar h=q(h)\), which covers all of \(Q\) since \(q\) is surjective,
\[
  \bar\alpha(gg',q(h))=q(\alpha(gg',h))=q(\alpha(g,\alpha(g',h)))
  =\bar\alpha(g,q(\alpha(g',h)))=\bar\alpha(g,\bar\alpha(g',q(h))),
\]
so the group law holds. For an action descended from \(\alpha\), surjectivity of
\(q\) also gives \(\bar\alpha(e,q(h))=q(\alpha(e,h))=q(h)\), so the induced map
is a group action. The readout is \(G\)-equivariant exactly when
\[
  r(\alpha(g,h))=\rho(g,r(h))
  \quad
  \forall g\in G,\ h\in H.
\]
If \(\bar\alpha(g,q(h))=q(\alpha(g,h))\) for all \(g,h\) and \(q(h)=q(h')\),
then \(q(\alpha(g,h))=\bar\alpha(g,q h)=\bar\alpha(g,q h')=q(\alpha(g,h'))\),
so \(\mathcal O_\alpha=\emptyset\); conversely a descended action requires
\(\mathcal O_\alpha=\emptyset\). If \(\mathcal O_\alpha=\emptyset\), the induced
action is well defined by the first equivalence and is a descended action, so
it preserves the group law by the computation above; the group law is
therefore not a separate source of anomaly, and an anomaly, as defined in the
statement, is exactly the failure of one of the four conditions for the
induced action. For the induced action, action descent and the group law
hold. In the claimed, unbroken case, an anomaly is therefore exactly a failure
of domain descent or readout preservation. An unclaimed or explicitly broken
symmetry is outside this anomaly designation.
\end{proof}

\paragraph{Nonclaims.} The theorem does not classify anomalies
into carrier-specific forms and does not assert anything about
non-symmetric gates (maps on the carrier, in the sense of
\Cref{sec:apparatus}, that do not commute with the action).  The standard quantum-field-theory treatment of
anomalies, including the Wess--Zumino consistency condition and the
descent-equation classification, is collected in
\citet{Bertlmann2000Anomalies}.

\subsection*{Layer instantiations}\label{layer-inst:F40}

\begingroup\small
\begin{description}
\item[\emph{{IEEE} 754 floating-point non-associativity (numerical
computing).}] \textit{Analogy.}
Associativity of real addition fails in IEEE 754 arithmetic, although addition of finite operands in a fixed format and rounding mode is commutative: \((a+b)+c\) generically differs from \(a+(b+c)\) because each operation rounds. Rebracketing is not a group action; read loosely, with rebracketings of a sum as the would-be symmetry and the rounded value as the readout, the failure is one of ReadoutPreserved; numerical
libraries declare explicit accumulation orders (Kahan or
pairwise summation) as the apparatus extension that mitigates the
anomaly
\citep{IEEE7542019FloatingPoint}.

\end{description}
\endgroup


\subsection{Duality Equivalence [\texorpdfstring{\Fref{F41}}{F41}]}\label{sec:symmetry:duality-equivalence}

Let \(\mathcal L_A\) and \(\mathcal L_B\) be formed closures with
declared symmetry presentations.  Each presentation has an object
quotient, transport data, declared observables, role registry,
residual registry, and status discipline.  A proposed duality is not
just a dictionary between symbols.  It is a pair of translations
between the two formed packages, with the claim that the translations
preserve the roles being compared.

\paragraph{Reading the statement.}
\(\mathcal K_A\) and \(\mathcal K_B\) are the claims of the two
packages, and \(\operatorname{status}_A\), \(\operatorname{status}_B\)
are their declared status assignments, which give each claim its
status in the package's status discipline. The last condition says
that translating a claim carries its status to the corresponding
status of the other package under \(\varsigma\). A \emph{role-preserving
duality equivalence} is defined as the conjunction of the four displayed
conditions; the role readouts live on the quotients, so the role condition
constrains the duality maps themselves. The theorem records this definition
together with its consequences: the inverse map \(\bar G\) also carries routes
to routes and roles back, it is the only two-sided inverse of \(\bar F\), the
role condition fixes \(\rho\) on the image of \(r_A\), and \(F\) and \(G\)
preserve and reflect the quotient classes. These consequences use only the first three conditions together with the descent equations: uniqueness of \(\bar G\) and the fiber clause use the first, the route clause the first and third, and the role clauses the first and second; the status condition, and with it \(\kappa\), \(\varsigma\) and the status maps, enters only the definition.

\begin{theorem}[Duality Equivalence]\label{thm:F41}
Let \(\mathcal C_A,\mathcal C_B\) be sets (in the intended reading, formed closures) with quotients
\(q_A:\mathcal C_A\to Q_A\) and \(q_B:\mathcal C_B\to Q_B\).  Let
\(F:\mathcal C_A\to\mathcal C_B\) and
\(G:\mathcal C_B\to\mathcal C_A\) descend to
\(\bar F:Q_A\to Q_B\) and \(\bar G:Q_B\to Q_A\), meaning explicitly
\(\bar F\circ q_A=q_B\circ F\) and
\(\bar G\circ q_B=q_A\circ G\).  Let
\(r_A:Q_A\to S_A\), \(r_B:Q_B\to S_B\) be role readouts on the quotients,
let \(\rho:S_A\to S_B\) be the role-value map, and let
\(\phi:\Gamma_A\to\Gamma_B\) and
\(\kappa:\mathcal K_A\to\mathcal K_B\) translate routes and claims, let
\(\operatorname{status}_A:\mathcal K_A\to\Sigma_A\) and
\(\operatorname{status}_B:\mathcal K_B\to\Sigma_B\) be the declared
claim statuses, and let \(\varsigma:\Sigma_A\to\Sigma_B\) be the
declared status-value map,
where a route \(\gamma\in\Gamma_A\) acts on \(Q_A\) by
\(\mathrm{act}_A(\gamma):Q_A\to Q_A\) and a route
\(\delta\in\Gamma_B\) acts on \(Q_B\) by
\(\mathrm{act}_B(\delta):Q_B\to Q_B\).
The data form a \emph{role-preserving duality equivalence} when
\begin{align*}
  \bar G\circ\bar F&=\mathrm{id}_{Q_A}
  \quad\text{and}\quad
  \bar F\circ\bar G=\mathrm{id}_{Q_B},\\
  r_B\circ\bar F&=\rho\circ r_A,\\
  \bar F\circ \mathrm{act}_A(\gamma)
  &=\mathrm{act}_B(\phi(\gamma))\circ\bar F
  \quad\forall \gamma\in\Gamma_A,\\
  \operatorname{status}_B(\kappa(k))
  &=\varsigma(\operatorname{status}_A(k))
  \quad\forall k\in\mathcal K_A.
\end{align*}
Then, for a role-preserving duality equivalence, the inverse also
intertwines the routes,
\(\bar G\circ\mathrm{act}_B(\phi(\gamma))=\mathrm{act}_A(\gamma)\circ\bar G\)
for every \(\gamma\in\Gamma_A\), and it transports the roles back:
\(r_B=\rho\circ r_A\circ\bar G\). Moreover \(\bar G\) is the only two-sided inverse of
\(\bar F\), and \(\rho(r_A(a))=r_B(\bar F(a))\) for every \(a\in Q_A\), so if \(r_A\)
is surjective then \(\rho\) is determined by \(\bar F\), \(r_A\) and \(r_B\).
Consequently \(q_A(x)=q_A(x')\iff q_B(Fx)=q_B(Fx')\) for all
\(x,x'\in\mathcal C_A\), and \(q_B(y)=q_B(y')\iff q_A(Gy)=q_A(Gy')\) for all
\(y,y'\in\mathcal C_B\).
\end{theorem}
\begin{proof}
The four displayed conditions define a role-preserving duality equivalence:
the first line makes \(Q_A\) and \(Q_B\) equivalent quotient presentations,
the second transports role observables, the third transports routes, and the
fourth transports claim statuses. We prove the consequences.

For the routes, let \(y\in Q_B\). Using \(\bar F\bar G=\mathrm{id}\),
then the route condition, then \(\bar G\bar F=\mathrm{id}\):
\[
  \bar G\,\mathrm{act}_B(\phi\gamma)\,y
  =\bar G\,\mathrm{act}_B(\phi\gamma)\,\bar F\bar G\,y
  =\bar G\bar F\,\mathrm{act}_A(\gamma)\,\bar G\,y
  =\mathrm{act}_A(\gamma)\,\bar G\,y .
\]
For the roles, \(r_B=r_B\circ\bar F\circ\bar G=\rho\circ r_A\circ\bar G\).
If \(\bar G'\) is another two-sided inverse of \(\bar F\), then \(\bar G'=\bar G'\circ\bar F\circ\bar G=\bar G\). The identity \(\rho(r_A(a))=r_B(\bar F(a))\) is the role condition read pointwise, and it fixes \(\rho\) on the image of \(r_A\).
For the fiber clause: since \(q_B\circ F=\bar F\circ q_A\) and \(\bar F\) is a
bijection, \(q_B(Fx)=q_B(Fx')\) exactly when \(q_A(x)=q_A(x')\); likewise for
\(G\).
\end{proof}

\paragraph{Nonclaims.} The theorem does not classify
symmetry presentations and does not assert a unique role presentation
beyond observable equivalence.

\subsection*{Layer instantiations}\label{layer-inst:F41}

\begingroup\small
\begin{description}
\item[\emph{Pontryagin duality (harmonic analysis).}] \textit{Analogy.}
For a locally compact abelian group \(G\), the Pontryagin dual
\(\widehat G\) of continuous characters \(G\to S^1\) satisfies the
Pontryagin duality theorem
\(G\cong\widehat{\widehat G}\), with the Fourier transform
\(L^2(G)\to L^2(\widehat G)\) a unitary isomorphism.  Pontryagin dualization is a contravariant equivalence: for closed \(H\), the
dual of \(G/H\) is the annihilator subgroup \(H^\perp\subseteq\widehat G\), and
double dualization recovers \(G\). For functions with defined convolution,
Fourier transform exchanges convolution and pointwise multiplication; this is
an analogy with reversible translation, with quotient variance and route data
specified separately \citep{Pontryagin1966TopologicalGroups}.

\end{description}
\endgroup


\subsection*{Symmetry at a glance}\label{sec:symmetry:summary}\addcontentsline{toc}{subsection}{Symmetry at a glance}
\begin{table}[H]
  \centering
  \footnotesize
  \begin{tabularx}{\textwidth}{@{}
      >{\raggedright\arraybackslash}p{0.24\textwidth}
      >{\raggedright\arraybackslash}X
      >{\raggedright\arraybackslash}p{0.12\textwidth}@{}}
    \toprule
    Name & One-line claim & Substrate cite \\
    \midrule
    Duality-Fixity Theorem (\Fref{F14}) &
    For a readout that is strongly measurable and square-integrable when the weight is a measure, and a positive cone satisfying (C1)--(C3) (monotone trace, positive trace-zero elements vanish, trace identity), domination by bounds \(B_n\) with \(\inf_n\operatorname{tr}B_n\le0\) forces the anti-invariant
    ledger to vanish and, for a separating readout, confines visible mass
    to \(\Fixinv\) up to a null set; such bounds exist exactly when \(\psi^-=0\) \(\mu\)-almost everywhere.  &
    \PaperSDTC, \PaperRH \\
    Selection-Obstruction Theorem (\Fref{F15a}) &
    For the orbit quotient with its induced action, an equivariant selector
    takes only fixed points; selectedness descends through the quotient exactly
    when its obstruction is empty.  &
    \textemdash \\
    Spontaneous Symmetry Breaking (\Fref{F15b}) &
    A branch-selection record exists exactly when some ground state is
    not invariant; no symmetric selector produces its branch.  &
    \PaperCSL\ (parallel) \\
    Anomaly Theorem (\Fref{F40}) &
    For a claimed symmetry not recorded as explicitly broken, an anomaly
    is the failure of action descent, domain descent or readout
    preservation; for a surjective layer quotient, a descended action is a
    group action (it satisfies the group law and the identity acts
    trivially).  &
    \textemdash \\
    Duality-Equivalence Theorem (\Fref{F41}) &
    In a role-preserving duality equivalence, the inverse also
    intertwines the routes and \(r_B=\rho\circ r_A\circ\bar G\).  &
    \textemdash \\
    \bottomrule
  \end{tabularx}
  \caption{Symmetry at a glance: the five laws of \Cref{sec:symmetry}. Duality Fixity has restricted alignment.}
  \label{tab:axis-symmetry-summary}
\end{table}

\FloatBarrier
 % sec:symmetry
\section{Status, Records, Coherence}\label{sec:status-records-coherence}

This chapter treats records as part of the mathematical object rather
than as commentary about it.  A formed fact has a claim record, a
status, and a place in a coherence ledger that records how it
persists, transports, composes and supports explanation; the laws of
the chapter govern that bookkeeping.  In the running example of
\Cref{sec:intro}, the record of the split pair \((1,2)\), with its
status and the repair chosen, is part of the object these laws
describe.

The fourteen laws run from shared objects to counterfactual use.
Objectivity as Common Quotient identifies when several presentations
share an object; Lawhood Descent tests when a law factors through the
available quotient; Object Persistence tests when an object survives
transport; and Fact / Record Formation says what it means for a fact
or record to form.  Reflexive Nonclosure, Interface Mediation,
Probability as Stable Unresolved Fiber and Causality as Stable
Intervention Transport describe how records behave under self-audit,
interfaces, unresolved fibers and interventions.  The remaining laws
concern coherence: Moduli / Contingency separates stable parameters
from contingent choices, Conservation as Orbit Descent reads
conservation as a quotient law on orbits, Composability and
Presentation Invariance say when records survive composition and change
of presentation, Explanation as Compression with Descent treats
explanation as compression that still factors through the declared
quotient, and Counterfactual Legitimacy says when a counterfactual
claim is lawful rather than an overread.

The chapter uses only the vocabulary of \Cref{sec:apparatus}.  The
summary table at the end groups the laws into three reading groups,
which do not follow document order: status (\Fref{F17}, \Fref{F18},
\Fref{F22}, \Fref{F23}), records (\Fref{F19}, \Fref{F20}, \Fref{F21},
\Fref{F25}) and coherence (\Fref{F26}--\Fref{F31}).  The chapter's axes
are status, records and interfaces (\Cref{subsec:reading-guide});
coherence is a reading group, not an axis.

\subsection{Objectivity as Common Quotient [\texorpdfstring{\Fref{F17}}{F17}]}\label{sec:status-records-coherence:objectivity}

\paragraph{Reading the statement.} The symbols are those of \Cref{thm:F17} below.
Here ``\(r\) descends through \(q\)'' means that \(r=\bar r\circ q\) for
some map \(\bar r\), as in \Cref{sec:apparatus}.
\(\operatorname{CommonAcrossAccesses}(q,r)\) says that \(r\) descends
through every access quotient \(q_i\), and
\(\operatorname{PublicDeterminesObject}(p,r)\) says that \(r\) descends
through the public quotient \(p\).
\(\operatorname{AccessObjectivityDefectEmpty}(q_i,r)\) says that no two
histories with the same \(q_i\)-class have different \(r\)-values.
\(\operatorname{PublicToObjectDefect}(p,r)\) is the set of pairs with
the same \(p\)-class and different \(r\)-values, and
\(\operatorname{ObstructionEmpty}\) says that this set is empty.
\(\operatorname{CommonQuotientVisible}(q,p)\) says that \(p\) itself
descends through every \(q_i\). The objective extension
\(\operatorname{objectiveExtension}(p,r)\) is the pairing
\(h\mapsto(p(h),r(h))\); it \emph{retains} \(p\), and \emph{carries}
\(r\), when \(p\), respectively \(r\), descends through it.
The status partition
\(\operatorname{ExactlyOneObjectivityStatus}(p,r)\) says that exactly
one of four statuses holds, according to whether \(r\) descends through
\(p\) and whether \(p\) descends through \(r\): exact common quotient when
both hold, coarse objective object when only \(r\) descends through \(p\),
private overrefinement when only \(p\) descends through \(r\), and
incomparable object claim when neither does.

\begin{theorem}[Objectivity as Common Quotient]\label{thm:F17}
Let \(H\) be a carrier with a family of access quotients
\(q_i:H\to Q_i\) indexed by \(i\in I\), a public quotient
\(p:H\to P\), and a target readout \(r:H\to R\).  Assume each
\(q_i\) is surjective, \(p\) is surjective, and the family satisfies
the common-quotient universal property
\(\operatorname{CommonQuotientUniversal}(q,p)\): \(p\) factors
through every \(q_i\), and any map \(s:H\to S\) that factors through every
\(q_i\) also factors through \(p\).  Then
\begin{align*}
  &\operatorname{CommonAcrossAccesses}(q,r)
    \iff \operatorname{PublicDeterminesObject}(p,r),\\
  &\operatorname{CommonAcrossAccesses}(q,r)\\
  &\qquad{}\iff
    \forall i\in I,\,
      \operatorname{AccessObjectivityDefectEmpty}(q_i,r),\\
  &\operatorname{PublicDeterminesObject}(p,r)\\
  &\qquad{}\iff
    \operatorname{ObstructionEmpty}
      \bigl(\operatorname{PublicToObjectDefect}(p,r)\bigr),
\end{align*}
together with
\begin{align*}
  &\operatorname{CommonQuotientVisible}(q,p),\\
  &\operatorname{RetainsPublicQuotient}
      \bigl(p,\operatorname{objectiveExtension}(p,r)\bigr),\\
  &\operatorname{CarriesObject}
      \bigl(\operatorname{objectiveExtension}(p,r),r\bigr),
\end{align*}
and the status partition
\[
  \operatorname{ExactlyOneObjectivityStatus}(p,r).
\]
In ordinary terms: \(p\) factors through every \(q_i\); the pairing
\(\langle p,r\rangle\) satisfies \(p=\mathrm{pr}_1\circ\langle p,r\rangle\) and
\(r=\mathrm{pr}_2\circ\langle p,r\rangle\); and exactly one of the following holds:
both \(r\) factors through \(p\) and \(p\) through \(r\) (exact common quotient);
only the first (coarse objective object); only the second (private
overrefinement); neither (incomparable).
\end{theorem}
\begin{proof}
Assume the surjectivity of each \(q_i\), the surjectivity of \(p\),
and \(\operatorname{CommonQuotientUniversal}(q,p)\).  We prove the
seven conjuncts in turn.

\emph{(1) Common across accesses iff public determines object.}
The universal property supplies comparison maps
\[
  p=\pi_i\circ q_i\quad(i\in I),
  \qquad
  s=\lambda\circ p
  \quad\text{whenever } s=\mu_i\circ q_i \text{ for every } i\in I .
\]
Thus \(p\) is the finest quotient that factors through every \(q_i\).
If \(r\) is common across the access quotients, then \(r\) factors through
every \(q_i\) by definition, hence through
\(p\), by the universal property applied to the map \(r\).  Conversely, if \(r\) factors through \(p\), it factors through
every \(q_i\), since \(p=\pi_i\circ q_i\).  So a readout \(r\) is common
across the access quotients exactly when it is constant on the common
public fibers:
\[
  \forall h,h'\in H,\quad
  p(h)=p(h')\Longrightarrow r(h)=r(h').
\]
Since \(p\) is surjective, this is equivalent to the existence of
\(\bar r:P\to R\) with
\[
  r=\bar r\circ p,
\]
which is precisely \(\operatorname{PublicDeterminesObject}(p,r)\).

\emph{(2) Common across accesses iff every access defect is empty.}
For \(i\in I\), the access objectivity defect is the set
\[
  \Delta_i(q_i,r)
  =
  \{(h,h')\in H\times H:
      q_i(h)=q_i(h')\ \text{and}\ r(h)\neq r(h')\}.
\]
Since \(q_i\) is surjective, the descent criterion of
\Cref{sec:apparatus} shows that \(\Delta_i(q_i,r)=\emptyset\) is
equivalent to \(r=\bar r_i\circ q_i\) for some \(\bar r_i\); requiring
this for every \(i\in I\) is \(\operatorname{CommonAcrossAccesses}(q,r)\),
giving the second equivalence.

\emph{(3) Public determines object iff the public-to-object
obstruction is empty.}
The public-to-object defect is
\[
  \Delta_P(p,r)
  =
  \{(h,h')\in H\times H:
      p(h)=p(h')\ \text{and}\ r(h)\neq r(h')\}.
\]
By the descent criterion, since \(p\) is surjective, its emptiness is
equivalent to \(r=\bar r\circ p\) for some \(\bar r:P\to R\), that is, to
\(\operatorname{PublicDeterminesObject}(p,r)\).

\emph{(4) Common quotient visible.}
The same universal property exhibits \(p\) as a declared common
quotient of the family \((q_i)_{i\in I}\):
\[
  \operatorname{CommonQuotientUniversal}(q,p)
  \Longrightarrow
  \operatorname{CommonQuotientVisible}(q,p).
\]
Here visibility is the first half of the universal property: \(p\)
descends through every \(q_i\), \(p=\pi_i\circ q_i\).

\emph{(5) The objective extension retains the public quotient.}
Let
\[
  \operatorname{objectiveExtension}(p,r)
  :=
  \langle p,r\rangle:H\longrightarrow P\times R,
  \qquad
  h\longmapsto (p(h),r(h)).
\]
Projection to the first coordinate recovers the public quotient:
\[
  \operatorname{pr}_P\circ\langle p,r\rangle=p.
\]
Therefore the objective extension retains \(p\).

\emph{(6) The objective extension carries the object.}
Projection to the second coordinate recovers the object readout:
\[
  \operatorname{pr}_R\circ\langle p,r\rangle=r.
\]
Thus the same extension carries the object \(r\).

\emph{(7) Exactly one objectivity status.}
The four status predicates compare the object \(r\) with the public
quotient \(p\) in both directions: whether \(r\) descends through \(p\)
(the public quotient determines the object) and whether \(p\) descends
through \(r\) (the object determines the public quotient).  The cases
are
\[
\begin{array}{c|c|c}
\text{status} & r\ \text{descends through}\ p
  & p\ \text{descends through}\ r\\
\hline
\text{exact common quotient} & \text{yes} & \text{yes}\\
\text{coarse objective object} & \text{yes} & \text{no}\\
\text{private overrefinement} & \text{no} & \text{yes}\\
\text{incomparable object claim} & \text{no} & \text{no}
\end{array}
\]
Every configuration of the two Boolean tests falls into one row, and
no configuration falls into two rows.  Hence
\(\operatorname{ExactlyOneObjectivityStatus}(p,r)\).
\end{proof}

\paragraph{Nonclaims.} The theorem does not determine which presentations are declarable and
does not enumerate the kinds of objects that may be shared.  The
closest physical analogy is quantum Darwinism, in which
objectivity emerges as redundant agreement of the environment's
records across many access channels \citep{Zurek2009QuantumDarwinism};
the layer-agnostic statement above is independent of that
substrate.

\subsection*{Layer instantiations}\label{layer-inst:F17}

\begingroup\small
\begin{description}
\item[\emph{Fixed-effect meta-analysis (statistics).}] \textit{Analogy.}
Under a declared fixed-common-effect model, a fixed-effect
meta-analysis combines independent study-level effect-size
estimates by inverse-variance weighting to produce a pooled
consensus estimate computed from the study-level estimates; when the
model holds, every study estimates the same common effect, and the
pooled estimate is its public estimator.  Cochran's \(Q\) (and the later \(I^2\) index of Higgins and Thompson) measures the
\Fref{F17} access-objectivity defect: heterogeneity is the witness that
the declared common-effect model is violated
\citep{HedgesOlkin1985MetaAnalysis}.

\end{description}
\endgroup


\subsection{Lawhood Descent [\texorpdfstring{\Fref{F18}}{F18}]}\label{sec:status-records-coherence:lawhood-descent}

\paragraph{Reading the statement.}
A candidate law here is a predicate or readout \(P:H\to V\). It is
\emph{lawful} at the layer when it is invariant under the symmetries of the
access quotient: \(P\circ\sigma=P\) for every bijection \(\sigma:H\to H\)
with \(\accessQ\circ\sigma=\accessQ\). For a relation or a transition rule,
lawfulness means factoring through \(\accessQ\) in each argument; invariance
under such \(\sigma\) is weaker there, since every such \(\sigma\) preserves
the equality relation on \(H\), which need not factor.

\begin{theorem}[Lawhood Descent]\label{thm:F18}
Let \(\accessQ:H\to Q\) be a surjective access quotient (the available access of
a formed closure in the intended reading).  A declared candidate law \(P:H\to V\) satisfies \(P\circ\sigma=P\) for
every bijection \(\sigma:H\to H\) with \(\accessQ\circ\sigma=\accessQ\) (that is,
\(P\) is lawful at the layer) if and only if \(P\) factors through \(\accessQ\),
that is, \(P(h)=P(h')\) whenever \(\accessQ(h)=\accessQ(h')\).
\end{theorem}
\begin{proof}
If \(P=\bar P\circ\accessQ\) and \(\accessQ\circ\sigma=\accessQ\), then
\(P\circ\sigma=\bar P\circ\accessQ\circ\sigma=\bar P\circ\accessQ=P\).
Conversely, let \(P\) be lawful and \(\accessQ(h)=\accessQ(h')\) with
\(h\ne h'\). The transposition \(\sigma\) exchanging \(h\) and \(h'\) is a
bijection with \(\accessQ\circ\sigma=\accessQ\), so
\(P(h')=P(\sigma(h))=P(h)\). Hence \(P\) is constant on access classes;
define \(\bar P(\accessQ(h))=P(h)\), which surjectivity defines on all of
\(Q\).  Constancy makes this well-defined, and
\(P=\bar P\circ\accessQ\).  Thus lawhood is precisely quotient descent.
\end{proof}

\paragraph{Nonclaims.} The theorem does not enumerate lawful candidates and does not
construct the access quotient; that is the role of Quotientality.

\subsection*{Layer instantiations}\label{layer-inst:F18}

\begingroup\small
\begin{description}
\item[\emph{Wilson loops in gauge theory (physics).}] \textit{Special case.}
In Yang--Mills theory, a candidate physical observable is a
lawful observable at the gauge-quotient iff it is
gauge-invariant.  Wilson loops --- traces of the holonomy around
closed loops --- are the canonical gauge-invariant
observables, descending through the gauge-orbit projection
\(\mathcal A\to\mathcal A/\mathcal G\)
\citep{Wilson1974ConfinementQuarks}.

\item[\emph{Parnas information hiding (software engineering).}] \textit{Analogy.}
Parnas's principle of information hiding declares that a module
exposes a public interface and hides its implementation behind
that interface.  A client predicate is a lawful use of the
module iff it depends only on the declared public operations;
predicates that read hidden state are overreads.  The interface
is the \Fref{F18} access quotient through which lawful client claims
descend
\citep{Parnas1972InformationHiding}.

\item[\emph{Bisimulation invariance in process calculi
(concurrency theory).}] \textit{Special case.}
In Milner--Park bisimulation theory, every Hennessy--Milner
modal-logic formula is bisimulation-invariant, and on image-finite
labelled transition systems two processes are bisimilar iff they
satisfy the same formulas.  A process property is lawful iff it is
bisimulation-invariant, and the bisimulation quotient on the labelled
transition system is the \Fref{F18} access quotient through which
lawful process-level laws, the Hennessy--Milner formulas among them,
descend
\citep{Milner1989CommunicationConcurrency}.

\item[\emph{Congruences and quotient algebras (universal
algebra).}] \textit{Special case.}
A predicate on an algebra \(A\) descends to a well-defined predicate
on the quotient \(A/{\sim}\) iff it is \(\sim\)-invariant, and the
operations of \(A\) descend to \(A/{\sim}\) iff \(\sim\) is a
congruence (compatible with all operations); for a congruence
\(\sim\), laws over the quotient algebra correspond exactly to
\(\sim\)-invariant predicates on the original algebra
\citep{Gratzer1979UniversalAlgebra}.

\end{description}
\endgroup


\subsection{Object Persistence [\texorpdfstring{\Fref{F19}}{F19}]}\label{sec:status-records-coherence:object-persistence}

\paragraph{Reading the statement.} The symbols are those of \Cref{thm:F19} below.
\(\operatorname{Persists}(q_i,\gamma,q_j)\) says that the later class
\(q_j(\gamma h)\) depends only on the earlier class \(q_i(h)\), that
is, \(q_j\circ\gamma\) descends through \(q_i\). The persistence
obstruction, also called the split obstruction, is the set of pairs
with the same \(q_i\)-class whose transports have different
\(q_j\)-classes; the \emph{ObstructionEmpty} predicates say that the
named set is empty. \(\operatorname{ViablePersistence}\) says that every
transported class is viable, and the dissolution obstruction is the set
of histories whose transported class is not viable.
\(\operatorname{RelationValuedPersistence}(q_i,q_j,C)\) says that
whether \(C(h,k)\) holds depends only on the classes \(q_i(h)\) and
\(q_j(k)\); its obstruction is the set of pairs of pairs with the same
classes on which \(C\) differs.
\(\operatorname{IdentityHolonomyCoherent}\) says that
\(\operatorname{compBar}=\bar\eta\circ\bar\gamma\), and the holonomy
defect is the set of classes where the two differ. The comparison
status compares the split obstruction with the merge obstruction, the
pairs merged by the transport that \(q_i\) separates.  The persistence
extension is the pairing \(h\mapsto(q_i(h),q_j(\gamma h))\); it
\emph{retains} the old identity \(q_i\), and \emph{carries} the later
identity \(q_j\circ\gamma\), when that map descends through it.

\begin{theorem}[Object Persistence]\label{thm:F19}
Let \(H_i,H_j\) be carriers with access quotients
\(q_i:H_i\to Q_i\) and \(q_j:H_j\to Q_j\), and let
\(\gamma:H_i\to H_j\) be a declared transport.  Assume \(q_i\) is surjective.  Let
\(\operatorname{viable}:Q_j\to\mathrm{Prop}\) be a viability
predicate, \(C:H_i\times H_j\to\mathrm{Prop}\) a relation predicate,
let \(Q_k\) be a declared target set, and let
\[
  \bar\gamma:Q_i\to Q_j,\qquad
  \bar\eta:Q_j\to Q_k,\qquad
  \operatorname{compBar}:Q_i\to Q_k
\]
be arbitrary holonomy data (no compatibility between
\(\operatorname{compBar}\) and \(\bar\eta\circ\bar\gamma\) is
assumed).  Then
\begin{align*}
  &\operatorname{Persists}(q_i,\gamma,q_j)
    \iff
    \operatorname{PersistenceObstructionEmpty}(q_i,\gamma,q_j)\\
  &\qquad{}\iff
    \operatorname{SplitObstructionEmpty}(q_i,\gamma,q_j),\\
  &\operatorname{ViablePersistence}(\gamma,q_j,\operatorname{viable})
    \iff
    \operatorname{DissolutionObstructionEmpty}
      (\gamma,q_j,\operatorname{viable}),\\
  &\operatorname{RelationValuedPersistence}(q_i,q_j,C)
    \iff
    \operatorname{RelationPersistenceObstructionEmpty}(q_i,q_j,C),\\
  &\operatorname{IdentityHolonomyCoherent}
      (\bar\gamma,\bar\eta,\operatorname{compBar})\\
  &\qquad{}\iff
    \operatorname{IdentityHolonomyDefectEmpty}
      (\bar\gamma,\bar\eta,\operatorname{compBar}),
\end{align*}
where \(\operatorname{ViablePersistence}:\Longleftrightarrow\forall h,\
\operatorname{viable}(q_j(\gamma h))\);
\(\operatorname{RelationValuedPersistence}:\Longleftrightarrow C=\bar C\circ(q_i\times
q_j)\) for some \(\bar C\); and
\(\operatorname{IdentityHolonomyCoherent}:\Longleftrightarrow\operatorname{compBar}=\bar\eta\circ\bar\gamma\);
together with the comparison-status partition
\[
  \operatorname{ExactlyOnePersistenceComparisonStatus}
    (q_i,\gamma,q_j),
\]
and the retention and carrying conditions
\begin{align*}
  &\operatorname{RetainsOldIdentity}
      \bigl(q_i,\operatorname{persistenceExtension}(q_i,\gamma,q_j)\bigr),\\
  &\operatorname{CarriesLaterIdentity}
      \bigl(\operatorname{persistenceExtension}(q_i,\gamma,q_j),
        \gamma,q_j\bigr).
\end{align*}
If moreover \(\eta:H_j\to H_k\) and \(q_k:H_k\to Q_k\) satisfy
\(\bar\gamma\circ q_i=q_j\circ\gamma\), \(\bar\eta\circ q_j=q_k\circ\eta\) and
\(\operatorname{compBar}\circ q_i=q_k\circ\eta\circ\gamma\), then
\(\operatorname{IdentityHolonomyCoherent}(\bar\gamma,\bar\eta,\operatorname{compBar})\)
holds; hence, when the holonomy defect is nonempty, one of the three
displayed commutation equations fails.
Here the persistence and split obstructions are the same set
\(\Delta_{\mathrm{split}}=\{(h,h'):q_i(h)=q_i(h'),\ q_j(\gamma h)\ne q_j(\gamma h')\}\), and
\(\operatorname{Persists}(q_i,\gamma,q_j)\) says that \(q_j\circ\gamma\)
factors through \(q_i\).
Here \(\operatorname{ExactlyOnePersistenceComparisonStatus}\) says that exactly
one of exact endurance (\(\Delta_{\mathrm{split}}=\Delta_{\mathrm{merge}}=\emptyset\)),
persistent with merge (only \(\Delta_{\mathrm{merge}}\ne\emptyset\)), identity
branching (only \(\Delta_{\mathrm{split}}\ne\emptyset\)) and identity rewrite (both
nonempty) holds, where \(\Delta_{\mathrm{merge}}=\{(h,h'):q_j(\gamma h)=q_j(\gamma
h'),\ q_i(h)\ne q_i(h')\}\).
\end{theorem}
\begin{proof}
Assume the surjectivity of \(q_i\), the predicates
\(\operatorname{viable}\) and \(C\), and the holonomy data
\(\bar\gamma,\bar\eta,\operatorname{compBar}\).  We prove the eight
conjuncts in turn. Part (3) proves the comparison-status partition, which the statement lists after the holonomy line, and part (6) also proves the added holonomy clause.

\emph{(1) Persists iff the persistence obstruction is empty.}
The persistence obstruction is
\[
  \Delta_{\mathrm{pers}}(q_i,\gamma,q_j)
  =
  \{(h,h')\in H_i\times H_i:
      q_i(h)=q_i(h')\ \text{and}\
      q_j(\gamma h)\ne q_j(\gamma h')\}.
\]
Since \(q_i\) is surjective, the descent criterion of
\Cref{sec:apparatus} applied to \(q_j\circ\gamma\) shows that its
emptiness is equivalent to the existence of a map
\(\gamma^{\flat}:Q_i\to Q_j\) satisfying
\[
  \gamma^{\flat}\circ q_i=q_j\circ\gamma,
\]
which is \(\operatorname{Persists}(q_i,\gamma,q_j)\).

\emph{(2) Split-obstruction empty iff persists.}
The split obstruction is the same pair set written in the comparison
vocabulary:
\[
  \Delta_{\mathrm{split}}(q_i,\gamma,q_j)
  =
  \Delta_{\mathrm{pers}}(q_i,\gamma,q_j).
\]
Thus \(\Delta_{\mathrm{split}}=\emptyset\) iff the persistence
obstruction is empty, and by (1) this is equivalent to persistence.

\emph{(3) Exactly-one persistence comparison status.}
The comparison statuses are decided by two split-pair defects on the
declared continuation \(\gamma\): the persistence/split obstruction
\(\Delta_{\mathrm{split}}\) of case~(1), and the merge obstruction
\[
  \Delta_{\mathrm{merge}}(q_i,\gamma,q_j)
  =
  \{(h,h')\in H_i\times H_i:
      q_j(\gamma h)=q_j(\gamma h')\ \text{and}\ q_i(h)\ne q_i(h')\}.
\]
The four statuses partition the configurations by the joint
\(\Delta_{\mathrm{split}}\)\textendash{}\(\Delta_{\mathrm{merge}}\)
emptiness pattern:
\[
\begin{array}{c|c|c}
\text{status} & \Delta_{\mathrm{split}}=\emptyset
  & \Delta_{\mathrm{merge}}=\emptyset\\
\hline
\text{exact endurance} & \text{yes} & \text{yes}\\
\text{persistent with merge} & \text{yes} & \text{no}\\
\text{identity branching} & \text{no} & \text{yes}\\
\text{identity rewrite} & \text{no} & \text{no}
\end{array}
\]
Every configuration of the two Boolean defect tests falls into
exactly one row.  Existence is by case enumeration; exclusivity is by
definitional disjointness of the four status predicates.  The
dissolution and relation-valued obstructions of cases~(4) and~(5) are
separate conditions and do not enter this partition.

\emph{(4) Viable persistence iff the dissolution obstruction is
empty.}
The dissolution obstruction is
\[
  \Delta_{\mathrm{diss}}(\gamma,q_j,\operatorname{viable})
  =
  \{h\in H_i:
      \neg\,\operatorname{viable}(q_j(\gamma h))\}.
\]
Viable persistence says that every transported \(q_j\)-class lies in
the viable region:
\[
  \forall h\in H_i,\quad
  \operatorname{viable}(q_j(\gamma h)).
\]
This is precisely \(\Delta_{\mathrm{diss}}=\emptyset\).

\emph{(5) Relation-valued persistence iff its obstruction is empty.}
The relation-persistence obstruction is
\[
  \Delta_C(q_i,q_j,C)
  =
  \left\{
  \begin{array}{l}
    ((h,k),(h',k'))\in(H_i\times H_j)^2:\\
    q_i(h)=q_i(h'),\ q_j(k)=q_j(k'),\
    C(h,k)\not\iff C(h',k')
  \end{array}
  \right\}.
\]
Emptiness of \(\Delta_C\) says that the truth of \(C(h,k)\) depends only
on the pair of classes \((q_i(h),q_j(k))\).  Setting \(\bar C:=\bot\) off the image of \(q_i\times q_j\), this holds exactly when \(C\) descends to a
well-defined relation on \(Q_i\times Q_j\), which is relation-valued
persistence; this part uses no surjectivity, and surjectivity of \(q_j\) is not assumed.

\emph{(6) Identity holonomy coherent iff the holonomy defect is
empty.}
The holonomy defect is
\[
  \Delta_{\mathrm{hol}}(\bar\gamma,\bar\eta,\operatorname{compBar})
  =
  \{a\in Q_i:
      \operatorname{compBar}(a)\ne\bar\eta(\bar\gamma(a))\}.
\]
Thus \(\Delta_{\mathrm{hol}}=\emptyset\) iff
\[
  \forall a\in Q_i,\quad
  \operatorname{compBar}(a)=\bar\eta(\bar\gamma(a)),
\]
which is the pointwise form of
\(\operatorname{compBar}=\bar\eta\circ\bar\gamma\), i.e. identity
holonomy coherence. Under the descent equations of the added clause,
\(\operatorname{compBar}(q_ih)=q_k\eta\gamma h=\bar\eta(q_j\gamma h)
=\bar\eta\bar\gamma(q_ih)\) for every \(h\), and \(q_i\) is surjective, so
the defect is empty.

\emph{(7) The persistence extension retains the old identity.}
Define the persistence extension as the joint quotient
\[
  \operatorname{persistenceExtension}(q_i,\gamma,q_j)
  :=
  \langle q_i,q_j\circ\gamma\rangle:H_i\to Q_i\times Q_j.
\]
Projection to the first coordinate recovers the old identity:
\[
  \operatorname{pr}_{Q_i}\circ
  \langle q_i,q_j\circ\gamma\rangle=q_i.
\]
Therefore the persistence extension retains \(q_i\).

\emph{(8) The persistence extension carries the later identity.}
Projection to the second coordinate recovers the transported later
identity:
\[
  \operatorname{pr}_{Q_j}\circ
  \langle q_i,q_j\circ\gamma\rangle=q_j\circ\gamma.
\]
Thus the same extension carries \(q_j\) along the declared transport
\(\gamma\).
\end{proof}

\paragraph{Nonclaims.} The theorem does not classify persistent objects or say which
transports should be declared.

\subsection*{Layer instantiations}\label{layer-inst:F19}

\begingroup\small
\begin{description}
\item[\emph{Mendelian inheritance (genetics).}] \textit{Analogy.}
Through meiotic segregation and fertilisation, parental allele
identity at a locus survives into the offspring's diploid
genotype.  The \Fref{F19} reading: allele identity persists when no
mutation or non-disjunction breaks the descent; mutation and
non-disjunction events are the persistence obstruction
\citep{Mendel1866VersuchePflanzenHybriden}.

\item[\emph{Pure-state identity under unitary evolution (quantum
mechanics).}] \textit{Special case.}
For a closed quantum system, time evolution is the unitary
\(U(t)=\exp(-iHt/\hbar)\); the projective ray
\([|\psi\rangle]\) (state modulo global phase) evolves through the
descended projective evolution.  Pure-state identity persists:
\(U(t)(e^{i\phi}\psi)=e^{i\phi}U(t)\psi\), so the persistence
obstruction is empty.  For decoherence, take unit state vectors, fix an initial environment state
and a unitary evolution of the combined system. With the earlier quotient
the system's projective ray and the later quotient its reduced density
operator, corestricted to the reduced states it attains, global phase cancels, so the persistence obstruction is empty.
When the reduced state is mixed, it fails
\(\operatorname{viable}(\rho):\Leftrightarrow\rho^2=\rho\) and contributes to
the dissolution obstruction \citep{vonNeumann1932MathematischeGrundlagen}.

\end{description}
\endgroup


\subsection{Fact / Record Formation [\texorpdfstring{\Fref{F20}}{F20}]}\label{sec:status-records-coherence:fact-record}

\begin{definition}[Record formation]\label{def:record-formation}
Let \(H_0,H_1\) be history sets (in the intended reading, carriers of a
formed closure \(\closureForm(A)\)), with continuation
\(\gamma:H_0\to H_1\) and record readout \(\rho_1:H_1\to R_1\). The
\emph{record} of a history \(h\in H_0\) is its record signature
\(s(h)=\rho_1(\gamma(h))\), and an event readout \(\varepsilon:H_0\to E\) is a
\emph{recorded fact} when \(\varepsilon=\bar\varepsilon\circ s\) for some
\(\bar\varepsilon:R_1\to E\).
\end{definition}

\paragraph{Reading the statement.} The theorem makes
\Cref{def:record-formation} precise for readouts: the record is the
signature \(\rho_1\circ\gamma\).
\(\operatorname{EventRecorded}(\varepsilon,\gamma,\rho_1)\) says that
the event \(\varepsilon\) descends through the record signature
\(\rho_1\circ\gamma\), so the record determines the event.
\(\operatorname{RecordStable}(\rho_s,\theta,\rho_t)\) says that
\(\rho_t\circ\theta\) descends through \(\rho_s\), so the transport
respects records. \(\operatorname{FactValueStable}\) says that
\(d_t\circ\bar\theta=u\circ d_s\), so the fact value moves by the
declared update. \(\operatorname{StableFact}\) is the conjunction of these
three conditions.  They describe the transport of one recorded fact when
\(H_s=H_1\), \(\rho_s=\rho_1\), \(E_s=E\) and \(d_s=\bar\varepsilon\), the
descended event readout of \Cref{def:record-formation}; the last clause
of the theorem treats this case.  \(\operatorname{EventRecorded}(\varepsilon,\gamma,\rho_1)\)
holds exactly when \(\varepsilon\) is a recorded fact in the sense of
\Cref{def:record-formation}.
\(\operatorname{RetainsRecord}\) and \(\operatorname{CarriesEvent}\) say
that the record signature, respectively the event, descends through
the record-event extension. \(\operatorname{IdempotentPackaging}(C)\)
says that packaging an already packaged history changes nothing. The
packaging clause is independent of the other data: it is the standard fact
that an idempotent map is a retraction onto its fixed set, included because a
packaged record must be unchanged by repackaging; \(H\) denotes the auxiliary
carrier only in this clause.

\begin{theorem}[Fact / Record Formation]\label{thm:F20}
Let \(H_0,H_1,H_s,H_t\) be history spaces with event readout
\(\varepsilon:H_0\to E\), continuation \(\gamma:H_0\to H_1\),
record readouts
\[
  \rho_1:H_1\to R_1,\qquad
  \rho_s:H_s\to R_s,\qquad
  \rho_t:H_t\to R_t,
\]
transport \(\theta:H_s\to H_t\), a declared record-level transport
\(\bar\theta:R_s\to R_t\), fact-value maps
\(d_s:R_s\to E_s\), \(d_t:R_t\to E_t\), update
\(u:E_s\to E_t\), and a packaging map \(C:H\to H\) on an auxiliary
carrier \(H\).  Define the record signature
\[
  s:=\rho_1\circ\gamma:H_0\to R_1,
\]
and assume that \(s\) and \(\rho_s\) are surjective.  Define
\begin{align*}
  \mathcal O_{\mathrm{ev}}
  &:=
  \{(h,h')\in H_0^2:
      \rho_1(\gamma h)=\rho_1(\gamma h')
      \text{ and } \varepsilon(h)\ne\varepsilon(h')\},\\
  \mathcal O_{\mathrm{red}}
  &:=
  \{(h,h')\in H_0^2:
      \varepsilon(h)=\varepsilon(h')
      \text{ and } \rho_1(\gamma h)\ne\rho_1(\gamma h')\},\\
  \mathcal O_{\mathrm{rec}}
  &:=
  \{(x,y)\in H_s^2:
      \rho_s(x)=\rho_s(y)
      \text{ and } \rho_t(\theta x)\ne\rho_t(\theta y)\},\\
  \mathcal O_{\mathrm{val}}
  &:=
  \{r\in R_s:
      d_t(\bar\theta r)\ne u(d_s r)\}.
\end{align*}
Then
\begin{align*}
  \operatorname{EventRecorded}(\varepsilon,\gamma,\rho_1)
  &\iff \mathcal O_{\mathrm{ev}}=\emptyset,\\
  \operatorname{RecordStable}(\rho_s,\theta,\rho_t)
  &\iff \mathcal O_{\mathrm{rec}}=\emptyset,\\
  \operatorname{FactValueStable}(\bar\theta,d_s,d_t,u)
  &\iff \mathcal O_{\mathrm{val}}=\emptyset,\\
  \operatorname{StableFact}
  &\iff
  \mathcal O_{\mathrm{ev}}=\emptyset
  \wedge
  \mathcal O_{\mathrm{rec}}=\emptyset
  \wedge
  \mathcal O_{\mathrm{val}}=\emptyset,
\end{align*}
together with the record-comparison status partition
\[
  \operatorname{ExactlyOneRecordComparisonStatus}
    (\varepsilon,\gamma,\rho_1),
\]
the retention and carrying conditions for the record-event extension
\(\langle\rho_1\circ\gamma,\varepsilon\rangle:H_0\to R_1\times E\),
\begin{align*}
  &\operatorname{RetainsRecord}
      \bigl(\rho_1\circ\gamma,
        \langle\rho_1\circ\gamma,\varepsilon\rangle\bigr),\\
  &\operatorname{CarriesEvent}
      \bigl(\langle\rho_1\circ\gamma,\varepsilon\rangle,
        \varepsilon\bigr),
\end{align*}
and the characterization of idempotent packaging by its image,
\[
  \operatorname{IdempotentPackaging}(C)
  \iff
  C(H)=\operatorname{Fix}(C),
\]
where \(\operatorname{IdempotentPackaging}(C):\Longleftrightarrow\forall h\in H,\ C(C(h))=C(h)\)
and \(\operatorname{Fix}(C)=\{h:C(h)=h\}\). Here, with \(s=\rho_1\circ\gamma\),
\(\operatorname{EventRecorded}(\varepsilon,\gamma,\rho_1):\Longleftrightarrow\exists\,\bar\varepsilon:R_1\to E,\ \bar\varepsilon\circ s=\varepsilon\);
\(\operatorname{RecordStable}(\rho_s,\theta,\rho_t):\Longleftrightarrow\exists\,\theta^\flat:R_s\to R_t,\ \theta^\flat\circ\rho_s=\rho_t\circ\theta\);
\(\operatorname{FactValueStable}(\bar\theta,d_s,d_t,u):\Longleftrightarrow d_t\circ\bar\theta=u\circ d_s\);
and \(\operatorname{StableFact}\) is the conjunction of these three.
The record-comparison status takes one of four values
(\emph{exact}, \emph{redundant}, \emph{ambiguous}, or
\emph{incomparable}) determined by the joint emptiness pattern of
\(\mathcal O_{\mathrm{ev}}\) and \(\mathcal O_{\mathrm{red}}\): exact when
both are empty, redundant when only \(\mathcal O_{\mathrm{ev}}\) is empty,
ambiguous when only \(\mathcal O_{\mathrm{red}}\) is empty, and incomparable
when neither is.
If moreover \(\bar\theta\) is a descended transport,
\(\bar\theta\circ\rho_s=\rho_t\circ\theta\), then fact-value stability is the
raw fact-value equation:
\[
  \operatorname{FactValueStable}(\bar\theta,d_s,d_t,u)
  \iff
  \forall x\in H_s,\ d_t(\rho_t(\theta x))=u(d_s(\rho_s x)).
\]
If moreover \(H_s=H_1\), \(\rho_s=\rho_1\), \(E_s=E\),
\(d_s\circ\rho_1\circ\gamma=\varepsilon\) and
\(\bar\theta\circ\rho_1=\rho_t\circ\theta\), then
\(\operatorname{FactValueStable}(\bar\theta,d_s,d_t,u)\) implies
\(d_t\circ\rho_t\circ\theta\circ\gamma=u\circ\varepsilon\): a stable fact
transports the event it records.
\end{theorem}
\begin{proof}
Assume the displayed maps, the surjectivity of \(s=\rho_1\circ\gamma\),
and the surjectivity of \(\rho_s\).  We prove the conjuncts in
turn.

\emph{(1) Event recorded iff event-record obstruction empty.}
This is the descent criterion of \Cref{sec:apparatus} for the
surjective record signature \(s\) and the readout \(\varepsilon\): the
descended event readout is \(\bar\varepsilon(s(h)):=\varepsilon(h)\).

\emph{(2) Record stable iff record-stability obstruction empty.}
The record-stability equivalence is the same quotient argument for
\(\rho_s\).  The condition \(\mathcal O_{\mathrm{rec}}=\emptyset\) is
exactly the independence needed to define a map
\[
  \theta^{\flat}(\rho_s x):=\rho_t(\theta x),
  \qquad\text{so that}\qquad
  \theta^{\flat}\circ\rho_s=\rho_t\circ\theta.
\]
The fact-value clause below concerns the given map \(\bar\theta\);
when record stability holds, the natural choice of \(\bar\theta\) is
this descended transport \(\theta^{\flat}\). For the last clause, if
\(\bar\theta\circ\rho_s=\rho_t\circ\theta\), then
\(d_t(\rho_t(\theta x))=d_t(\bar\theta(\rho_s x))\), so the raw equation at
\(x\) is the fact-value equation at \(\rho_s x\); since \(\rho_s\) is
surjective, every \(r\in R_s\) is of this form.

\emph{(3) Fact-value stable iff the fact-value defect is empty.}
The fact-value defect at \(r\in R_s\) is the failure of the
commutative square
\[
  d_t\circ\bar\theta=u\circ d_s.
\]
Equivalently,
\[
  \mathcal O_{\mathrm{val}}=\emptyset
  \quad\Longleftrightarrow\quad
  \forall r\in R_s,\ d_t(\bar\theta r)=u(d_s r),
\]
which is precisely \(\operatorname{FactValueStable}(\bar\theta,d_s,d_t,u)\).

\emph{(4) Stable fact iff the three obstructions are empty.}
By definition, a stable fact requires event recording, record
stability, and fact-value stability.  Combining (1)--(3) gives
\[
  \operatorname{StableFact}
  \iff
  \mathcal O_{\mathrm{ev}}=\emptyset
  \wedge
  \mathcal O_{\mathrm{rec}}=\emptyset
  \wedge
  \mathcal O_{\mathrm{val}}=\emptyset.
\]

\emph{(5) Exactly-one record-comparison status.}
The four statuses are decided by the joint
\((\mathcal O_{\mathrm{ev}},\mathcal O_{\mathrm{red}})\) emptiness
pattern:
\[
\begin{array}{c|c|c}
\text{status} & \mathcal O_{\mathrm{ev}}=\emptyset
  & \mathcal O_{\mathrm{red}}=\emptyset\\
\hline
\text{exact record} & \text{yes} & \text{yes}\\
\text{redundant record} & \text{yes} & \text{no}\\
\text{ambiguous record} & \text{no} & \text{yes}\\
\text{incomparable record} & \text{no} & \text{no}
\end{array}
\]
Every configuration of the two Boolean defect tests falls into one
row, and no configuration falls into two rows.  Existence is by case
enumeration; exclusivity is by definitional disjointness of the four
status predicates.

\emph{(6) Retains record.}
Define the canonical record-event extension
\[
  \langle\rho_1\circ\gamma,\varepsilon\rangle:
  H_0\longrightarrow R_1\times E,
  \qquad
  h\longmapsto(\rho_1(\gamma h),\varepsilon(h)).
\]
Projection to the first coordinate recovers the record signature:
\[
  \operatorname{pr}_{R_1}
  \circ\langle\rho_1\circ\gamma,\varepsilon\rangle
  =
  \rho_1\circ\gamma.
\]

\emph{(7) Carries event.}
Projection to the second coordinate recovers the event readout:
\[
  \operatorname{pr}_{E}
  \circ\langle\rho_1\circ\gamma,\varepsilon\rangle
  =
  \varepsilon.
\]

\emph{(8) Idempotent packaging is image equal to fixed locus.}
Always \(\operatorname{Fix}(C)\subseteq C(H)\), since \(h=C(h)\) for a fixed
point. If \(C(C(h))=C(h)\) for every \(h\), then each \(C(h)\) is fixed, so
\(C(H)\subseteq\operatorname{Fix}(C)\). Conversely, if
\(C(H)\subseteq\operatorname{Fix}(C)\), then \(C(h)\) is fixed for every
\(h\), that is, \(C(C(h))=C(h)\).

\emph{(9) A stable fact transports its event.} Under the additional
identifications, for every \(h\in H_0\),
\(d_t\rho_t\theta\gamma h=d_t\bar\theta\rho_1\gamma h=u\,d_s\rho_1\gamma h
=u\,\varepsilon h\), using the descended transport, fact-value stability and
the decoding identity in turn.
\end{proof}

\paragraph{Nonclaims.} The theorem does not determine which candidates count as evidence and
does not identify the fact with the record.

\subsection*{Layer instantiations}\label{layer-inst:F20}

\begingroup\small
\begin{description}
\item[\emph{{FAIR} scientific-data principles (research
methodology).}] \textit{Analogy.}
The FAIR Guiding Principles --- Findable, Accessible,
Interoperable, Reusable --- are guidelines for the findability and
reuse of scientific data.  Read through \Fref{F20}, data form stable
facts when generation events are recorded as provenance
metadata, records persist under storage / migration transport via
persistent identifiers and access protocols, and derived analyses
preserve fact-value consistency through declared processing
pipelines.  Each obstruction (unrecorded provenance, broken DOI,
undocumented pipeline change) is an \Fref{F20}-style stable-fact failure
\citep{WilkinsonDumontier2016FAIRPrinciples}.

\end{description}
\endgroup


\subsection{Reflexive Nonclosure [\texorpdfstring{\Fref{F21}}{F21}]}\label{sec:status-records-coherence:reflexive-nonclosure}

\paragraph{Reading the statement.} The symbols are those of \Cref{thm:F21} below.
\(\operatorname{CompleteSelfSoundness}_{\pi}(c)\) says that the claim
\(c\) asserts the complete soundness of the policy at its own layer; a
scoped lower audit is an audit of a more limited claim below that
assertion. \(\operatorname{InScope}_{\pi}\) and
\(\operatorname{SameLayerCycle}_{\pi}\) mark the claims the policy
covers and those that depend on themselves at the same layer. Rules 2 and 3 fix the statuses of outside-scope and in-scope cyclic claims, and the rule on \(\mu\)
says that the meta layer does not accept a meta-claim about its own
soundness. Rule 4 is the premise through which same-layer circularity enters;
the blocking of self-certification is derived from rules 2--4, each needed by
the sharpness clause, and the meta-layer clause is the contrapositive of the
rule on \(\mu\). The scoped-lower-audit clause follows from the blocking of self-certification.

\begin{theorem}[Reflexive Nonclosure]\label{thm:F21}
Let \(\mathcal C\) be a space of claims and let \(\pi\) be a same-layer
audit policy on \(\mathcal C\): predicates
\(\operatorname{InScope}_{\pi}\),
\(\operatorname{CompleteSelfSoundness}_{\pi}\),
\(\operatorname{ScopedLowerAudit}_{\pi}\) and
\(\operatorname{SameLayerCycle}_{\pi}\) on \(\mathcal C\), and a status
map \(\operatorname{st}_{\pi}\) from \(\mathcal C\) to a set of statuses
containing three distinct statuses \(\mathsf{accepted}\),
\(\mathsf{outsideScope}\) and \(\mathsf{undefinedCircular}\), subject to rules 2--4 below (rule 1 describes the intended policies and is not assumed):
\begin{enumerate}
\item (intended, not assumed) a scoped lower audit is not a complete self-soundness claim;
\item a claim outside the scope receives \(\mathsf{outsideScope}\);
\item a claim in scope with a same-layer cycle receives
  \(\mathsf{undefinedCircular}\);
\item a complete self-soundness claim in scope has a same-layer cycle.
\end{enumerate}
Let \(\mu\) be a meta-layer audit into a space \(\mathcal M\) of
meta-claims: a relation \(\operatorname{Audits}_{\mu}(m,c)\) and
predicates \(\operatorname{Accepted}_{\mu}\) and
\(\operatorname{MetaSelfSound}_{\mu}\) on \(\mathcal M\), such that \(\operatorname{MetaSelfSound}_{\mu}(m)\Rightarrow\neg\operatorname{Accepted}_{\mu}(m)\) for every \(m\in\mathcal M\): no
meta-claim asserting its own soundness is accepted by the meta layer.
Write
\begin{align*}
  \operatorname{SelfCertificationBlocked}_{\pi}(c)
  &:\iff
  \bigl(\operatorname{CompleteSelfSoundness}_{\pi}(c)
    \Rightarrow \operatorname{st}_{\pi}(c)\ne\mathsf{accepted}\bigr),\\
  \operatorname{AcceptedScopedLowerAudit}_{\pi}(c)
  &:\iff
  \operatorname{ScopedLowerAudit}_{\pi}(c)
    \wedge\operatorname{st}_{\pi}(c)=\mathsf{accepted},\\
  \operatorname{HasMetaLayerCertificate}_{\mu}(c)
  &:\iff
  \exists m\in\mathcal M,\
    \operatorname{Audits}_{\mu}(m,c)\wedge\operatorname{Accepted}_{\mu}(m).
\end{align*}
Then for every \(c\in\mathcal C\),
\begin{align*}
  &\operatorname{SelfCertificationBlocked}_{\pi}(c),\\
  &\operatorname{AcceptedScopedLowerAudit}_{\pi}(c)
    \Longrightarrow
    \neg\operatorname{CompleteSelfSoundness}_{\pi}(c),\\
  &\operatorname{HasMetaLayerCertificate}_{\mu}(c)\\
  &\qquad{}\Longrightarrow
    \exists m\in\mathcal M,\,
      \operatorname{Audits}_{\mu}(m,c)
      \wedge\operatorname{Accepted}_{\mu}(m)
      \wedge
      \neg\operatorname{MetaSelfSound}_{\mu}(m),\\
  &\operatorname{CompleteSelfSoundness}_{\pi}(c)\\
  &\qquad{}\Longrightarrow
    \bigl(\neg\operatorname{InScope}_{\pi}(c)
      \wedge\operatorname{st}_{\pi}(c)=\mathsf{outsideScope}\bigr)\\
  &\hspace{5em}\vee
    \bigl(\operatorname{InScope}_{\pi}(c)
      \wedge\operatorname{SameLayerCycle}_{\pi}(c)
      \wedge\operatorname{st}_{\pi}(c)=\mathsf{undefinedCircular}\bigr).
\end{align*}
Rules~2, 3 and~4 cannot be dropped: there is a policy on two claims (without rule 4) satisfying
rules 1--3 in which \(\operatorname{SelfCertificationBlocked}_{\pi}\) fails,
a policy on one claim (without rule 2) satisfying rules 1, 3 and 4 in which it fails, and a
policy on one claim (without rule 3) satisfying rules 1, 2 and 4 in which it fails. Rule~1 is not assumed: the second line follows from the first, since a claim with status \(\mathsf{accepted}\) is not a complete self-soundness claim.
\end{theorem}
\begin{proof}
Assume the same-layer policy \(\pi\) and the meta-layer audit
\(\mu\), and fix \(c\in\mathcal C\).  Suppose
\(\operatorname{CompleteSelfSoundness}_{\pi}(c)\).  There are two
cases.  If \(c\) is outside the scope of \(\pi\), rule~2 gives
\(\operatorname{st}_{\pi}(c)=\mathsf{outsideScope}\).  If \(c\) is in
scope, rule~4 gives a same-layer cycle, and rule~3 then gives
\(\operatorname{st}_{\pi}(c)=\mathsf{undefinedCircular}\).  This proves
the fourth line, and in both cases
\(\operatorname{st}_{\pi}(c)\ne\mathsf{accepted}\), which is
\(\operatorname{SelfCertificationBlocked}_{\pi}(c)\): a claim certifying
its own complete soundness at the same layer is never accepted.

If \(\operatorname{AcceptedScopedLowerAudit}_{\pi}(c)\) holds, then
\(\operatorname{st}_{\pi}(c)=\mathsf{accepted}\), so
\(\neg\operatorname{CompleteSelfSoundness}_{\pi}(c)\) by the first line: the audit is
accepted for its scoped lower claim, not for the full assertion of its
own soundness.

If \(\operatorname{HasMetaLayerCertificate}_{\mu}(c)\) holds, choose
\(m\) with \(\operatorname{Audits}_{\mu}(m,c)\) and
\(\operatorname{Accepted}_{\mu}(m)\).  If
\(\operatorname{MetaSelfSound}_{\mu}(m)\) held, the rule on \(\mu\)
would give \(\neg\operatorname{Accepted}_{\mu}(m)\); hence
\(\neg\operatorname{MetaSelfSound}_{\mu}(m)\).  The certificate is
supplied by a meta-claim that does not certify its own soundness, so
the audit has moved to a strict meta-layer.

For the sharpness claim, let \(\mathcal C=\{a,b\}\), put both claims in scope,
let \(b\) be the only complete self-soundness claim and \(a\) the only scoped
lower audit, let no claim have a same-layer cycle, and let
\(\operatorname{st}_{\pi}\) assign \(\mathsf{accepted}\) to both. Rule~1 holds
because \(a\) is not a complete self-soundness claim, and rules~2 and~3 hold
vacuously. Rule~4 fails at \(b\), and \(b\) is a complete self-soundness claim
with status \(\mathsf{accepted}\), so
\(\operatorname{SelfCertificationBlocked}_{\pi}(b)\) fails. For rule~2, let
\(\mathcal C=\{c\}\) with \(c\) a complete self-soundness claim outside the
scope of \(\pi\), with no scoped lower audit, no same-layer cycle, and status \(\mathsf{accepted}\).
Rule~1 holds because there is no scoped lower audit, rules~3 and~4 because
\(c\) is not in scope, rule~2 fails at \(c\), and
\(\operatorname{SelfCertificationBlocked}_{\pi}(c)\) fails. For rule~3, let
\(\mathcal C=\{c\}\) with \(c\) a complete self-soundness claim in scope, with a
same-layer cycle, no scoped lower audit, and status \(\mathsf{accepted}\).
Rule~1 holds because there is no scoped lower audit, rule~2 holds vacuously,
rule~4 holds because \(c\) has a same-layer cycle, rule~3 fails at \(c\), and
\(\operatorname{SelfCertificationBlocked}_{\pi}(c)\) fails.
\end{proof}

\paragraph{Nonclaims.} The theorem does not forbid scoped self-audit; it forbids unrestricted
same-layer self-closure.  A proof-theoretic analogue is G\"odel's second
incompleteness theorem \citep{Godel1931Unentscheidbare}, where unprovability
of self-consistency is derived from the arithmetic of the theory; here
non-acceptance of complete self-soundness follows from rules 2--4 and the
distinctness of the three statuses, which
the policy supplies. The modal-logic reformulation in
terms of a provability operator is developed in
\citet{Boolos1993LogicProvability}.

\subsection*{Layer instantiations}\label{layer-inst:F21}

\begingroup\small
\begin{description}
\item[\emph{{ISO}/{IEC} 17025 laboratory accreditation
(industrial metrology).}] \textit{Analogy.}
Under ISO/IEC 17025, a testing or calibration laboratory
cannot self-certify its measurement competence: an
accredited laboratory must be assessed by a recognised
accreditation body (an ILAC member such as A2LA, UKAS, or
DAkkS), which acts as the meta-layer.  Internal
quality-management records are a scoped lower audit,
sufficient
for internal corrective action but not for legal or
commercial recognition of competence; the meta-layer
accreditation certificate is the legitimate reflexive
audit witness \citep{ISOIEC17025_2017}.

\end{description}
\endgroup


\subsection{Interface Mediation [\texorpdfstring{\Fref{F22}}{F22}]}\label{sec:status-records-coherence:interface-mediation}

\paragraph{Reading the statement.} The symbols are those of \Cref{thm:F22} below.
\(\operatorname{InterfaceMediated}(q,\beta,U,q')\) says that the target
class \(q'(U(x,y))\) depends only on the source class \(q(x)\) and the
boundary value \(\beta(y)\). \(\operatorname{BoundaryLeakageAbsent}\)
says that, for each fixed \(x\), two inputs \(y,y'\) with the same
boundary value give the same target class;
\(\operatorname{InternalInsufficiencyAbsent}\) says that, for each
fixed \(y\), two states \(x,x'\) with the same source class give the
same target class. \(\operatorname{ExternallyInert}(\bar U)\) says that
\(\bar U(a,b)\) does not depend on the boundary argument \(b\), and
\(\operatorname{ActiveMediatedInfluence}(\bar U)\) says that it does
for some \(a\).

\begin{theorem}[Interface Mediation]\label{thm:F22}
Let \(q:X\to Q\) be a source quotient, let
\(\beta:Y\to B\) be a boundary readout, let
\(U:X\times Y\to X'\) be an interface update, and let
\(q':X'\to Q'\) be the target quotient.  Assume \(q\) and \(\beta\) are surjective.  Define
\[
  \mathcal O_U
  :=
  \{((x,y),(x',y')) :
    q(x)=q(x'),\ \beta(y)=\beta(y'),\
    q'(U(x,y))\ne q'(U(x',y'))\}.
\]
Then
\begin{align*}
  \operatorname{InterfaceMediated}(q,\beta,U,q')
  &\iff \mathcal O_U=\emptyset\\
  &\iff
  \exists!\,\bar U:Q\times B\to Q'\ \text{with}\
    \bar U(q(x),\beta(y))=q'(U(x,y))\ \forall x,y,\\
  \mathcal O_U=\emptyset
  &\iff
  \operatorname{BoundaryLeakageAbsent}\wedge
  \operatorname{InternalInsufficiencyAbsent},
\end{align*}
and, by the definitions of active mediated influence and external inertness,
for every \(\bar U:Q\times B\to Q'\), in particular the descended \(\bar U\) of the
second line,
\[
  \operatorname{ActiveMediatedInfluence}(\bar U)
  \iff \neg\operatorname{ExternallyInert}(\bar U).
\]
\end{theorem}
\begin{proof}
Assume the source quotient \(q\), boundary readout \(\beta\),
interface update \(U\), and target quotient \(q'\).  If
\(\operatorname{InterfaceMediated}(q,\beta,U,q')\) holds, then
changing representatives inside the same \(q\)-class and
\(\beta\)-class cannot change the target \(q'\)-class; hence
\(\mathcal O_U=\emptyset\).

Conversely, suppose \(\mathcal O_U=\emptyset\).  Define
\[
  \bar U(q(x),\beta(y)):=q'(U(x,y)).
\]
If \(q(x)=q(x')\) and \(\beta(y)=\beta(y')\), the emptiness of
\(\mathcal O_U\) gives
\[
  q'(U(x,y))=q'(U(x',y')),
\]
so \(\bar U\) is independent of representatives.  Thus the
interface update descends through \(Q\times B\).  The absence of
\(\mathcal O_U\) also rules out the two possible mediation failures:
boundary-equivalent inputs cannot leak distinct target classes, and
source-equivalent internal states cannot create an undeclared target
difference.  Conversely, assume both absences, and let \(q(x)=q(x')\) and
\(\beta(y)=\beta(y')\).  Boundary leakage is absent at the fixed state
\(x\), and internal insufficiency is absent at the fixed input \(y'\), so
\[
  q'(U(x,y))=q'(U(x,y'))=q'(U(x',y')),
\]
and \(\mathcal O_U=\emptyset\).  Finally, \(\bar U\) has active mediated influence
exactly when it is not constant in the external argument, which is
the negation of external inertness.
The descended map \(\bar U\) is unique because \(q\times\beta\) is surjective.
\end{proof}

\paragraph{Nonclaims.} The theorem does not determine the interface's internal structure and
does not apply to undeclared interfaces.

\subsection*{Layer instantiations}\label{layer-inst:F22}

\begingroup\small
\begin{description}
\item[\emph{Abstract-data-type encapsulation (software
engineering).}] \textit{Special case.}
In an abstract data type, the carrier \(X\) is the set of
concrete representations satisfying the representation invariant (e.g., a stack realised as a
linked list, a dynamic array, or a deque); the source
quotient \(q\) is the abstraction function, which sends each
representation to the abstract value it represents. Take \(Y\) to be the set of public method calls with concretely encoded arguments, \(\beta\) the map sending a call to its method name and the abstract values of its arguments, \(U(x,y)\) the implementation's post-state and return value, and \(q'\) the abstraction of the post-state paired with the abstract return value. Require argument representations to satisfy their invariants and \(U\) to be total with invariant-satisfying post-states, including declared error outcomes. Take \(Q=q(X)\), \(B=\beta(Y)\) and \(Q'=q'(X')\).  Encapsulation is the
requirement that two representations in the same source
class respond identically (modulo the target quotient)
to the same public-method call; private-state exposure
through the public interface, where two representations of one abstract
value answer the same public call with different abstract results, is the
internal-insufficiency failure mode; boundary leakage would be
two calls at a fixed concrete representation, with the same
boundary value \(\beta(y)=\beta(y')\), that give different
abstract results \citep{LiskovZilles1974ADT}.

\item[\emph{Black-box security reduction
(theoretical cryptography).}] \textit{Analogy.}
A standard security proof shows that any adversary
breaking a scheme can be turned into an algorithm
solving an underlying hard problem; in the black-box
reduction discipline, the reduction accesses the
adversary only through oracle queries.  The carrier is
the adversary's internal state, the source quotient
collapses adversaries with the same query/response
distribution, the boundary is the oracle transcript,
and the update is the adversary's next step.  The
reduction's success thereby depends only on the
boundary I/O profile, not on the adversary's internal
code.  The random-oracle methodology, in which the hash
function is itself modelled as a uniformly random
oracle, is one influential setting within this
discipline \citep{BellareRogaway1993RandomOracles}.

\end{description}
\endgroup


\subsection{Probability as Stable Unresolved Fiber [\texorpdfstring{\Fref{F23}}{F23}]}\label{sec:status-records-coherence:probability}

\paragraph{Reading the statement.} The symbols are those of \Cref{thm:F23} below.
\(\operatorname{Unresolved}(q,r)\) says that some \(q\)-fiber contains
two readout values. \(\operatorname{Stable}\) says that each ledger has its
support contained in its own fiber, at the source and at the target, and that
the two declared maps \(\operatorname{push}\) and \(\operatorname{mix}\),
the transported ledger and the declared mixture of target ledgers, agree on
every source class; \(\operatorname{push}\) reads the source ledger through
\(\Phi\), and \(\operatorname{mix}\) reads the ledger of the continued class
through \(\Psi\).
\(\operatorname{StableProbability}(q,r)\) says that the readout is
unresolved and the declared ledgers satisfy \(\operatorname{Stable}\).
\(\operatorname{Stable}\) constrains \(\lambda\), \(\lambda'\), their supports and
the declared maps \(\Phi\), \(\Psi\), \(\bar\kappa\); it does not involve the
weight assignment \(m\), so the values \(\operatorname{Pr}_{q,r}\) enter only
clauses (ii) and (iv)--(vi). With \(S\) and \(S'\) empty and \(\mathcal T\) a
one-point set, \(\operatorname{Stable}\) holds, and by (iii)
\(\operatorname{StableProbability}(q,r)\) is then equivalent to
\(\operatorname{Unresolved}(q,r)\).  Clause (iii) unfolds the definition
of \(\operatorname{StableProbability}\); no normalization or additivity
of the weights is assumed beyond what clauses (iv)--(vi) state. The relation \(S\) is declared separately from the weights: it constrains \(m(\lambda(a))\) only through the hypothesis of clause (v). Under \(\operatorname{Stable}\), a ledger shared by two classes has empty support, since its support lies in both fibers; clause (vi), which does not use \(S\), allows such sharing.

\begin{theorem}[Probability as Stable Unresolved Fiber]\label{thm:F23}
Let \(q:H\to Q\) be a surjective quotient of histories, let
\(\lambda:Q\to L\) be a ledger map with support relation
\(S\subseteq L\times H\), let
\(m:L\to (H\to\mathrm{Prop})\to W\) be a weight assignment, and let
\(r:H\to Z\) be a readout.  Let a declared continuation carry a target
quotient \(q':H'\to Q'\) with ledger \(\lambda':Q'\to L'\) and support
relation \(S'\subseteq L'\times H'\), and let
\(\bar\kappa:Q\to Q'\) be the class-level continuation, let \(\mathcal T\) be a
set with a declared ledger transport \(\Phi:L\to\mathcal T\) and target
mixture \(\Psi:L'\to\mathcal T\), and put
\(\operatorname{push}:=\Phi\circ\lambda\) and
\(\operatorname{mix}:=\Psi\circ\lambda'\circ\bar\kappa\); only their pointwise
equality in \(\operatorname{Stable}\) is used. The continuation data enter only
through \(\operatorname{Stable}\), which enters clause (iii), the last sentence
of clause (v), and the final equivalence of clause (vi). Apart from these uses
of \(\operatorname{Stable}\), clauses (i), (ii) and (iv)--(vi) depend only on
\(q\), \(\lambda\), \(S\), \(m\), \(r\) and their clause-specific hypotheses.  For \(a\in Q\) and \(v\in Z\), set
\[
  \operatorname{Pr}_{q,r}(a,v)
  :=
  m(\lambda(a))\bigl(\{h\in H:q(h)=a\ \text{and}\ r(h)=v\}\bigr),
\]
and write
\begin{align*}
  \operatorname{Stable}
  :\iff{}&
  \bigl(\forall a\in Q,\forall h\in H,\
    S(\lambda(a),h)\Rightarrow q(h)=a\bigr)\\
  &{}\wedge
  \bigl(\forall a'\in Q',\forall h'\in H',\
    S'(\lambda'(a'),h')\Rightarrow q'(h')=a'\bigr)\\
  &{}\wedge
  \bigl(\forall a\in Q,\
    \operatorname{push}(a)=\operatorname{mix}(a)\bigr),\\
  \operatorname{Unresolved}(q,r)
  :\iff{}&
  \exists h,h'\in H,\ q(h)=q(h')\ \text{and}\ r(h)\ne r(h'),\\
  \operatorname{StableProbability}(q,r)
  :\iff{}&
  \operatorname{Unresolved}(q,r)
  \wedge
  \operatorname{Stable}.
\end{align*}
Then:
\begin{enumerate}[label=(\roman*)]
\item \(r\) descends through \(q\) if and only if
  \(q(h)=q(h')\Rightarrow r(h)=r(h')\) for all \(h,h'\in H\), and exactly
  one of ``\(r\) descends through \(q\)'' and \(\operatorname{Unresolved}(q,r)\)
  holds;
\item if \(r=\bar r\circ q\), then for every \(a\in Q\) and \(v\in Z\) the
  event \(\{h:q(h)=a\ \text{and}\ r(h)=v\}\) is the fiber \(q^{-1}(a)\) when
  \(v=\bar r(a)\) and is empty otherwise, so
  \(\operatorname{Pr}_{q,r}(a,v)=m(\lambda(a))(q^{-1}(a))\) for
  \(v=\bar r(a)\) and \(\operatorname{Pr}_{q,r}(a,v)=m(\lambda(a))(\emptyset)\)
  for \(v\ne\bar r(a)\);
\item \(\operatorname{StableProbability}(q,r)\) holds if and only if
  \(\operatorname{Stable}\) holds and \(r\) does not descend through \(q\);
\item if the weights form a commutative monoid, \(Z\) is finite and each
  \(m(\lambda(a))\) vanishes on the empty event and is additive on disjoint
  events, then
  \(\sum_{v\in Z}\operatorname{Pr}_{q,r}(a,v)=m(\lambda(a))(q^{-1}(a))\) for
  every \(a\in Q\): the readout distributes the weight of the fiber;
\item if the support conjunct of \(\operatorname{Stable}\) holds at \(a\), that
  is, \(S(\lambda(a),h)\Rightarrow q(h)=a\), and there are an operation \(\oplus\) on \(W\) and
  an element \(0\in W\) (no monoid laws are needed) with
  \(m(\lambda(a))(A\cup B)=m(\lambda(a))(A)\oplus m(\lambda(a))(B)\) for disjoint
  events \(A,B\) and \(m(\lambda(a))(E)=0\) for every event \(E\) disjoint from
  \(\{h:S(\lambda(a),h)\}\), then
  \(\operatorname{Pr}_{q,r}(a,v)=m(\lambda(a))(\{h:r(h)=v\})\) for every
  \(v\in Z\): the ledger of a class is concentrated on its fiber; in
  particular, under \(\operatorname{Stable}\) the support hypothesis holds at
  every \(a\in Q\), so this conclusion holds at every \(a\) at which the two
  weight conditions hold;
\item if \(W=\mathbb R\), \(H\) is finite and for each \(a\in Q\) there are weights
  \(w_a:q^{-1}(a)\to(0,1]\) with \(\sum_{h\in q^{-1}(a)}w_a(h)=1\) and
  \(m(\lambda(a))(E)=\sum_{h\in E}w_a(h)\) for every event \(E\subseteq q^{-1}(a)\)
  (no condition is imposed on other events, so distinct classes may share a
  ledger), then \(r\) descends through \(q\) if and only if every
  \(\operatorname{Pr}_{q,r}(a,\cdot)\) is a point mass, that is,
  \(\operatorname{Pr}_{q,r}(a,v)=1\) for some \(v\); hence
  \(\operatorname{Unresolved}(q,r)\) holds exactly when some
  \(\operatorname{Pr}_{q,r}(a,\cdot)\) charges two values; consequently,
  under these weights, \(\operatorname{StableProbability}(q,r)\) holds if and
  only if \(\operatorname{Stable}\) holds and some
  \(\operatorname{Pr}_{q,r}(a,\cdot)\) charges two values.
\end{enumerate}
\end{theorem}
\begin{proof}
Assume the surjective quotient \(q\), the ledgers, supports and
weights, the continuation data, and the readout \(r\).  (i) The first equivalence is the descent criterion of
\Cref{sec:apparatus}, since \(q\) is surjective.
\(\operatorname{Unresolved}(q,r)\) is the negation of the constancy
condition, so exactly one of the two holds.

(ii) Let \(r=\bar r\circ q\) and \(a\in Q\).  Every \(h\) with \(q(h)=a\)
has \(r(h)=\bar r(a)\).  So the event \(\{h:q(h)=a\ \text{and}\ r(h)=v\}\)
is the whole fiber \(q^{-1}(a)\) when \(v=\bar r(a)\) and is empty
otherwise, and the formula for \(\operatorname{Pr}_{q,r}\) follows from its
definition: the readout is deterministic on each fiber.

(iii) \(\operatorname{StableProbability}(q,r)\) is
\(\operatorname{Unresolved}(q,r)\wedge\operatorname{Stable}\), and by (i)
\(\operatorname{Unresolved}(q,r)\) holds exactly when \(r\) does not descend
through \(q\).
Unresolvedness alone is not probability: a stable probability also requires
that each ledger have its support contained in its own fiber, at the source and at the
target, and that transporting the ledger of each source class give the
declared mixture of target ledgers, which is \(\operatorname{Stable}\).

(iv) For fixed \(a\), the events \(\{h:q(h)=a\ \text{and}\ r(h)=v\}\),
\(v\in Z\), are pairwise disjoint and their union is \(q^{-1}(a)\). Finite
additivity of \(m(\lambda(a))\) gives the stated sum.

(v) The event \(\{h:r(h)=v\}\) is the disjoint union of
\(\{h:q(h)=a\ \text{and}\ r(h)=v\}\) and \(\{h:q(h)\ne a\ \text{and}\ r(h)=v\}\).
By the support conjunct, the second is disjoint from
\(\{h:S(\lambda(a),h)\}\), so \(m(\lambda(a))\) vanishes on it; so does the
empty event. Write \(E_{a,v}=\{h:q(h)=a\ \text{and}\ r(h)=v\}\) and \(\oplus\)
for the addition of weights. Additivity on \(E_{a,v}\) and \(\emptyset\) gives
\(\operatorname{Pr}_{q,r}(a,v)=m(\lambda(a))(E_{a,v})\oplus 0\), and additivity
on \(E_{a,v}\) and \(\{h:q(h)\ne a\ \text{and}\ r(h)=v\}\) gives
\(m(\lambda(a))(\{h:r(h)=v\})=m(\lambda(a))(E_{a,v})\oplus 0\), where \(0\) is
the common value on events disjoint from the support. Hence
\(m(\lambda(a))(\{h:r(h)=v\})=\operatorname{Pr}_{q,r}(a,v)\).

(vi) The event \(\{h:q(h)=a,\ r(h)=v\}\) lies in \(q^{-1}(a)\), and each \(h\in
q^{-1}(a)\) has \(w_a(h)>0\), so \(\operatorname{Pr}_{q,r}(a,v)=\sum_{h\in q^{-1}(a),\,r(h)=v}w_a(h)\),
these values sum to \(1\) over \(v\), and \(\operatorname{Pr}_{q,r}(a,v)>0\)
exactly when some \(h\in q^{-1}(a)\) has \(r(h)=v\). Hence
\(\operatorname{Pr}_{q,r}(a,\cdot)\) is a point mass exactly when \(r\) takes
one value on \(q^{-1}(a)\), and by (i) this holds for every \(a\) exactly when
\(r\) descends through \(q\).
\end{proof}

\paragraph{Nonclaims.} The theorem does not determine numerical probability values and does
not apply to resolved fibers.  The measure-theoretic foundation of
probability is \citet{Kolmogorov1933Grundbegriffe}; the branch /
decoherent-histories interpretation that comes closest to a stable
unresolved-fiber reading appears in \citet{Griffiths1984ConsistentHistories}.

\subsection*{Layer instantiations}\label{layer-inst:F23}

\begingroup\small
\begin{description}
\item[\emph{Design-based survey sampling (social
statistics).}] \textit{Analogy.}
For a fixed finite population, the carrier is the set of
possible sample-draw realisations under the chosen
randomisation scheme; the source quotient collapses
draws sharing the same population-level summary; the
readout is the realised sample statistic.  The same
population-level summary supports many distinct sample
realisations.  The sampling design's inclusion-probability
measure is the ledger; the design-induced probability
law on draws is the weight assignment.  The
Horvitz-Thompson estimator and the design-based variance
estimator are the stable probability records derived
from this design measure
\citep{Cochran1977SamplingTechniques}.

\item[\emph{Chain-ladder loss-reserving in non-life
insurance (actuarial science).}] \textit{Analogy.}
The carrier is the space of individual
claim-development histories within an accident year;
the source quotient collapses histories sharing the
same observed cumulative-payment pattern through the
current development cell; the readout is the ultimate
loss, which is unresolved over the observed
development-pattern fibre.  Mack's distribution-free
assumptions on the cumulative payments constitute the
ledger and supply the first-and-second-moment structure
of incremental development factors as the weight
assignment.  The chain-ladder reserve estimator and
Mack's prediction-error formula are the stable
probability records used for reserve setting and
solvency capital \citep{Mack1993ChainLadder}.

\item[\emph{Wright-Fisher stochastic population genetics
(theoretical biology).}] \textit{Analogy.}
The carrier is the set of gamete-sampling realisations
within a generation; the source quotient collapses
realisations to the current-generation allele frequency;
the readout is the realised next-generation allele
count.  Two distinct gamete-sampling outcomes sharing
the same current frequency can produce different
next-generation counts, so the readout is unresolved
over the current-frequency fibre.  The ledger records
the effective population size and any
selection/mutation parameters; the weight assignment is
the Wright-Fisher binomial (or diffusion) transition
kernel they induce.  The resulting transition law is
the stable probability record that descends to the
population-level frequency dynamics
\citep{Wright1931Evolution}.

\item[\emph{Shannon noisy-channel coding (information
theory).}] \textit{Analogy.}
The carrier is the set of (input, output) trajectories
of a noisy memoryless channel; the source quotient
collapses trajectories sharing the same observed output
sequence; the readout is the transmitted input.  The
receiver cannot uniquely recover the input from the
output (unresolved fibre).  The channel input
distribution supplies the ledger; the channel law
supplies the weight assignment; the channel capacity and the noisy-channel coding
theorem are results about this channel that \Fref{F23}
neither states nor implies; \Fref{F23} classifies the input readout as unresolved exactly
when at least one output fibre contains trajectories with
different transmitted inputs
\citep{Shannon1948Mathematical}.

\item[\emph{Weibull life-distribution model in
reliability engineering (engineering reliability).}] \textit{Analogy.}
The carrier is the set of individual component
life trajectories under specified operating conditions;
the source quotient collapses trajectories sharing the
same operating-condition class; the readout is the
realised time-to-failure of a particular component.
Individual times-to-failure are unresolved over the
operating-condition fibre.  The Weibull distribution
function \(F(t)=1-\exp(-(t/\eta)^{\beta})\) is the
standard ledger and supplies a stable two-parameter
hazard model as the weight assignment; fitted shape
\(\beta\) and scale \(\eta\) give the stable
probability record used for population-level
reliability prediction, warranty pricing, and burn-in
design \citep{Weibull1951Statistical}.
\end{description}
\endgroup

\subsection{Causality as Stable Intervention Transport [\texorpdfstring{\Fref{F25}}{F25}]}\label{sec:status-records-coherence:causality}

\paragraph{Reading the statement.} The symbols are those of \Cref{thm:F25} below.
\(\operatorname{ITD}\) and \(\operatorname{PITD}\) say that the target
readout, respectively the kernel outcome, depends only on the context
class and the intervention class. \(\operatorname{Caus}(c,u,v)\) says
that interventions of classes \(u\) and \(v\) at context class \(c\)
have a stable causal effect: both are admissible, the quotient-level
transport gives different results for them, and the residue is stable
in the declared sense; \(\operatorname{ResStable}\) is declared data.
\(\operatorname{ProbCaus}\) is the same for kernel outcomes.

\begin{theorem}[Causality as Stable Intervention Transport]\label{thm:F25}
Let \(q:H\to Q\) be a surjective context quotient, let
\(\alpha:A\to\bar A\) be a surjective intervention abstraction, let
\(\Gamma:H\times A\to H'\) be an intervention transport, let
\(r:H'\to R\) be a target readout, and let
\(K:H\times A\to K_0\) be a kernel outcome (in the intended reading \(K_0\) is a set of probability laws on outcomes and \(K(h,a)\) is the outcome law of intervention \(a\) in context \(h\); the theorem uses \(K\) only as a map).  Let
\(\operatorname{Adm}:Q\times\bar A\to\mathrm{Prop}\) be admissibility,
let
\(\operatorname{ResStable},\operatorname{ProbResStable}:
Q\times\bar A\times\bar A\to\mathrm{Prop}\) be residue-stability
predicates, and let \(\bar\Gamma:Q\times\bar A\to R\) and
\(\bar K:Q\times\bar A\to K_0\) be quotient-level transports.  Define
\[
  \mathcal O_{\Gamma}
  :=
  \{((h,a),(h',a')) :
    q(h)=q(h'),\ \alpha(a)=\alpha(a'),\
    r(\Gamma(h,a))\ne r(\Gamma(h',a'))\},
\]
and \(\mathcal O_K\) in the same way with \(K\) in place of
\(r\circ\Gamma\).  Define intervention-transport descent
\(\operatorname{ITD}\), probabilistic intervention-transport descent
\(\operatorname{PITD}\), and, for \(c\in Q\) and \(u,v\in\bar A\), a stable
causal effect \(\operatorname{Caus}(c,u,v)\) and a stable probabilistic
causal effect \(\operatorname{ProbCaus}(c,u,v)\):
\begin{align*}
  \operatorname{ITD}(q,\alpha,\Gamma,r)
  &:\Longleftrightarrow
  \exists\,g:Q\times\bar A\to R,\
    g(q(h),\alpha(a))=r(\Gamma(h,a))\ \forall h,a,\\
  \operatorname{PITD}(q,\alpha,K)
  &:\Longleftrightarrow
  \exists\,k:Q\times\bar A\to K_0,\
    k(q(h),\alpha(a))=K(h,a)\ \forall h,a,\\
  \operatorname{Caus}(c,u,v)
  &:\Longleftrightarrow
  \operatorname{Adm}(c,u)\wedge\operatorname{Adm}(c,v)
  \wedge\bar\Gamma(c,u)\ne\bar\Gamma(c,v)
  \wedge\operatorname{ResStable}(c,u,v),\\
  \operatorname{ProbCaus}(c,u,v)
  &:\Longleftrightarrow
  \operatorname{Adm}(c,u)\wedge\operatorname{Adm}(c,v)
  \wedge\bar K(c,u)\ne\bar K(c,v)
  \wedge\operatorname{ProbResStable}(c,u,v).
\end{align*}
Then
\begin{align*}
  \operatorname{ITD}(q,\alpha,\Gamma,r)
  &\iff \mathcal O_{\Gamma}=\emptyset
  \iff
  \exists!\,g:Q\times\bar A\to R,\
    g(q(h),\alpha(a))=r(\Gamma(h,a))\ \forall h,a,\\
  \operatorname{PITD}(q,\alpha,K)
  &\iff \mathcal O_{K}=\emptyset
  \iff
  \exists!\,k:Q\times\bar A\to K_0,\
    k(q(h),\alpha(a))=K(h,a)\ \forall h,a.
\end{align*}
Stable causal effects have raw witnesses, each under its own descent
equation. If \(\bar\Gamma\) is a descended transport,
\(\bar\Gamma(q(h),\alpha(a))=r(\Gamma(h,a))\) for all \(h,a\), then
\(\operatorname{Caus}(c,u,v)\) implies that there are \(h\in H\) and
\(a_1,a_2\in A\) with \(q(h)=c\), \(\alpha(a_1)=u\), \(\alpha(a_2)=v\) and
\(r(\Gamma(h,a_1))\ne r(\Gamma(h,a_2))\). If \(\bar K\) is a descended
transport, \(\bar K(q(h),\alpha(a))=K(h,a)\) for all \(h,a\), then
\(\operatorname{ProbCaus}(c,u,v)\) implies the same with \(K\) in place of
\(r\circ\Gamma\).
\end{theorem}
\begin{proof}
Assume the surjective quotients \(q\) and \(\alpha\), the transport
\(\Gamma\), the target readout \(r\), and the kernel outcome \(K\).
The first line is the descent criterion of \Cref{sec:apparatus} for
the surjective source \(q\times\alpha:H\times A\to Q\times\bar A\) and the
readout \((h,a)\mapsto r(\Gamma(h,a))\), with
\(g(q(h),\alpha(a)):=r(\Gamma(h,a))\); the probabilistic statement is
identical with \(K\) in place of \(r\circ\Gamma\).

A causal effect at context class \(c\) compares two intervention
classes \(u\) and \(v\).  It requires both to be admissible at \(c\),
the quotient-level transport to give different results on them (a
nonzero residue), and the residue to be stable; this is the definition of
\(\operatorname{Caus}\), and \(\operatorname{ProbCaus}\) is the same with
\(\bar K\) and \(\operatorname{ProbResStable}\).  A zero residue,
\(\bar\Gamma(c,u)=\bar\Gamma(c,v)\), is a finding of no effect, not a
failure of transport: a transport can descend with every residue zero.

For the last clause, choose by surjectivity \(h\) with \(q(h)=c\) and
\(a_1,a_2\) with \(\alpha(a_1)=u\), \(\alpha(a_2)=v\). Then
\(r(\Gamma(h,a_1))=\bar\Gamma(c,u)\ne\bar\Gamma(c,v)=r(\Gamma(h,a_2))\), using
only the descent equation of \(\bar\Gamma\); the same computation applies to
\(K\) and \(\bar K\), using only the descent equation of \(\bar K\).
Uniqueness of \(g\) and \(k\) holds because \(q\times\alpha:H\times A\to Q\times\bar A\) is surjective.
\end{proof}

\paragraph{Nonclaims.} The theorem does not identify causes with raw interventions and does
not apply to undeclared alternatives.  The intervention-counterfactual
hierarchy that this layer-agnostic formulation echoes is developed
in detail in \citet{Pearl2009Causality}.

\subsection*{Layer instantiations}\label{layer-inst:F25}

\begingroup\small
\begin{description}
\item[\emph{Randomised controlled trial / Fisher
randomisation (experimental design).}] \textit{Analogy.}
The carrier is the population of experimental units
(patients, field plots); the context quotient collapses
units to their pre-randomisation baseline stratum; the
intervention space is the set of treatment arms; the
abstraction collapses operational delivery details to
the trial-arm class; the transport is the trial's
delivery-and-follow-up process; the target readout is
the primary outcome.  Individual outcomes do not
deterministically descend through (baseline stratum,
arm); the \Fref{F25} reading is its probabilistic form:
randomised assignment supplies admissibility, and the
stable residue is the arm-specific outcome distribution
or the within-stratum average treatment effect, secured
by pre-registration, blinding, and intention-to-treat
analysis \citep{Fisher1935DesignExperiments}.

\end{description}
\endgroup


\subsection{Moduli / Contingency [\texorpdfstring{\Fref{F26}}{F26}]}\label{sec:status-records-coherence:moduli}

\paragraph{Reading the statement.} The symbols are those of \Cref{thm:F26} below.
Lawful closures are grouped into moduli by the equivalence \(\sim\), and
\(p\) sends a lawful closure to its modulus. A readout can fail to be a
function of the modulus (it is presentation-dependent), or it can be a
function of the modulus and then either take one value on every lawful
closure (necessary) or take different values on two inequivalent lawful
closures (contingent). The definitions come first; the theorem proves that
these alternatives are exhaustive and exclusive, and that descent through
\(p\) is the same as independence of presentation. The predicate
\(\operatorname{NecessaryBackground}\) of \Cref{thm:F50} is necessity in the
present sense for a background already defined on moduli.

\begin{theorem}[Moduli / Contingency]\label{thm:F26}
Let \(\mathcal C\) be a closure space with lawfulness predicate
\(\operatorname{Lawful}:\mathcal C\to\mathrm{Prop}\), equivalence relation
\(\sim\) on lawful closures (an equivalence relation, by the kernel condition), moduli quotient \(p:\mathcal C\to M\) with
\(p(c)=p(d)\iff c\sim d\) for lawful \(c,d\), and readout
\(r:\mathcal C\to V\). Assume that every \(m\in M\) is \(p(c)\) for some
lawful \(c\). Define, with \(c,d\) ranging over lawful closures,
\begin{align*}
  \operatorname{ClosureObstructed}
  &:\Longleftrightarrow
  \text{no closure is lawful},\\
  \operatorname{NontrivialModuli}
  &:\Longleftrightarrow
  \exists c,d:\ \neg(c\sim d),\\
  \operatorname{ModuliReadout}(r)
  &:\Longleftrightarrow
  \forall c,d:\ c\sim d\Rightarrow r(c)=r(d),\\
  \operatorname{PresentationDependent}(r)
  &:\Longleftrightarrow
  \exists c,d:\ c\sim d\ \wedge\ r(c)\ne r(d),\\
  \operatorname{NecessaryReadout}(r)
  &:\Longleftrightarrow
  \forall c,d:\ r(c)=r(d),\\
  \operatorname{ContingentReadout}(r)
  &:\Longleftrightarrow
  \exists c,d:\ \neg(c\sim d)\ \wedge\ r(c)\ne r(d).
\end{align*}
Then:
\begin{enumerate}[label=(\roman*)]
\item \(\operatorname{ModuliReadout}(r)\) holds if and only if there is
  \(\bar r:M\to V\) with \(\bar r(p(c))=r(c)\) for every lawful \(c\), and
  \(\operatorname{PresentationDependent}(r)\) holds if and only if
  \(\operatorname{ModuliReadout}(r)\) fails;
\item if \(\operatorname{ModuliReadout}(r)\) holds, then exactly one of
  \(\operatorname{NecessaryReadout}(r)\) and
  \(\operatorname{ContingentReadout}(r)\) holds;
\item if \(\operatorname{NontrivialModuli}\) fails, every moduli readout is
  necessary; if \(\operatorname{ClosureObstructed}\) holds, every readout is
  necessary and none is contingent or presentation-dependent, since these
  three predicates quantify over lawful closures only; this part does not use
  the standing assumption. Under the standing assumption,
  \(\operatorname{ClosureObstructed}\) moreover forces \(M=\emptyset\), hence
  \(\mathcal C=\emptyset\), because \(p\) is defined on all of \(\mathcal C\);
\item for a map \(T:\mathcal C\to\mathcal C'\) into a second closure space
  with lawfulness \(\operatorname{Lawful}'\) and equivalence \(\sim'\), taking
  lawful closures to lawful closures, \(T\) respects moduli
  (\(c\sim d\Rightarrow T(c)\sim' T(d)\) for lawful \(c,d\)) if and only if no
  pair of lawful closures \(c\sim d\) has \(\neg(T(c)\sim' T(d))\).
\end{enumerate}
\end{theorem}
\begin{proof}
(i) If \(\bar r\) exists and \(c\sim d\), then \(p(c)=p(d)\), so
\(r(c)=\bar r(p(c))=\bar r(p(d))=r(d)\). Conversely, let \(r\) be a moduli
readout. For \(m\in M\) choose a lawful \(c_m\) with \(p(c_m)=m\) and set
\(\bar r(m):=r(c_m)\). For lawful \(c\), \(p(c_{p(c)})=p(c)\) gives
\(c_{p(c)}\sim c\), so \(\bar r(p(c))=r(c_{p(c)})=r(c)\). The second
statement is the negation of the definition of a moduli readout.

(ii) Let \(r\) be a moduli readout. If \(r\) is necessary, no two lawful
closures have different readouts, so \(r\) is not contingent. If \(r\) is not
necessary, there are lawful \(c,d\) with \(r(c)\ne r(d)\); since \(r\) is a
moduli readout, \(c\not\sim d\), so \(r\) is contingent.

(iii) If \(\operatorname{NontrivialModuli}\) fails, all lawful closures are
equivalent, so a moduli readout takes one value on them. If no closure is
lawful, the universal conditions hold vacuously and the existential ones fail.

(iv) The obstruction pairs are exactly the failures of the implication that
defines respecting moduli.
\end{proof}

\paragraph{Nonclaims.} The theorem does not enumerate a carrier's moduli and does not
collapse stable parameters with contingent choices.

\subsection*{Layer instantiations}\label{layer-inst:F26}

\begingroup\small
\begin{description}
\item[\emph{Moduli space of elliptic curves (algebraic
geometry).}] \textit{Special case.}
The carrier is the set of Weierstrass-form presentations
\(y^2=x^3+ax+b\) over an algebraically closed base field of characteristic other
than 2 and 3; the lawfulness
predicate restricts to non-singular cubics
\((4a^3+27b^2\ne 0)\); the equivalence is isomorphism of
elliptic curves; the moduli quotient sends a
non-singular presentation to its isomorphism class and every singular
presentation to the class of \(y^2=x^3-x\) (\Fref{F26}'s kernel condition is imposed only on lawful
closures). Take \(M\) to be the isomorphism classes of
nonsingular presentations, and extend \(j\) to singular
presentations by \(j(a,b)=1728\). This total readout is a moduli readout: distinct Weierstrass
presentations of isomorphic elliptic curves share the
same j-invariant, and the descended map
\(\bar\jmath:M\to\mathbb A^1\) is a bijection onto the
coarse moduli space.  It is contingent, since non-isomorphic
curves have different j-invariants.  The individual coefficients
\((a,b)\) are presentation-dependent: \((a,b)\) and
\((\lambda^4a,\lambda^6b)\) present isomorphic curves.  The genus, equal
to one on every lawful closure, is a necessary readout
\citep{Mumford1965GIT}.

\item[\emph{Crystallographic space groups and Bravais
lattices (solid-state physics).}] \textit{Special case.}
The carrier is the set of three-dimensional crystalline
arrangements; the lawfulness
predicate requires invariance under a rank-3 lattice of
translations; the equivalence is conjugacy of space groups under
orientation-preserving affine maps (change of origin and
right-handed basis); the moduli quotient sends a crystal to its
space-group label among the 230 space-group types.  Abstract
isomorphism, equivalently conjugacy under all affine maps, merges
the 11 enantiomorphic pairs and gives 219 classes.  The space-group label is a
moduli readout, and a contingent one, since lawful crystals with different
space groups exist (rock salt, \(Fm\bar3m\); \(\alpha\)-quartz,
\(P3_121\)).  The lattice metric (unit-cell lengths
\(a,b,c\), angles \(\alpha,\beta,\gamma\),
atomic-displacement parameters) is presentation-dependent for this
equivalence: lawful crystals with the same space group can have different
lattice metrics \citep{HahnITC2016}.

\item[\emph{Standard Model free parameters as
renormalisation moduli (particle physics).}] \textit{Analogy.}
The carrier is the space of Standard-Model Lagrangians
with the fixed gauge group
SU(3)\(\times\)SU(2)\(\times\)U(1), three-generation
fermion content, and Higgs doublet built into the
lawfulness predicate (renormalisability, gauge
invariance, anomaly cancellation, observed particle
content); the equivalence identifies Lagrangians giving
identical low-energy phenomenology under the
renormalisation group; the moduli quotient sends a
Lagrangian to its renormalised-parameter point.  The 19
parameters of the Standard Model with massless neutrinos (three
gauge couplings, nine charged-fermion masses, four CKM-matrix
parameters, the Higgs mass and vacuum expectation value, and the QCD
vacuum angle), extended by the neutrino-sector parameters, are
contingent moduli readouts; in dimensionless form the fermion
masses are replaced by Yukawa couplings and the Higgs mass by the
quartic self-coupling; predicted phenomenology descends
through the renormalised-parameter point
\citep{Workman2022Review}.

\end{description}
\endgroup


\subsection{Conservation as Orbit Descent [\texorpdfstring{\Fref{F27}}{F27}]}\label{sec:status-records-coherence:conservation}

\paragraph{Reading the statement.}
A declared route or symmetry structure on \(X\) is a set \(\Gamma\) of partial
steps \(\gamma:\operatorname{dom}\gamma\to X\) with
\(\operatorname{dom}\gamma\subseteq X\); a family of maps \(X\to X\) is the case
\(\operatorname{dom}\gamma=X\). A declared relation \(\to\) on \(X\) is the case
\(\Gamma=\{\gamma_{x,x'}:x\to x'\}\) with \(\operatorname{dom}\gamma_{x,x'}=\{x\}\) and
\(\gamma_{x,x'}(x)=x'\); its orbits are the classes of the equivalence generated
by \(\to\), and conservation says \(C(x)=C(x')\) whenever \(x\to x'\). The reaction,
Reidemeister and reduction instances below are declared in this form. The \emph{orbits} are the classes of the
equivalence relation generated by \(x\sim\gamma(x)\) for \(\gamma\in\Gamma\)
and \(x\in\operatorname{dom}\gamma\), and \(q_{\mathrm{orb}}\) is the quotient
map onto them. A quantity \(C:X\to V\) is \emph{conserved} when
\(C(\gamma(x))=C(x)\) for every \(\gamma\in\Gamma\) and
\(x\in\operatorname{dom}\gamma\).

\begin{theorem}[Conservation as Orbit Descent]\label{thm:F27}
Let \(\Gamma\) be a set of partial maps \(\gamma:\operatorname{dom}\gamma\to X\)
with \(\operatorname{dom}\gamma\subseteq X\), let \(q_{\mathrm{orb}}\) be the
quotient map of the equivalence relation generated by \(x\sim\gamma(x)\)
(\(\gamma\in\Gamma\), \(x\in\operatorname{dom}\gamma\)), and let \(C:X\to V\).
Then \(C\) is conserved, \(C(\gamma(x))=C(x)\) for all \(\gamma\in\Gamma\) and
\(x\in\operatorname{dom}\gamma\), if and only if it descends through the orbit
quotient, \(C=\bar C\circ q_{\mathrm{orb}}\) for some \(\bar C\); that is,
\(C(x)=C(x')\) whenever \(x\) and \(x'\) lie in the same orbit.
\end{theorem}
\begin{proof}
If \(C\) descends through the orbit quotient, then, since \(x\) and
\(\gamma(x)\) lie in the same orbit for \(x\in\operatorname{dom}\gamma\), \(C(\gamma(x))=C(x)\), so \(C\) is
conserved.  Conversely, let \(C\) be conserved. The relation
\(C(x)=C(x')\) is an equivalence relation containing every generating pair
\((x,\gamma(x))\), so it contains the generated equivalence: \(C\) is constant
on orbits.  Define the quotient readout by \(\bar C([x])=C(x)\).  Orbit constancy
makes the definition independent of representative, and
\(C=\bar C\circ q_{\mathrm{orb}}\).  Thus conservation is orbit
descent.
\end{proof}

For the conserved quantity to be lawful at the current layer in the sense of \Cref{thm:F18}, \(C\) must also factor through the access quotient \(\accessQ\) of that layer; descent through \(q_{\mathrm{orb}}\) alone does not give this.

\paragraph{Nonclaims.} The theorem does not enumerate conserved quantities and does not
apply to undeclared routes or symmetries.  Noether's theorem
\citep{Noether1918Invariante} produces conserved quantities from continuous
variational symmetries; the theorem above reads each such quantity as a
function on the orbits of the flow.

\subsection*{Layer instantiations}\label{layer-inst:F27}

\begingroup\small
\begin{description}
\item[\emph{Lavoisier conservation of mass in chemical
reactions (chemistry).}] \textit{Special case.}
The carrier is the set of finite multisets of molecules, each molecule a finite multiset of atoms of declared elements \(e\) with masses \(m_e\); the declared route structure
is the set of elementary chemical reactions that
preserve atomic identity; the orbit quotient refines the map sending a
configuration to its multiset of atoms.  Total mass
\(C(x)=\sum_e n_e(x)\,m_e\) is constant across orbits
of chemical rearrangement and so descends through the
orbit quotient.  Lavoisier's stoichiometric balancing
of chemical equations is the founding expression of
this orbit-descent reading \citep{Lavoisier1789Traite}.

\item[\emph{Knot invariants under Reidemeister moves
(combinatorial knot theory).}] \textit{Special case.}
The carrier is the set of planar knot diagrams; the
declared route structure is the Reidemeister moves (R1
twist, R2 poke, R3 slide) together with planar isotopy;
the orbit quotient sends a diagram to its
ambient-isotopy class.  Polynomial invariants such as
the Jones polynomial or the Conway-normalized Alexander polynomial are unchanged under
every Reidemeister move and so descend through the
orbit quotient to functions on knots.  Reidemeister's
theorem certifies that ambient isotopy of knots is
exactly generated by the three planar moves
\citep{Reidemeister1927Knotentheorie}.

\item[\emph{Subject reduction in operational semantics
(programming-language theory).}] \textit{Special case.}
The carrier is the set of typable program expressions
under a fixed typing context \(\Gamma\) and a type
system in which each typable expression has a unique type (as in the
Church-style simply typed \(\lambda\)-calculus); the declared route structure is the small-step
reduction relation \(\to\) with its reachability
preorder.  Subject reduction (type preservation) is the
conservation statement: if
\(\Gamma\vdash e:\tau\) and \(e\to e'\) then
\(\Gamma\vdash e':\tau\).  On the typable carrier every reduction step therefore links
two expressions of the same type, so the type readout is constant on
each class of the equivalence generated by \(\to\) and descends
through this orbit quotient; no subject expansion is needed, because
both ends of each step are typable.  Together with progress, it yields
the Wright-Felleisen formulation of type soundness
\citep{WrightFelleisen1994SyntacticTypeSoundness}.

\item[\emph{Meselson-Stahl semi-conservative {DNA}
replication (molecular biology).}] \textit{Analogy.}
The carrier is the population of DNA duplexes in a
dividing cell; the declared route structure is one
round of replication, in which each parental duplex
yields two daughter duplexes that each contain one
parental strand and one newly synthesised strand; the
orbit quotient sends a daughter duplex to its parental
template-strand identity.  The conserved quantity is
the template-strand sequence content carried by each
daughter molecule, which is constant across the
replication orbit (modulo replication-error
corrections).  The Meselson-Stahl density-gradient
experiment is the canonical demonstration,
distinguishing semi-conservative replication from
conservative and dispersive alternatives
\citep{MeselsonStahl1958Replication}.

\item[\emph{Martingale property under conditional
expectation (probability theory).}] \textit{Special case.}
For a stochastic process \((M_t)_{t\ge 0}\) with \(\mathbb E|M_t|<\infty\) for every \(t\), adapted to a filtration
\((\mathcal F_t)\) and a fixed time \(s\), the carrier is the time set
\([s,\infty)\); the declared route structure is time advance, whose
orbit quotient collapses \([s,\infty)\) to a single class; the
quantity is \(C_s(t)=\mathbb E[M_t\mid\mathcal F_s]\), valued in
integrable random variables modulo almost-sure equality.  The
martingale property is this conservation statement for every \(s\):
\(\mathbb E[M_t\mid\mathcal F_s]=M_s\) almost surely for all
\(t\ge s\), so \(C_s\) is constant along
the orbit, equal to its value \(M_s\) at \(t=s\).  The
\(\mathcal F_s\)-measurability of \(\mathbb E[M_t\mid\mathcal F_s]\)
holds for every integrable process and is not the content.  Optional
stopping, Doob's maximal inequalities, and martingale convergence
are built on this conservation property
\citep{Doob1953Stochastic}.
\end{description}
\endgroup

\subsection{Composability [\texorpdfstring{\Fref{F28}}{F28}]}\label{sec:status-records-coherence:composability}

\paragraph{Reading the statement.} The symbols are those of \Cref{thm:F28} below.
\(\operatorname{COD}\), \(\operatorname{CRD}\) and
\(\operatorname{AD}\) say that the composite outcome, the cross residual
and admissibility depend only on the parts' classes and interface
values \(\kappa(a,b)\). \(\operatorname{Composable}\) is exact
composability, the conjunction of the three. \(\operatorname{Assoc}\)
and \(\operatorname{Id}\) are the coherence equations named in the
display.

\begin{theorem}[Composability]\label{thm:F28}
Let \(X_A,X_B\) be component spaces with quotients
\(q_A:X_A\to Q_A\), \(q_B:X_B\to Q_B\), interface readouts
\(\beta_A:X_A\to I_A\), \(\beta_B:X_B\to I_B\), composition
\(U:X_A\times X_B\to X_{AB}\), composite quotient
\(q_{AB}:X_{AB}\to Q_{AB}\), admissibility predicate
\(\operatorname{Adm}:X_A\times X_B\to\mathrm{Prop}\), and residual
\(\delta:X_A\times X_B\to D\).  Write
\[
  \kappa(a,b):=\bigl((q_A(a),\beta_A(a)),(q_B(b),\beta_B(b))\bigr)
\]
for the part/interface quotient and assume \(\kappa\) is surjective.
For a map \(f\) on \(X_A\times X_B\), let
\[
  \mathcal O(f)
  :=
  \{((a,b),(a',b')) :
     \kappa(a,b)=\kappa(a',b'),\ f(a,b)\ne f(a',b')\},
\]
and set \(\mathcal O_U:=\mathcal O(q_{AB}\circ U)\),
\(\mathcal O_{\delta}:=\mathcal O(\delta)\) and
\(\mathcal O_{\operatorname{Adm}}:=\mathcal O(\operatorname{Adm})\).
Write \(\operatorname{COD}\), \(\operatorname{CRD}\) and
\(\operatorname{AD}\) for descent through \(\kappa\) of the composite
outcome, the cross residual and admissibility, and
\(\operatorname{Composable}\) for exact composability.  Then
\begin{align*}
  \operatorname{COD}
  &\iff
  \mathcal O_U=\emptyset
  \iff
    \exists!\,\bar U:(Q_A\times I_A)\times(Q_B\times I_B)\to Q_{AB}\\
  &\hspace{6em}\text{with}\ \bar U(\kappa(a,b))=q_{AB}(U(a,b)),\\
  \operatorname{CRD}
  &\iff
  \mathcal O_{\delta}=\emptyset,
  \qquad
  \operatorname{AD}
  \iff
  \mathcal O_{\operatorname{Adm}}=\emptyset,\\
  \operatorname{Composable}
  &\iff
  \mathcal O_{\operatorname{Adm}}=\emptyset
  \wedge\mathcal O_U=\emptyset
  \wedge\mathcal O_{\delta}=\emptyset.
\end{align*}
The maps through which the cross residual and admissibility descend are likewise unique.
Associativity and identity are separate coherence conditions: for the
two bracketings \(\ell,\varrho:Y\to Y'\) of a declared quotient-level triple
composite, define \(\operatorname{Assoc}:\Longleftrightarrow\ell=\varrho\), and
for a declared composite \(c:Z\times E\to Z\) and a declared \(e\in E\) define
\(\operatorname{Id}:\Longleftrightarrow\forall z\in Z,\ c(z,e)=z\); then
\(\operatorname{Assoc}\iff\{y:\ell(y)\ne\varrho(y)\}=\emptyset\) and
\(\operatorname{Id}\iff\{z\in Z:c(z,e)\ne z\}=\emptyset\). If moreover \(\ell\circ\kappa_3=q'\circ L\) and \(\varrho\circ\kappa_3=q'\circ P\) for raw bracketings \(L,P:X_3\to X'\), a surjection \(\kappa_3:X_3\to Y\) and a map \(q':X'\to Y'\), then \(\operatorname{Assoc}\iff q'\circ L=q'\circ P\); in particular raw associativity \(L=P\) gives \(\operatorname{Assoc}\). Likewise, if \(c(\kappa_Z x,e)=\kappa_Z(\tilde c(x))\) for a surjection \(\kappa_Z\) and a map \(\tilde c\), then \(\operatorname{Id}\iff\kappa_Z\circ\tilde c=\kappa_Z\).
\end{theorem}
\begin{proof}
Assume the component quotients \(q_A,q_B\), interface readouts
\(\beta_A,\beta_B\), composition \(U\), composite quotient
\(q_{AB}\), admissibility predicate \(\operatorname{Adm}\), residual
\(\delta\), and surjectivity of \(\kappa\).  The composite outcome
descends exactly when it is independent of representatives of the
\(\kappa\)-classes, that is, of the parts' classes and interface values.
Thus, if \(\mathcal O_U=\emptyset\), the definition
\[
  \bar U(\kappa(a,b)):=q_{AB}(U(a,b))
\]
is well-defined, and surjectivity of \(\kappa\) defines it everywhere; two such maps agree on the image of the surjection \(\kappa\), so they are equal;
if \(\mathcal O_U\ne\emptyset\), a pair in \(\mathcal O_U\) is precisely
a split-pair obstruction.  The same argument applied to \(\delta\) and
to \(\operatorname{Adm}\) gives the second line, and exact
composability is the conjunction of the three descent conditions.
Associativity and identity are pointwise equations, so each holds
exactly when its set of failures is empty. For the raw clauses, precompose with the surjections: \(\ell=\varrho\) iff \(\ell\circ\kappa_3=\varrho\circ\kappa_3\) iff \(q'\circ L=q'\circ P\), and \(c(\cdot,e)=\mathrm{id}_Z\) iff \(c(\kappa_Z x,e)=\kappa_Z x\) for all \(x\) iff \(\kappa_Z\circ\tilde c=\kappa_Z\).
\end{proof}

The three descent conditions characterize exact quotient-level
composability. Associativity and identity are additional equations for
declared composites.

\paragraph{Nonclaims.} The theorem does not enumerate composable records and does not apply
to undeclared compositions.

\subsection*{Layer instantiations}\label{layer-inst:F28}

\begingroup\small
\begin{description}
\item[\emph{Universally composable cryptographic
protocols (theoretical computer science).}] \textit{Analogy.}
The carriers are implementations of protocol components;
the source quotients collapse implementations to their
UC ideal-functionality class; the interface readouts
record input-output behaviour visible to the
environment; the composition combines components by
message-passing; the composite quotient collapses the
composite to its ideal-functionality class.  The UC
composition theorem yields composite-outcome descent;
the residual records simulator output and descends
through the same component quotient.  Admissibility,
associativity, and identity correspond to syntactic
well-formedness of the protocol-composition operation
\citep{Canetti2001UC}.

\end{description}
\endgroup


\subsection{Presentation Invariance [\texorpdfstring{\Fref{F29}}{F29}]}\label{sec:status-records-coherence:presentation-invariance}

\paragraph{Reading the statement.}
Each predicate on the left is defined by the first condition displayed on
its right, as the symbol \(:\Longleftrightarrow\) marks: invariance under change of presentation within a
\(\pi\)-class, invariance under the declared compatible relation \(\sim\),
invariance of the relation \(R\) under the tuple quotient, and
covariance of \(T\) under the group.

\begin{theorem}[Presentation Invariance]\label{thm:F29}
Let \(P\) be a presentation space with quotient
\(\pi:P\to\Pi\), let \(\sim\) be a relation compatible
with \(\pi\), that is, \(p\sim p'\Rightarrow\pi(p)=\pi(p')\), let \(s:\Pi\to P\) be a section with
\(\pi(s(x))=x\), and let \(T:P\to V\) be a readout.  Let
\(\tau:\operatorname{Tuple}(P)\to\operatorname{Tuple}(\Pi)\) be the
tuple quotient on a declared carrier \(\operatorname{Tuple}(P)\) of presentation
tuples, for instance \(P^n\) with \(\tau=\pi^n\), let \(R:\operatorname{Tuple}(P)\to\mathrm{Prop}\) be a
relation, and let a group \(G\) act on \(P\) and \(V\).  Define
\begin{align*}
  \operatorname{PresentationInvariant}(\pi,T)
  &:\Longleftrightarrow
  \forall p,p'\in P,\ \pi(p)=\pi(p')\Rightarrow T(p)=T(p'),\\
  \operatorname{EquivalentPresentationInvariant}(\sim,T)
  &:\Longleftrightarrow
  \forall p,p'\in P,\ p\sim p'\Rightarrow T(p)=T(p'),\\
  \operatorname{PresentationRelationInvariant}(\tau,R)
  &:\Longleftrightarrow
  \forall x,y,\ \tau(x)=\tau(y)\Rightarrow (R(x)\leftrightarrow R(y)),\\
  \operatorname{GroupEquivariant}(T)
  &:\Longleftrightarrow
  \forall g\in G,\forall p\in P,\ T(g\cdot p)=g\cdot T(p).
\end{align*}
Then
\[
  \operatorname{PresentationInvariant}(\pi,T)
  \iff
  \exists\,\bar T:\Pi\to V\ \text{with}\ \bar T\circ\pi=T,
\]
\(\operatorname{PresentationInvariant}(\pi,T)\) implies
\(\operatorname{EquivalentPresentationInvariant}(\sim,T)\); when \(\sim\) is
an equivalence relation, the converse holds for every \(T\) into a
two-element \(V\) exactly when \(p\sim p'\iff\pi(p)=\pi(p')\); and
\(\operatorname{PresentationInvariant}(\pi,T)\) holds exactly when
\(T(p)=T(s(\pi(p)))\) for every \(p\in P\). More generally, if \(r\) is any
relation on \(P\) whose equivalence closure is the kernel of \(\pi\), that is,
\(\pi(p)=\pi(p')\) exactly when \(p\) and \(p'\) are joined by a finite chain
of \(r\)-steps taken in either direction, then
\(\operatorname{PresentationInvariant}(\pi,T)\) holds exactly when
\(T(p)=T(p')\) whenever \(r(p,p')\). For every \(\tau\), \(\operatorname{PresentationRelationInvariant}(\tau,R)\) holds if and only if
\(R=\bar R\circ\tau\) for some \(\bar R:\operatorname{Tuple}(\Pi)\to\mathrm{Prop}\).
\(\operatorname{GroupEquivariant}(T)\) is the displayed intertwining condition;
presentation invariance does not imply it.
\end{theorem}
\begin{proof}
Assume the quotient \(\pi\), compatible relation \(\sim\), section
\(s\), readout \(T\), tuple quotient \(\tau\), relation \(R\), and
group action of \(G\).  If \(T\) is presentation-invariant, it is
constant on \(\pi\)-fibers.  Conversely, if \(T\) is constant on
\(\pi\)-fibers, define
\[
  \bar T(x):=T(s(x)).
\]
For \(p\in P\), the presentations \(p\) and \(s(\pi(p))\) lie in the same
\(\pi\)-fiber, because \(\pi(s(\pi(p)))=\pi(p)\); since \(T\) is constant on
\(\pi\)-fibers, \(\bar T(\pi(p))=T(s(\pi(p)))=T(p)\), that is,
\(\bar T\circ\pi=T\). Since \(p\sim p'\) implies \(\pi(p)=\pi(p')\), a readout constant on
\(\pi\)-fibers is invariant under \(\sim\). If \(\sim\) is the kernel of
\(\pi\), the converse is immediate. If \(\sim\) is an equivalence relation
strictly smaller than the kernel, pick \(p\not\sim p'\) with
\(\pi(p)=\pi(p')\); the indicator of the \(\sim\)-class of \(p\) is
\(\sim\)-invariant but not presentation-invariant. If \(T\) is
presentation-invariant, then \(T(p)=T(s(\pi(p)))\) because \(p\) and
\(s(\pi(p))\) lie in one \(\pi\)-fiber; conversely, if \(T(p)=T(s(\pi(p)))\) for
every \(p\) and \(\pi(p)=\pi(p')\), then \(T(p)=T(s(\pi(p)))=T(s(\pi(p')))=T(p')\).

The relation statement holds for every \(\tau\): given invariance, put \(\bar R(z):\Leftrightarrow\exists x,\ \tau(x)=z\wedge R(x)\); then \(\bar R(\tau(y))\Leftrightarrow R(y)\). The converse is immediate.  The group-action statement is a separate
intertwining condition:
\[
  T(g\cdot p)=g\cdot T(p).
\]
For a generating relation \(r\): if \(T\) is presentation-invariant, then
\(r(p,p')\) gives \(\pi(p)=\pi(p')\) and so \(T(p)=T(p')\); conversely, if
\(T\) is constant along \(r\)-steps, it is constant along every finite chain of
steps taken in either direction, by induction on the length, hence on every
\(\pi\)-fiber.
For the non-implication, take \(P=\Pi=V=\{0,1\}\) and \(\pi=s=T=\mathrm{id}\),
and let \(C_2\) swap the points of \(P\) and act trivially on \(V\). Then \(T\) is
presentation-invariant, but \(T(g\cdot0)=1\ne0=g\cdot T(0)\) for the generator
\(g\).
\end{proof}

Readouts and relations satisfying the stated descent conditions are
independent of representatives of their respective presentation quotients.
The group-action clause characterizes equivariance; presentation independence
of a formed record requires descent of each relevant field.

\paragraph{Nonclaims.} The theorem does not enumerate presentations and does not collapse
the presentation/object distinction.

\subsection*{Layer instantiations}\label{layer-inst:F29}

\begingroup\small
\begin{description}
\item[\emph{Coordinate-free formulation of general
relativity (mathematical physics).}] \textit{Special case.}
The presentation space is the set of
coordinate-component representations of a geometric
tensor over a spacetime manifold; the quotient
collapses representations to the underlying
coordinate-free tensor field.  The diffeomorphism
group acts on representations by pushforward. The
readout \(T\) sending a coordinate representation to its tensor field is
presentation-invariant; it is also diffeomorphism-equivariant,
\(T(\phi\cdot p)=\phi_*T(p)\), which \Fref{F29} treats as a separate
condition.  General
covariance is the requirement that physically
meaningful readouts be presentation-invariant in this
sense \citep{MisnerThorneWheeler1973Gravitation}.

\item[\emph{Gr\"obner basis as canonical
polynomial-ideal presentation (computer algebra).}] \textit{Special case.}
The presentation space \(P\) is the set of finite subsets of
\(k[x_1,\dots,x_n]\), and \(\pi\) sends a set to the ideal it generates.
Under a fixed monomial ordering, the reduced Gr\"obner basis supplies a
canonical section: every ideal has a unique reduced Gr\"obner basis.
Membership of a fixed polynomial is a presentation-invariant readout.
Equality of generated ideals is an invariant relation on \(P^2\); for any
fixed ideal \(J\), so is the predicate \(I(A)\cap I(B)=J\), fitting the
relation clause with \(\tau=\pi^2\). Membership and equality can be checked
using Gr\"obner bases and normal forms, while an intersection can be computed
using an elimination Gr\"obner basis. Buchberger's algorithm constructs the
required bases \citep{Buchberger1965ThesisGroebner}.

\item[\emph{Jordan canonical form and matrix similarity
(linear algebra).}] \textit{Special case.}
For a finite-dimensional vector space \(V\) over an
algebraically closed field \(k\), the presentation
space is the set of matrix representatives of linear
endomorphisms; the quotient sends a matrix to its
similarity class.  Under a fixed eigenvalue and
block ordering, the Jordan canonical form supplies a
canonical section.  Trace, determinant, characteristic
polynomial, minimal polynomial, eigenvalue multiset
with algebraic and geometric multiplicities, rank, and
nullity are presentation-invariant readouts.  The
general linear group \(\mathrm{GL}_n(k)\) acts on
matrices by similarity; the listed invariants are
\(\mathrm{GL}_n(k)\)-invariant readouts, but they do not
separate similarity classes (nilpotent \(7\times 7\)
matrices with Jordan blocks of sizes \(3,3,1\) and
\(3,2,2\) agree on all of them), while the Jordan form,
read through the canonical section, is a complete
invariant \citep{HoffmanKunze1971LinearAlgebra}.

\item[\emph{De Bruijn indices for alpha-equivalence
(programming-language theory).}] \textit{Special case.}
The presentation space is the set of lambda terms with
named bound variables; the quotient sends a term to its
alpha-equivalence class.  The de Bruijn-index encoding
(each bound variable replaced by the distance to its
binder) is a complete invariant: two terms have the
same encoding iff they are alpha-equivalent.  Renaming
each bound variable by the depth of its binder, from a
reserved supply of names disjoint from the free
variables, gives a canonical named representative of
each class, which is the section \(s\).  The set of alpha-equivalence classes of one-step beta-reducts,
and the alpha-equivalence class of the beta-normal form when
one exists (with a marker otherwise), are presentation-invariant
readouts. Named-term outputs need not be presentation-invariant,
unless a canonical representative convention is imposed; the de Bruijn encoding
solves the variable-capture problem by construction
\citep{deBruijn1972LambdaNotation}.

\item[\emph{Lewin transformational music theory and the
chromatic group (music theory).}] \textit{Special case.}
The presentation space is the set of pitch-class
assignments to a musical motif; the chromatic group
\(G=\mathbb Z/12\mathbb Z\) acts by transposition, and
the dihedral group \(D_{12}\) of order 24 extends this with
inversion.  Some motif features, such as the pitch-class set, transform equivariantly,
\(T(g\cdot p)=g\cdot T(p)\); others are invariant.  Pitch-class invariants
such as the interval-class content (a canonical
\(G\)-invariant readout) descend through the orbit
quotient of the \(D_{12}\)-action and are
presentation-invariant musical features
\citep{Lewin1987GMIT}.
\end{description}
\endgroup

\subsection{Explanation as Compression with Descent [\texorpdfstring{\Fref{F30}}{F30}]}\label{sec:status-records-coherence:explanation}

\paragraph{Reading the statement.} The symbols are those of \Cref{thm:F30} below.
\(\operatorname{ExactExplanation}(b,e,t)\) says that the explanation
\(e\) descends through the base \(b\) and the target \(t\) descends
through \(e\): the base determines the explanation, and the explanation
determines the target. \(\operatorname{StrictCompression}(b,e)\) says
that \(e\) descends through \(b\) but \(b\) does not descend through
\(e\), so the explanation forgets some distinction of the base.

\begin{theorem}[Explanation as Compression with Descent]\label{thm:F30}
Let \(H\) be a history space with base readout \(b:H\to B\),
explanation readout \(e:H\to E\), target readout \(t:H\to T\),
continuation \(\kappa:H\to H'\), later explanation
\(e':H'\to E'\), and later target \(t':H'\to T'\).  Assume \(b\)
and \(e\) are surjective.  Define the obstruction sets
\begin{align*}
  \mathcal O_{\mathrm{src}}
  &:=
  \{(h,h')\in H^2:b(h)=b(h')\ \text{and}\ e(h)\ne e(h')\},\\
  \mathcal O_{\mathrm{tgt}}
  &:=
  \{(h,h')\in H^2:e(h)=e(h')\ \text{and}\ t(h)\ne t(h')\},\\
  \mathcal O_{\mathrm{cmp}}
  &:=
  \{(h,h')\in H^2:e(h)=e(h')\ \text{and}\ b(h)\ne b(h')\},
\end{align*}
and the predicates
\begin{align*}
  \operatorname{ExactExplanation}(b,e,t)
  &:\Longleftrightarrow
  \exists\,\alpha:B\to E,\ \exists\,\beta:E\to T,\
    \alpha\circ b=e\ \text{and}\ \beta\circ e=t,\\
  \operatorname{StrictCompression}(b,e)
  &:\Longleftrightarrow
  \exists\,\alpha:B\to E,\ \alpha\circ b=e
  \ \text{and}\ 
  \neg\exists\,\delta:E\to B,\ \delta\circ e=b,\\
  \operatorname{DynamicExplanation}(e,\kappa,e',t')
  &:\Longleftrightarrow
  \exists\,\gamma:E\to E',\ \gamma\circ e=e'\circ\kappa
  \ \text{and}\
  \exists\,\theta:E'\to T',\ \theta\circ e'=t'.
\end{align*}
Then:
\begin{enumerate}[label=(\roman*)]
\item \(\operatorname{ExactExplanation}(b,e,t)\) holds if and only if
  \(\mathcal O_{\mathrm{src}}=\emptyset\) and
  \(\mathcal O_{\mathrm{tgt}}=\emptyset\);
\item \(\operatorname{StrictCompression}(b,e)\) holds if and only if
  \(\mathcal O_{\mathrm{src}}=\emptyset\) and
  \(\mathcal O_{\mathrm{cmp}}\ne\emptyset\): the explanation is determined by
  the base and merges two histories that the base separates;
\item if \(\operatorname{DynamicExplanation}(e,\kappa,e',t')\), then the later
  target \(t'\circ\kappa\) descends through \(e\): the explanation now
  determines the target after the continuation.
\end{enumerate}
\end{theorem}
\begin{proof}
Assume the base readout \(b\), explanation readout \(e\), target
readout \(t\), and the surjectivity of \(b\) and \(e\).

(i) By the descent criterion of \Cref{sec:apparatus}, since \(b\) and
\(e\) are surjective, \(\alpha\) exists exactly when
\(\mathcal O_{\mathrm{src}}=\emptyset\), and \(\beta\) exactly when
\(\mathcal O_{\mathrm{tgt}}=\emptyset\).

(ii) By (i), \(\alpha\) exists exactly when
\(\mathcal O_{\mathrm{src}}=\emptyset\), and by the descent criterion for
the surjective \(e\), a map \(\delta:E\to B\) with \(\delta\circ e=b\)
exists exactly when \(\mathcal O_{\mathrm{cmp}}=\emptyset\).  So strict compression is
\(\mathcal O_{\mathrm{src}}=\emptyset\) together with
\(\mathcal O_{\mathrm{cmp}}\ne\emptyset\).

(iii) Given \(\gamma\) and \(\theta\),
\(t'\circ\kappa=\theta\circ e'\circ\kappa=\theta\circ\gamma\circ e\), so
\(t'\circ\kappa\) factors through \(e\).
\end{proof}

\paragraph{Nonclaims.} The theorem does not rank explanations and does not identify
explanation with the fact explained.  The Kolmogorov-complexity and
minimum-description-length traditions develop the
shortest-program / shortest-code reading of explanation in detail
\citep{LiVitanyi2008Kolmogorov}.

\subsection*{Layer instantiations}\label{layer-inst:F30}

\begingroup\small
\begin{description}
\item[\emph{Newtonian mechanics explaining
Kepler's laws (mathematical physics).}] \textit{Analogy.}
The history space is the set of solar-system
trajectories; the base records observed positions and
velocities.  The explanation is Newton's gravitational
dynamics: masses of the bodies, the inverse-square
gravitational law, and initial conditions.  The target
is the set of Kepler-type empirical regularities
(elliptical orbits, equal areas in equal times,
\(T^2\propto a^3\)) and the ephemeris of future
planetary positions.  The descended map identifies
masses and initial conditions from kinematics, and
the target map solves the two-body problem to derive
Kepler's laws; because the dynamics is deterministic,
integrating the equations of motion recovers the trajectory from the
explanation, so this compression is exact but not strict in the
\Fref{F30} sense; its gain is that finitely many parameters and one
law replace the continuous trajectory \citep{Newton1687Principia}.

\item[\emph{Principal component analysis as
dimensionality-reducing explanation (statistics).}] \textit{Special case.}
The history space is observed records
\(x\in\mathbb R^p\); the base is the full record.
The explanation is the top-\(k\) principal-component
representation, computed as eigenvectors of the
empirical covariance matrix.  The target is the
low-rank reconstruction or rank-\(k\) orthogonal
projection visualisation, both exactly determined by
the principal-component representation through a
linear map.  Taking the history space to be all of
\(\mathbb R^p\), the compression is strict whenever \(k<p\), because
the rank-\(k\) projection is not injective; on a finite data set the
projection can be injective, and non-zero discarded variance then
measures reconstruction error rather than strictness.  The
explained-variance ratio measures the compression fidelity \citep{Pearson1901PCA}.

\end{description}
\endgroup


\subsection{Counterfactual Legitimacy [\texorpdfstring{\Fref{F31}}{F31}]}\label{sec:status-records-coherence:counterfactual}

\paragraph{Reading the statement.} The symbols are those of \Cref{thm:F31} below.
\(\operatorname{CFL}\) says that admissibility is constant on
\((q,\alpha)\)-classes and that target truth is constant on the
admissible pairs of each class, so that a counterfactual claim depends
only on the declared context and intervention classes. The target
quotient \(q'\) is recorded with the data, but
neither obstruction set uses it.

\begin{theorem}[Counterfactual Legitimacy]\label{thm:F31}
Let \(q:H\to Q\) be a context quotient, let \(\alpha:A\to\bar A\) be an intervention abstraction, let \(\operatorname{Adm}:H\times A\to\mathrm{Prop}\) be admissibility, let
\(\Gamma:H\times A\to H'\) be a counterfactual route, let
\(q':H'\to R\) be a target quotient, and let
\(\tau:H'\to\mathrm{Bool}\) be target truth.  Define
\[
  \mathcal O_{\mathrm{adm}}
  :=
  \{((h,a),(h',a')) :
    q(h)=q(h'),\ \alpha(a)=\alpha(a'),\
    \operatorname{Adm}(h,a)\not\leftrightarrow
    \operatorname{Adm}(h',a')\},
\]
and
\[
  \mathcal O_{\mathrm{cf}}
  :=
  \left\{
  \begin{array}{l}
    ((h,a),(h',a')) :
    q(h)=q(h'),\ \alpha(a)=\alpha(a'),\\
    \operatorname{Adm}(h,a),\ \operatorname{Adm}(h',a'),\\
    \tau(\Gamma(h,a))\ne\tau(\Gamma(h',a'))
  \end{array}
  \right\}.
\]
Define the counterfactual-legitimacy predicate by
\(\operatorname{CFL}(q,\alpha,\operatorname{Adm},\Gamma,q',\tau):\Longleftrightarrow\)
for all \(h,h'\in H\) and \(a,a'\in A\) with \(q(h)=q(h')\) and
\(\alpha(a)=\alpha(a')\), \(\operatorname{Adm}(h,a)\leftrightarrow\operatorname{Adm}(h',a')\),
and if moreover \(\operatorname{Adm}(h,a)\) and \(\operatorname{Adm}(h',a')\), then
\(\tau(\Gamma(h,a))=\tau(\Gamma(h',a'))\); the target quotient \(q'\) is not used
in this definition.  Then
\begin{align*}
  \operatorname{CFL}(q,\alpha,\operatorname{Adm},\Gamma,q',\tau)
  &\iff
  \mathcal O_{\mathrm{adm}}=\emptyset
  \wedge
  \mathcal O_{\mathrm{cf}}=\emptyset\\
  &\iff
  \begin{gathered}
    \exists\,\overline{\operatorname{Adm}}:Q\times\bar A\to\mathrm{Prop},\
    \exists\,\bar\tau:Q\times\bar A\to\mathrm{Bool}
    \ \text{such that}\\
    \overline{\operatorname{Adm}}(q(h),\alpha(a))
      \leftrightarrow\operatorname{Adm}(h,a)
    \ \text{for all }(h,a)\ \text{and}\\
    \bar\tau(q(h),\alpha(a))=\tau(\Gamma(h,a))
    \ \text{for all admissible }(h,a).
  \end{gathered}
\end{align*}
\end{theorem}
\begin{proof}
Assume the context quotient \(q\), intervention abstraction
\(\alpha\), admissibility predicate \(\operatorname{Adm}\), route
\(\Gamma\), target quotient \(q'\), and target truth \(\tau\).  If a
counterfactual is legitimate, admissibility is constant on
\((q,\alpha)\)-classes, so
\(\mathcal O_{\mathrm{adm}}=\emptyset\).  Its target truth is also
constant on admissible \((q,\alpha)\)-classes, so
\(\mathcal O_{\mathrm{cf}}=\emptyset\).

Conversely, assume both obstruction sets are empty.  Define
\[
  \overline{\operatorname{Adm}}(q(h),\alpha(a)):=\operatorname{Adm}(h,a)
\]
and, for admissible pairs,
\[
  \bar\tau(q(h),\alpha(a)):=\tau(\Gamma(h,a)).
\]
The emptiness of \(\mathcal O_{\mathrm{adm}}\) makes
\(\overline{\operatorname{Adm}}\) well-defined, it is set false off the image of \(q\times\alpha\), \(\bar\tau\) is set to false at classes with no admissible representative, and the emptiness of
\(\mathcal O_{\mathrm{cf}}\) makes \(\bar\tau\) independent of the
admissible representative.  Conversely again, given
\(\overline{\operatorname{Adm}}\) and \(\bar\tau\), two pairs with the
same classes have the same admissibility and, when admissible, the
same target truth, so both obstruction sets are empty.  Thus a pair in either obstruction set is exactly a witness that
admissibility or admissible target truth fails to descend to
\(Q\times\bar A\).
\end{proof}

\paragraph{Nonclaims.} The theorem does not determine the truth value of counterfactuals and
does not apply to arbitrary undeclared possible worlds.  The
possible-worlds semantics in the background is the classical reference
\citep{Lewis1973Counterfactuals}.

\subsection*{Layer instantiations}\label{layer-inst:F31}

\begingroup\small
\begin{description}
\item[\emph{{MIL-STD-810} environmental-qualification
bench testing (engineering qualification).}] \textit{Analogy.}
US Department of Defense MIL-STD-810 prescribes a
catalogue of environmental qualification test methods
(thermal cycling, vibration, shock, humidity, salt fog,
etc.).  The context quotient collapses specimens to
their design baseline (model number, build
configuration); the intervention abstraction collapses
test profiles to the standardised method class;
admissibility requires the profile to lie within the
standard envelope; the counterfactual route is the
laboratory test execution under the prescribed
profile; the target truth records pass/fail according
to the acceptance criterion.  The legitimacy of a
qualification-test counterfactual --- ``this specimen
would pass under method X'' --- is the \Fref{F31} statement:
admissibility and pass/fail descend through
(design-baseline, test-method) classes
\citep{MILSTD810H2019}.

\end{description}
\endgroup


\subsection*{Status, Records, Coherence at a glance}\label{sec:status-records-coherence:summary}\addcontentsline{toc}{subsection}{Status, Records, Coherence at a glance}
\begin{center}
  \centering
  \footnotesize
  \setlength{\tabcolsep}{4pt}
  \medskip\noindent\begin{minipage}{\textwidth}\centering
  \textbf{Status}\par\smallskip
  \begin{tabularx}{\textwidth}{@{}
      >{\raggedright\arraybackslash}p{0.24\textwidth}
      >{\raggedright\arraybackslash}X
      >{\raggedright\arraybackslash}p{0.10\textwidth}@{}}
    \toprule
    Name & One-line claim & Substrate cite \\
    \midrule
    Objectivity as Common Quotient (\Fref{F17}) &
    For surjective access and public quotients, when the public quotient is the common quotient of the accesses (it factors through each, and every map factoring through each factors through it), a readout factors through every access exactly when the public quotient determines it, and exactly when no access identifies two histories with different readout values.  &
    \textemdash \\
    Lawhood-Descent Theorem (\Fref{F18}) &
    For a surjective access quotient, a candidate law is lawful exactly
    when it factors through it.  &
    \textemdash \\
    Interface-Mediation Theorem (\Fref{F22}) &
    For surjective source quotient and boundary readout, an update is mediated iff its target class descends, uniquely, through (source
    class, boundary value); failure splits into boundary leakage and internal
    insufficiency.  &
    \textemdash \\
    Probability as Stable Unresolved Fiber (\Fref{F23}) &
    For a surjective quotient, a readout either descends or is unresolved, never both; stable
    probability is unresolvedness together with the declared ledger
    conditions Stable, which do not involve the weights and hold trivially for empty supports and a one-point target set; on a finite carrier, with the weight assignment given on fiber events by real fiber
    weights, positive and summing to one on each fiber, descent holds iff every fiber chance is a point mass (value 1 at one readout value).  &
    \textemdash \\
    \bottomrule
  \end{tabularx}
  \end{minipage}\par

  \medskip\noindent\begin{minipage}{\textwidth}\centering
  \textbf{Records}\par\smallskip
  \begin{tabularx}{\textwidth}{@{}
      >{\raggedright\arraybackslash}p{0.24\textwidth}
      >{\raggedright\arraybackslash}X
      >{\raggedright\arraybackslash}p{0.10\textwidth}@{}}
    \toprule
    Name & One-line claim & Substrate cite \\
    \midrule
    Object-Persistence Theorem (\Fref{F19}) &
    For a surjective source access quotient \(q_i\), an object persists exactly when its split obstruction is empty; exactly one of four comparison statuses holds; the pairing \(h\mapsto(q_i(h),q_j(\gamma h))\) retains \(q_i\) and carries \(q_j\circ\gamma\); viability, relations and holonomy hold exactly when their violation sets are empty.  &
    \textemdash \\
    Fact--Record Theorem (\Fref{F20}) &
    For surjective record signatures \(s\) and \(\rho_s\), an event is recorded, a record is stable, and a fact value is stable
    exactly when the corresponding obstruction is empty; the record
    comparison has exactly one of four statuses.  &
    \textemdash \\
    Reflexive-Nonclosure Theorem (\Fref{F21}) &
    Under rules 2--4 of a same-layer audit policy, each of which is needed, a complete self-soundness claim is never accepted, and, when the meta layer accepts no self-sound meta-claim, a meta-layer certificate comes from an accepted
    meta-claim that is not itself self-sound.  &
    \textemdash \\
    Causality Theorem (\Fref{F25}) &
    For surjective context and intervention quotients, intervention
    transport descends to (context class, intervention
    class) exactly when its obstruction is empty, and then uniquely;
    under a descended transport every stable causal effect has a witness
    at one raw context with two raw interventions.  &
    \textemdash \\
    \bottomrule
  \end{tabularx}
  \end{minipage}\par

  \medskip\noindent\begin{minipage}{\textwidth}\centering
  \textbf{Coherence}\par\smallskip
  \begin{tabularx}{\textwidth}{@{}
      >{\raggedright\arraybackslash}p{0.24\textwidth}
      >{\raggedright\arraybackslash}X
      >{\raggedright\arraybackslash}p{0.10\textwidth}@{}}
    \toprule
    Name & One-line claim & Substrate cite \\
    \midrule
    Moduli--Contingency Theorem (\Fref{F26}) &
    When the moduli map identifies exactly the equivalent lawful closures and every modulus is attained by a lawful closure, a readout agrees
    on lawful closures with a map on the moduli set
    exactly when it is presentation-independent on lawful closures; it is then either
    necessary or contingent, never both.  &
    \textemdash \\
    Conservation Theorem (\Fref{F27}) &
    Conserved quantities are exactly readouts that factor through
    route-orbit quotients.  &
    \textemdash \\
    Composability Theorem (\Fref{F28}) &
    For a surjective part/interface quotient \(\kappa\), the composite outcome, cross residual and admissibility each descend through \(\kappa\), uniquely, exactly when their split-pair sets are empty; exact composability is emptiness of all three.  &
    \textemdash \\
    Presentation-Invariance Theorem (\Fref{F29}) &
    For a quotient with a section, a readout constant on presentation classes is exactly one that factors through the quotient, agrees with its value at the section representative, or is constant along any relation generating the kernel; relation invariance under any tuple quotient is descent; presentation invariance does not imply equivariance.  &
    \textemdash \\
    Explanation-as-Compression Theorem (\Fref{F30}) &
    For surjective base and explanation readouts, an explanation is exact
    exactly when its source and target obstructions are empty, and is a
    strict compression exactly when it is determined by the base and merges
    two histories the base separates; a dynamic explanation makes the later
    target descend through the explanation.  &
    \textemdash \\
    Counterfactual-Legitimacy Theorem (\Fref{F31}) &
    A counterfactual is legitimate exactly when both obstruction sets are empty, exactly when admissibility, and target truth on admissible pairs, descend to (context class, intervention class).  &
    \textemdash \\
    \bottomrule
  \end{tabularx}
  \end{minipage}\par

  \captionof{table}{Status, Records, Coherence at a glance: the fourteen laws of \Cref{sec:status-records-coherence}, grouped into reading groups.}
  \label{tab:axis-status-records-coherence-summary}
\end{center}

\FloatBarrier
          % sec:status-records-coherence
\section{Self-Reference and Limits}\label{sec:self-reference-limits}

The Self-Reference and Limits chapter collects the laws that mark where
closure formation reaches a boundary.  Some boundaries are horizons:
boundary data can determine a declared target even when the ambient
carrier is larger.  Others concern information loss, singular limits,
locality, topology and completeness: what is visible, what is glued
globally, what survives a limit, and what cannot be established from
inside the same closure without additional audit structure.  The
self-audit law, Reflexive Nonclosure (\Fref{F21}), is stated in the
Status, Records, Coherence chapter because its data are claim records;
here the law that treats self-reference directly is the Closure
Completeness Boundary (\Fref{F52}).  In the running example of
\Cref{sec:intro}, the identification of \(1\) and \(2\) made by \(q\)
loses information that \(F\) records; in the language of the
Information Loss Normal Form (\Fref{F34}), with the identity as
transport, \(q\) as recovery quotient and \(F\) as source information,
\((1,2)\) lies in the loss obstruction.

The eighteen laws proceed from local sufficiency to completeness.
Horizon Sufficiency / Holographic Compression (\Fref{F33}) opens the
chapter: a target depends on a history only through its boundary
readout and exterior class exactly when it descends through their
pairing.  It is followed by Information Loss, Scale Descent /
Renormalization, Singular Limit Completion, Complementarity as
Non-Joint Access, Arrow from Record Monotonicity and Entropy as Fiber
Volume.  The middle of the chapter treats Irreducibility / No Short
Closed Description, Continuum Emergence, Criticality / Phase Boundary,
Locality / Domain of Dependence, Law--State Split and Fine-Tuning as
Small Selector Region.  The last part gives Topology as Global Gluing
Invariant, Common-Source / Nonlocal Correlation, Vacuum / Background
Selection, Unification as Common Refinement and the Closure
Completeness Boundary, which is distinct from the closure--reality
boundary of \Cref{sec:anchor:closure-reality-boundary}.  The chapter
depends only on \Cref{sec:apparatus}, and the summary table at its end
groups the laws into limits, emergence and invariants.

\subsection{Horizon Sufficiency / Holographic Compression [\texorpdfstring{\Fref{F33}}{F33}]}\label{sec:self-reference-limits:horizon-sufficiency}

\begin{definition}[Horizon-bounded support]\label{def:horizon-bounded-support}
Let \(\beta:I\to B\) be a boundary readout and \(\varepsilon:E\to Q\) an
exterior classifier on an interior--exterior decomposition \(I\times E\). A
target readout \(t:I\times E\to T\) is \emph{horizon-bounded} for
\((\beta,\varepsilon)\) when \(t(i,e)=t(i',e')\) whenever \(\beta(i)=\beta(i')\)
and \(\varepsilon(e)=\varepsilon(e')\); the pair \((\beta,\varepsilon)\) is its
horizon record.
\end{definition}

\begin{theorem}[Horizon Sufficiency / Holographic Compression]\label{thm:F33}
Let \(I\) and \(E\) be sets, the interior and exterior factors (a set \(H\) of histories with a decomposition \(H\cong I\times E\) is identified with \(I\times E\)), and define
\[
  \pi : I \times E \longrightarrow B \times Q,
  \qquad
  \pi(i,e)=(\beta(i),\varepsilon(e)),
\]
where \(\beta:I\to B\) is the boundary readout and
\(\varepsilon:E\to Q\) is the exterior classifier. Assume \(\pi\) is
surjective, and let \(t:I\times E\to T\) be a declared target readout.
Then
\begin{align*}
  t\ \text{horizon-bounded for }(\beta,\varepsilon)
  &\iff
  t \text{ descends through }\pi\\
  &\iff
  \forall (i,e),(i',e'),\
    \pi(i,e)=\pi(i',e')\Rightarrow t(i,e)=t(i',e')\\
  &\iff
  \exists\,\bar t:B\times Q\to T\ \text{ such that }\
    \bar t\circ\pi=t .
\end{align*}
Here horizon-boundedness is that of \Cref{def:horizon-bounded-support}. Call \(t\) \emph{horizon-sufficient}
when these hold. For a family \((t_\sigma)_{\sigma\in\Sigma}\), the bundled
readout \(x\mapsto(t_\sigma(x))_{\sigma}\) is horizon-sufficient if and only if
every \(t_\sigma\) is.
\end{theorem}
\begin{proof}
By \Cref{def:horizon-bounded-support}, \(t\) is horizon-bounded exactly
when it is constant on the fibers of \(\pi\).  Since \(\pi\) is
surjective, the descent criterion of \Cref{sec:apparatus} gives the
equivalence with descent through \(\pi\) and with the existence of
\(\bar t\).  Two points have the same bundled value exactly when they
have the same value under every \(t_\sigma\), so the bundled readout is
constant on the fibers of \(\pi\) exactly when every \(t_\sigma\) is.
\end{proof}

The motivating substrate-physics analogy is the holographic principle, in which
the bulk information of a region is encoded on its boundary
\citep{Maldacena1997LargeN,Bousso2002Holographic}.  The
layer-agnostic statement here is the descent-through-boundary form
of that pattern; it does not assert the AdS/CFT correspondence or
any particular gravitational realization.

\subsection*{Layer instantiations}\label{layer-inst:F33}

\begingroup\small
\begin{description}
\item[\emph{Boundary-value problem for a linear elliptic
{PDE} (mathematical analysis).}] \textit{Special case.}
For a fixed bounded Lipschitz domain \(\Omega\subseteq\mathbb R^n\), the
first factor \(I\) of the decomposition (its name in \Cref{thm:F33} is only a slot name: here the first factor carries the boundary data and the second factor carries the bulk data \(f\) and \(A\)) is the set of admissible Dirichlet data \(g\); the boundary readout records \(g\) in a fixed
function space (e.g., \(H^{1/2}(\partial\Omega)\)).
The second factor \(E\) parameterises sources \(f\in H^{-1}(\Omega)\) and bounded measurable, uniformly elliptic
divergence-form coefficient fields; the exterior classifier records the
source-and-coefficient data in their function-space
equivalence class.  The target readout is the entire
weak solution \(u\in H^1(\Omega)\).  By the Lax-Milgram
theorem, applied, after lifting \(g\) to \(H^1(\Omega)\), to the bilinear form
\(\int_\Omega A\nabla u\cdot\nabla v\), coercive on \(H^1_0(\Omega)\) by the Poincar\'e inequality, the target descends through the product
quotient: same boundary data plus same
source-coefficient class yield the same solution.  For Neumann data over a connected \(\Omega\), put
\(g\in H^{-1/2}(\partial\Omega)\) in the first factor \(I\) and
\(f\in(H^1(\Omega))'\), with the coefficient field, in the second factor \(E\),
and let the target be the weak solution modulo constants when
\(\langle f,1\rangle+\langle g,1\rangle_{\partial\Omega}=0\) and a marker
\(\bot\) otherwise. The compatibility test and, on compatible data,
coercivity on \(H^1(\Omega)/\mathbb R\) make this target a function of
\((g,f,A)\), so it again descends through the product quotient
\citep{Evans2010PartialDifferential}.

\item[\emph{Boundary representation in solid {CAD}
modelling (engineering).}] \textit{Analogy.}
A 3D rigid solid \(B\subset\mathbb R^3\) is represented
by its boundary surface together with adjacency and
orientation data: a B-rep records faces, edges, and
vertices with incidence relations.  The interior is the
open solid; the boundary readout records the B-rep with
face-orientation conventions; the exterior classifier
records the combinatorial topology (genus, handles,
through-holes).  Geometric and topological readouts ---
volume via the divergence theorem, mass distribution,
collision detection, boolean operations --- descend
through the product quotient
\citep{Mortenson2006GeometricModeling}.

\end{description}
\endgroup


\subsection{Information Loss Normal Form [\texorpdfstring{\Fref{F34}}{F34}]}\label{sec:self-reference-limits:information-loss}

\paragraph{Reading the statement.} The symbols are those of \Cref{thm:F34} below.
\(\operatorname{Recoverable}(\sigma,\tau,\rho)\) says that the source
information can be read back after transport: \(\sigma(h)\) depends only
on the later recovery class \(\rho(\tau h)\), that is, \(\sigma\)
descends through \(\rho\circ\tau\).

\begin{theorem}[Information Loss Normal Form]\label{thm:F34}
Let \(H_0,H_1\) be source and later carriers with transport
\(\tau:H_0\to H_1\). Let \(\sigma:H_0\to I\) be source information,
and \(\rho:H_1\to R\) a recovery quotient. Assume that
\(\rho\circ\tau\) is surjective onto \(R\). Define
\[
  \mathcal O_\rho
  =\{(h,h'):\rho(\tau h)=\rho(\tau h')\ \text{and}\
      \sigma(h)\ne\sigma(h')\}.
\]
Then
\begin{align*}
  \operatorname{Recoverable}(\sigma,\tau,\rho)
  &\iff \mathcal O_\rho=\emptyset\\
  &\iff
  \exists\,\bar\sigma:R\to I\ \text{ such that }\
    \bar\sigma\circ\rho\circ\tau=\sigma .
\end{align*}
The later quotient preserves the declared source information exactly
when every identification made by \(\rho\circ\tau\) is invisible to
\(\sigma\).
\end{theorem}
\begin{proof}
\(\operatorname{Recoverable}(\sigma,\tau,\rho)\) says that \(\sigma\)
descends through \(\rho\circ\tau\), and \(\mathcal O_\rho\) is the
split-pair obstruction of \(\rho\circ\tau\) and \(\sigma\).  Since
\(\rho\circ\tau\) is surjective onto \(R\), both equivalences are the
descent criterion of \Cref{sec:apparatus}.  A pair in
\(\mathcal O_\rho\) has the same recovered class but different source
information, which is precisely loss across a recovery class.
\end{proof}

A physical analogue is the black-hole information paradox, posed by Hawking \citep{Hawking1976Breakdown} after his thermodynamic analysis of black holes \citep{Hawking1976BHThermodynamics}, where the loss of information is argued
from semiclassical gravity; the statement above records only the descent
obstruction of a declared horizon quotient.

\subsection*{Layer instantiations}\label{layer-inst:F34}

\begingroup\small
\begin{description}
\item[\emph{Huffman entropy coding (information theory).}] \textit{Special case.}
The source carrier is the set of source-symbol sequences
over an alphabet \(\Sigma\); the source information is
the original symbol sequence.  The transport encodes
into a prefix-free binary codeword sequence; the later carrier is the set \(\tau(H_0)\) of encoded streams, and the recovery quotient is the prefix decoder on it.  Distinct
source sequences produce distinct codeword streams; the
obstruction set is empty and the source is fully
recoverable.  Compression achieves the entropy bound to
within one bit per symbol
\citep{Huffman1952Method}.

\item[\emph{Reed-Solomon forward error correction
(coding theory).}] \textit{Special case.}
For a Reed-Solomon \((n,k)\) code over a finite field
\(\mathbb F_q\), the source carrier is the set of pairs
\((m,\varepsilon)\) of a \(k\)-symbol message and an error
vector of Hamming weight at most
\(t=\lfloor(n-k)/2\rfloor\), with \(\sigma(m,\varepsilon)=m\).
Let the later carrier be the union of radius-\(t\) Hamming
balls around the codewords. The transport sends
\((m,\varepsilon)\) to \(c(m)+\varepsilon\), where \(c(m)\)
evaluates the message polynomial at \(n\) distinct field
points. Let \(\rho\) be bounded-distance decoding from this
carrier to the message set, using syndrome decoding and error
correction. Disjoint radius-\(t\) balls make
\(\rho(c(m)+\varepsilon)=m\), so the obstruction is empty;
taking \(\varepsilon=0\) for each message shows that
\(\rho\circ\tau\) is surjective, and the source message is
recoverable; Reed-Solomon codes underlie CD/DVD, QR
codes, and deep-space communication
\citep{ReedSolomon1960Codes}.

\item[\emph{Cryptographic one-way hash function
(cryptography).}] \textit{Special case.}
For a cryptographic hash function
\(H:\{0,1\}^{<2^{64}}\to\{0,1\}^{256}\) (SHA-256), the
source information is the original bitstring.  The
transport maps bitstrings to fixed-length digests;
the recovery quotient is the identity on the set \(H(\{0,1\}^{<2^{64}})\) of attained digests.
Since the domain has more than \(2^{256}\) elements, distinct source
bitstrings collide to the same digest by pigeonhole, so the obstruction set is non-empty.  The
cryptographic hash is the canonical non-recoverable
instance; its cryptographic use rests on collisions
being hard to find, not absent
\citep{NIST2015FIPS180}.

\item[\emph{Database write-ahead logging and {ARIES}
recovery (database systems).}] \textit{Analogy.}
The source information is the committed-state effects
of database transactions.  The transport writes log
records (before-image, after-image, LSN, transaction
ID) to stable storage before any modified page is
written.  The recovery quotient is the ARIES
three-pass algorithm (analysis, redo, undo) applied
to the on-disk log and data pages after a crash.
Within the failure model, the obstruction set is
empty and the committed-state effects are recoverable
\citep{Mohan1992ARIES}.

\item[\emph{Magnetic-resonance image reconstruction
(medical imaging).}] \textit{Analogy.}
The source information is the tissue proton-density
map, supported within the field of view and represented at the
declared resolution.  The transport
samples k-space (the Fourier dual of physical space)
under controlled magnetic-field gradient sequences;
the later carrier is the set of sampled k-space
datasets.  The recovery quotient is the inverse
Fourier transform, in clinical practice combined with
parallel-imaging unfolding and partial-Fourier
filling.  When the k-space sampling step is at most the reciprocal of the
object's spatial extent, the obstruction set is empty and the
proton-density map is recoverable up to the resolution set by the
sampled k-space extent; with a coarser step, aliasing folds the
object onto itself and makes the loss illegitimate
\citep{Lauterbur1973MRI}.
\end{description}
\endgroup

\subsection{Scale Descent / Renormalization [\texorpdfstring{\Fref{F35}}{F35}]}\label{sec:self-reference-limits:scale-descent}

\begin{theorem}[Scale Descent / Renormalization]\label{thm:F35}
Let \(H\) be a carrier with consecutive scale quotients
\(q_n:H\to Q_n\) and \(q_{n+1}:H\to Q_{n+1}\), with \(q_n\) surjective. Let
\(c:Q_n\to Q_{n+1}\) be a surjective coarsening with
\[
  c\circ q_n=q_{n+1},
\]
and let \(t:H\to T\) be a target law. Put
\[
  \mathcal O_n(t)=
  \{(h,h'):q_n(h)=q_n(h')\ \text{and}\ t(h)\ne t(h')\}.
\]
Then
\begin{align*}
  t \text{ descends to } Q_n
  &\iff \mathcal O_n(t)=\emptyset\\
  &\iff
  \exists\,\bar t_n:Q_n\to T\ \text{ such that }\
    \bar t_n\circ q_n=t .
\end{align*}
When \(t\) descends to \(Q_n\), the map \(\bar t_n\) is unique, and a
renormalized law at scale \(Q_{n+1}\), a map \(\bar t_{n+1}:Q_{n+1}\to T\) with
\(\bar t_{n+1}\circ c=\bar t_n\), exists exactly when \(\bar t_n\) is constant
on the fibers of \(c\), and then \(t=\bar t_{n+1}\circ q_{n+1}\), and \(\bar t_{n+1}\) is unique since \(c\) is surjective. Consequently
\(t\) descends to \(Q_{n+1}\) if and only if it descends to \(Q_n\) and
\(\bar t_n\) is constant on the fibers of \(c\).
\end{theorem}
\begin{proof}
Since \(q_n\) is surjective, the first two equivalences and the
uniqueness of \(\bar t_n\) are the descent criterion of
\Cref{sec:apparatus} for \(q_n\) and \(t\). For the next scale, if \(\bar t_{n+1}\circ c=\bar t_n\),
then \(c(x)=c(x')\) gives \(\bar t_n(x)=\bar t_n(x')\), so \(\bar t_n\) is
constant on the fibers of \(c\). Conversely, if \(\bar t_n\) is constant on
the fibers of \(c\), then \(\bar t_{n+1}(c(x)):=\bar t_n(x)\) is independent of
the choice of \(x\), and surjectivity of \(c\) defines it on all of
\(Q_{n+1}\). Then \(t=\bar t_n\circ q_n=\bar t_{n+1}\circ c\circ q_n=\bar t_{n+1}\circ q_{n+1}\).
Conversely, if \(t=s\circ q_{n+1}=s\circ c\circ q_n\), then \(t\) descends to
\(Q_n\), uniqueness gives \(\bar t_n=s\circ c\), and \(\bar t_n\) is constant on
the fibers of \(c\). If two renormalized laws \(\bar t_{n+1},\bar t'_{n+1}\)
satisfy \(\bar t_{n+1}\circ c=\bar t_n=\bar t'_{n+1}\circ c\), then they agree on
the image of \(c\), which is all of \(Q_{n+1}\). Thus scale descent is quotient
descent iterated along the tower.
\end{proof}

The Wilsonian renormalization group of statistical and quantum field
theory \citep{Wilson1975RG} is the standard substrate analogy, in
which scale-by-scale quotient descent is the central calculational
device.

\subsection*{Layer instantiations}\label{layer-inst:F35}

\begingroup\small
\begin{description}
\item[\emph{Hydrodynamic limit via Chapman-Enskog
expansion (kinetic theory).}] \textit{Analogy.}
For a rarefied gas, the carrier is the space of
phase-space distribution functions \(f(x,v,t)\).  The
scale quotient at the kinetic scale records the
one-particle distribution at Knudsen number
\(\mathrm{Kn}\sim O(1)\); the coarsening sends \(f\)
to its hydrodynamic moments (density, momentum
density, energy density) in the limit
\(\mathrm{Kn}\to 0\).  The target law at the kinetic
scale is the Boltzmann equation; the renormalised law
at the hydrodynamic scale is the Navier-Stokes (or
Euler) system, derived by the Chapman-Enskog
asymptotic expansion.  Scale descent succeeds in
\(\mathrm{Kn}\ll 1\) regimes and fails near shocks
and boundary layers \citep{Cercignani1988Boltzmann}.

\item[\emph{Burt-Adelson Laplacian image pyramid
(computer vision and signal processing).}] \textit{Special case.}
For a digital image, the carrier is the space of
full-resolution images.  The scale quotient
sends an image to its Gaussian-pyramid level at
resolution \(2^{-n}\) (Gaussian filter followed by
downsampling), corestricted to the levels it attains; the coarsening sends level \(n\) to
level \(n+1\) by repeating the filter-and-downsample
step.  The band-pass Laplacian layer
\(L_n=G_n-\mathrm{up}(G_{n+1})\) is the target law at
level \(n\) and descends through the scale quotient.
The Laplacian pyramid is a complete representation:
the original image \(G_0\) is recovered exactly from the coarsest level \(G_N\) and the residues \(L_0,\dots,L_{N-1}\) by \(G_n=L_n+\mathrm{up}(G_{n+1})\), \(n=N-1,\dots,0\)
\citep{BurtAdelson1983LaplacianPyramid}.

\end{description}
\endgroup


\subsection{Singular Limit Completion [\texorpdfstring{\Fref{F36}}{F36}]}\label{sec:self-reference-limits:singular-limit}

\paragraph{Reading the statement.} In \Cref{thm:F36} below, a
limiting process \(p\) has a completed object \(\ell(p)\), and
\(\operatorname{Completed}(u,\ell)\) says that its target-limit value
depends only on that object.  \(\operatorname{PathDependent}(u,\ell)\)
says that two processes with the same completed object have different
target-limit values.  \(\operatorname{Removable}(j,\ell,\sigma,p_0)\)
says that completion removes the singularity at \(p_0\): the completed
object \(\ell(p_0)\) is not the image of a current object under \(j\),
and the signature-limit readout \(\sigma\) descends through \(\ell\).

\begin{theorem}[Singular Limit Completion]\label{thm:F36}
Let \(P\) be a set of limiting processes, let \(\ell:P\to\widehat Q\)
be a surjective limit quotient, let \(j:Q\to\widehat Q\) be a
completion map, let \(u:P\to\widehat T\) be the target-limit readout,
let \(\sigma:P\to\widehat S\) be a declared signature-limit readout,
let \(u^{\mathrm{ren}}:P\to\widehat T^{\mathrm{ren}}\) be a declared
renormalized target-limit readout, and let \(p_0\in P\) be a selected
process.  Define the singular target obstruction
\[
  \mathcal O_{\ell}(u)
  =\{(p_1,p_2)\in P^2:
      \ell(p_1)=\ell(p_2)\ \text{and}\ u(p_1)\ne u(p_2)\}
\]
and the corresponding signature obstruction
\(\mathcal O_{\ell}(\sigma)\).  Write
\[
  \rho_{u,\ell}:P\to\widehat Q\times\widehat T,
  \qquad
  \rho_{u,\ell}(p):=(\ell(p),u(p)),
\]
for the joint \((\ell,u)\) refinement, and let
\(\operatorname{Current}_{j}(\ell,p_0)\) abbreviate
\(\exists c\in Q,\ j(c)=\ell(p_0)\), and define
\[
  \operatorname{Removable}(j,\ell,\sigma,p_0)
  :\Longleftrightarrow
  \neg\operatorname{Current}_{j}(\ell,p_0)
  \wedge
  \operatorname{Completed}(\sigma,\ell).
\]
Here \(\operatorname{Completed}(v,\kappa)\), for a readout \(v\) and a map
\(\kappa\) on \(P\), says that \(v(p)\) depends only on \(\kappa(p)\),
\(\operatorname{Refines}(\kappa,\ell)\) says that \(\kappa(p_1)=\kappa(p_2)\)
implies \(\ell(p_1)=\ell(p_2)\), and
\(\operatorname{PathDependent}(u,\ell)\) says that
\(\mathcal O_\ell(u)\ne\emptyset\).
Then
\begin{align*}
  \operatorname{Completed}(u,\ell)
  &\iff \mathcal O_{\ell}(u)=\emptyset,\\
  \operatorname{Completed}(u,\ell)
  &\iff
  \exists\,\bar u:\widehat Q\to\widehat T,\ \bar u\circ\ell=u,\\
  \operatorname{Completed}(\sigma,\ell)
  &\iff \mathcal O_{\ell}(\sigma)=\emptyset,\\
  \operatorname{Removable}(j,\ell,\sigma,p_0)
  &\iff \neg\operatorname{Current}_{j}(\ell,p_0)\wedge\mathcal O_{\ell}(\sigma)=\emptyset,\\
  \operatorname{PathDependent}(u,\ell)
  &\iff \neg\operatorname{Completed}(u,\ell),\\
  &\operatorname{Refines}(\rho_{u,\ell},\ell),\\
  &\operatorname{Completed}(u,\rho_{u,\ell}),\\
  \operatorname{Completed}(u^{\mathrm{ren}},\ell)
  &\iff
  \mathcal O_{\ell}(u^{\mathrm{ren}})=\emptyset.
\end{align*}
Equivalently, completion is descent of \(u\) through \(\ell\); the
canonical path-memory repair \(\rho_{u,\ell}\) is a refinement of
\(\ell\) that completes \(u\); a renormalized target is tested by the
same descent criterion as its raw target, and if \(u^{\mathrm{ren}}=R\circ u\) for a
map \(R\), then \(\operatorname{Completed}(u,\ell)\) implies
\(\operatorname{Completed}(u^{\mathrm{ren}},\ell)\), with the converse when \(R\) is
injective.
\end{theorem}
\begin{proof}
We prove the eight conjuncts. Parts (4) and (5) prove the PathDependent and Removable lines, which the statement lists in the opposite order; part (8) also proves the clause on \(u^{\mathrm{ren}}=R\circ u\).

\emph{(1)--(3) Completion and obstructions.}
Since \(\ell\) is surjective, the descent criterion of
\Cref{sec:apparatus}, applied to \(\ell\) and \(u\), gives (1) and (2),
and applied to \(\ell\) and \(\sigma\) gives (3).

\emph{(4) Path-dependent iff not completed.}
Path dependence is witnessed by a split pair:
\[
  \operatorname{PathDependent}(u,\ell)
  \iff
  \exists(p_1,p_2)\in\mathcal O_{\ell}(u).
\]
By (1), \(\mathcal O_{\ell}(u)\ne\emptyset\) is equivalent to
\(\neg\operatorname{Completed}(u,\ell)\).

\emph{(5) Removable.}
By definition \(\operatorname{Removable}(j,\ell,\sigma,p_0)\) is
\(\neg\operatorname{Current}_{j}(\ell,p_0)\wedge\operatorname{Completed}(\sigma,\ell)\),
and (3) replaces the second conjunct by \(\mathcal O_{\ell}(\sigma)=\emptyset\).

\emph{(6) The path repair refines \(\ell\).}
The canonical path-memory repair is
\[
  \rho_{u,\ell}(p)=(\ell(p),u(p)).
\]
Projection to the first coordinate recovers the original limit:
\[
  \operatorname{pr}_{\widehat Q}\circ\rho_{u,\ell}=\ell.
\]
Thus each \(\rho_{u,\ell}\)-fiber lies inside an \(\ell\)-fiber, so
\(\rho_{u,\ell}\) is a source refinement of \(\ell\).

\emph{(7) After the path repair, the target is completed.}
Projection to the second coordinate recovers the target-limit readout:
\[
  \operatorname{pr}_{\widehat T}\circ\rho_{u,\ell}=u.
\]
Hence \(u\) factors through the repaired quotient by the explicit
factor \(\operatorname{pr}_{\widehat T}\), and therefore
\[
  \operatorname{Completed}(u,\rho_{u,\ell}).
\]

\emph{(8) Renormalized completed iff its obstruction is empty.}
The descent criterion applied to \(\ell\) and \(u^{\mathrm{ren}}\) gives
\[
  \operatorname{Completed}(u^{\mathrm{ren}},\ell)
  \iff
  \mathcal O_{\ell}(u^{\mathrm{ren}})=\emptyset,
\]
so a renormalized target is subject to the same completion criterion
as a raw target. If \(u^{\mathrm{ren}}=R\circ u\) and \(u=\bar u\circ\ell\), then
\(u^{\mathrm{ren}}=(R\circ\bar u)\circ\ell\); conversely, assume \(R\) is injective
and \(\operatorname{Completed}(u^{\mathrm{ren}},\ell)\). For \(\ell(p_1)=\ell(p_2)\),
this completion gives \(R(u(p_1))=R(u(p_2))\); injectivity gives
\(u(p_1)=u(p_2)\). Thus \(\mathcal O_\ell(u)=\emptyset\), and the first
equivalence gives \(\operatorname{Completed}(u,\ell)\).
\end{proof}

\subsection*{Layer instantiations}\label{layer-inst:F36}

\begingroup\small
\begin{description}
\item[\emph{Riemann removable singularity in complex
analysis (complex analysis).}] \textit{Special case.}
For a function \(f\) holomorphic on a punctured disk,
the set of limiting processes is the family of
approach paths to the puncture along which \(f\) has a
limit in the Riemann sphere; the limit quotient
\(\ell\) is constant, since every path approaches the same puncture; the
target readout is the path-limit value of \(f\).  By
Riemann's removable-singularity theorem, if \(f\) is
bounded on the punctured disk then the limit is
independent of approach path and the obstruction set
is empty.  At a pole
\(|f|\to\infty\) along every approach path, so the limit is
path-independent on the Riemann sphere.  An essential
singularity can give the path-dependent case: \(e^{1/z}\) tends
to \(\infty\) along the positive real axis and to \(0\) along the
negative real axis
\citep{Ahlfors1979ComplexAnalysis}.

\item[\emph{L'H\^opital's rule for indeterminate
forms (real analysis).}] \textit{Analogy.}
For functions \(f,g\) with \(f(a)=g(a)=0\), the
literal substitution \(\lim_{x\to a}f(x)/g(x)\)
yields the indeterminate form \(0/0\), which by
itself does not descend through the limit quotient.
L'H\^opital's rule provides the completion: if \(f\)
and \(g\) are differentiable with \(g'(a)\ne 0\),
the limit equals \(f'(a)/g'(a)\).  The path-memory
repair records the order of vanishing of \(f\) and
\(g\) at \(a\), completing the limit through the
derivative-ratio readout
\citep{Cauchy1821CoursAnalyse}.

\item[\emph{Riemann integral as limit of Riemann sums
(classical analysis).}] \textit{Special case.}
For a bounded \(f:[a,b]\to\mathbb R\), the set of
limiting processes is the family of sequences of
tagged partitions with mesh tending to zero along
which the Riemann sums converge; the target readout
is the limit of the Riemann sums; the limit quotient \(\ell\) is the constant
map onto a one-point set.  The function is Riemann-integrable precisely when
\(\mathcal O_\ell(u)=\emptyset\) (for bounded \(f\), tags approximating the interval suprema
and infima give sequences converging to the upper and lower
Darboux integrals, which differ otherwise), i.e., when the limit of Riemann sums
is independent of the choice of partition
refinement and tag selection.  The completed map is
the Riemann integral
\(\int_a^b f\,dx\); non-Riemann-integrable functions
(e.g., the Dirichlet function) instantiate the
path-dependent case
\citep{Riemann1854Habilitationsschrift}.
\end{description}
\endgroup


\subsection{Complementarity as Non-Joint Access [\texorpdfstring{\Fref{F37}}{F37}]}\label{sec:self-reference-limits:complementarity}

Position and momentum \citep{Bohr1928Quantum} motivate a restricted choice
of \(\mathcal J\): if an admissible carrier must arise from a single joint
quantum observable with sharp position and sharp momentum as its marginals, no
such carrier exists. The pair of statistical readouts still exists as a map.
Merely requiring a joint probability distribution with the correct marginals
would not give non-joint access, since the product of those marginals is one.
The statement below abstracts non-joint access relative to the declared class
\(\mathcal J\).

\paragraph{Reading the statement.}
The class \(\mathcal J\) and the formedness of each access are declared
data: they fix which joint carriers count as admissible at the stated
scope. \(q_A\perp_{\mathrm{acc}}q_B\) says that each access is lawful
on its own but no admissible joint carrier determines both.

\begin{theorem}[Complementarity as Non-Joint Access]\label{thm:F37}
Let \(H\) be a nonempty carrier with two access quotients
\(q_A:H\to Q_A\) and \(q_B:H\to Q_B\) and target readouts
\(t_A:H\to T_A\), \(t_B:H\to T_B\).  The access \(q_A\) is
\emph{separately admissible} when \(t_A\) descends through \(q_A\)
(formedness of the access is part of the setting and adds no condition), and
likewise for \(q_B\).  Let \(K\) be a declared
joint-carrier type and let \(\mathcal J\) be a declared class of
admissible joint carriers \(j:H\to K\).  A joint access in
\(\mathcal J\) is a carrier \(j\in\mathcal J\) together with maps
\(a:K\to Q_A\) and \(b:K\to Q_B\) satisfying
\[
  a\circ j=q_A,\qquad b\circ j=q_B .
\]
Define
\begin{align*}
  q_A \perp_{\mathrm{acc}} q_B
  \;:\Longleftrightarrow\;&
  q_A,\,q_B\ \text{separately admissible}\\
  &\wedge\ \neg\exists\, j\in\mathcal J,\ \exists\,a:K\to Q_A,\ \exists\,b:K\to Q_B,\
  a\circ j=q_A\ \wedge\ b\circ j=q_B .
\end{align*}
Then \(q_A \perp_{\mathrm{acc}} q_B\) holds if and only if \(q_A\) and \(q_B\) are separately admissible and every
\(j\in\mathcal J\) identifies two histories that one of the accesses
separates: for each \(j\in\mathcal J\) there are \(h,h'\in H\) with
\(j(h)=j(h')\) and (\(q_A(h)\neq q_A(h')\) or \(q_B(h)\neq q_B(h')\)).
Thus complementarity is precisely separate admissibility together with the presence, in every admissible joint
carrier, of a split pair against one of the two accesses. Equivalently, \(q_A\perp_{\mathrm{acc}}q_B\) holds if and only if \(q_A\) and \(q_B\) are separately admissible and no \(j\in\mathcal J\) has
\(\langle q_A,q_B\rangle=g\circ j\) for a map \(g:K\to Q_A\times Q_B\); in
particular, if \(K=Q_A\times Q_B\) and \(\langle q_A,q_B\rangle\in\mathcal J\),
then \(q_A\perp_{\mathrm{acc}}q_B\) fails.
\end{theorem}
\begin{proof}
Fix
\(j\in\mathcal J\). If \(a\circ j=q_A\), then \(j(h)=j(h')\) gives
\(q_A(h)=q_A(h')\). (For the factorization form, take \(g=\langle a,b\rangle\), and
conversely \(a=\mathrm{pr}_1\circ g\), \(b=\mathrm{pr}_2\circ g\); if the pairing
itself lies in \(\mathcal J\), take \(g=\mathrm{id}\).) Conversely, if \(q_A\) is constant on the fibers of
\(j\), set \(a(j(h)):=q_A(h)\) on the image of \(j\) and give \(a\) any value
in \(Q_A\) elsewhere, which is possible since \(H\), and hence \(Q_A\), is
nonempty; then \(a\circ j=q_A\). So a map \(a\) exists exactly when no pair
identified by \(j\) is separated by \(q_A\), and likewise for \(b\) and
\(q_B\). Hence \(j\) is a joint access exactly when it identifies no pair that
\(q_A\) or \(q_B\) separates, and \(\mathcal J\) contains no joint access
exactly when every \(j\in\mathcal J\) identifies such a pair. Together with
separate admissibility this is complementarity. The product map
\(h\mapsto(q_A(h),q_B(h))\) always carries both accesses, and it factors
through every carrier of both accesses; complementarity is
therefore a statement about the declared class \(\mathcal J\), which
decides whether such a carrier counts as an admissible joint access.
\end{proof}

\subsection*{Layer instantiations}\label{layer-inst:F37}

\begingroup\small
\begin{description}
\item[\emph{Time-frequency Gabor uncertainty (signal
processing).}] \textit{Analogy.}
For a square-integrable signal \(s\in L^2(\mathbb R)\),
the access quotient \(q_A\) records its time-marginal
\(|s(t)|^2\) and \(q_B\) records its
frequency-marginal \(|\hat s(f)|^2\).  Within the
admissibility class of Cohen-class bilinear
time-frequency representations satisfying the correct
marginals, the covariance properties of bilinear
distributions, and non-negativity (probability-density
interpretation), no joint quotient exists: the
Wigner-Ville distribution satisfies marginals and
covariance but is not non-negative, and no member of
Cohen's class is jointly compatible.  The access
quotients are \Fref{F37}-complementary
\citep{Cohen1995TimeFrequency}.

\end{description}
\endgroup


\subsection{Arrow from Record Monotonicity [\texorpdfstring{\Fref{F38}}{F38}]}\label{sec:self-reference-limits:arrow}

The phrase ``arrow of time'' is due to \citet{Eddington1928Nature}, who tied it to the increase of entropy; the formulation below takes the arrow instead from monotone growth of records, without committing to a specific physical realization.

\paragraph{Reading the statement.} The symbols are those of \Cref{thm:F38} below.
\(\operatorname{Arrow}(\kappa)\) is defined at the level of histories:
the continuation respects records, and record status does not decrease along
it. The theorem says that this is the same as a map \(\bar\kappa\) on records
along which status does not decrease.

\begin{theorem}[Arrow from Record Monotonicity]\label{thm:F38}
Let \(H_s,H_t\) be source and target history sets with continuation
\(\kappa:H_s\to H_t\). Let \(r_s:H_s\to R_s\) be a surjective record quotient
and \(r_t:H_t\to R_t\) a target record map, and let
\(s_s:R_s\to S\), \(s_t:R_t\to S\) be status readouts into a set \(S\) with a relation \(\le_S\) (no order axioms are used). Define
\[
  \operatorname{Arrow}(\kappa)
  :\Longleftrightarrow
  r_t\circ\kappa\ \text{descends through}\ r_s
  \ \wedge\
  \forall h\in H_s,\ s_s(r_s h)\le_S s_t(r_t(\kappa h)).
\]
Then
\begin{align*}
  \bigl(\forall h,h'\in H_s,\ r_s(h)=r_s(h')\Rightarrow r_t(\kappa h)=r_t(\kappa h')\bigr)
  &\iff
  \exists\,\bar\kappa:R_s\to R_t\ \text{ such that }\
    \bar\kappa\circ r_s=r_t\circ\kappa,\\
  \operatorname{Arrow}(\kappa)
  &\iff
  \exists\,\bar\kappa:R_s\to R_t,\
  \bar\kappa\circ r_s=r_t\circ\kappa\\
  &\hphantom{{}\iff{}}\text{and } s_s(x)\le_S s_t(\bar\kappa x)\
    \forall x\in R_s.
\end{align*}
The arrow is the monotone direction of the descended record transport.
\end{theorem}
\begin{proof}
Since \(r_s\) is surjective, the first equivalence is the descent
criterion of \Cref{sec:apparatus} for \(r_s\) and \(r_t\circ\kappa\),
and the descended map \(\bar\kappa\) is unique.  The arrow condition
therefore splits into descent of \(r_t\circ\kappa\) and a status
inequality.  For \(x=r_s(h)\), \(s_t(\bar\kappa x)=s_t(r_t(\kappa h))\); since \(r_s\) is
surjective, the history-level and record-level inequalities are the same
condition, which gives the second equivalence.
\end{proof}

\subsection*{Layer instantiations}\label{layer-inst:F38}

\begingroup\small
\begin{description}
\item[\emph{Git commit {DAG} and version-control history
(software engineering).}] \textit{Analogy.}
The source history set is a project-state at a given
commit; the target is the project-state after a child
commit; \(\kappa\) records the commit operation.  The
record quotient sends a state to its commit hash, with
status readout equal to the DAG-ancestry generation
(the longest ancestor chain length).  The hash of a
new commit is the cryptographic digest of its content
together with parent-commit references, so the
descended map carries each parent hash to the child
hash.  The status readout is strictly monotone: each
child commit has generation greater than each of its
parents \citep{ChaconStraub2014ProGit}.

\end{description}
\endgroup


\subsection{Entropy as Fiber Volume [\texorpdfstring{\Fref{F39}}{F39}]}\label{sec:self-reference-limits:entropy}

The fiber-volume reading of entropy combines two classical substrate
sources: the statistical-mechanical interpretation
\(S=k\log W\) developed in \citet{Boltzmann1877Beziehung} and the
communication-theoretic entropy of \citet{Shannon1948Mathematical}.
The layer-agnostic statement below records the structural form
common to both.

\paragraph{Reading the statement.} In \Cref{thm:F39} below,
\(\operatorname{TargetEntropyZero}(q,t)\) is defined combinatorially,
as \(\mathcal O_q(t)=\emptyset\): no \(q\)-fiber contains two different
target values.  Its identification with a vanishing entropy value is
the added clause on \(m(\langle q,t\rangle,q)\).  The fiber-volume
functional of the last clause is the conditional Shannon entropy of a
uniformly distributed history given its class, and \(m(q_F,q_C)\) is
then the entropy of the \(q_F\)-class given the \(q_C\)-class, so
\(E(q_C)=E(q_F)+m(q_F,q_C)\) is the chain rule of conditional entropy;
the theorem evaluates \(m\) only on refining pairs.  In this instance
\(m(\langle q,t\rangle,q)\) is the conditional entropy of the target
given the \(q\)-class, and the last clause identifies
\(\operatorname{TargetEntropyZero}(q,t)\) with its vanishing.

\begin{theorem}[Entropy as Fiber Volume]\label{thm:F39}
Let \(q:H\to Q\) be a surjective access quotient with target readout
\(t:H\to T\).  Let \(q_F:H\to Q_F\) be a finer quotient and
\(q_C:H\to Q_C\) a coarsening, related by a factorization
\(c:Q_F\to Q_C\) with
\[
  c\circ q_F=q_C.
\]
Let \(\mathcal E\) be a set of entropy values with a relation \(\le_E\) (no order axioms are used), and let
\[
  E:\{\,q':H\to Q'\ :\ Q'\text{ a set of a fixed universe}\,\}\longrightarrow \mathcal E
\]
be an entropy functional assigning a fiber-volume value to every
declared quotient.  Suppose the apparatus declares a mixing
functional \(m\) and a combiner
\(\odot:\mathcal E\times\mathcal E\to\mathcal E\) satisfying the
chain-rule identity
\[
  E(q_C)=E(q_F)\odot m(q_F,q_C)
\]
and the combine-monotonicity condition at this pair
\[
  E(q_F)\le_E E(q_F)\odot m(q_F,q_C)
\]
(which holds when \(a\le_E a\odot b\) for all \(a,b\), and in the fiber-volume
instance because \(m\ge0\)).
Suppose also that \(\mathcal E\) has a declared zero \(0\) such that
\(E(q')=0\) exactly when \(q'\) is injective, that is, every nonempty fiber of \(q'\) is a singleton.
Define the raw split-pair and target obstructions
\begin{align*}
  \mathcal O_q^{\mathrm{raw}}
  &:=
  \{(h,h')\in H^2:q(h)=q(h')\ \text{and}\ h\ne h'\},\\
  \mathcal O_q(t)
  &:=
  \{(h,h')\in H^2:q(h)=q(h')\ \text{and}\ t(h)\ne t(h')\},
\end{align*}
and put \(\operatorname{TargetEntropyZero}(q,t):\Longleftrightarrow\mathcal O_q(t)=\emptyset\)
and \(\operatorname{TargetUnresolved}(q,t):\Longleftrightarrow\mathcal O_q(t)\neq\emptyset\);
the entropy identification is the clause on \(m(\langle q,t\rangle,q)\) below.
Of the displayed lines, only the first and the last involve \(E\); the target predicates involve \(E\) only through the mixing clause below.
Then
\begin{align*}
  E(q)=0
  &\iff
  \mathcal O_q^{\mathrm{raw}}=\emptyset,\\
  \operatorname{TargetEntropyZero}(q,t)
  &\iff
  t\ \text{descends through}\ q,\\
  \operatorname{TargetUnresolved}(q,t)
  &\iff
  \neg(t\ \text{descends through}\ q),\\
  \forall b\in Q_C,\ \forall h\in H,\quad
  q_C(h)=b
  &\iff
  \exists a\in Q_F,\ c(a)=b\ \wedge\ q_F(h)=a,\\
  E(q_F)
  &\le_E
  E(q_C).
\end{align*}
If moreover \(m(\langle q,t\rangle,q)=0\) exactly when \(\langle q,t\rangle\)
and \(q\) have the same fibers (as holds when \(m(q',q'')=0\) exactly when
\(q'\) and \(q''\) have the same fibers, for every \(q'\) refining \(q''\)), then
the target entropy is a mixing value:
\[
  m(\langle q,t\rangle,q)=0
  \iff t\ \text{descends through}\ q
  \iff \operatorname{TargetEntropyZero}(q,t).
\]
For finite nonempty \(H\), the fiber-volume functional
\(E(q')=|H|^{-1}\sum_{h\in H}\log|q'^{-1}(q'(h))|\) with
\(m(q',q'')=E(q'')-E(q')\) for \(q'\) refining \(q''\), with values in
\(\mathbb R\) under addition and zero \(0\), satisfies the zero criterion, and the chain-rule identity and
combine-monotonicity at every refining pair,
\(m\ge0\), and \(m(q',q'')=0\) exactly when \(q'\) and \(q''\) have the same
fibers; hence all conclusions above hold for it, in particular
\(E(q_F)\le E(q_C)\) and
\(\operatorname{TargetEntropyZero}(q,t)\iff m(\langle q,t\rangle,q)=0\).
\end{theorem}
\begin{proof}
Assume \(q\) is surjective, \(c\circ q_F=q_C\), the chain-rule
identity \(E(q_C)=E(q_F)\odot m(q_F,q_C)\), and combine-monotonicity
\(E(q_F)\le_E E(q_F)\odot m(q_F,q_C)\).  We prove the five conjuncts and then
the added clause.

\emph{(1) Raw entropy zero iff no raw split pair.}
By the declared zero, \(E(q)=0\) exactly when every \(q\)-fiber is a
singleton, that is, when no \(q\)-fiber contains distinct elements,
which is \(\mathcal O_q^{\mathrm{raw}}=\emptyset\).

\emph{(2) Target entropy zero iff target descends.}
Target entropy is measured by the target obstruction
\(\mathcal O_q(t)\).  Its emptiness is exactly the quotient-descent
condition for \(t\) through \(q\):
\[
  \mathcal O_q(t)=\emptyset
  \iff
  \exists\,\bar t:Q\to T,\ \bar t\circ q=t.
\]
This is \(\operatorname{TargetEntropyZero}(q,t)\) iff target descent.

\emph{(3) Target unresolved iff target does not descend.}
The target-unresolved predicate is witnessed by a split pair
\((h,h')\in\mathcal O_q(t)\).  By (2), nonemptiness of
\(\mathcal O_q(t)\) is the negation of target descent.

\emph{(4) Coarse fibers decompose into fine fibers.}
Let \(b\in Q_C\) and \(h\in H\).  If \(q_C(h)=b\), set
\(a:=q_F(h)\).  The factorization \(c\circ q_F=q_C\) gives
\[
  c(a)=c(q_F(h))=q_C(h)=b,
\]
and by construction \(q_F(h)=a\).  Conversely, if there exists
\(a\in Q_F\) with \(c(a)=b\) and \(q_F(h)=a\), then the same
factorization gives
\[
  q_C(h)=c(q_F(h))=c(a)=b.
\]

\emph{(5) Entropy is monotone under coarsening.}
Apply the chain-rule identity to the factorized coarsening pair
\((q_F,q_C)\):
\[
  E(q_C)=E(q_F)\odot m(q_F,q_C).
\]
Combine-monotonicity with \(a:=E(q_F)\) and \(b:=m(q_F,q_C)\) gives
\[
  E(q_F)\le_E E(q_F)\odot m(q_F,q_C)=E(q_C).
\]
Thus \(E(q_F)\le_E E(q_C)\).

\emph{(6) Target entropy as a mixing value.} Under the added hypothesis,
\(m(\langle q,t\rangle,q)=0\) exactly when the two quotients have the same
fibers, that is, when \(q(x)=q(y)\) implies
\((q(x),t(x))=(q(y),t(y))\), the converse being automatic; that is, exactly
when \(t\) is constant on \(q\)-fibers, which by (2) is descent of \(t\) and
\(\operatorname{TargetEntropyZero}(q,t)\). For the fiber-volume functional,
\(m(q',q'')=|H|^{-1}\sum_h\log\bigl(|q''^{-1}(q''h)|/|q'^{-1}(q'h)|\bigr)\) is
a sum of nonnegative terms, since each \(q'\)-fiber lies in a \(q''\)-fiber, and
it vanishes exactly when every term does, that is, when the fibers coincide.
\end{proof}

\subsection*{Layer instantiations}\label{layer-inst:F39}

\begingroup\small
\begin{description}
\item[\emph{Simpson concentration index (ecology).}] \textit{Analogy.}
For a community of \(N\) individuals classified into
species, the access quotient sends each of the \(N\) individuals to its species.  The
Simpson concentration index
\(\lambda=\sum_i p_i^2\) (with \(p_i\) the abundance of
species \(i\)) is the probability that two individuals drawn with
replacement lie in one species fiber; it is not the fiber-volume functional
\(E\) of \Cref{thm:F39} (for the species quotient \(E=\sum_i p_i\log(Np_i)\)),
but, like \(E\), it does not decrease under coarsening.  Taxonomic coarsening (genus,
family) combines species classes; since
\(p_g=\sum_{i\in g}p_i\) and the cross terms of \(p_g^2\) are
nonnegative, \(\sum_g p_g^2\ge\sum_i p_i^2\), so the index is
monotone non-decreasing under coarsening
\citep{Simpson1949Diversity}.

\item[\emph{Burnside orbit-counting in finite group
theory (group theory).}] \textit{Special case.}
For a finite group \(G\) acting on a finite nonempty set \(X\),
the access quotient sends each element to its
\(G\)-orbit. Take any target readout \(t\), for instance \(t(x)=|Gx|\).  The fiber over an orbit \(O\) has
cardinality \(|O|\). For every quotient \(q'\) of the finite nonempty set
\(X\), take \(\mathcal E=[0,\infty)\), \(E(q')=\log(|X|/|\operatorname{im}q'|)\)
and \(\odot=+\); for a coarsening \(q_C=c\circ q_F\), set
\(m(q_F,q_C)=\log(|\operatorname{im}q_F|/|\operatorname{im}q_C|)\). Then
\(m\ge0\), the chain rule and combine-monotonicity of \Fref{F39} hold, and
\(E(q')=0\) exactly when \(q'\) is injective.  This \(E\) is not the
fiber-volume functional of the theorem's last clause; the two agree
when all fibers have the same size. The orbit quotient has entropy
\(\log(|X|/|X/G|)\), with \(|X/G|\) given by
Burnside's lemma as the average number of fixed
points of group elements.  Coarsening to a supergroup
\(G'\supseteq G\) merges orbits into larger fibers,
non-decreasing the average orbit size
\citep{Burnside1897Theory}.

\end{description}
\endgroup


\subsection{Irreducibility / No Short Closed Description [\texorpdfstring{\Fref{F42}}{F42}]}\label{sec:self-reference-limits:irreducibility}

The notion of an object's shortest closed description is the central
object of Kolmogorov complexity and algorithmic information theory;
the standard reference is \citet{LiVitanyi2008Kolmogorov}.

\paragraph{Reading the statement.} The symbols are those of \Cref{thm:F42} below.
\(\mathcal O_t(e,h)\) is the set of pairs in the \(e\)-class of \(h\)
with different target values, and \(\mathcal O_\gamma(e,h)\) the set of
pairs in that class whose compressed states differ after the route
\(r_\gamma\). \(\operatorname{Irreducible}_e(h)\) says that no closed
compression \(e'\) in the admissible class \(\mathcal A\) describes \(h\) more
cheaply than \(e\) does. The predicate does not require \(e\) itself to lie in
\(\mathcal A\) or to be closed for \(h\); it is meant for \(e\in\mathcal A\) with
\(\operatorname{ClosedCompression}(e,h)\), as produced by the existence clause. The class matters: closedness at \(h\) looks only at
the class of \(h\), so if the admissible maps can give \(h\) any value on a
class of its own (for instance, if every map \(H\to E\) is admissible and
\(E\) has at least two elements), irreducibility reduces to minimality of \(\operatorname{cost}(e(h))\) over all of \(E\).

\begin{theorem}[Irreducibility / No Short Closed Description]\label{thm:F42}
\emph{Pointwise form.} Let \(H\) be a set of objects with compression
map \(e:H\to E\), target readout \(t:H\to T\), and route updates
\(r_\gamma:H\to H\) for \(\gamma\in\Gamma\). Let \(\mathcal A\) be a declared
class of admissible compressions \(H\to E\), and let
\(\operatorname{cost}:E\to C\) take values in a set \(C\) with a relation \(<\) (no order axioms are used; the existence clause uses only well-foundedness). For \(h\in H\) put
\(\mathcal O_t(e,h):=\{(x,x')\in H^2:e(x)=e(x')=e(h),\ t(x)\ne t(x')\}\) and
\(\mathcal O_\gamma(e,h):=\{(x,x')\in H^2:e(x)=e(x')=e(h),\ e(r_\gamma x)\ne
e(r_\gamma x')\}\). A compression \(e\) is closed for \(h\), written
\(\operatorname{ClosedCompression}(e,h)\), when \(t\) and every
\(e\circ r_\gamma\) are constant on the class \(e^{-1}(e(h))\). Define
\[
  \operatorname{Irreducible}_e(h)
  :\Longleftrightarrow
  \neg\exists e'\in\mathcal A,\
  \operatorname{ClosedCompression}(e',h)
  \ \wedge\
  \operatorname{cost}(e'(h))<\operatorname{cost}(e(h)).
\]
Then
\[
  \operatorname{ClosedCompression}(e,h)
  \iff
  \mathcal O_t(e,h)=\emptyset
  \ \wedge\
  \forall\gamma\in\Gamma,\ \mathcal O_\gamma(e,h)=\emptyset.
\]
By the displayed equivalence, \(\operatorname{Irreducible}_e(h)\) holds
exactly when every \(e'\in\mathcal A\) with
\(\operatorname{cost}(e'(h))<\operatorname{cost}(e(h))\) has
\(\mathcal O_t(e',h)\ne\emptyset\) or \(\mathcal O_\gamma(e',h)\ne\emptyset\)
for some \(\gamma\). Moreover, if \(<\) is well founded on \(C\) and some compression
in \(\mathcal A\) is closed for \(h\), then there is \(e^\ast\in\mathcal A\)
that is closed for \(h\) and satisfies
\(\operatorname{Irreducible}_{e^\ast}(h)\). If for every \(v\in E\) the class \(\mathcal A\) contains a map \(e'\) with
\(e'(h)=v\) and \(e'(x)\ne v\) for every \(x\ne h\) (for instance, if
\(\mathcal A\) contains every map \(H\to E\) and \(E\) has at least two
elements), then
\(\operatorname{Irreducible}_e(h)\) holds exactly when no \(v\in E\) has
\(\operatorname{cost}(v)<\operatorname{cost}(e(h))\). For such a class
\(\mathcal A\) the pointwise notion depends on neither \(t\) nor the routes
\(r_\gamma\); it constrains the target only when \(\mathcal A\) excludes maps
isolating \(h\). The global notion below requires closure on every class, so an
isolating map is globally closed only when \(t\) and every \(r_\gamma\) respect
all its classes.
\emph{Global form.} Let \(b:H\to B\) be a map. Separately declare an admissible
class \(\mathcal A_{\mathrm{glob}}\) of descriptions \(e:H\to E'\), allowing the
codomain \(E'\) to vary, and a cost \(c\) on these descriptions and on \(b\).
Call \(e\) \emph{globally closed over} \(b\) when \(e\) descends through \(b\),
\(t\) descends through \(e\), and \(e\circ r_\gamma\) descends through \(e\) for
every \(\gamma\). For surjective \(b\) and \(e\), this holds exactly when the sets
\(\{(x,x'):b(x)=b(x'),\ e(x)\ne e(x')\}\) and
\(\{(x,x'):e(x)=e(x'),\ t(x)\ne t(x')\}\) and, for every \(\gamma\),
\(\{(x,x'):e(x)=e(x'),\ e(r_\gamma x)\ne e(r_\gamma x')\}\) are empty. Define
\(b\) to be \emph{globally irreducible} when no \(e\in\mathcal A_{\mathrm{glob}}\)
globally closed over \(b\) satisfies \(c(e)<c(b)\). If \(c\) takes values in a
well-founded strict order and some description in \(\mathcal A_{\mathrm{glob}}\)
is globally closed over \(b\), then some \(e^\ast\in\mathcal A_{\mathrm{glob}}\)
globally closed over \(b\) has no description in \(\mathcal A_{\mathrm{glob}}\)
globally closed over \(b\) of smaller cost, and \(e^\ast\) is globally
irreducible (with \(e^\ast\) in place of \(b\)).
\end{theorem}
\begin{proof}
Assume the named maps and the declared cost relation. A compression is
closed exactly when no target or route update distinguishes two
representatives of the same compressed description. In symbols,
\[
  \mathcal O_t(e,h)=
  \{(x,x'):e(x)=e(x')=e(h),\ t(x)\ne t(x')\},
\]
and similarly \(\mathcal O_\gamma(e,h)\) is formed using
\(e(r_\gamma x)\). Hence closedness is precisely the simultaneous
emptiness of these obstruction sets. Substituting this for each cheaper \(e'\in\mathcal A\) in the definition of
\(\operatorname{Irreducible}_e(h)\) gives the reformulation after the display:
a shorter closed \(e'\) carries all declared readouts with smaller cost, so
\(h\) is reducible exactly when some cheaper admissible candidate has all its
obstruction sets empty.
If \(\operatorname{Irreducible}_e(h)\) holds, no admissible closed compression
has cost below \(\operatorname{cost}(e(h))\). If moreover \(e\in\mathcal A\)
is closed at \(h\), then \(e\) is minimal among admissible closed compressions.
For the existence clause, the costs \(\operatorname{cost}(e'(h))\) of the
closed compressions \(e'\in\mathcal A\) at \(h\) form a nonempty subset of
\(C\). Since \(<\) is well founded, this subset has a \(<\)-minimal element
\(\operatorname{cost}(e^\ast(h))\) with \(e^\ast\in\mathcal A\) closed at
\(h\), and no closed compression in \(\mathcal A\) at \(h\) has a cost below
it, which is \(\operatorname{Irreducible}_{e^\ast}(h)\).
For the last clause, let \(v\in E\) and take \(e'\in\mathcal A\) with
\(e'(h)=v\) and \(e'(x)\ne v\) for \(x\ne h\); when \(\mathcal A\) contains
every map and \(E\) has two elements, choose \(w\ne v\) and set
\(e'(h)=v\) and \(e'(x)=w\) for \(x\ne h\). The map \(e'\) has
\(\{h\}\) as the \(e'\)-class of \(h\), so both obstruction sets are empty and
\(e'\) is closed for \(h\). Hence a \(v\) with
\(\operatorname{cost}(v)<\operatorname{cost}(e(h))\) gives a cheaper closed
\(e'\in\mathcal A\); conversely, a cheaper closed \(e'\) gives such a \(v\),
namely \(e'(h)\). For the global clause, each descent condition is the
emptiness of the corresponding set by the descent criterion, using
surjectivity of \(b\) and \(e\). An isolating map is globally closed only if
\(t\) and every \(r_\gamma\) respect its classes on all of \(H\). For the global
minimum, take a closed description of minimal cost; any description globally
closed over \(e^\ast\) descends through \(e^\ast\), hence through \(b\), so it is
closed over \(b\) and cannot cost less.
\end{proof}

\subsection*{Layer instantiations}\label{layer-inst:F42}

\begingroup\small
\begin{description}
\item[\emph{Abel-Ruffini-Galois unsolvability of the
quintic (algebra).}] \textit{Analogy.}
For a quintic over \(\mathbb Q\) with Galois group \(S_5\), no
root is a nested-radical expression in its coefficients, by Galois's solvability criterion; Abel proved the corresponding statement for the general quintic \citep{Abel1824Quintic,Galois1846Memoire}. To instantiate \Fref{F42}, one must
specify \(H,E,t,\mathcal A\) and establish that every compression
in \(\mathcal A\) closed at the chosen polynomial would yield such
a radical expression for the target root.

\end{description}
\endgroup


\subsection{Continuum Emergence [\texorpdfstring{\Fref{F43}}{F43}]}\label{sec:self-reference-limits:continuum-emergence}

\paragraph{Reading the statement.} \Cref{thm:F43} below does not
construct a limit: the candidate limit \(L\), the projections \(p_i\) and
the limit quotient \(\ell\) are given, and no topology or convergence
beyond the declared predicate \(\operatorname{Vanishes}\) on the
residuals is used.  \(\operatorname{Cofinal}(I)\) says that every
declared finite test is covered at some index; it is a covering
condition, not order-theoretic cofinality.  Coherence, cofinality and
vanishing enter \(\operatorname{StableLimitQuotient}(L)\) only as
assumed conjuncts and are carried unchanged into the conclusion.  The
content of the theorem is the existence and uniqueness of a readout
\(\bar\ell\) on the limit quotient that agrees with every stage readout:
the equations \(\phi_{ij}\circ p_j=p_i\) and \(r_i\circ\phi_{ij}=r_j\),
together with directedness, make the stage readouts agree, and
\(\bar\ell\) exists when their common value descends through \(\ell\).

\begin{theorem}[Continuum Emergence]\label{thm:F43}
Let \((I,\le)\) be a nonempty directed index set, let \(A_i\) be a discrete
approximation at index \(i\), and let
\(\phi_{ij}:A_j\to A_i\) be comparison maps for \(i\le j\). Let
\(p_i:L\to A_i\) be projections from a candidate limit \(L\), let
\(\xi_i\in R_i\) be residuals, let
\(\ell:L\to Q_\infty\) be a surjective limit quotient, let
\(r_i:A_i\to L_{\mathrm{rec}}\) be stage readouts, and let
\(t:L\to L_{\mathrm{rec}}\) be a target readout. Let \(\mathcal T\) be a set of
declared finite tests with a covering relation
\(\operatorname{cov}\subseteq I\times\mathcal T\), write
\(\operatorname{Cofinal}(I)\) (a covering condition, not order-theoretic
cofinality) \(:\Longleftrightarrow\forall\tau\in\mathcal T\ \exists i\in I\ \operatorname{cov}(i,\tau)\),
and let \(\operatorname{Vanishes}\) be a declared predicate on residual
families \((\xi_i)_{i\in I}\). Write \(\operatorname{TowerCoherent}(\phi_{ij})\)
when \(\phi_{ii}=\mathrm{id}\) and \(\phi_{ik}=\phi_{ij}\circ\phi_{jk}\) for
\(i\le j\le k\), and call \(\bar\ell:Q_\infty\to L_{\mathrm{rec}}\) compatible with
all \(p_i\) when \(\bar\ell(\ell(x))=r_i(p_i(x))\) for all \(i\in I\) and
\(x\in L\). Write
\begin{align*}
  \operatorname{StableLimitQuotient}(L)
  :\iff\
  &\operatorname{TowerCoherent}(\phi_{ij})
  \ \wedge\
  \operatorname{Cofinal}(I)\\
  &\wedge\
  \operatorname{Vanishes}(\xi_i)
  \ \wedge\
  \exists\,\bar\ell:Q_\infty\to L_{\mathrm{rec}}\ \text{ compatible with all }p_i .
\end{align*}
Suppose that the tower is coherent and cofinal, that the residuals vanish,
that \(\phi_{ij}\circ p_j=p_i\) and \(r_i\circ\phi_{ij}=r_j\) for \(i\le j\),
and that for some index \(b\) the target readout satisfies
\(t=r_b\circ p_b\) and is constant on the fibers of \(\ell\). Then there is exactly one \(\bar\ell:Q_\infty\to L_{\mathrm{rec}}\) compatible
with all \(p_i\), namely \(\bar\ell(\ell(x))=t(x)\); with the assumed coherence,
cofinality and vanishing, \(\operatorname{StableLimitQuotient}(L)\) holds. More generally, assume only that \(\phi_{ij}\circ p_j=p_i\) and
\(r_i\circ\phi_{ij}=r_j\) for \(i\le j\). Then \(r_i\circ p_i=r_j\circ p_j\) for all
\(i,j\in I\); for every index \(b\), a map \(\bar\ell\) compatible with all \(p_i\)
exists if and only if \(r_b\circ p_b\) is constant on the fibers of \(\ell\), and
it is then unique; hence \(\operatorname{StableLimitQuotient}(L)\) holds if and
only if, in addition, the tower is coherent and cofinal and the residuals
vanish.
\end{theorem}
\begin{proof}
Fix \(x\in L\) and \(i\in I\). Since \(I\) is directed, choose \(k\) with
\(b\le k\) and \(i\le k\). Using \(p_b=\phi_{bk}\circ p_k\) and
\(r_b\circ\phi_{bk}=r_k\),
\[
  t(x)=r_b(p_b(x))=r_b(\phi_{bk}(p_k(x)))=r_k(p_k(x)),
\]
and in the same way
\(r_i(p_i(x))=r_i(\phi_{ik}(p_k(x)))=r_k(p_k(x))\). Hence
\(r_i(p_i(x))=t(x)\) for every \(i\). Since \(t\) is constant on the fibers
of \(\ell\) and \(\ell\) is surjective, \(\bar\ell(\ell(x)):=t(x)\) defines a
map \(\bar\ell:Q_\infty\to L_{\mathrm{rec}}\), and
\(\bar\ell(\ell(x))=r_i(p_i(x))\) for all \(i\) and \(x\), so \(\bar\ell\) is
compatible with all \(p_i\). Together with the coherence, cofinality and
vanishing hypotheses, the four conditions defining
\(\operatorname{StableLimitQuotient}(L)\) hold.
For the general clause, the same computation with a common upper bound \(k\)
of \(i\) and \(j\) gives \(r_i\circ p_i=r_k\circ p_k=r_j\circ p_j\). If
\(\bar\ell\) is compatible with all \(p_i\), then
\(r_b(p_b(x))=\bar\ell(\ell(x))\) is constant on the fibers of \(\ell\);
conversely, if \(r_b\circ p_b\) is constant on those fibers, the argument above
with \(t:=r_b\circ p_b\) gives a compatible \(\bar\ell\). Since \(\ell\) is
surjective, \(\bar\ell\) is determined by \(\bar\ell(\ell(x))=r_b(p_b(x))\).
\end{proof}

The result supplies a readout on the specified quotient of a candidate limit
that agrees with every stage readout.

\subsection*{Layer instantiations}\label{layer-inst:F43}

\begingroup\small
\begin{description}
\item[\emph{Dedekind cut construction of the reals
(analysis).}] \textit{Analogy.}
The discrete approximations are rational intervals
indexed by precision; comparison maps are
inclusions of finer intervals into coarser ones.
A Dedekind cut is a non-empty proper downward-closed
subset \(A\subsetneq\mathbb Q\) without greatest element.  The candidate limit is the set of cuts;
projections send a cut to its rational interval
approximant; the surjective limit quotient sends
a cut to its real value, and tower coherence is
nested-approximation consistency.  The analogy
is informal: the cut construction is not
presented here as a tower satisfying the
hypotheses of \Fref{F43} \citep{Dedekind1872Stetigkeit}.

\end{description}
\endgroup


\subsection{Criticality / Phase Boundary [\texorpdfstring{\Fref{F44}}{F44}]}\label{sec:self-reference-limits:criticality}

The motivating substrate analogy is the theory of phase transitions and
critical phenomena; the standard textbook treatment is
\citet{Stanley1971PhaseTransitions}.  The layer-agnostic statement
below records the descent-and-margin shape that the substrate
realizes.

\paragraph{Reading the statement.} In \Cref{thm:F44} below, the three
displayed lines define phase transition, phase boundary and criticality from
the declared data, and clause (i) is the ordinary quotient-descent criterion. Neighborhoods are indexed by \(W\in\mathcal W\), and \(U\) is the non-uniform-response
predicate. For a general relation \(N\) the predicate is only as informative
as \(N\): if \(\mathcal W\) is empty, every formed control is a phase boundary,
and an index \(W\) with no \(y\) satisfying \(N(x,W,y)\) makes \(x\) locally
phase-constant. The relation of (iv) excludes both cases. Clause (iii) asks that, at each step and for every \(W\), some
neighbour of \(x_k\) have the phase of \(x_{k+1}\); with the neighbourhood
relation of (iv) this says that \(x_k\) lies in the closure of the phase class
of \(x_{k+1}\), and (iv) is the form that applies to continuous paths.
The relation of (iv) is chosen so that the empty open set does not
witness local constancy, as it would for the relation
\(x\in W\wedge y\in W\).

\begin{theorem}[Criticality / Phase Boundary]\label{thm:F44}
Let \(\mathcal C\) be a control space with surjective quotient
\(c:\mathcal C\to C\), phase readout \(\varphi:\mathcal C\to\Phi\),
formedness predicate \(F:\mathcal C\to\mathrm{Prop}\), and
neighborhood relation \(N\subseteq\mathcal C\times\mathcal W\times\mathcal C\) over a
declared set \(\mathcal W\) of neighbourhood indices (in (iv), \(\mathcal W\) is
the set of open subsets of \(\mathcal C\)). Let
\(M,U,R:\mathcal C\to\mathrm{Prop}\) denote, respectively,
margin-vanishing, non-uniform response, and
residuals-statused. For a readout \(r:\mathcal C\to V\) let
\(\mathcal O_c(r)=\{(x,y):c(x)=c(y)\ \text{and}\ r(x)\ne r(y)\}\), and for
\(x,x_-,x_+\in\mathcal C\) define
\begin{align*}
  \operatorname{PhaseTransition}(x_-,x_+)
  &:\Longleftrightarrow \varphi(x_-)\ne\varphi(x_+),\\
  \operatorname{PhaseBoundary}(x)
  &:\Longleftrightarrow F(x)\ \wedge\
    \neg\exists W,\ \forall y,\ N(x,W,y)\Rightarrow \varphi(y)=\varphi(x),\\
  \operatorname{Critical}(x)
  &:\Longleftrightarrow \operatorname{PhaseBoundary}(x)\wedge (M(x)\vee U(x))\wedge R(x).
\end{align*}
Then:
\begin{enumerate}[label=(\roman*)]
\item \(r\) descends through \(c\), that is \(r=\bar r\circ c\) for some
  \(\bar r:C\to V\), if and only if \(\mathcal O_c(r)=\emptyset\);
\item if \(\varphi\), \(F\), \(M\), \(U\) and \(R\) are constant on the fibers
  of \(c\), and \(N\) respects \(c\) (\(N(x,W,y)\iff N(x',W,y)\) whenever
  \(c(x)=c(x')\)), then \(\operatorname{PhaseBoundary}\) and
  \(\operatorname{Critical}\) are constant on the fibers of \(c\), so by (i)
  they descend through \(c\): whether a control is critical depends only on
  its class;
\item if \(x_0,\dots,x_n\) is a finite list of controls with \(F(x_k)\) for
  every \(k\) and, for every \(k<n\) and every \(W\), some \(y\) with
  \(N(x_k,W,y)\) has \(\varphi(y)=\varphi(x_{k+1})\) (for instance,
  \(N(x_k,W,x_{k+1})\) for every \(W\), as when \(N\) is a neighbourhood graph
  whose relation does not depend on \(W\); with the topological relation \(N(x,W,y):\Leftrightarrow(x\in W\Rightarrow y\in W)\), \(W\) open, used in (iv), the condition says exactly that \(x_k\) lies in the closure of the phase class of \(x_{k+1}\)), and
  \(\varphi(x_0)\ne\varphi(x_n)\), then some \(k<n\) has
  \(\varphi(x_k)\ne\varphi(x_{k+1})\) and
  \(\operatorname{PhaseBoundary}(x_k)\): a phase transition along a path of
  neighbouring formed controls crosses a phase boundary;
\item if \(\mathcal C\) is a topological space, \(W\) ranges over its open
  sets and \(N(x,W,y)\) means that \(y\in W\) whenever \(x\in W\), and
  \(\gamma:[0,1]\to\mathcal C\) is continuous with \(F(\gamma(s))\) for every
  \(s\) and \(\varphi(\gamma(0))\ne\varphi(\gamma(1))\), then
  \(\operatorname{PhaseBoundary}(\gamma(s))\) for some \(s\in[0,1]\): a phase
  transition along a continuous path of formed controls crosses a phase
  boundary. With this \(N\), \(\operatorname{PhaseBoundary}(x)\) holds exactly
  when \(F(x)\) holds, \(\varphi\) is not constant on \(\mathcal C\), and
  \(\varphi\) is constant on no open neighbourhood of \(x\).
\end{enumerate}
\end{theorem}
\begin{proof}
Since \(c\) is surjective, (i) is the descent criterion of
\Cref{sec:apparatus}.  The defined predicates sort a formed control \(x\) into two cases:
\[
  \begin{cases}
  \exists W\,\forall y,\ N(x,W,y)\Rightarrow\varphi(y)=\varphi(x),
    &\text{\(x\) is locally phase-constant,}\\
  \neg\exists W\,\forall y,\ N(x,W,y)\Rightarrow\varphi(y)=\varphi(x),
    &\text{\(x\) lies on a phase boundary,}
  \end{cases}
\]
and a critical control is a boundary control at which the separating margin
vanishes or the response is non-uniform, with its residuals statused.
For (ii), let \(c(x)=c(x')\). Then \(F(x)=F(x')\) and
\(\varphi(x)=\varphi(x')\), and for every \(W\) and \(y\),
\(N(x,W,y)\iff N(x',W,y)\); so \(x\) is locally phase-constant exactly when
\(x'\) is, and \(\operatorname{PhaseBoundary}(x)\iff
\operatorname{PhaseBoundary}(x')\). With \(M(x)=M(x')\), \(U(x)=U(x')\) and
\(R(x)=R(x')\), also \(\operatorname{Critical}(x)\iff
\operatorname{Critical}(x')\).
(iii) Since \(\varphi(x_0)\ne\varphi(x_n)\), some \(k<n\) has
\(\varphi(x_{k+1})\ne\varphi(x_k)\). For every \(W\), the hypothesis gives
\(y\) with \(N(x_k,W,y)\) and \(\varphi(y)=\varphi(x_{k+1})\ne\varphi(x_k)\), so
no \(W\) makes \(\varphi\) constant at \(x_k\); with \(F(x_k)\) this is
\(\operatorname{PhaseBoundary}(x_k)\).
(iv) Suppose no \(\gamma(s)\) is a boundary point. Since every
\(\gamma(s)\) is formed, for each \(s\) there is an open \(W_s\) with
\(\varphi(y)=\varphi(\gamma(s))\) for every \(y\) such that
\(\gamma(s)\in W_s\Rightarrow y\in W_s\). If \(\gamma(s)\notin W_s\), every
\(y\) qualifies, so \(\varphi\) is constant on \(\mathcal C\), contradicting
\(\varphi(\gamma(0))\ne\varphi(\gamma(1))\). Hence \(\gamma(s)\in W_s\) and
\(\varphi\) is constant on \(W_s\), so \(\varphi\circ\gamma\) is constant on
the open neighbourhood \(\gamma^{-1}(W_s)\) of \(s\). A locally constant map
on the connected interval \([0,1]\) is constant, so
\(\varphi(\gamma(0))=\varphi(\gamma(1))\), a contradiction.
For the last sentence of (iv): an open \(W\not\ni x\) relates \(x\) to every
\(y\), so it witnesses local constancy exactly when \(\varphi\) is constant on
\(\mathcal C\); an open \(W\ni x\) does so exactly when \(\varphi\) is constant on
\(W\). Since \(\mathcal C\) is an open neighbourhood of \(x\),
\(\operatorname{PhaseBoundary}(x)\) holds exactly when \(F(x)\) holds and
\(\varphi\) is constant on no open neighbourhood of \(x\), which implies that
\(\varphi\) is not constant on \(\mathcal C\).
\end{proof}

\subsection*{Layer instantiations}\label{layer-inst:F44}

\begingroup\small
\begin{description}
\item[\emph{Allee effect threshold (population
ecology).}] \textit{Analogy.}
For a species population \(\dot N=N\,r(N)\), the
control space is population size; the phase readout
records ``extinction'' vs ``persistence''.  The Allee
effect makes the per-capita growth \(r(N)\) negative
below an Allee threshold \(N_A\), producing a phase
boundary at \(N=N_A\): the per-capita growth rate
\(r(N_A)\) vanishes, response to environmental perturbations is
non-uniform, and the residuals-statused predicate \(R\) would have
to be supplied as a declared status of demographic
noise, which the analogy does not provide
\citep{Allee1931Aggregations}.

\item[\emph{Thom cusp catastrophe (catastrophe
theory).}] \textit{Analogy.}
The cusp catastrophe is the family of potentials
\(V_a(x)=x^4/4+a_1 x^2/2+a_2 x\) parameterised by
\((a_1,a_2)\in\mathbb R^2\); the phase readout
records the number of stable equilibria.  The
bifurcation set \(4a_1^3+27a_2^2=0\) is the phase boundary: on its
fold branches one stable equilibrium merges with the unstable one,
all three equilibria merge at the cusp point \(a=0\), the energy
margin vanishes, the response is
non-uniform, and the residual smooth-perturbation
status is classified by the elementary catastrophe
normal form \citep{Thom1972Stabilite}.

\item[\emph{{NBER} business-cycle peak/trough
(macroeconomics).}] \textit{Analogy.}
The control space is macroeconomic time series of
GDP, employment, industrial production, and personal
income; the phase readout records expansion vs
recession.  NBER-defined peaks and troughs are
phase boundaries: leading-indicator marginal growth
vanishes, sectoral responses are non-uniform, and
the residual status is handled by the NBER Business
Cycle Dating Committee's revision protocol over
noisy indicators \citep{BurnsMitchell1946Measuring}.

\item[\emph{Hopf bifurcation onset of oscillation
(nonlinear dynamics).}] \textit{Analogy.}
For a parameter-dependent ODE
\(\dot{\mathbf x}=f(\mathbf x;\mu)\), the control
space is \(\mu\); the phase readout records ``stable
fixed point'' vs ``limit cycle''.  A Hopf bifurcation
at \(\mu=\mu_H\) occurs when a complex-conjugate
eigenvalue pair of the Jacobian crosses the
imaginary axis: the real part vanishes, in the supercritical case
oscillation amplitude grows as \(\sqrt{\mu-\mu_H}\) (a subcritical
bifurcation instead produces a jump to a distant attractor), and the
residual perturbation is statused by smooth
normal-form coefficients \citep{Hopf1942Abzweigung}.
\end{description}
\endgroup


\subsection{Locality / Domain of Dependence [\texorpdfstring{\Fref{F45}}{F45}]}\label{sec:self-reference-limits:locality}

The classical relativistic analogy for a local quotient with a
domain-of-dependence structure is the causal structure of
Lorentzian spacetime; the standard reference for the
domain-of-dependence formalism is \citet[Chapter~8]{Wald1984GR}.

\begin{theorem}[Locality / Domain of Dependence]\label{thm:F45}
Let \(H\) be a history set with local quotient \(\ell:H\to L\),
target readout \(t:H\to T\), signature \(\sigma:H\to\Sigma\), update
map \(u:H\to H'\), and target quotient \(q':H'\to Q'\). Assume
\(\ell\) is surjective. Let \(<\) be a declared strict-smaller relation on a
declared class of local quotients of \(H\) that contains \(\ell\); in
\(\operatorname{MinimalDomain}\), \(\ell'\) ranges over that class. Write
\(\operatorname{LocalToDomain}(\ell,t)\) when \(t\) factors through \(\ell\),
that is \(t=\bar t\circ\ell\) for some \(\bar t:L\to T\);
\(\operatorname{TargetFamilyLocal}(\ell,\sigma)\) when \(\sigma\) factors
through \(\ell\); and \(\operatorname{DynamicLocality}(\ell,u,q')\) when
\(q'\circ u\) factors through \(\ell\). Define
\[
  \mathcal O_\ell(t)=
  \{(h,h'):\ell(h)=\ell(h')\ \text{and}\ t(h)\ne t(h')\},
\]
and define
\[
  \operatorname{MinimalDomain}(\ell,\sigma)
  :\Longleftrightarrow \operatorname{TargetFamilyLocal}(\ell,\sigma)\ \wedge\
  \neg\exists \ell',\ \ell'\!<\ell\ \wedge\
  \operatorname{TargetFamilyLocal}(\ell',\sigma).
\]
Then
\begin{align*}
  \operatorname{LocalToDomain}(\ell,t)
  &\iff \mathcal O_\ell(t)=\emptyset,\\
  \operatorname{TargetFamilyLocal}(\ell,\sigma)
  &\iff
  \forall a,b\in H,\ \ell(a)=\ell(b)\Rightarrow \sigma(a)=\sigma(b),\\
  \operatorname{DynamicLocality}(\ell,u,q')
  &\iff
  \forall a,b\in H,\ \ell(a)=\ell(b)\Rightarrow q'(u(a))=q'(u(b)).
\end{align*}
Moreover, the local repair \(\langle\ell,t\rangle:H\to L\times T\) satisfies
\(\operatorname{LocalToDomain}(\langle\ell,t\rangle,t)\). If \(\ell'<\ell\)
means that \(\ell'\) factors through \(\ell\) and \(\ell\) does not factor
through \(\ell'\), the declared class contains \(\sigma\), and \(H\) is
nonempty, then \(\operatorname{MinimalDomain}(\ell,\sigma)\) holds if and only if \(\ell\) and
\(\sigma\) have the same fibers.
\end{theorem}
\begin{proof}
Since \(\ell\) is surjective, the three displayed equivalences are the
descent criterion of \Cref{sec:apparatus} for \(\ell\) and the readouts
\(t\), \(\sigma\) and \(q'\circ u\).  For the local repair, \(t=\mathrm{pr}_2\circ\langle\ell,t\rangle\), so \(t\)
factors through \(\langle\ell,t\rangle\). For the last clause, note that for
nonempty \(H\) two maps on \(H\) with the same fibers factor through each other
(for empty \(H\) a map into an empty codomain need not exist). Suppose
\(\operatorname{MinimalDomain}(\ell,\sigma)\). Then \(\sigma\) factors through
\(\ell\); if \(\ell\) did not factor through \(\sigma\), then \(\sigma<\ell\)
and \(\sigma\) factors through itself, contradicting minimality, so \(\ell\)
and \(\sigma\) have the same fibers. Conversely, if they have the same fibers
and \(\sigma\) factors through some \(\ell'\) with \(\ell'<\ell\), then
\(\ell\), which factors through \(\sigma\), factors through \(\ell'\),
contradicting \(\ell'<\ell\).
\end{proof}

\subsection*{Layer instantiations}\label{layer-inst:F45}

\begingroup\small
\begin{description}
\item[\emph{Markov property in stochastic processes
(probability theory).}] \textit{Special case.}
For a time-homogeneous Markov chain \((X_t)\) on a countable
state space \(S\) with transition matrix \(P\), let \(H\) be the set of
state sequences \(\omega=(x_0,\dots,x_k)\) of positive probability,
\(\ell(\omega)=x_k\) and \(L=\ell(H)\), so \(\ell\) is surjective. For an
integer \(\Delta\ge1\), let \(t(\omega)\) be the law of \(X_{k+\Delta}\) given
\((X_0,\dots,X_k)=\omega\). The Markov property says
\(t(\omega)=P^{\Delta}(x_k,\cdot)\), which implies
\(\operatorname{LocalToDomain}(\ell,t)\) with \(\bar t(x)=P^{\Delta}(x,\cdot)\). Take \(\sigma=t\), all quotients of \(H\)
as the declared class, and strict coarsening as \(<\). Then \(\ell\)
is sufficient and is minimal exactly when distinct states in \(L\)
have distinct \(\Delta\)-step rows; otherwise equality of those rows
gives a coarser sufficient quotient. This descent condition concerns the \(\Delta\)-step law of the unlumped state; lumpability asks that the lumped process itself be Markov, and the two conditions can coincide \citep{Markov1906Investigation}.

\item[\emph{Sheaf locality on a topological space
(algebraic geometry).}] \textit{Special case.}
For a sheaf \(\mathcal F\) on a topological space
\(X\), the history set consists of pairs
\((U,s)\) of an open subset and a section; the
local quotient sends \((U,s)\) to \(U\) together with its family of
germs \((s_x)_{x\in U}\), \(s_x\in\mathcal F_x\), corestricted to its image so that \(\ell\) is surjective.  The target
readout is the section.  The sheaf condition
states that two sections agreeing on all germs
coincide, so the section descends through the germ-family quotient (gluing is not needed, since every history is an actual section).  Take \(\sigma=t\), all quotients of \(H\) as
the declared class and strict coarsening as \(<\). Since the family of
germs and the section determine each other, \(\ell\) and the readout have
the same fibers, so \(\ell\) is a minimal domain by the last clause of
\Cref{thm:F45} \citep{Grothendieck1957Tohoku}.

\item[\emph{Cellular automaton local-update rule
(theoretical computer science).}] \textit{Special case.}
For a cellular automaton on \(\mathbb Z^d\) with
finite state space, the history set is global
configurations; the local quotient sends a
configuration to its restriction to a radius-\(r\)
neighbourhood of a cell.  The local update rule
computes the next state from the radius-\(r\)
neighbourhood; the next-state readout descends
through the neighbourhood quotient, so the
radius-\(r\) neighbourhood is a sufficient domain.
Take \(\sigma\) the next-state readout, all quotients of the history set as the declared class and strict coarsening as \(<\); by the last clause of \Cref{thm:F45}, a local quotient is minimal exactly when it has the fibers of \(\sigma\). The minimal quotient identifies neighbourhood
patterns with the same next state; for an XOR rule it
is the parity of the neighbourhood
\citep{vonNeumann1966SelfReproducing}.

\item[\emph{Lamport happens-before relation
(distributed systems).}] \textit{Special case.}
For a distributed system of asynchronous processes exchanging
messages, fix an event set containing at least three events. The history set consists of ordered pairs \((e_1,e_2)\) of distinct events; the local quotient \(\ell\) sends such a pair to its pair of causal pasts \((\downarrow e_1,\downarrow e_2)\), corestricted to its image, with \(\downarrow e:=\{f:f\to e\}\cup\{e\}\), where \(\to\) is Lamport's relation, the smallest transitive relation in which events within one process are ordered by occurrence and every message-send precedes its matching receive. Let \(\sigma(e_1,e_2)\) be the verdict \(e_1\to e_2\), \(e_2\to e_1\) or \(e_1\parallel e_2\) (concurrent). Since \(e_1\to e_2\) exactly when \(e_1\in\downarrow e_2\), the verdict descends through \(\ell\). Take all quotients of the history set as the declared class and strict coarsening as \(<\). Because \(e\) is the greatest element of \(\downarrow e\), \(\ell\) is injective, while \(\sigma\) takes at most three values on at least six pairs; so by the last clause of \Cref{thm:F45} \(\ell\) is sufficient but not minimal, and the minimal domain is the verdict itself \citep{Lamport1978HappensBefore}.

\item[\emph{Voronoi cell as nearest-neighbour
domain (computational geometry).}] \textit{Special case.}
For a nonempty finite indexed set of distinct seed points
\(P=\{p_1,\ldots,p_m\}\subset\mathbb R^d\), the history set is the ambient
space; the local quotient sends each point to the
index of its nearest seed, ties broken by the smallest index.  The local
domain is the cell \(\{x:\text{the tie-broken nearest seed is }p_i\}\),
contained in the Voronoi cell
\(V_i=\{x:\|x-p_i\|\le\|x-p_j\|\,\forall j\}\), and
the nearest-neighbour classification descends
through this quotient.  On interior points of these cells the verdict is
robust to small perturbations. At distance ties the declared rule gives a
unique verdict, although it need not be robust to perturbations.
Take \(\sigma\) the verdict readout, all quotients of the history set as the declared class and strict coarsening as \(<\); by the last clause of \Cref{thm:F45}, a local quotient is minimal exactly when it has the fibers of \(\sigma\). The minimal domain is the partition into the tie-broken cells
\citep{Voronoi1908Nouvelles}.
\end{description}
\endgroup

\subsection{Law--State Split [\texorpdfstring{\Fref{F46}}{F46}]}\label{sec:self-reference-limits:law-state-split}

\paragraph{Reading the statement.} The symbols are those of \Cref{thm:F46} below.
By clause (ii), for a readout that descends the split does not depend on
\(\mu\): a law readout is a readout constant on \(\mathcal X\), and \(\mu\) enters
only through descent, uniqueness of \(\bar r\), and clauses (iii) and (iv).
Taking every closure lawful and \(c\sim d:\Longleftrightarrow\mu(c)=\mu(d)\)
in \Cref{thm:F26}, a readout that descends through \(\mu\) is a moduli
readout, \(\operatorname{LawReadout}\) is \(\operatorname{NecessaryReadout}\),
and for a descending readout \(\operatorname{StateDependent}\) is
\(\operatorname{ContingentReadout}\). Clause~(i) below is the first half of
clause~(i) of \Cref{thm:F26} with uniqueness of \(\bar r\) added; clause~(ii)
contains clause~(ii) there; clause~(iv) is clause~(iv) there with \(\sim'\)
the kernel of \(\mu_2\); and clause~(iii) evaluates the descended readout at
the selected modulus \(m_\ast\).

\begin{theorem}[Law--State Split]\label{thm:F46}
Let \(\mathcal X\) be a closure space with surjective moduli quotient
\(\mu:\mathcal X\to M\) and readout \(r:\mathcal X\to V\). Let
\(m_\ast\in M\) be a selected modulus, and let
\(\tau:\mathcal X\to\mathcal X_2\) be a transport into a second closure
space with moduli quotient \(\mu_2:\mathcal X_2\to M_2\). The readout
\(r\) \emph{descends} through \(\mu\) when \(r=\bar r\circ\mu\) for some
\(\bar r:M\to V\), its \emph{descended readout}. For a readout that descends,
define
\begin{align*}
  \operatorname{LawReadout}(r)
  &:\Longleftrightarrow
  \forall m,m'\in M,\ \bar r(m)=\bar r(m'),\\
  \operatorname{StateDependent}(r)
  &:\Longleftrightarrow \neg\operatorname{LawReadout}(r).
\end{align*}
Then:
\begin{enumerate}[label=(\roman*)]
\item \(r\) descends through \(\mu\) if and only if no pair \(x,y\) has
  \(\mu(x)=\mu(y)\) and \(r(x)\ne r(y)\), and then its descended readout
  \(\bar r\) is unique;
\item if \(r\) descends, \(\operatorname{LawReadout}(r)\) holds if and only if no
  pair \(m,m'\in M\) has \(\bar r(m)\ne\bar r(m')\), if and only if \(r\) is
  constant on \(\mathcal X\), and exactly one of
  \(\operatorname{LawReadout}(r)\) and \(\operatorname{StateDependent}(r)\)
  holds; moreover \(\operatorname{StateDependent}(r)\) holds if and only if
  some \(x,y\in\mathcal X\) have \(\mu(x)\ne\mu(y)\) and \(r(x)\ne r(y)\);
\item if \(r\) descends, every \(x\) with \(\mu(x)=m_\ast\) has
  \(r(x)=\bar r(m_\ast)\): the selected state fixes the readout;
\item \(\mu_2\circ\tau\) descends through \(\mu\) if and only if no pair
  \(x,y\) has \(\mu(x)=\mu(y)\) and \(\mu_2(\tau x)\ne\mu_2(\tau y)\).
\end{enumerate}
\end{theorem}
\begin{proof}
(i) and (iv), with the uniqueness of \(\bar r\), are the descent
criterion of \Cref{sec:apparatus} for the surjective quotient \(\mu\),
applied to \(r\) and to \(\mu_2\circ\tau\). (ii) The first equivalence restates
the definition with the order of quantifier and negation exchanged. For the
second, if \(\bar r\) is constant then \(r=\bar r\circ\mu\) is constant; if
\(r\) is constant and \(m=\mu(x)\), \(m'=\mu(y)\), then
\(\bar r(m)=r(x)=r(y)=\bar r(m')\). Either \(\bar r\) is constant on \(M\), and
\(r\) is a law-readout, or there exist \(m,m'\) with
\(\bar r(m)\ne\bar r(m')\), and the value depends on the selected moduli
class; \(\operatorname{StateDependent}(r)\) is defined as the second case, so
exactly one holds. If \(\bar r(m)\ne\bar r(m')\), choose \(x,y\) with
\(\mu(x)=m\), \(\mu(y)=m'\); then \(r(x)\ne r(y)\) and \(\mu(x)\ne\mu(y)\).
Conversely such \(x,y\) give \(\bar r(\mu x)=r(x)\ne r(y)=\bar r(\mu y)\). (iii) If \(\mu(x)=m_\ast\), then
\(r(x)=\bar r(\mu(x))=\bar r(m_\ast)\).
\end{proof}

\subsection*{Layer instantiations}\label{layer-inst:F46}

\begingroup\small
\noindent In each item the constancy of the readout is the cited classical
result, supplied as input; \Fref{F46} then classifies that readout as a law
readout and the listed coordinates as state-dependent. It does not reprove the
cited law. In each item \(M\) is the image of the moduli quotient, so that \(\mu\) is surjective as \Cref{thm:F46} requires and the descended readout is unique.
\begin{description}
\item[\emph{Clapeyron ideal gas law (thermodynamics).}] \textit{Special case.}
For the ideal-gas model, the closure space is the
set of states \((P,V,T,n)\) with \(P,V,T,n>0\) and \(PV=nRT\), where
\(R\approx8.314\;\mathrm{J\,mol^{-1}\,K^{-1}}\) is the molar gas constant; the
moduli quotient sends a state to its \((P,V,T,n)\) coordinates; this map is
injective, so every readout descends, and the item illustrates clause (ii) only.  The readout
\(PV/(nT)\) equals \(R\) on every state of the model, so the descended readout is
constant on the moduli space and is a law-readout
in the \Fref{F46} sense.  The state-dependent readouts
(individual \(P\), \(V\), \(T\), \(n\)) vary across
moduli \citep{Clapeyron1834Memoire}.

\item[\emph{Hardy-Weinberg principle (population
genetics).}] \textit{Special case.}
For a randomly mating diploid population with two
alleles at frequencies \(p,q\) under the
Hardy-Weinberg assumptions, the closure space is
the set of equilibrium populations and the moduli
quotient collapses populations to their
allele-frequency pair \((p,q)\).  The residual readout
\((f_{AA}-p^2,\,f_{Aa}-2pq,\,f_{aa}-q^2)\) equals
\((0,0,0)\) on every equilibrium population: the
descended readout is constant and is a law-readout.
The individual genotype frequencies
\(f_{AA},f_{Aa},f_{aa}\) vary across moduli and are
state-dependent \citep{Hardy1908Mendelian}.

\item[\emph{Pythagorean theorem (Euclidean
geometry).}] \textit{Special case.}
For right triangles in Euclidean \(\mathbb R^2\),
the moduli quotient sends a triangle to its leg
lengths and hypotenuse \((a,b,c)\).  The Pythagorean
readout \(c^2-a^2-b^2\) equals \(0\) on every right
triangle; the descended readout is constant and is
a law-readout.  Triangle area \(ab/2\), the two acute
angles, and perimeter \(a+b+c\) are state-dependent
\citep{Euclid300BCElements}.

\item[\emph{Coulomb's law (electromagnetism).}] \textit{Special case.}
For a system of nonzero point charges in vacuum, the
moduli quotient sends configurations to charge values and pairwise
distances.  For each pair, with \(F\) the signed pairwise force
between charges 1 and 2 (positive for repulsion), the readout
\(F r_{12}^2/(q_1 q_2)\) equals
\(k_e=1/(4\pi\varepsilon_0)\approx8.988\times 10^9\;\mathrm{N\,m^2\,C^{-2}}\)
on every configuration of the electrostatic model of nonzero point charges at
rest in vacuum; the
descended readout is constant and is a law-readout.
Individual force magnitudes, charge values, and
distances are state-dependent
\citep{Coulomb1785Memoires}.

\item[\emph{Bayes' rule for probability updating
(probability theory).}] \textit{Special case.}
Let the closure space be the set of quadruples
\((P(H),P(D\mid H),P(D),P(H\mid D))\in[0,1]^4\) with \(P(H)>0\) and \(P(D)>0\) arising from
events \(H,D\) of some probability space, and let the moduli quotient send a
quadruple to its first three coordinates (prior, likelihood, data).  The Bayes residual
readout \(P(H\mid D)\,P(D)-P(D\mid H)\,P(H)\)
equals zero on every valid tuple; the descended
readout is constant and is a law-readout.  The
specific posterior, prior, and data values are
state-dependent \citep{Bayes1763Essay}.
\end{description}
\endgroup

\subsection{Fine-Tuning as Small Selector Region [\texorpdfstring{\Fref{F47}}{F47}]}\label{sec:self-reference-limits:fine-tuning}

The canonical physical analogy for a small target-compatible region
inside a moduli space is Weinberg's cosmological-constant analysis
\citep{Weinberg1989CosmologicalConstant}.  The layer-agnostic
statement below abstracts the smallness test from that substrate
discussion.

\paragraph{Reading the statement.} The symbols are those of \Cref{thm:F47} below.
\(\operatorname{TargetImpossible}(\Sigma(r,T))\) says that the selector
region is empty: no modulus reaches the target region. The declarations
of moduli space, target, measure and threshold are part of the setting and
add no hypothesis.

\begin{theorem}[Fine-Tuning as Small Selector Region]\label{thm:F47}
Let \(M\) be a moduli space, \(r:M\to V\) a readout, and
\(T:V\to\mathrm{Prop}\) a target-region predicate.  Let
\(\mathcal M\) be the carrier type of measure-record values, let
\(\mu:\mathrm{Pow}(M)\to\mathcal M\) be a measure record on subsets
of \(M\), let \(\mu_{\mathrm{tot}}\in\mathcal M\) be the total
measure, let \(\theta\in\Theta\) be a threshold in a declared set \(\Theta\), and let
\(\odot:\Theta\times\mathcal M\to\mathcal M\) be the declared
threshold scaling.  Let \(\operatorname{Pos}:\mathcal M\to
\mathrm{Prop}\) and \({\le_{\mathcal M}}\subseteq
\mathcal M\times\mathcal M\) be the declared positivity and order
predicates on the measure record.  Let \(\mathrm{sel}\in M\) be a
selected modulus, let \(\mathsf{selDec}\) be a selector-declared
predicate.  Define the selector region
\(\Sigma(r,T):=\{m\in M:T(r(m))\}\) and
\begin{align*}
  \operatorname{Small}(\Sigma(r,T))
  &:\Longleftrightarrow
  \operatorname{Pos}(\mu(\Sigma(r,T)))
  \wedge
  \mu(\Sigma(r,T))\le_{\mathcal M}\theta\odot\mu_{\mathrm{tot}},\\
  \operatorname{Realized}(\Sigma(r,T),\mathrm{sel},\mathsf{selDec})
  &:\Longleftrightarrow
  \mathsf{selDec}
  \wedge
  \mathrm{sel}\in\Sigma(r,T)
  \wedge
  \operatorname{Small}(\Sigma(r,T)),\\
  \operatorname{TargetImpossible}(\Sigma(r,T))
  &:\Longleftrightarrow
  \Sigma(r,T)=\emptyset.
\end{align*}
Then \(m\in\Sigma(r,T)\) if and only if \(T(r(m))\), and
\[
  \operatorname{TargetImpossible}(\Sigma(r,T))
  \Longrightarrow
  \neg\operatorname{Realized}(\Sigma(r,T),\mathrm{sel},\mathsf{selDec}):
\]
fine-tuning cannot be realized in a selector region that no modulus
reaches, whatever the measure record says; this uses only the conjunct
\(\mathrm{sel}\in\Sigma(r,T)\) of \(\operatorname{Realized}\). Moreover, realized fine-tuning
gives \(\operatorname{Pos}(\mu(\Sigma(r,T)))\) and \(\Sigma(r,T)\ne\emptyset\);
and if \(\mu\) is monotone (\(\Sigma\subseteq\Sigma'\) implies
\(\mu(\Sigma)\le_{\mathcal M}\mu(\Sigma')\)), \(\le_{\mathcal M}\) is transitive
and \(T\) implies \(T'\), then
\(\operatorname{Pos}(\mu(\Sigma(r,T)))\) and \(\operatorname{Small}(\Sigma(r,T'))\)
give \(\operatorname{Small}(\Sigma(r,T))\): a stricter target keeps a small
selector region small. If moreover \(\mathcal M=\mathbb R\) with its usual order,
\(\mu(M)=\mu(\Sigma(r,T))+\mu(\Sigma(r,\neg T))=\mu_{\mathrm{tot}}>0\),
\(\operatorname{Pos}(x):\Leftrightarrow x>0\), \(\Theta\subseteq\mathbb R\), \(\odot\) is multiplication and
\(\theta<\tfrac12\), then \(\operatorname{Small}(\Sigma(r,T))\) and
\(\operatorname{Small}(\Sigma(r,\neg T))\) do not both hold.
\end{theorem}
\begin{proof}
The membership statement is the definition of
\(\Sigma(r,T)\) as the preimage \(r^{-1}(T)\).  If
\(\operatorname{TargetImpossible}(\Sigma(r,T))\), then no \(m\in M\) lies in
\(\Sigma(r,T)\); in particular \(\mathrm{sel}\notin\Sigma(r,T)\), which
contradicts the second conjunct of \(\operatorname{Realized}\).  Hence
\(\neg\operatorname{Realized}(\Sigma(r,T),\mathrm{sel},\mathsf{selDec})\).
Realized fine-tuning contains \(\operatorname{Small}(\Sigma(r,T))\), hence
\(\operatorname{Pos}(\mu(\Sigma(r,T)))\), and \(\mathrm{sel}\in\Sigma(r,T)\).
If \(T\) implies \(T'\), then \(\Sigma(r,T)\subseteq\Sigma(r,T')\), so
monotonicity gives \(\mu(\Sigma(r,T))\le_{\mathcal M}\mu(\Sigma(r,T'))\), and
\(\operatorname{Small}(\Sigma(r,T'))\) gives
\(\mu(\Sigma(r,T'))\le_{\mathcal M}\theta\odot\mu_{\mathrm{tot}}\);
transitivity and the assumed positivity give
\(\operatorname{Small}(\Sigma(r,T))\).
For the quantitative clause, the two regions partition \(M\), so both being
small would give \(\mu_{\mathrm{tot}}\le2\theta\mu_{\mathrm{tot}}<\mu_{\mathrm{tot}}\).
\end{proof}

\subsection*{Layer instantiations}\label{layer-inst:F47}

\begingroup\small
\begin{description}
\item[\emph{Circumstellar habitable zone (astrobiology).}] \textit{Analogy.}
The moduli space is the set of orbital and
atmospheric parameter values for a planet (semi-major
axis, eccentricity, atmospheric composition, greenhouse
mass column); the readout is the equilibrium surface
temperature; the target predicate restricts to the
liquid-water range.  The selector region is the circumstellar habitable zone.
Whether it occupies a small fraction of the parameter space
depends on a declared measure and threshold satisfying
\(\operatorname{Small}\). Earth meets the liquid-water target
criterion; realizing the fine-tuning predicate additionally
requires those measure and threshold conditions
\citep{Kasting1993Habitable}.

\end{description}
\endgroup


\subsection{Topology as Global Gluing Invariant [\texorpdfstring{\Fref{F48}}{F48}]}\label{sec:self-reference-limits:topology}

The motivating mathematical example is sheaf theory, in which
topological information emerges as the consistent global gluing of
locally defined data; the standard reference is
\citet{MacLaneMoerdijk1992Sheaves}.

\begin{theorem}[Topology as Global Gluing Invariant]\label{thm:F48}
Let \(G\) be a set of global states with gluing quotient
\(\gamma:G\to\Gamma\), readout \(r:G\to V\), transport
\(\tau:G\to G'\), and target gluing quotient \(\gamma':G'\to\Gamma'\).
Assume \(\gamma\) is surjective. Write
\(\operatorname{TopologicalReadout}(r)\) when \(r\) factors through \(\gamma\),
that is \(r=\bar r\circ\gamma\) for some \(\bar r:\Gamma\to V\), and define
\[
  \mathcal O_\gamma(r)=
  \{(g,g'):\gamma(g)=\gamma(g')\ \text{and}\ r(g)\ne r(g')\},
\]
and define
\begin{align*}
  \operatorname{CompleteInvariant}(r)
  &:\Longleftrightarrow \operatorname{TopologicalReadout}(r)\wedge
    \forall g,g',\ r(g)=r(g')\Rightarrow \gamma(g)=\gamma(g'),\\
  \operatorname{CoarseInvariant}(r)
  &:\Longleftrightarrow \operatorname{TopologicalReadout}(r)\wedge
    \exists g,g',\gamma(g)\ne\gamma(g')\wedge r(g)=r(g'),\\
  \operatorname{TransportDescends}(\tau)
  &:\Longleftrightarrow
    \exists\,\bar\tau:\Gamma\to\Gamma'\ \text{ with }\
    \bar\tau\circ\gamma=\gamma'\circ\tau .
\end{align*}
Then
\[
  \operatorname{TopologicalReadout}(r)\iff\mathcal O_\gamma(r)=\emptyset,
\]
every topological readout is exactly one of a complete and a coarse
invariant, and \(\operatorname{TransportDescends}(\tau)\) holds if and only if
\(\mathcal O_\gamma(\gamma'\circ\tau)=\emptyset\).
\end{theorem}
\begin{proof}
Since \(\gamma\) is surjective, the first equivalence is the descent
criterion of \Cref{sec:apparatus}.  If \(r\) also separates \(\Gamma\)-classes, it is complete; if two
distinct gluing classes have the same \(r\)-value, it is coarse. These
are the two cases
\[
  \forall g,g',\ r(g)=r(g')\Rightarrow\gamma(g)=\gamma(g')
  \quad\text{or}\quad
  \exists g,g',\gamma(g)\ne\gamma(g')\wedge r(g)=r(g'),
\]
and each is the negation of the other, so a topological readout satisfies
exactly one of them. The transport clause is the same descent criterion applied to
\(\gamma'\circ\tau\).
\end{proof}

\subsection*{Layer instantiations}\label{layer-inst:F48}

\begingroup\small
\begin{description}
\item[\emph{Euler characteristic (algebraic
topology).}] \textit{Special case.}
For a finite cell complex, the gluing quotient
collapses cell-complex presentations to their
underlying topological space.  The Euler
characteristic
\(\chi=V-E+F\) (alternating sum of cell counts in
higher dimensions) is invariant under
cell-decomposition refinement and homotopy, hence a
TopologicalReadout.  It is a coarse invariant: the
torus and the Klein bottle both have \(\chi=0\) but
the torus is orientable and the Klein bottle is
not.  For the transport clause take \(\tau\) the identity of presentations
and \(\gamma'\) the map to homotopy types; homeomorphic spaces are homotopy
equivalent, so \(\operatorname{TransportDescends}(\tau)\) holds
\citep{Euler1750Polyhedron}.

\item[\emph{Genus of a closed orientable surface
(topology).}] \textit{Special case.}
For closed connected orientable 2-manifolds, the gluing
quotient collapses polygon-with-edge-identification
presentations to homeomorphism classes.  The genus
\(g(\Sigma)\) — the number of handles, equivalently
\((2-\chi)/2\) — is a TopologicalReadout \citep{Riemann1857Theorie}, and the later
classification theorem of closed surfaces makes
genus a CompleteInvariant for closed connected
orientable surfaces.

\item[\emph{Fundamental group (algebraic topology).}] \textit{Special case.}
For based spaces drawn from a fixed set (for instance, finite based CW
complexes), the gluing quotient sends a space to its based homotopy type.
The isomorphism class of the fundamental group \(\pi_1(X,x_0)\) — homotopy
classes of based loops — is homotopy-invariant,
hence a TopologicalReadout.  It is coarse: the
2-sphere and 3-sphere both have trivial \(\pi_1\)
but are not homotopy-equivalent.  Continuous
basepoint-preserving maps induce \(\pi_1\)-compatible
group homomorphisms \citep{Poincare1895Analysis}.

\item[\emph{{DNA} linking number of closed circular
{DNA} (molecular biology).}] \textit{Special case.}
Model a closed circular DNA duplex as an ordered pair of disjoint smooth oriented embedded circles in \(\mathbb R^3\), its two strands; the gluing quotient collapses pairs related by smooth ambient isotopy of \(\mathbb R^3\) preserving the strand orientations (no strand breaks).  The linking
number \(\mathrm{Lk}\in\mathbb Z\) — the signed
Gauss linking integral counting strand wrappings
— is invariant under smooth ambient isotopy preserving the strand orientations, hence a
TopologicalReadout.  For a duplex represented by the centreline and offset edge of an embedded closed ribbon, with compatible orientations, the Călugăreanu-White formula
\(\mathrm{Lk}=\mathrm{Tw}+\mathrm{Wr}\) decomposes
\(\mathrm{Lk}\) into non-topological twist and
writhe; only their sum is topological.
Topoisomerase enzymes implement the transport that
changes \(\mathrm{Lk}\) by strand cleavage and
religation \citep{BatesMaxwell2005DNATopology}.

\item[\emph{Reeb graph from {M}orse-function level
sets (geometric topology).}] \textit{Analogy.}
For a smooth manifold with a Morse function \(f\),
the gluing quotient collapses configurations
continuously deformable along the Morse flow.  The
Reeb graph quotients the manifold by the
``connected component of \(f^{-1}(c)\)'' relation,
encoding topology changes of level sets at
critical values.  Within the declared
function-isotopy class, the Reeb graph is a
TopologicalReadout, and a coarse invariant
(distinct Morse functions can yield isomorphic
Reeb graphs) \citep{Reeb1946Points}.
\end{description}
\endgroup

\subsection{Common-Source / Nonlocal Correlation Normal Form [\texorpdfstring{\Fref{F49}}{F49}]}\label{sec:self-reference-limits:common-source}

The Einstein--Podolsky--Rosen / Bell situation \citep{Bell1964EPR}, of correlations that admit no local common-source explanation, is an analogy, not a special case: Bell's theorem concerns probabilistic local models, whereas \(\operatorname{CommonSource}\) below is a deterministic factorization; the layer-agnostic statement below
records the descent-to-common-source obstruction abstracted from
that substrate.

\paragraph{Reading the statement.} The symbols are those of \Cref{thm:F49} below.
The locality package \(\operatorname{Loc}\) is declared data: the
theorem uses it only through its failure, which is what makes a
correlation residual nonlocal. \(\operatorname{Corr}\), by contrast, is defined from the
data: the joint readout is not a function of the two local access classes. Surjectivity of \((\sigma,q_A,q_B)\), which the theorem does not assume but a model may impose, says that every source class occurs with every pair of local access classes; it is the deterministic
form of measurement independence, and it excludes
the trivial source: for injective \(\sigma\) (for example
\(\sigma=\mathrm{id}_H\)) every fiber of \((\sigma,q_A)\) is a point, so
\(\operatorname{CommonSource}(\sigma)\) holds automatically, and surjectivity
then forces \(L_A\) and \(L_B\) to be singletons; hence when \(L_A\) or \(L_B\)
has two elements, \(\sigma\) is not injective. The
theorem tests the declared \(\sigma\), \(\iota\) and \(\rho\); it does not
assert that no source quotient explains the correlation.

\begin{theorem}[Common-Source / Nonlocal Correlation Normal Form]\label{thm:F49}
Let \(H\) be a nonempty history space with source quotient
\(\sigma:H\to S\), shared interface \(\iota:H\to I\), local access
quotients
\(q_A:H\to L_A\), \(q_B:H\to L_B\), and readouts
\(r_A:H\to A\), \(r_B:H\to B\), with joint readout
\(r=(r_A,r_B):H\to A\times B\). Let a declared direct route consist of
two declared propositions, its admissibility record and its stability record, and a route quotient
\(\rho:H\to P\). Let \(\operatorname{Corr}:\Longleftrightarrow\exists h,h'\) with
\((q_A,q_B)(h)=(q_A,q_B)(h')\) and \(r(h)\ne r(h')\), that is, the joint
readout is not fixed by the two local access classes, and let
\(\operatorname{Loc}\) be the declared locality package, a proposition the
theorem does not analyse. \(\operatorname{CommonSource}(\sigma)\) below is a
deterministic factorization; probabilistic common-cause models are not
covered. No surjectivity is assumed. Write \(\operatorname{CommonSource}(\sigma)\) when \(r_A\) is
constant on the fibers of \((\sigma,q_A)\) and \(r_B\) is constant on the
fibers of \((\sigma,q_B)\); \(\operatorname{InterfaceFactorization}(\iota)\)
when \(r\) is constant on the fibers of \((\iota,q_A,q_B)\); and
\(\operatorname{DirectRoute}\) when both records hold and
\(r\) is constant on the fibers of \(\rho\), and define
\[
  \operatorname{LocallyExplainable}
  :\iff \operatorname{CommonSource}(\sigma)\vee\operatorname{DirectRoute}
    \vee \operatorname{InterfaceFactorization}(\iota).
\]
Then
\begin{align*}
  \operatorname{CommonSource}(\sigma)
  &\iff
  \exists\,\bar r_A:S\times L_A\to A,\ \bar r_B:S\times L_B\to B,\\
  &\hphantom{{}\iff{}}
    \bar r_A\circ(\sigma,q_A)=r_A\ \wedge\ \bar r_B\circ(\sigma,q_B)=r_B,\\
  \operatorname{InterfaceFactorization}(\iota)
  &\iff
  \exists\,\hat r:I\times L_A\times L_B\to A\times B,\
    \hat r\circ(\iota,q_A,q_B)=r,\\
  \operatorname{DirectRoute}
  &\iff
  \text{both records hold, and}\
    \exists\,\tilde r:P\to A\times B,\ \tilde r\circ\rho=r.
\end{align*}
Call \(\operatorname{Corr}\wedge\neg\operatorname{LocallyExplainable}\) a
\emph{correlation residual}, and a \emph{nonlocal residual} when moreover
\(\neg\operatorname{Loc}\). By the three equivalences, a correlation residual
is present exactly when \(\operatorname{Corr}\) holds, neither the common-source
nor the interface factorization exists, and it is not the case that the route
is admissible, stable and has some \(\tilde r\) with \(\tilde r\circ\rho=r\).
Moreover, \(\neg\operatorname{Corr}\) implies \(\operatorname{InterfaceFactorization}(\iota)\)
for every \(\iota\), so a correlation residual is exactly
\(\neg\operatorname{LocallyExplainable}\); and \(\operatorname{CommonSource}(\sigma)\)
implies \(\operatorname{InterfaceFactorization}\) for the interface \(\iota:=\sigma\).
\end{theorem}
\begin{proof}
Since \(H\) is nonempty, so are \(A\) and \(B\); a map on \(H\) that is constant on the fibers of a map \(p\) factors through \(p\), defined from representatives on the image of \(p\) and by a fixed value elsewhere. (Without nonemptiness the CommonSource equivalence fails: for \(H=L_B=A=\varnothing\) and \(S\times L_A\) a point, \(\operatorname{CommonSource}(\sigma)\) holds vacuously but no \(\bar r_A\) exists.) So a map on \(H\) with values in \(A\), \(B\) or \(A\times B\) factors through a map \(p\) exactly when it is constant on the fibers of \(p\). Hence \(r_A\) and \(r_B\) are constant on the
fibers of \((\sigma,q_A)\) and \((\sigma,q_B)\) exactly when response maps
\(\bar r_A\), \(\bar r_B\) as displayed exist: each local readout depends only
on the shared source class and its own access data. The
joint readout factors through \((\iota,q_A,q_B)\)
exactly when it is constant on its fibers, which gives the interface
case. The direct-route case is the same factorization argument for \(\rho\): the joint readout factors through the route
quotient exactly when it is constant on its fibers, and a declared route
explains the correlation when, in addition, its admissibility and
stability records hold. A correlation is
locally explained exactly when one of the three explanations is
supplied. If a correlation is present and none is supplied, it is
recorded as a residual, not as an explained relation; it is a nonlocal
residual only relative to a failed locality package.
If \(\neg\operatorname{Corr}\), then \(r\) is constant on the fibers of \((q_A,q_B)\), hence on the finer fibers of \((\iota,q_A,q_B)\). Under \(\operatorname{CommonSource}(\sigma)\), \(r_A\) and \(r_B\) are constant on the fibers of \((\sigma,q_A)\) and \((\sigma,q_B)\), so \(r=(r_A,r_B)\) is constant on the fibers of \((\sigma,q_A,q_B)\).
\end{proof}

\subsection*{Layer instantiations}\label{layer-inst:F49}

\begingroup\small
\begin{description}
\item[\emph{{MISTRA} twin studies (behavioral
genetics).}] \textit{Analogy.}
For the Minnesota Study of Twins Reared Apart, the
source quotient sends a twin pair to its
shared-genotype class; the readouts are conditional
trait
distributions
\(P(\mathrm{trait}\mid\sigma)\) under a declared
quantitative-genetic model.  The CommonSource
hypothesis holds when conditional distributions
descend through the shared-genotype quotient; the
interface captures shared rearing environment
(absent for the reared-apart sample).  The reared-apart monozygotic
correlation \(r_{MZA}\), which estimates heritability directly,
operationalises the \Fref{F49} LocallyExplainable decomposition; the
Falconer formula \(h^2=2(r_{MZ}-r_{DZ})\) is the counterpart for
twins reared together
\citep{Bouchard1990Sources}.

\end{description}
\endgroup


\subsection{Vacuum / Background Selection [\texorpdfstring{\Fref{F50}}{F50}]}\label{sec:self-reference-limits:vacuum}

The most discussed contemporary substrate analogy is the string
landscape of vacua surveyed in \citet{Susskind2003Landscape};
the layer-agnostic statement below records the selector-singleton
shape without committing to that physical realization.

\begin{theorem}[Vacuum / Background Selection]\label{thm:F50}
Let \(\mathcal X\) be a closure space with moduli quotient
\(\mu:\mathcal X\to M\), raw background \(b_0:\mathcal X\to B\),
a declared background \(b:M\to B\), target \(t:M\to T\), and vacuum
criterion \(V\subseteq M\), and a declared set \(T_{\mathrm{f}}\subseteq T\)
of formed targets. Assume \(\mu\) is surjective. Define
\begin{align*}
  \operatorname{NecessaryBackground}(b)
  &:\Longleftrightarrow \forall m,m'\in M,\ b(m)=b(m'),\\
  \operatorname{BackgroundModuli}(b)
  &:\Longleftrightarrow \exists m,m'\in M,\ b(m)\ne b(m'),\\
  \operatorname{Vacuum}(m_\ast)
  &:\Longleftrightarrow V=\{m_\ast\}\ \wedge\ t(m_\ast)\in T_{\mathrm{f}}.
\end{align*}
Then
\[
  \bigl(\forall x,y\in\mathcal X,\ \mu(x)=\mu(y)\Rightarrow b_0(x)=b_0(y)\bigr)
  \iff
  \exists\,\beta:M\to B\ \text{ with } \beta\circ\mu=b_0.
\]
The necessary/moduli alternatives are exhaustive and disjoint. If \(b\) is
a descended raw background, \(b\circ\mu=b_0\), then
\[
  \operatorname{NecessaryBackground}(b)
  \iff \forall x,y\in\mathcal X,\ b_0(x)=b_0(y).
\]
Exactly one of \(V=\emptyset\), \(V\) a singleton, and \(V\) containing two
distinct elements holds. Some \(m_\ast\) satisfies
\(\operatorname{Vacuum}(m_\ast)\) if and only if \(V=\{m\}\) for some \(m\)
with \(t(m)\in T_{\mathrm{f}}\), and such an \(m_\ast\) is unique. If
\(\operatorname{Vacuum}(m_\ast)\), then under \(\operatorname{BackgroundModuli}(b)\) some modulus outside \(V\) carries a background different from \(b(m_\ast)\); if moreover \(b\circ\mu=b_0\), then \(b_0(x)=b(m_\ast)\) for
every \(x\) with \(\mu(x)=m_\ast\).
\end{theorem}
\begin{proof}
Since \(\mu\) is surjective, the descent clause is the descent criterion
of \Cref{sec:apparatus} for \(\mu\) and \(b_0\).  For the declared background \(b\), either \(b\) is constant on \(M\)
or there exist \(m,m'\) with \(b(m)\ne b(m')\):
\[
  \bigl(\forall m,m', b(m)=b(m')\bigr)
  \quad\vee\quad
  \bigl(\exists m,m', b(m)\ne b(m')\bigr).
\]
These alternatives are mutually exclusive. For the vacuum criterion,
either \(V\) is empty, or it has an element \(m\); in the second case either
every element of \(V\) equals \(m\), so \(V=\{m\}\), or some element differs
from \(m\), so \(V\) contains two distinct elements. A singleton has no two
distinct elements and a nonempty set is not empty, so exactly one case holds.
By definition, \(\operatorname{Vacuum}(m_\ast)\) says that \(V=\{m_\ast\}\) and
\(t(m_\ast)\in T_{\mathrm{f}}\), which gives the existence equivalence; if
\(\operatorname{Vacuum}(m_\ast)\) and \(\operatorname{Vacuum}(m'_\ast)\), then
\(\{m_\ast\}=V=\{m'_\ast\}\), so \(m_\ast=m'_\ast\).
If \(b\circ\mu=b_0\), then \(b_0(x)=b(\mu x)\), so \(b_0\) is constant when
\(b\) is; conversely every \(m\in M\) is \(\mu(x)\) for some \(x\), so \(b\) is
constant when \(b_0\) is. For the last clause, \(b_0(x)=b(\mu x)=b(m_\ast)\) when
\(\mu(x)=m_\ast\); and if \(b(m)\ne b(m')\), one of \(m,m'\) has background
different from \(b(m_\ast)\), and it is not \(m_\ast\), the only member of \(V\).
\end{proof}

\subsection*{Layer instantiations}\label{layer-inst:F50}

\begingroup\small
\begin{description}
\item[\emph{Waddington epigenetic landscape
(developmental biology).}] \textit{Analogy.}
For a developing embryo, the moduli quotient
collapses stem-cell trajectory configurations
sharing the same cell-type endpoint; the
descended background records the cell-fate class.
Waddington's epigenetic landscape depicts cell-fate
decision as a ball rolling down a branching
landscape: each branching point is a vacuum-style
selection of one cell-fate from a set of accessible
attractors \citep{Waddington1957Strategy}.

\end{description}
\endgroup


\subsection{Unification as Common Refinement [\texorpdfstring{\Fref{F51}}{F51}]}\label{sec:self-reference-limits:unification}

The motivating particle-physics analogy is the electroweak unification
introduced in \citet{Weinberg1967Leptons}.  The statement below records a common-refinement shape; electroweak unification is an analogy, and no physical unification is asserted to satisfy its hypotheses.

\paragraph{Reading the statement.}
\emph{Exact unification} is defined as the conjunction displayed in the
theorem: formation of the parent closure, the three pairs of equations, role
preservation of the two child theories, and the status check of every cross
residual. Role preservation and the residual status check are declared
predicates of the apparatus. The content of the theorem is what follows from
the first pair of equations: the parent quotient is a common refinement of the
two children, in the sense of the displayed square, and, when \(q_U\) is
surjective, \(\langle\pi_A,\pi_B\rangle\) is injective exactly when the parent
carries no distinction beyond those of the two children.

\begin{theorem}[Unification as Common Refinement]\label{thm:F51}
Let \(\mathcal C_U,\mathcal C_A,\mathcal C_B\) be parent and child
closures with quotients \(q_U:\mathcal C_U\to Q_U\),
\(q_A:\mathcal C_A\to Q_A\), and \(q_B:\mathcal C_B\to Q_B\). Let
\(\iota_A:\mathcal C_U\to\mathcal C_A\) and
\(\iota_B:\mathcal C_U\to\mathcal C_B\) be interpretations, and let
\(\pi_A:Q_U\to Q_A\), \(\pi_B:Q_U\to Q_B\) be quotient projections.
Let \(R_U^A,R_U^B:Q_U\to Q_U\) and \(R_A:Q_A\to Q_A\),
\(R_B:Q_B\to Q_B\) be the declared route systems. For claim sets
\(K_A,K_B,K_U\), lifts \(\lambda_A:K_A\to K_U\),
\(\lambda_B:K_B\to K_U\), status maps \(s_A,s_B,s_U\), and status
translations \(\alpha_A,\alpha_B\) into the statuses of \(K_U\), \emph{exact
unification} is defined as the conjunction of the following: the
declared formedness proposition for the parent closure \(\mathcal C_U\) holds; the equations
\begin{align*}
  \pi_A\circ q_U=q_A\circ\iota_A
  &\quad\wedge\quad
  \pi_B\circ q_U=q_B\circ\iota_B,\\
  \pi_A\circ R_U^A=R_A\circ\pi_A
  &\quad\wedge\quad
  \pi_B\circ R_U^B=R_B\circ\pi_B,\\
  s_U(\lambda_A k)=\alpha_A(s_A k)
  &\quad\wedge\quad
  s_U(\lambda_B \ell)=\alpha_B(s_B \ell)
  \quad\text{for all }k\in K_A,\ell\in K_B
\end{align*}
hold; the declared role-preservation proposition holds; and the declared cross-residual status-check proposition holds. If
\(\pi_A\circ q_U=q_A\circ\iota_A\) and \(\pi_B\circ q_U=q_B\circ\iota_B\), in
particular under exact unification, then
\[
  \langle\pi_A,\pi_B\rangle\circ q_U
  =(q_A\times q_B)\circ\langle\iota_A,\iota_B\rangle
  :\mathcal C_U\to Q_A\times Q_B ,
\]
and the parent kernel is contained in both pulled-back child kernels: for all
\(x,x'\in\mathcal C_U\), \(q_U(x)=q_U(x')\) implies
\(q_A(\iota_A x)=q_A(\iota_A x')\) and \(q_B(\iota_B x)=q_B(\iota_B x')\).
If moreover \(q_U\) is surjective, then \(\langle\pi_A,\pi_B\rangle\) is injective
if and only if \(q_U(x)=q_U(x')\Leftrightarrow q_A(\iota_A x)=q_A(\iota_A x')
\wedge q_B(\iota_B x)=q_B(\iota_B x')\) for all \(x,x'\): the parent carries
exactly the distinctions of the two children.
\end{theorem}
\begin{proof}
The first pair of equations gives, for every
\(x\in\mathcal C_U\),
\[
  \pi_A(q_U x)=q_A(\iota_A x),
  \qquad
  \pi_B(q_U x)=q_B(\iota_B x),
\]
which is the displayed square componentwise, so \(Q_U\) is a common
refinement. If \(q_U(x)=q_U(x')\), then
\(q_A(\iota_A x)=\pi_A(q_U x)=\pi_A(q_U x')=q_A(\iota_A x')\), and likewise
for \(B\). For the last clause, let \(q_U\) be surjective. By the square,
\(q_A(\iota_A x)=q_A(\iota_A x')\wedge q_B(\iota_B x)=q_B(\iota_B x')\) says
\(\langle\pi_A,\pi_B\rangle(q_U x)=\langle\pi_A,\pi_B\rangle(q_U x')\). If the
pairing is injective, this gives \(q_U(x)=q_U(x')\). Conversely, every two
points of \(Q_U\) are \(q_U(x)\) and \(q_U(x')\) for some \(x,x'\), so the
equivalence gives injectivity.
\end{proof}

\subsection*{Layer instantiations}\label{layer-inst:F51}

\begingroup\small
\begin{description}
\item[\emph{Maxwell unification of electromagnetism
(classical physics).}] \textit{Analogy.}
The parent closure is Maxwell's electromagnetic
field theory; the children are pre-Maxwellian
electrostatics (Coulomb-Gauss) and magnetostatics
(Amp\`ere-Biot-Savart).  Setting time derivatives to
zero in Maxwell's equations yields each child as a
specialisation.  The unified parent predicts
electromagnetic waves at the speed of light,
identifying light as electromagnetic — a
prediction not derivable from either child alone
\citep{Maxwell1865Dynamical}.

\end{description}
\endgroup


\subsection{Closure Completeness Boundary [\texorpdfstring{\Fref{F52}}{F52}]}\label{sec:self-reference-limits:closure-completeness}

A proof-theoretic analogue of a closure-completeness boundary is
G\"odel's incompleteness theorem \citep{Godel1931Unentscheidbare}, where
incompleteness is derived from the arithmetic of the theory; the statement
below records a declared scope boundary instead. The modal-logic systematization of provability obstructions is
collected in \citet{Boolos1993LogicProvability}.  The layer-agnostic
meta-coherence statement below records the apparatus-side scope
discipline that complements those classical results.

\begin{definition}[Catalog boundary]\label{def:catalog-boundary}
Let \(D\) be a direction type with predicates \(\operatorname{claimed}\),
\(\operatorname{inScope}\) and \(\operatorname{namedOutside}\), as in
\Cref{thm:F52}. The \emph{catalog boundary} is
\(\partial D:=\{d:\operatorname{claimed}(d)\wedge\neg\operatorname{inScope}(d)
\wedge\operatorname{namedOutside}(d)\}\), the claimed directions explicitly
marked out of scope, for example because they require substrate-specific data
or stronger apparatus. The \emph{unmarked frontier} is
\(\{d:\operatorname{claimed}(d)\wedge\neg\operatorname{inScope}(d)\wedge
\neg\operatorname{namedOutside}(d)\}\).
\end{definition}

The theorem below states when a closure package reaches this boundary:
it is complete up to its declared scope, and every direction it claims
lies inside that scope or is named as outside it.

\paragraph{Reading the statement.}
\(\operatorname{formed}\) records that the closure package is formed in
the sense of \Cref{def:closure-formation}, and
\(\operatorname{scopeDeclared}\) records that its completeness scope has
been declared. Both are declared data of the package. Each displayed
line of the theorem unfolds the predicate on its left; together they
state what completeness up to a declared scope requires. The last sentence of the theorem, that each conjunct can fail while the other three hold, is proved by one-element examples.

\begin{theorem}[Closure Completeness Boundary]\label{thm:F52}
Let \(R\) be a residual type with predicates
\(\operatorname{hasStatus}:R\to\mathrm{Prop}\) and
\(\operatorname{adequate}:R\to\mathrm{Prop}\). Let \(D\) be a direction
type with predicates \(\operatorname{inScope}:D\to\mathrm{Prop}\),
\(\operatorname{claimed}:D\to\mathrm{Prop}\), and
\(\operatorname{namedOutside}:D\to\mathrm{Prop}\). Let
\(\operatorname{formed}\) and \(\operatorname{scopeDeclared}\) be
propositions, and let
\(v\) be a meta-audit verdict, drawn from a declared set of verdicts
containing \(\operatorname{accepted}\). Define
\begin{align*}
  \operatorname{ResidualCoverage}
  &:\Longleftrightarrow
  \forall r\in R,\ \operatorname{hasStatus}(r)\wedge\operatorname{adequate}(r),\\
  \operatorname{MetaCovered}(v)
  &:\Longleftrightarrow v=\operatorname{accepted},\\
  \operatorname{ScopedClosureComplete}
  &:\Longleftrightarrow
  \operatorname{formed}\wedge\operatorname{scopeDeclared}\wedge
  \operatorname{ResidualCoverage}\\
  &\qquad\wedge\
  \forall d\in D,\ \operatorname{claimed}(d)\Rightarrow
    \bigl(\operatorname{inScope}(d)\vee\operatorname{namedOutside}(d)\bigr),\\
  \operatorname{MetaComplete}
  &:\Longleftrightarrow \operatorname{ScopedClosureComplete}\wedge\operatorname{MetaCovered}(v).
\end{align*}
Then \(\operatorname{ScopedClosureComplete}\) holds if and only if
\(\operatorname{formed}\) and \(\operatorname{scopeDeclared}\) hold and the
three obstruction sets
\[
\begin{gathered}
  \{r:\neg\operatorname{hasStatus}(r)\},\qquad
  \{r:\neg\operatorname{adequate}(r)\},\\
  \{d:\neg\operatorname{inScope}(d)\wedge\operatorname{claimed}(d)
      \wedge\neg\operatorname{namedOutside}(d)\}
\end{gathered}
\]
are empty; the third set is the unmarked frontier of
\Cref{def:catalog-boundary}. Hence \(\operatorname{MetaComplete}\) fails
exactly when one of these conditions fails or \(v\ne\operatorname{accepted}\).
Each of the four conjuncts of \(\operatorname{ScopedClosureComplete}\) can fail
while the other three hold.
\end{theorem}
\begin{proof}
\(\operatorname{ResidualCoverage}\) says that every residual has a status and
that every residual is adequate, that is, that the first two sets are empty.
With \(\operatorname{hasStatus}(r):\Longleftrightarrow\exists\rho\in\mathrm{Rec},\ \operatorname{loc}(\rho)=r\),
its first conjunct is (RC) of \Cref{sec:apparatus}; adequacy is an additional
requirement.
For a direction \(d\) there are three cases:
\[
  \operatorname{inScope}(d),\qquad
  \operatorname{namedOutside}(d),\qquad
  \neg\operatorname{inScope}(d)\wedge\neg\operatorname{namedOutside}(d).
\]
The direction clause of \(\operatorname{ScopedClosureComplete}\) excludes
the third case for every claimed direction, which is emptiness of the third
set. The penultimate sentence follows from the definition of
\(\operatorname{MetaComplete}\). For the last, take \(R\) and \(D\) with one
element each. Making \(\operatorname{formed}\) false, or
\(\operatorname{scopeDeclared}\) false, or \(\operatorname{hasStatus}\) false
at the residual, or \(d\) claimed with \(\operatorname{inScope}(d)\) and
\(\operatorname{namedOutside}(d)\) false, while every other datum is true,
falsifies exactly one conjunct.
\end{proof}

For the catalog itself, take as \(D\) the candidate directions
considered for it, let \(\operatorname{inScope}(d)\) mean that \(d\) is a
law of the catalog or reduces to one, let
\(\operatorname{namedOutside}(d)\) mean that \(d\) is named in
\Cref{sec:sketches}, and let \(\operatorname{claimed}(d)\) hold for
every \(d\in D\).  The direction clause then requires every candidate to
be covered or named, and the residual clauses require a residual
register for the catalog, which is not given here; \Cref{sec:sketches}
discusses this application.

\subsection*{Layer instantiations}\label{layer-inst:F52}

\begingroup\small
\begin{description}
\item[\emph{{CONSORT} statement for randomised-trial
reporting (clinical research).}] \textit{Analogy.}
The residuals are the methodological items required
for adequate trial reporting; directions are
declarable claims and outcomes.  CONSORT 2010
requires 25 explicit reporting items, with primary
and secondary outcomes \emph{inScope} and
exploratory analyses \emph{namedOutside}.  Residual
coverage holds when every CONSORT item has status
and is adequate; the meta-audit is journal-editor
acceptance \citep{Schulz2010CONSORT}.

\end{description}
\endgroup


\clearpage
\subsection*{Self-Reference and Limits at a glance}\label{sec:self-reference-limits:summary}\addcontentsline{toc}{subsection}{Self-Reference and Limits at a glance}
\begin{center}
  \centering
  \captionof{table}{Self-Reference and Limits at a glance: the eighteen laws of \Cref{sec:self-reference-limits}, grouped into reading groups.}
  \label{tab:axis-self-reference-limits-summary}
  \scriptsize
  \setlength{\tabcolsep}{4pt}
  \medskip\noindent\begin{minipage}{\textwidth}\centering
  \textbf{Limits}\par\smallskip
  \begin{tabularx}{\textwidth}{@{}
      >{\raggedright\arraybackslash}p{0.24\textwidth}
      >{\raggedright\arraybackslash}X
      >{\raggedright\arraybackslash}p{0.10\textwidth}@{}}
    \toprule
    Name & One-line claim & Substrate cite \\
    \midrule
    Horizon-Sufficiency Theorem (\Fref{F33}) &
    When the pairing is surjective, a target depends on a history only
    through its boundary readout and
    exterior class exactly when it factors through their pairing; a
    bundled family does so exactly when each member does.  &
    \textemdash \\
    Information-Loss Theorem (\Fref{F34}) &
    When recovery after transport is surjective, source information is
    recoverable after transport exactly when it
    agrees on any two sources that transport and recovery identify.  &
    \textemdash \\
    Singular-Limit Theorem (\Fref{F36}) &
    For a surjective limit quotient, a target is completed exactly when no two processes with the same
    limit object have different target values, and is path-dependent
    otherwise; the pairing \(\rho_{u,\ell}\) refines \(\ell\), and \(u\) is completed
    through it; \(R\circ u\) is completed when \(u\) is, conversely for injective
    \(R\).  &
    \textemdash \\
    Locality Theorem (\Fref{F45}) &
    For a surjective local quotient, a target is local to a domain exactly
    when it is constant on the domain's classes. For strict coarsening, with the signature in the
    declared class and a nonempty carrier, a domain is minimal exactly
    when it has the signature's fibers.  &
    \textemdash \\
    \bottomrule
  \end{tabularx}
  \end{minipage}\par

  \medskip\noindent\begin{minipage}{\textwidth}\centering
  \textbf{Emergence}\par\smallskip
  \begin{tabularx}{\textwidth}{@{}
      >{\raggedright\arraybackslash}p{0.24\textwidth}
      >{\raggedright\arraybackslash}X
      >{\raggedright\arraybackslash}p{0.10\textwidth}@{}}
    \toprule
    Name & One-line claim & Substrate cite \\
    \midrule
    Scale-Descent Theorem (\Fref{F35}) &
    For a surjective scale quotient and coarsening, a law descends to the
    coarser scale exactly when it descends to the
    finer scale and its descended law is constant on the fibers of the
    coarsening; the renormalized law is then unique.  &
    \textemdash \\
    Continuum-Emergence Theorem (\Fref{F43}) &
    Over a nonempty directed index set with a surjective limit quotient,
    a coherent, cofinal tower with vanishing residuals, compatible
    projections and stage readouts, and a target readout equal to one
    stage readout after projection and constant on the limit fibers
    yields \(\operatorname{StableLimitQuotient}(L)\), witnessed by the unique readout on the quotient agreeing with every stage readout; coherence, cofinality and vanishing enter only as conjuncts of StableLimitQuotient, and the content is descent of the common stage readout through \(\ell\).  &
    \textemdash \\
    Fine-Tuning Theorem (\Fref{F47}) &
    If no modulus reaches the target region, fine-tuning is not
    realized; for a monotone measure record with a transitive order, a
    stricter target whose selector region has positive measure keeps a small
    selector region small.  &
    \textemdash \\
    Vacuum Theorem (\Fref{F50}) &
    For a surjective moduli quotient and a descended background, exactly
    one holds: it is necessary
    (equivalently, the raw background is constant) or it varies across
    moduli; a vacuum exists exactly when the vacuum criterion is a
    singleton with formed target, and is then unique.  &
    \textemdash \\
    Unification Theorem (\Fref{F51}) &
    Under exact unification the parent quotient is a common refinement of
    both children; when \(q_U\) is surjective, \(\langle\pi_A,\pi_B\rangle\) is
    injective iff the parent keeps exactly the children's distinctions.  &
    \textemdash \\
    \bottomrule
  \end{tabularx}
  \end{minipage}\par

  \medskip\noindent\begin{minipage}{\textwidth}\centering
  \textbf{Invariants}\par\smallskip
  \begin{tabularx}{\textwidth}{@{}
      >{\raggedright\arraybackslash}p{0.24\textwidth}
      >{\raggedright\arraybackslash}X
      >{\raggedright\arraybackslash}p{0.10\textwidth}@{}}
    \toprule
    Name & One-line claim & Substrate cite \\
    \midrule
    Complementarity Theorem (\Fref{F37}) &
    On a nonempty carrier, two separately admissible accesses are
    complementary exactly when
    every admissible joint carrier identifies two histories that one
    access separates, so that no admissible carrier has both accesses
    factoring through it.  &
    \textemdash \\
    Arrow Theorem (\Fref{F38}) &
    For a surjective record quotient, a continuation has an arrow exactly
    when it descends to a map of record classes that never lowers a
    record's status.  &
    \textemdash \\
    Entropy Theorem (\Fref{F39}) &
    For a surjective access quotient, the target obstruction is empty
    exactly when the target descends. For an entropy functional with a declared zero satisfying
    \(E(q')=0\) exactly when \(q'\) is injective,
    and with the chain rule and combine-monotonicity at a coarsening
    pair, \(E(q)=0\) exactly when \(q\) has no raw split pair, and entropy
    does not decrease under the coarsening. For finite nonempty \(H\), the fiber-volume functional
    \(E(q')=|H|^{-1}\sum_h\log|q'^{-1}(q'h)|\) satisfies these hypotheses and has \(m(\langle q,t\rangle,q)=0\) exactly
    when the target descends.  &
    \textemdash \\
    Irreducibility Theorem (\Fref{F42}) &
    Closedness at \(h\) is emptiness of the target and route split-pair sets on the class of \(h\); with well-founded costs and some admissible closed compression, an irreducible one exists; if for every \(v\in E\) the admissible class contains a map taking \(h\) to \(v\) and isolating \(h\), irreducibility is minimal cost over \(E\) and ignores the target.  &
    \textemdash \\
    Criticality Theorem (\Fref{F44}) &
    A phase change along a finite path of formed controls in which, for every neighbourhood index, each control has a neighbour in the phase of the next, or along a continuous path when neighbourhoods are the open sets, crosses a phase boundary; boundary and criticality
    descend through the control quotient when their data do.  &
    \textemdash \\
    Law--State Split Theorem (\Fref{F46}) &
    For a surjective moduli quotient, a readout descends exactly when no two closures with one modulus differ, and then uniquely; a descended readout is a law readout exactly when it is constant on the closure space and state-dependent exactly when two closures with different moduli differ; the selected modulus fixes the readout on its fiber.  &
    \textemdash \\
    Topology Theorem (\Fref{F48}) &
    For a surjective gluing quotient, a readout factors through it exactly when its split-pair set is empty; each such readout is exactly one of complete and coarse; a transport induces a map of gluing classes exactly when \(\gamma'\circ\tau\) has empty split-pair set.  &
    \textemdash \\
    Common-Source Theorem (\Fref{F49}) &
    On a nonempty history space, common-source and interface
    explanations are equivalent to their factorizations; direct-route
    explanation additionally requires admissibility and stability. A
    correlation residual is exactly failure of local explainability, and a
    common source supplies an interface factorization with \(\iota=\sigma\).  &
    \textemdash \\
    Closure Completeness Theorem (\Fref{F52}) &
    Scoped closure completeness holds exactly when the closure is formed
    and scope-declared, every residual is statused and adequate, and
    every claimed direction is in scope or named outside.  &
    \textemdash \\
    \bottomrule
  \end{tabularx}
  \end{minipage}\par

\end{center}

\FloatBarrier
  % sec:self-reference-limits
\section{The closure--reality boundary}\label{sec:anchor:closure-reality-boundary}

The final row of the catalog is not another internal normal form.  It
marks the boundary between formal closure-of-record and bare physical
realisability.  A closure-of-record (\Cref{def:closure-formation}) is a
formed closure with an audit record of its closure data.  Bare physical
realisability says only that some physical or operational system
instantiates a model in an ordinary sense.  The question is whether the
two notions can be identified; without further structure they cannot.
The running example of \Cref{sec:intro} is a formal closure-of-record:
its carrier, split pair and repairs are audited data, and whether a
physical system instantiates it is the separate question treated here.

We model a closure-of-record abstractly.  Let \(\mathcal C\) be a
declared class of carriers, each considered over a declared scope
\(\mathfrak S\) that records the targets, routes, instruments,
interfaces, residuals, budgets and nonclaims under discussion.  (In this
chapter a carrier is an element of \(\mathcal C\), such as a subset of
a record set in the finite models, not the underlying set of the
working vocabulary of \Cref{sec:apparatus}.)  Formal closure over
\(\mathfrak S\) is given by a declared operator
\(\mathrm{Cl}_{\mathfrak S}:\mathcal C\to\mathcal C\), and a carrier is
formally closed when it is a fixed point.  The naive
\emph{closure-equals-reality} (CER) principle would identify bare
physical realisability of \(C\) with \(C=\mathrm{Cl}_{\mathfrak S}(C)\).
Its direct direction says that physical realisability implies
closure-of-record, and its converse that closure-of-record implies
physical realisability.  The two no-go theorems below show that neither
direction is valid without further structure.  Their implications use
no closure axiom; the finite models satisfy all three closure axioms,
so neither direction follows from those axioms.

Three operators occur in this chapter: an arbitrary operator
\(\mathrm{Cl}_{\mathfrak S}\) in the two no-go theorems; the five-flag
saturation \(\operatorname{Sat}_5\), a finite formal model of zero
structural defect, in clauses (a), (c) and (e) of \Cref{prop:F53}; and
the inference reflector \(\ClSB\) of \PaperAOR\
\citep{Tsiokos2026AOR}, which reflects ledgers into inference-closed
ledgers while keeping their evidence unchanged, in clause (b) only.
Full accounting in that apparatus requires supplied witnesses in
addition to inference closure.  Neither construction identifies formal
closure with bare physical realisability.

\begin{definition}[Reflective closure operator]\label{def:reflective-closure}
This definition fixes two different operators: the retained five-flag saturation \(\operatorname{Sat}_5\), defined here, and the AOR inference reflector \(\ClSB\), recalled from \citet{Tsiokos2026AOR}.
In the five-flag encoding used in this paper (called the \emph{retained} encoding below: it retains a carrier's structural data and changes only its five flags), a carrier over a scope \(\mathfrak S\) is a pair
\(C=(s,(P_1(C),\ldots,P_5(C)))\) of structural data \(s\) in a fixed set
\(\mathcal D_{\mathfrak S}\) and five propositions, read as: declared records are
typed, cited routes are audited, residuals are statused, nonclaims are
registered, and meta-audit obligations are discharged; two carriers are equal
when their data are equal and their flags are pairwise equivalent.
Write \(\operatorname{Sat}_5(C)\) for the operator keeping the structural
data and replacing each flag by a true proposition, and
\(\Xi_5(C)=\{i\in\{1,\ldots,5\}:\neg P_i(C)\}\).
The retained operational predicate is defined by \(\Xi_5(C)=\varnothing\).
This is a predicate on the five flags, not a construction of witnesses.

The inference-closure operator \(\ClSB\) of \PaperAOR, called the AOR operator below, is a different operator. Only \(\operatorname{Sat}_5\), \(\Xi_5\) and the retained operational predicate are defined here; the AOR notions below are recalled from \citet{Tsiokos2026AOR} for clause (b) of \Cref{prop:F53}. It takes a ledger to its
least rule-closed extension, retaining the actual obligation meanings,
evidence carrier, validity relation and certificate decoder. The reflected ledger
contains the seed ledger (the unit of the reflection). Inference closure does not supply missing
certificates. Concretely, a ledger over an audit model is a set of facts of
four kinds about obligations \(o\): \(o\) is declared, \(o\) is registered, \(e\)
is a certificate for \(o\), and \(o\) is discharged. A ledger is \emph{sound}
when every recorded certificate is valid for its obligation and the
meaning of every discharged obligation holds,
\emph{inference-closed} when it is closed under the rules of \PaperAOR, and
\emph{fully discharged} when every declared obligation is discharged. For a kernel state \(S\), \(\operatorname{Accounted}(S)\)
means that every reachable record has a supplied status and payload satisfying
its original readout. The \emph{canonical reflected presentation} of \(S\) is \(\ClSB\) applied to the ledger that presents \(S\) over its audit model; clause (b) of \Cref{prop:F53} therefore concerns the AOR reflector, not \(\operatorname{Sat}_5\). The relation between full discharge of the canonical reflected presentation and \(\operatorname{Accounted}(S)\) is the imported equivalence stated in clause (b) of \Cref{prop:F53} \citep{Tsiokos2026AOR}. Kernel states, their reachable records, supplied
statuses and payloads, and the canonical reflected presentation are the
notions of the AOR kernel of \citet{Tsiokos2026AOR}; they are used only in
clause (b) of \Cref{prop:F53}, whose equivalence is imported. Only the second sentence of clause (b), whose ledger is the scalar-budget ledger in the proof of \Cref{prop:F53}, is checked here, with the inference rules of \citet{Tsiokos2026AOR} stated in its proof; the first sentence is read in the notation of \citet{Tsiokos2026AOR}. In this paper
\(\ClSB\) always denotes that inference reflector, and no identification
\(\operatorname{Sat}_5=\ClSB\) is assumed.
\end{definition}

The two no-go theorems below derive each failure from a hypothesis that
supplies it; their general content is the clause of
\Cref{thm:CER-converse-no-go} for an arbitrary extensive operator: for an
extensive operator (in particular a closure operator) on the subsets of a set,
there exists a bare-realisability predicate, stable under every reduct
satisfying \(\operatorname{Red}C\subseteq C\), for which both CER implications
fail if and only if the operator is not the identity.

The \emph{reduct-obstruction hypothesis} of the next theorem concerns a
declared reduct map \(\operatorname{Red}:\mathcal C\to\mathcal C\), which forgets
a closure-relevant datum \(d\), and consists of three items:
\begin{enumerate}[label=(\roman*),topsep=2pt,itemsep=1pt]
  \item a carrier \(C^{+}\) with \(\operatorname{BareReal}_{\mathfrak S}(C^{+})\);
  \item \(\operatorname{BareReal}_{\mathfrak S}(C)\Rightarrow
    \operatorname{BareReal}_{\mathfrak S}(\operatorname{Red}C)\) for every
    \(C\in\mathcal C\);
  \item \(\mathrm{Cl}_{\mathfrak S}(\operatorname{Red}C^{+})\ne\operatorname{Red}C^{+}\).
\end{enumerate}
In the intended reading, \(\operatorname{Red}C^{+}\) keeps a target claim whose
audit uses \(d\), and (iii) holds because closure must add \(d\) back. The
proof uses all three items.

\begin{theorem}[Direct CER no-go]\label{thm:CER-direct-no-go}
Let \(\mathfrak S\) be a declared scope, let \(\mathcal C\) be a
class of carriers, let
\(\operatorname{BareReal}_{\mathfrak S}:\mathcal C\to\mathrm{Prop}\)
be bare physical realisability, and let
\(\mathrm{Cl}_{\mathfrak S}:\mathcal C\to\mathcal C\) be any operator on
carriers, the formal closure over \(\mathfrak S\).
Assume the reduct-obstruction hypothesis. Then \(C:=\operatorname{Red}C^{+}\) satisfies
\[
  \operatorname{BareReal}_{\mathfrak S}(C)
  \quad\text{and}\quad
  C\ne \mathrm{Cl}_{\mathfrak S}(C),
\]
so the direct CER implication
\[
  \operatorname{BareReal}_{\mathfrak S}(C)
  \Longrightarrow
  C=\mathrm{Cl}_{\mathfrak S}(C)
\]
fails at \(C\); by definition it fails exactly when some bare-realised carrier is not fixed by \(\mathrm{Cl}_{\mathfrak S}\). Moreover, the hypothesis is satisfiable by a closure
operator: on the subsets of \(\{d,e,f\}\), with
\(\mathrm{Cl}_{\mathfrak S}(C)=C\cup\{d\}\) if \(e\in C\) and
\(\mathrm{Cl}_{\mathfrak S}(C)=C\) otherwise, \(\operatorname{Red}C=C\setminus\{d\}\)
and \(\operatorname{BareReal}_{\mathfrak S}(C):\Longleftrightarrow C\subseteq\{d,e\}\),
the operator \(\mathrm{Cl}_{\mathfrak S}\) is extensive, monotone and
idempotent, \(\operatorname{BareReal}_{\mathfrak S}\) is stable under
\(\operatorname{Red}\), items (i)--(iii) hold with \(C^{+}=\{d,e\}\), and the
direct implication fails at \(\{e\}\). For every operator
\(\mathrm{Cl}_{\mathfrak S}\ne\mathrm{id}\) on \(\mathcal P(X)\), extensive or not (in particular every closure operator other than the identity), choose \(C_0\subseteq X\)
with \(\mathrm{Cl}_{\mathfrak S}(C_0)\ne C_0\). Defining
\(\operatorname{BareReal}_0(C):\Longleftrightarrow C\subseteq C_0\) gives a predicate
stable under every reduct satisfying \(\operatorname{Red}C\subseteq C\), for which
the direct implication fails at \(C_0\) (proved at the end of the proof below).
\end{theorem}
\begin{proof}
Assume the scope \(\mathfrak S\), the predicate
\(\operatorname{BareReal}_{\mathfrak S}\), the operator
\(\mathrm{Cl}_{\mathfrak S}\), and the reduct-obstruction
hypothesis.  Let \(C:=\operatorname{Red}C^{+}\). By (i) and (ii),
\[
  \operatorname{BareReal}_{\mathfrak S}(C),
\]
and by (iii),
\[
  C\ne \mathrm{Cl}_{\mathfrak S}(C).
\]
In the intended reading, the retained claim has a native residual whose
required record was the forgotten datum \(d\), and closure adds it back.
This \(C\) is the displayed witness to failure of the direct CER
implication. The finite-model clause is checked in the paragraph that follows; for the clause on arbitrary nonidentity operators, \(\operatorname{Red}C\subseteq C\subseteq C_0\) gives \(\operatorname{BareReal}_0(\operatorname{Red}C)\), and \(C_0\) satisfies \(\operatorname{BareReal}_0\) while \(\mathrm{Cl}_{\mathfrak S}(C_0)\neq C_0\).
\end{proof}

\paragraph{A finite model.} The hypothesis is satisfiable. Let the carriers
be the subsets of a record set \(\{d,e,f\}\), where \(e\) is a target claim
whose audit uses the datum \(d\); let \(\mathrm{Cl}_{\mathfrak S}(C)\) add \(d\)
to \(C\) whenever \(e\in C\), and let \(\operatorname{Red}C=C\setminus\{d\}\). Let \(\operatorname{BareReal}_{\mathfrak S}(C)\) mean
\(C\subseteq\{d,e\}\), the records of a realised system; it is preserved by
reducts. With \(C^{+}=\{d,e\}\), items (i)--(iii) hold, and the witness of the
theorem is \(C=\{e\}\): it is bare-realised, and
\(\mathrm{Cl}_{\mathfrak S}(C)=\{d,e\}\ne C\). The theorem states what any
carrier satisfying the hypothesis implies; the model shows that such carriers
exist. The operator \(\mathrm{Cl}_{\mathfrak S}\) of the model is extensive,
monotone and idempotent, so the model also shows that the direct CER
implication does not follow from the axioms of a closure operator together
with stability of bare realisability under reducts.

\subsection*{Layer instantiations}\label{layer-inst:CER-direct}

\begingroup\small
\begin{description}
\item[\emph{Common Rule informed-consent record
(medical research ethics).}] \textit{Analogy.}
The carrier \(C\) is the clinical procedure as
performed; the scope \(\mathfrak S\) imposes the
Common Rule audit (protocol, IRB approval,
informed consent, adverse-event reporting).  The
closure-relevant datum \(d\) is the signed
consent form with its IRB version history.  The
procedure can be bare-realised (physically
performed) while \(d\) is missing, so the reduct
satisfies \(\operatorname{BareReal}_{\mathfrak S}\)
but is not a closure-of-record: the retained
claim ``subject enrolled under IRB-approved
protocol'' has a residual whose required record
is the missing consent \citep{DHHS2018CommonRule}.

\end{description}
\endgroup

The \emph{closure-instantiation-gap hypothesis} of the next theorem
consists of two items:
\begin{enumerate}[label=(\roman*),topsep=2pt,itemsep=1pt]
  \item a formal carrier \(C\) whose records, residual statuses, bridges
    and nonclaims are closed over \(\mathfrak S\), so that
    \(C=\mathrm{Cl}_{\mathfrak S}(C)\);
  \item a witness that \(C\) has no bare physical instantiation, so that
    \(\neg\operatorname{BareReal}_{\mathfrak S}(C)\).
\end{enumerate}
The proof takes both items from the hypothesis.

\begin{theorem}[Converse CER no-go]\label{thm:CER-converse-no-go}
Let \(\mathfrak S\), \(\mathcal C\),
\(\operatorname{BareReal}_{\mathfrak S}\), and
\(\mathrm{Cl}_{\mathfrak S}\) be as above.  The converse CER implication
\[
  C=\mathrm{Cl}_{\mathfrak S}(C)
  \Longrightarrow
  \operatorname{BareReal}_{\mathfrak S}(C)
\]
fails for some carrier if and only if some fixed point of
\(\mathrm{Cl}_{\mathfrak S}\) is not bare-realised. In the model of
\Cref{thm:CER-direct-no-go}, whose closure operator is extensive, monotone and
idempotent and whose bare realisability is stable under reducts, the
closure-instantiation-gap hypothesis holds at \(C=\{f\}\) and the converse
implication fails there, so neither CER implication follows from the axioms of
a closure operator together with stability of bare realisability under
reducts. More generally, let \(\mathrm{Cl}_{\mathfrak S}\) be any
operator on the subsets of a set \(X\), and choose
\(C_0\subseteq X\) with \(\mathrm{Cl}_{\mathfrak S}(C_0)\ne C_0\), which exists
whenever \(\mathrm{Cl}_{\mathfrak S}\ne\mathrm{id}\). For
the predicate \(B_0(C):\Leftrightarrow C\subseteq C_0\) (the predicate \(\operatorname{BareReal}_0\) of \Cref{thm:CER-direct-no-go}) in place of
\(\operatorname{BareReal}_{\mathfrak S}\), \(B_0\) is stable under every reduct with
\(\operatorname{Red}C\subseteq C\) and the direct implication fails at \(C_0\); no other property of \(\mathrm{Cl}_{\mathfrak S}\) is used. If \(\mathrm{Cl}_{\mathfrak S}\) is extensive (in particular a closure operator), then \(\mathrm{Cl}_{\mathfrak S}(X)=X\) and the converse implication fails at \(X\). For every extensive operator,
including the identity, the always-false predicate, which is stable under
every reduct, makes the converse implication fail at \(X\). For every operator on \(\mathcal P(X)\), extensive or not, some reduct-stable predicate makes the direct implication fail if
and only if \(\mathrm{Cl}_{\mathfrak S}\ne\mathrm{id}\), since for the identity every
carrier is fixed. Hence, for every extensive operator, both CER implications fail for some reduct-stable
bare-realisability predicate if and only if \(\mathrm{Cl}_{\mathfrak S}\ne\mathrm{id}\).
The closure-instantiation-gap hypothesis at a carrier \(C\) is, item by item,
the statement that \(C\) witnesses the failure of the converse implication;
the content beyond this reformulation is the finite model and the clause on
extensive operators. In particular, under the closure-instantiation-gap
hypothesis there exists a carrier \(C\in\mathcal C\) such that
\[
  C=\mathrm{Cl}_{\mathfrak S}(C)
  \quad\text{and}\quad
  \neg\operatorname{BareReal}_{\mathfrak S}(C).
\]
\end{theorem}
\begin{proof}
Assume the scope \(\mathfrak S\), the predicate
\(\operatorname{BareReal}_{\mathfrak S}\), the operator
\(\mathrm{Cl}_{\mathfrak S}\), and the closure-instantiation-gap
hypothesis.  The hypothesis supplies a formal carrier \(C\) whose
records, residual statuses, bridges, and nonclaims are internally
closed:
\[
  C=\mathrm{Cl}_{\mathfrak S}(C).
\]
It also supplies the missing-instantiation witness:
\[
  \neg\operatorname{BareReal}_{\mathfrak S}(C).
\]
The point is that formal closure only checks the declared audit data
inside \(\mathfrak S\).  It does not by itself construct an external
physical system realising \(C\).  Therefore the displayed carrier
\(C\) is closed over \(\mathfrak S\) but not bare-realised, which is
the desired witness to failure of the converse CER implication. For the
last sentence of the theorem: the converse implication fails at a carrier
\(C\) exactly when its hypothesis \(C=\mathrm{Cl}_{\mathfrak S}(C)\) holds and
its conclusion \(\operatorname{BareReal}_{\mathfrak S}(C)\) does not. The
clause on the finite model is checked in the paragraph that follows. For the
general clause, \(\operatorname{Red}C\subseteq C\subseteq C_0\) gives stability;
\(C_0\) is bare-realised and not closed, so the direct implication fails there;
extensivity gives \(\mathrm{Cl}_{\mathfrak S}(X)=X\), and \(X\not\subseteq C_0\),
since \(X=C_0\) would make \(C_0\) closed, so the converse implication fails
at \(X\).
For the identity case, \(\mathrm{id}\) fixes every carrier, so the direct implication holds for every predicate; \(\bot\) is stable under every reduct and fails at the fixed carrier \(X\).
\end{proof}

\paragraph{A finite model.} In the model after \Cref{thm:CER-direct-no-go},
the carrier \(C=\{f\}\) contains no claim whose audit needs another record, so
\(\mathrm{Cl}_{\mathfrak S}(C)=C\), and \(f\notin\{d,e\}\), so
\(\neg\operatorname{BareReal}_{\mathfrak S}(C)\). Both items of the hypothesis
hold, and \(C\) witnesses the failure of the converse implication; in the same
model the direct implication fails at \(\{e\}\). So neither CER direction follows
from the closure-operator axioms together with reduct stability.

\subsection*{Layer instantiations}\label{layer-inst:CER-converse}

\begingroup\small
\begin{description}
\item[\emph{Construction documents for a
never-constructed building (architectural
design).}] \textit{Analogy.}
The carrier \(C\) is the building considered as a
closure-of-record: sealed construction documents,
stamped structural calculations, life-safety
analysis, code-review approvals, and issued
permit.  Under \(\mathfrak S\) the
architectural-design audit is closed at the
construction-document set.  When the owner
declines to break
ground or funding lapses, no physical building
instantiates \(C\), so \(C=\mathrm{Cl}_{\mathfrak
S}(C)\) yet \(\neg\operatorname{BareReal}
_{\mathfrak S}(C)\)
\citep{AllenIano2019Fundamentals}.

\end{description}
\endgroup


\paragraph{Reading the statement.}
The positive fixed-point equivalence below is proved for the retained
five-flag saturation. The AOR inference reflector has a separate
accounting criterion. An inference-closed ledger can still have an unpaid
obligation. Therefore the fixed-point equivalence for the flags cannot be
imported as an equivalence between inference closure and payment.
The proposition keeps both statements at their actual scopes.

\begin{proposition}[Closure--Reality Boundary resolution]\label{prop:F53}
For the retained encoding and the AOR kernel of
\Cref{def:reflective-closure}, the following statements hold.
\begin{enumerate}[label=(\alph*),topsep=2pt,itemsep=2pt]
\item A carrier is fixed by \(\operatorname{Sat}_5\) if and only if
  \(\Xi_5(C)=\varnothing\), equivalently all five flags hold.
\item (The accounting equivalence is imported; the kernel notions
  are those of \citet{Tsiokos2026AOR}, and the counterexample
  below is checked here.) For an AOR kernel state \(S\), its canonical reflected
  presentation is fully discharged if and only if
  \(\operatorname{Accounted}(S)\). Inference closure alone does not imply
  full discharge: a sound, inference-closed ledger over the scalar-budget
  evidence model of \PaperAOR, with all records declared, leaves the obligation \(5\le2\) unpaid.
\item The retained operational predicate, defined by
  \(\Xi_5(C)=\varnothing\), is consequently equivalent to the fixed-point
  predicate for \(\operatorname{Sat}_5\), and by (d) it is the unique
  predicate \(B\) for which both CER implications hold with
  \(\mathrm{Cl}=\operatorname{Sat}_5\) and \(B\) in place of bare
  realisability. No implication to or from bare
  physical realisability is asserted, and no identification with
  AOR inference closure is made.
\item For any operator \(\mathrm{Cl}\) and predicate \(B\), both CER
  implications hold for every carrier exactly when
  \(B(C)\iff C=\mathrm{Cl}(C)\) for every \(C\). Thus the fixed-point
  predicate is the unique predicate satisfying both implications.
\item If the set \(\mathcal D_{\mathfrak S}\) of structural data is nonempty, let
  \(D(C)=\{i\in\{1,\dots,5\}:P_i(C)\}\). For every proper subset
  \(D_0\subsetneq\{1,\dots,5\}\), both CER implications fail for
  \(\operatorname{Sat}_5\) and the predicate \(D(C)\subseteq D_0\). That
  predicate is stable under every reduct \(\operatorname{Red}\) with
  \(D(\operatorname{Red}C)\subseteq D(C)\) for all \(C\).
\end{enumerate}
Clause (b)'s accounting equivalence is imported from the AOR
kernel; its counterexample is checked using that same substrate. The other
clauses are proved for the retained encoding or for an arbitrary operator.
\end{proposition}
\begin{proof}
For (a), saturation keeps the structural data and sets all five flags to true.
Equality with the original carrier holds exactly when every original flag is
true, by propositional extensionality. This is precisely
\(\Xi_5(C)=\varnothing\). Clause (c) is (a) read through the definition of the retained predicate, together with (d) for \(\mathrm{Cl}=\operatorname{Sat}_5\).

For (b), \PaperAOR\ proves that a canonical discharged fact is derivable
exactly when its original record has a supplied valid status and payload.
Applying this equivalence to every reachable record yields the stated full
accounting equivalence. For the counterexample, take the scalar-budget model
with meaning \((c,b)\mapsto c\le b\), valid evidence a natural slack
\(e\) satisfying \(c+e=b\), and no dependency edges. Its interpreted
ledger declares and registers every obligation, stores exactly valid
certificates, and discharges exactly true inequalities. The inference rules of \citet{Tsiokos2026AOR} are four: from \(o\) declared
infer \(o\) registered; from a certificate \(e\) for \(o\) infer \(o\)
discharged; from \(o\) discharged infer \(o\) registered; and from \(o\)
declared and a dependency edge \(o\to d\) infer \(d\) declared. The
interpreted ledger contains every declaration and registration; a stored
certificate \(e\) for \((c,b)\) is valid, so \(c+e=b\) and \(c\le b\), and the
discharge the second rule concludes is already present; the model has no
dependency edges. Hence every rule preserves the ledger, and it is inference
closed. The declared obligation \((5,2)\)
is not discharged since \(5\nleq2\). Thus even a sound inference fixed point
need not be fully discharged.

For (d), the two implications are exactly the two directions of the stated
equivalence. Predicates with the same truth values are equal. For (e), keeping
any carrier's structural data and clearing its flags gives \(D(C)=\varnothing\),
which lies below \(D_0\) but is not fixed. Its saturation is fixed and has
all five flags, so it does not lie below the proper subset \(D_0\).
A reduct removing discharged flags only shrinks \(D(C)\).

The generic no-go theorems do not depend on identifying the two operators.
The retained saturation explains a formal fixed-point predicate. The
AOR construction additionally tracks the original witnesses required for
accounting; a consumer policy may require accepted delivery beyond accounting.
No bare physical existence claim follows from either construction.
\end{proof}

\subsection*{Layer instantiations}\label{layer-inst:F53}

\begingroup\small
\begin{description}
\item[\emph{Good Manufacturing Practice
(pharmaceutical regulatory affairs).}] \textit{Analogy.}
The raw carrier is a manufactured drug substance and the declared scope
includes batch records, validated processes, traceability and quality-unit
release. Closing the documentation's inference rules does not establish
that a missing validation was performed. Actual certificates must satisfy
those original obligations; release may additionally require the consumer's
acceptance policy. The distinction illustrates inference closure versus
accounting, without deriving a regulatory guarantee \citep{ICH2016Q7}.

\end{description}
\endgroup


This chapter closes the catalog by separating formal fixed points, audited
accounting and bare physical realisability. Neither CER direction follows
from the closure-operator axioms with reduct stability. The retained positive
equivalence concerns \(\operatorname{Sat}_5\); the AOR reflector
requires the separate source-accounting condition. The chapter imports that
condition's kernel theorem and does not claim that inference closure creates
its witnesses. This chapter does not re-prove the AOR apparatus or extend its
source-accounting equivalence to non-AOR reflective frameworks. The broader philosophical setting is the structural-realism
debate launched by \citet{Worrall1989StructuralRealism}; this chapter records
the formal closure / realisability gap without taking a stand in that debate.
         % sec:anchor:closure-reality-boundary
\section{Substrate instances at a glance}\label{sec:substrate-instances}

This chapter introduces no new claim.  It is a map for readers arriving
from a substrate paper: it says which law of the catalog is
instantiated by, or motivated by, which paper of the series.

No-Needles Stability (\Cref{thm:F6}) is the layer-agnostic form of the
layer-dissolving membrane theorem of \PaperNeedleKiller\
\citep{Tsiokos2026NeedleKiller}, which proves it stagewise for positive
semidefinite audit forms with Moore--Penrose pseudoinverses;
\PaperNS\ \citep{Tsiokos2026NavierStokes} applies that theorem to shell
readouts of the Navier--Stokes equations.

The Strict Extension lemma (\Cref{lem:F8}) is proved here; the
Schur-calculus picture of residuals composing along refinement chains,
from \PaperAR\ \citep{Tsiokos2026XiCompanion}, is context, not a
premise.

The Hiddenness Normal Form (\Cref{thm:F13a}) is the current/predictive
quotient statement associated with \PaperPvNP\
\citep{Tsiokos2026PvNP}.  Its source hypothesis \((\mathsf S)\) is
motivated by the saturated-SAT construction of that paper, but no
interface of that construction is shown here to satisfy it.

The CSL Formed-Closure Law (\Cref{thm:F13b}) is motivated by \PaperCSL\
\citep{Tsiokos2026CSL}.  It applies when the declared state readout,
together with current access, determines the predictive class; that
paper supplies an internal-state readout and a loop and exposure
diagnosis, but does not establish this pairwise hypothesis.

Duality Fixity (\Cref{thm:F14}) restates, for an arbitrary involutive
ledger, the confinement squeeze of \PaperSDTC\ \citep{Tsiokos2026SDTC},
which \PaperRH\ \citep{Tsiokos2026RH} applies to the zeros of the
Riemann zeta function.  The domination bounds are a hypothesis, and by
the last clause of the theorem they exist exactly when the anti-invariant
readout vanishes almost everywhere; the source papers likewise take them
as input.  In the second paper, the hypothesis supplying them is
equivalent to the Riemann hypothesis on the bare shell of all
nontrivial zeros, and in its concrete version on finite windows of zeros
closed under \(s\mapsto1-\bar s\) with unit weights that eventually
contain every zero.

Spontaneous Symmetry Breaking / Branch-Selection Formation
(\Cref{thm:F15b}) is proved here and runs parallel to the CSL diagnosis
pattern: a branch-selection record exists exactly when some ground
state is not invariant under the symmetry.

The Closure--Reality Boundary resolution (\Cref{prop:F53}) separates
the five-flag saturation from the inference reflector of \PaperAOR\
\citep{Tsiokos2026AOR}.  That paper supplies the inference reflection
and the source-accounting equivalence of clause (b), which is imported;
the fixed-point criterion, the no-go theorems and the counterexample are
proved here.

\begin{table}[t]
  \centering
  \footnotesize
  \begin{tabularx}{\textwidth}{@{}
      >{\raggedright\arraybackslash}p{0.18\textwidth}
      >{\raggedright\arraybackslash}p{0.18\textwidth}
      >{\raggedright\arraybackslash}X
      >{\raggedright\arraybackslash}p{0.22\textwidth}@{}}
    \toprule
    Law & Substrate paper(s) & Substrate result & Treatment here \\
    \midrule
    \Fref{F6} No-Needles &
    \citet{Tsiokos2026NeedleKiller, Tsiokos2026NavierStokes} &
    Layer-dissolving membrane theorem and its Navier--Stokes
    application. &
    anchored; proved here \\
    \Fref{F8} Strict Extension &
    \citet{Tsiokos2026XiCompanion} &
    Schur-calculus picture of residuals along refinement chains. &
    proved here; the source supplies context \\
    \Fref{F13a} Hiddenness &
    \citet{Tsiokos2026PvNP} &
    Saturated-SAT construction motivating the source hypothesis, which is not verified for it. &
    motivated; conditional on \((\mathsf S)\) \\
    \Fref{F13b} CSL Formed-Closure &
    \citet{Tsiokos2026CSL} &
    Internal-state readout and exposure diagnosis; pairwise determination is an additional hypothesis. &
    motivated; proved here under the determining hypothesis \\
    \Fref{F14} Duality Fixity &
    \citet{Tsiokos2026SDTC, Tsiokos2026RH} &
    Confinement theorem and its application to the zeros of the zeta function. &
    anchored; proved here, with the domination bounds as a hypothesis \\
    \Fref{F15b} Spontaneous SSB &
    \citet{Tsiokos2026CSL} (parallel) &
    Parallel to the CSL diagnosis. &
    parallel; proved here \\
    \Fref{F53} Closure--Reality Boundary &
    \citet{Tsiokos2026AOR} &
    Inference reflector and its source-accounting equivalence (clause (b) of \Cref{prop:F53}). &
    anchored; clause (b) imported, the rest proved here \\
    \bottomrule
  \end{tabularx}
  \caption{Laws of the catalog with substrate papers.}
  \label{tab:substrate-instances}
\end{table}

        % sec:substrate-instances
\section{Sketches and future work}\label{sec:sketches}

This chapter records candidates, not laws.  Several structural
candidates were considered for the catalog but are not laws of it,
because their recognition source has not been identified.  A
recognition source is the substrate, apparatus layer or carrier
discipline that makes a named structure precise enough to state in the
template of \Cref{def:normal-form}: a setup, a theorem and a proof.
Without one, a phrase may be suggestive and still fail to determine
these.  Some of the candidates below may turn out to belong to a
substrate paper only, because the structure is carrier-specific; others
may become laws once a recognition source is supplied.

\paragraph{Dimension as stable rank of access.}  The sketch asks
whether dimension can be read, at the layer-agnostic level, as a
presentation-invariant rank of the access quotient: roughly, the stable
number of independent directions that survive the declared observation
discipline.  Such a reading would explain why dimension is neither raw
carrier data nor a coordinate convention.  To state it as a law one
must know which rank notion is meant, which class of quotient
presentations preserves it, and which audit certifies that the rank is
observed rather than imposed.  These data are not part of the present
apparatus.

\paragraph{Smoothness as coherent linearization.}  The sketch asks
whether smooth structure can be characterized by compatible local
linearizations across declared overlaps: local linear approximations
should globalize exactly when they descend through the relevant access
and overlap quotients.  This would connect the local--global
obstruction of \Cref{sec:transport:local-global} with a
smoothness-specific substrate, but it needs a substrate-independent
account of what counts as a linearization, which overlaps preserve it,
and how the associated coherence ledger is audited.

\paragraph{Stability of nothing.}  A ``nothing'' state at a layer is
rarely absolute absence; it is more naturally a terminal, initial, null
or background class relative to a quotient and an audit.  The catalog
already contains Vacuum / Background Selection (\Cref{thm:F50}): a
vacuum is the single element of the declared vacuum criterion, with
formed target.  The sketch asks the stronger question of when the
stability of the null class forces a nonempty residual structure, which
needs a recognition source for the null class and for the residual
supposed to survive it.

\paragraph{Other candidates.}  Boundary compression, topology,
continuum behavior, critical loci, uncertainty and inaccessible
determining state are treated by laws of the catalog: Horizon Sufficiency, Information Loss, Continuum Emergence,
Criticality, Complementarity as Non-Joint Access, Topology as Global
Gluing Invariant and Hiddenness.  Other formulations on these six
topics are not stated here; in the application of \Cref{thm:F52} below
they count as named under those topics.

\paragraph{Scope of the catalog.}  The eight axes are an organizing
scheme that makes the catalog navigable, not a structured object, and
the catalog is not claimed to be complete.  The Closure Completeness
Boundary (\Cref{thm:F52}) characterizes scoped completeness for any
declared residuals and directions.  Take as \(D\) the candidate
directions considered for the catalog, let
\(\operatorname{inScope}(d)\) mean that \(d\) is a law of the catalog
or reduces to one, let \(\operatorname{namedOutside}(d)\) mean that
\(d\) is named in this chapter, and let \(\operatorname{claimed}(d)\)
hold for every \(d\in D\).  The direction clause of \Cref{thm:F52} then
requires every candidate to be covered or named.  Since the list of
candidates is not reproduced here, the paper does not verify that
clause; it records only the conditional statement that if every
candidate is a law, reduces to one, or is named here, then the direction
clause holds and the unmarked frontier is empty.  The residual clauses
of \Cref{thm:F52} would require a residual register for the catalog,
which is not given.  The candidates named here that neither belong to
nor reduce to a law form the catalog boundary \(\partial D\) of
\Cref{def:catalog-boundary}.  \Cref{thm:F52} is not an absolute
completeness theorem: a new recognition source may move a sketch out of
\(\partial D\) and into the catalog.

A candidate that comes with a recognition source can be tested in the
same way as the laws of this paper: identify the setup, state the
quotient or closure normal form, give the proof, record the cases and
write the nonclaims.  If it passes without special carrier features, it
belongs in the catalog; if it depends on special analytic, geometric,
arithmetic or physical structure, it belongs in the corresponding
substrate paper.
        % sec:sketches
\section{Nonclaims and caveats}\label{sec:nonclaims}

This chapter collects the scope limits of \FIV.  The chapters state the
local caveats where they matter; here the scope of the catalog, and
answers to common questions about it, are gathered in one place.

\begin{enumerate}[label=\textbf{N\arabic*.}, leftmargin=*]
\item \textbf{\FIV\ does not re-prove substrate theorems.}
The paper states layer-agnostic laws; it does not re-derive the
carrier-specific theorems of the substrate papers.  Where a law is
anchored to or motivated by a substrate paper, the body says so, and
\Cref{tab:substrate-instances} collects these cases.

\item \textbf{\FIV\ is not a foundational reset.}
The apparatus inherited from \FoundI, \FoundII\ and \FoundIII\ is used
as background after the review in \Cref{sec:apparatus}.  Closure
formation, admissibility and quotient packaging are inputs here, not
new foundations; the \(\BirdInt\) judgment is used by no hypothesis or
proof.

\item \textbf{\FIV\ is not a unification claim.}
The axes are an organizing device for recurring obstruction types.  The
paper does not assert that they form a categorical, fibrational,
monadic or otherwise privileged structure.

\item \textbf{\FIV\ is not a complete catalog.}
The \NumLaws\ laws are the scope of this paper, not an exhaustive list.
Candidates whose recognition source, in the sense of
\Cref{sec:sketches}, has not been identified appear in
\Cref{sec:sketches}.

\item \textbf{The conditional law is conditional.}
The Hiddenness Normal Form (\Fref{F13a}) assumes its source hypothesis
\((\mathsf S)\), which is not verified here for any interface derived
from the saturated-SAT construction that motivates it.

\item \textbf{\FIV\ does not establish the substrate endpoints.}
The catalog does not prove the self-dual trace confinement hypothesis,
the Riemann hypothesis, any separation of complexity classes, or the
hiddenness results of its substrate papers
\citep{Tsiokos2026SDTC, Tsiokos2026RH, Tsiokos2026PvNP, Tsiokos2026CSL};
it states layer-agnostic laws whose instances those papers study.

\item \textbf{The catalog discipline is the contribution.}
The contribution is the uniform treatment of \NumLaws\ layer-agnostic
laws, each with its hypotheses, proof, provenance and nonclaims.  The
substrate proofs live in the cited substrate papers.

\item \textbf{The eight-axis scheme is not minimal or canonical.}
The scheme is an editorial choice, not a theorem.  Other groupings are
legitimate if they preserve the same statements and exclusions.
\end{enumerate}

The same scope answers the common questions about the catalog.

\begin{enumerate}[label=\textbf{O\arabic*.}, leftmargin=*]
\item \textbf{Why call this ``Foundations IV'' if it is a catalog?}
The title names the position in the series, not a reset.  The
contribution is the uniform statement and proof of \NumLaws\
layer-agnostic laws.

\item \textbf{Are the anchored laws merely restating substrate theorems?}
\Fref{F6} is proved here, over the reals, for any block data
satisfying its identities and budgets; its real substrate instances
use Moore--Penrose pseudoinverses of positive semidefinite forms, which
satisfy these identities.  For
\Fref{F53}, the retained saturation/zero-flag equivalence is proved
here; the source-accounting equivalence of the audited-realisability
paper is imported at its stated scope, and the counterexample of clause
(b) of \Cref{prop:F53}, a sound inference-closed ledger leaving the
obligation \(5\le2\) undischarged, is checked here, together with the
general no-go theorems and finite models.

\item \textbf{Why is the Hiddenness Normal Form stated conditionally?}
Its first four conjuncts concern declared predicates of an interface
whose intended source is the saturated-SAT construction of
\citet{Tsiokos2026PvNP}.  No such interface is constructed there or
here, so these conjuncts enter as the hypothesis \((\mathsf S)\); the
fifth conjunct and the equivalence of surplus with non-collapse are
proved.

\item \textbf{Is the hypothesis \((\mathsf S)\) consistent?}
Yes: it holds for small explicit interfaces, including one with a
predictive surplus (\Cref{sec:access:hiddenness}).  This shows that
part (i) of \Cref{thm:F13a} is not vacuous; it says nothing about the
intended saturated-SAT interface.

\item \textbf{Why a single paper rather than one per axis?}
A single paper states the apparatus once and lets laws of different
axes be compared in one register.

\item \textbf{Why these \NumLaws\ laws and not others?}
They are the recurring structural obstructions and closures of the
Six Birds series.  The catalog is not exhaustive; \Cref{sec:sketches}
records candidates that remain outside it.

\item \textbf{Why do the Status/Records/Coherence and
Self-Reference/Limits chapters hold so many laws?}
Those axes cover the broadest range of record, coherence, information,
limit and self-reference phenomena.  Short per-law subsections and the
summary tables keep them navigable.

\item \textbf{Do the apparatus restatements lose nuance from the
foundations papers?}
Yes; the review in \Cref{sec:apparatus} is intentionally compact.  It
suffices for reading this paper, and the cited foundations papers carry
the operational detail.
\end{enumerate}
         % sec:nonclaims

\appendix
\crefalias{section}{appendix}
\crefalias{subsection}{appendix}
\crefname{appendix}{appendix}{appendices}
\Crefname{appendix}{Appendix}{Appendices}
\section{Lean formalization}\label{app:formalization}

\subsection{Scope}\label{subsec:lean-wording-discipline}

The laws of this paper are formalized in Lean~4 with Mathlib, in the
modules \texttt{SixBirdsMetaMath.FoundationsIV} of the accompanying
repository.  The modules and everything they import build without
\texttt{sorry}, and the development declares no axioms: the declarations
listed below depend only on Lean's standard axioms \texttt{propext},
\texttt{Classical.choice} and \texttt{Quot.sound}.  Each law is listed
in \Cref{tab:formalization-coverage} with its declarations and one of
three alignments.
\begin{description}
\item[Faithful.] The listed declarations together prove every clause of
  the paper statement, each from at most the paper's hypotheses for that
  clause; they may prove more, or state the data more generally.  A
  listed declaration marked \emph{companion} uses a stronger hypothesis
  and is not the covering declaration.
\item[Restricted.] The declarations cover less than the paper
  statement; this applies to \Fref{F14} only.
\item[Imported clause.] One clause applies a theorem of the
  formalization of a substrate paper; this applies to clause (b) of
  \Fref{F53} only.
\end{description}
Fifty laws are faithful, \Fref{F14} is restricted, and \Fref{F53} is
faithful apart from its imported clause.  A faithful formalization of a
conditional law, such as \Fref{F13a}, proves the implication from its
hypotheses, not the hypotheses.

\subsection{Notes on particular laws}\label{subsec:narrowing-disclosure}

\paragraph{F13a.} The source hypothesis \((\mathsf S)\) is the
predicate \texttt{HiddennessSourceObligation}: surjectivity of the
current map together with conjuncts 1--4.  Part (i) of \Cref{thm:F13a}
is \path|hiddenness_normal_form_with_exposure_collapse_surplus|, which
takes \((\mathsf S)\) as an explicit hypothesis;
\path|hiddenness_normal_form_of_source_obligation| states the same for
every interface satisfying a source predicate that implies
\((\mathsf S)\).  Part (ii) is
\path|predictive_surplus_iff_not_collapse_iff| together with
\path|hidden_iff_collapse_iff_of_surplus_iff_not_collapse|.  The
satisfiability remark of the reading paragraph is
\path|hiddenness_source_obligation_with_surplus_satisfiable|.  No
interface derived from the saturated-SAT construction of
\citet{Tsiokos2026PvNP} is formalized.  In the namespace of
\Fref{F13b}, the predicate \(\operatorname{HiddenStateSurplus}\) is
\texttt{HiddenFromCurrent}, with \texttt{HiddenStateSurplus} an
abbreviation for it; the \Fref{F13a} predicate \texttt{HiddenFromCurrent}
is a different definition.

\paragraph{F6.} \path|Stability.NoNeedlesReal.no_needles_stability_real|
proves \Cref{thm:F6} for real matrices, with the Loewner preorder
defined by quadratic forms, from exactly the displayed identities and
budgets; the decomposition \(K_{DD}=\widehat E K_{LL}\widehat E^{\top}+\Xi\)
and monotonicity under congruence are separate lemmas.  Companion
declarations prove the endpoint over the rational matrices of the
substrate development, for arbitrary rational blocks and for supplied
symmetric pseudoinverses.

\paragraph{F14.} The formalization proves the direct squeeze for an
abstract ledger satisfying the trace laws (C1)--(C3) and the
almost-everywhere laws used in the proof, with an additive involution \(U\)
preserving the declared squared norm in place of the unitary involution;
the proof of the squeeze does not use \(U\).  The equivalence under the
condition on \(U\), and the converse on ladders, are proved for the
abstract ledger under two explicit laws that the ledger structure does
not include (the squared norm of \(0\) is \(0\), and an
almost-everywhere-zero nonnegative integrand has integral \(0\)); a
one-point ledger with \(U=-\mathrm{id}\) shows that the equivalence fails
without the condition on \(U\).  The laws are instantiated on finite
ledgers valued in \(1\times1\) positive semidefinite rational matrices
with \(U=\mathrm{id}\).  An instance for Bochner integrals over a Hilbert
space is not formalized here; the companion formalization of
\citet{Tsiokos2026SDTC} treats real Hilbert spaces.

\paragraph{F29.} The first listed declaration assumes that \(\sim\) is
the kernel of \(\pi\), whereas \Cref{thm:F29} assumes only
\(p\sim p'\Rightarrow\pi(p)=\pi(p')\).  The clauses not mentioning
\(\sim\) follow from it with \(\sim\) the kernel of \(\pi\); the clauses on
\(\sim\) are \texttt{presentation\_invariance\_general\_relation}, which
assumes one-way compatibility only.

\paragraph{F39.} The fiber-volume instance is formalized over \(\mathbb
R\), with the chain identity, monotonicity and the zero criterion proved
directly; the abstract theorem, with combine-monotonicity at the pair,
applies to it.

\paragraph{F43.} Stage carriers and residual sets may vary with the
index, as in the statement; the candidate limit carrier is supplied, not
constructed, in both.

\paragraph{F53 and the closure--reality theorems.} The direct no-go
theorem is \path|ClosureRealityBoundary.cer_direct_reduct_witness|,
which takes items (i)--(iii) as hypotheses over an abstract carrier
type; \path|finite_cer_direct_instance| and
\path|finite_cer_paper_bundle| instantiate the finite model, and
\path|cer_direct_no_go_of_nonidentity_operator| proves the clause on
arbitrary nonidentity operators.  The converse theorem is
\path|cer_converse_closed_witness|, with the characterization
\path|cer_converse_failure_iff_closed_not_bare| and the clauses on
extensive operators \path|cer_no_go_of_nonidentity_extensive_operator|
and \path|cer_no_go_iff_nonidentity_extensive_operator|.  For
\Cref{prop:F53}, the retained fixed-point equivalence, the generic
uniqueness statement and the five-flag counterexamples of clause (e) are
proved directly.  The equivalence of clause (b) is
\path|current_aor_accounted_completion_iff|, which applies the kernel
theorem of the formalization of \citet{Tsiokos2026AOR}, and
\path|current_aor_inference_closed_not_fully_discharged| checks the
scalar-budget counterexample with the ledger definitions and inference
rules of that development; the proof in the body is self-contained given
the four rules it states.  The names
\texttt{ClSB} and \texttt{SixBirdsClosed}, and the prefix \texttt{aor}
in \texttt{closed\_iff\_aor}, \texttt{aor\_saturation\_cer\_both\_fail}
and \texttt{aor\_realisability\_unique\_fixed\_point\_predicate}, denote
\(\operatorname{Sat}_5\) and the five-flag encoding, not the AOR
inference reflector \(\ClSB\) of this paper; only the two
\texttt{current\_aor} declarations concern the AOR kernel.  The wrapper
\path|closure_reality_boundary| uses opaque obstruction fields and does
not instantiate them; it is not listed in \Cref{tab:formalization-coverage},
and no listed declaration uses it.

\subsection{Declarations by law}\label{subsec:formalization-coverage}

In \Cref{tab:formalization-coverage}, declaration names omit the prefix
\texttt{SixBirdsMetaMath.FoundationsIV}.

\begingroup
\scriptsize
\renewcommand{\arraystretch}{0.95}
\begin{longtable}{@{}
      >{\raggedright\arraybackslash}p{0.06\linewidth}
      >{\raggedright\arraybackslash}p{0.13\linewidth}
      >{\raggedright\arraybackslash}p{0.75\linewidth}@{}}
\caption{Lean declarations for the laws of the catalog.}\label{tab:formalization-coverage}\\
    \toprule
    Law & Alignment & Lean declarations \\
    \midrule
\endfirsthead
\caption[]{Lean declarations for the laws of the catalog (continued).}\\
    \toprule
    Law & Alignment & Lean declarations \\
    \midrule
\endhead
    \bottomrule
\endfoot
    \multicolumn{3}{@{}l}{\textbf{Access}} \\
    \Fref{F10} & faithful & \path|Access.Quotientality.quotientality| \\
    \Fref{F11} & faithful & \path|Access.Adequacy.paper_adequacy_three_parts| \\
    \Fref{F12} & faithful & \path|Access.NoFreeDistinction.no_free_distinction_paper| \\
    \Fref{F13a} & faithful & \path|Access.HiddennessNormalForm.hiddenness_normal_form_with_exposure_collapse_surplus|; \path|Access.HiddennessNormalForm.hiddenness_normal_form_of_source_obligation|; \path|Access.HiddennessNormalForm.predictive_surplus_iff_not_collapse_iff|; \path|Access.HiddennessNormalForm.hidden_iff_collapse_iff_of_surplus_iff_not_collapse|; \path|Access.HiddennessNormalForm.hiddenness_source_obligation_with_surplus_satisfiable| \\
    \Fref{F13b} & faithful & \path|Access.CSLFormedClosure.csl_formed_closure_direct_full| \\
    \Fref{F24} & faithful & \path|Access.LayerMultiplicity.paper_layer_multiplicity_strict_iff_corestriction| \\
    \addlinespace
    \multicolumn{3}{@{}l}{\textbf{Stability}} \\
    \Fref{F6} & faithful & \path|Stability.NoNeedlesReal.no_needles_stability_real|; \path|Stability.NoNeedlesReal.no_needles_stability_ordered_field|; \path|Stability.NoNeedlesReal.block_schur_decomposition|; \path|Stability.NoNeedlesReal.loewnerLE_congruence|; companions over the rationals: \path|Stability.NoNeedles.no_needles_stability|; \path|Stability.NoNeedles.no_needles_stability_general_blocks|; \path|Stability.NoNeedles.no_needles_stability_supplied_pseudoinverse| \\
    \Fref{F7} & faithful & \path|Stability.SufficiencyClosure.sufficiency_closure_probe_family| \\
    \Fref{F8} & faithful & \path|Stability.StrictExtension.strict_extension|; \path|Stability.StrictExtension.strict_extension_status_partition| \\
    \Fref{F9} & faithful & \path|Stability.NoUnstatusedResidual.no_unstatused_residual_paper_scoped| \\
    \addlinespace
    \multicolumn{3}{@{}l}{\textbf{Transport}} \\
    \Fref{F2} & faithful & \path|Transport.DescentRepair.paper_descent_repair_normal_form_maps| \\
    \Fref{F3} & faithful & \path|Transport.HolonomyMemoryRepair.paper_holonomy_memory_repair| \\
    \Fref{F4} & faithful & \path|Transport.LocalGlobalObstruction.local_global_obstruction_paper_full| \\
    \Fref{F5} & faithful & \path|Transport.RecombinationComparison.recombination_comparison| \\
    \addlinespace
    \multicolumn{3}{@{}l}{\textbf{Symmetry}} \\
    \Fref{F14} & restricted & \path|Symmetry.DualityFixity.duality_fixity_direct|; \path|Symmetry.DualityFixity.duality_fixity_direct_without_anti_invariance|; \path|Symmetry.DualityFixity.rational_psd_duality_fixity|; \path|Symmetry.DualityFixity.rational_trace_weighted_rank_one|; \path|Symmetry.DualityFixity.rational_psd_ledger_nonzero|; \path|Symmetry.DualityFixity.duality_fixity_identity_involution|; \path|Symmetry.DualityFixity.duality_fixity_no_minus_one_eigen|; \path|Symmetry.DualityFixity.duality_fixity_minus_one_eigen_counterexample|; \path|Symmetry.DualityFixity.finite_duality_fixity_identity_involution|; \path|Symmetry.DualityFixity.duality_ladder_exists_iff_readout_vanishes| \\
    \Fref{F15a} & faithful & \path|Symmetry.SelectionObstruction.selection_obstruction_orbit_conditional_paper|; \path|Symmetry.SelectionObstruction.selection_obstruction_orbit_strengthened|; \path|Symmetry.SelectionObstruction.selection_obstruction_equivariant_maps_characterization| \\
    \Fref{F15b} & faithful & \path|Symmetry.SpontaneousSSB.spontaneous_ssb_orbit_constant_paper| \\
    \Fref{F40} & faithful & \path|Symmetry.Anomaly.anomaly_symmetry_obstruction_full_paper| \\
    \Fref{F41} & faithful & \path|Symmetry.DualityEquivalence.paper_duality_equivalence_quotient_roles_full| \\
    \addlinespace
    \multicolumn{3}{@{}l}{\textbf{Status / Records / Coherence}} \\
    \Fref{F17} & faithful & \path|StatusRecordsCoherence.ObjectivityAsCommonQuotient.paper_objectivity_from_factorization_minimality| \\
    \Fref{F18} & faithful & \path|StatusRecordsCoherence.LawhoodDescent.readout_invariant_under_fiber_bijections_iff_descends| \\
    \Fref{F19} & faithful & \path|StatusRecordsCoherence.ObjectPersistence.object_persistence_strengthened_of_source_surjective| \\
    \Fref{F20} & faithful & \path|StatusRecordsCoherence.FactRecordFormation.paper_fact_record_formation_conditional_raw_transport_with_event| \\
    \Fref{F21} & faithful & \path|StatusRecordsCoherence.ReflexiveNonclosure.paper_reflexive_nonclosure_with_sharpness_rules_2_3_4_without_rule_1|; \path|StatusRecordsCoherence.ReflexiveNonclosure.paper_reflexive_nonclosure_with_sharpness_rules_2_3_4| (companion; also assumes rule 1) \\
    \Fref{F22} & faithful & \path|StatusRecordsCoherence.InterfaceMediation.paper_interface_mediation_of_source_boundary_surjective|; \path|StatusRecordsCoherence.InterfaceMediation.paper_interface_mediation_split_full| (companion; also assumes \(q'\) surjective) \\
    \Fref{F23} & faithful & \path|StatusRecordsCoherence.ProbabilityAsStableUnresolvedFiber.paper_probability_fiber_split_per_clause_stable| (clauses (i)--(v), any weight set); \path|StatusRecordsCoherence.ProbabilityAsStableUnresolvedFiber.paper_probability_fiber_split_per_clause_point_mass| (companion; clause (vi) under global weights); \path|StatusRecordsCoherence.ProbabilityAsStableUnresolvedFiber.finite_positive_fiber_chance_criterion_on_fiber|; \path|StatusRecordsCoherence.ProbabilityAsStableUnresolvedFiber.paper_probability_fiber_split_per_clause_stable_strengthened| (clause (vi) as stated) ; \path|StatusRecordsCoherence.ProbabilityAsStableUnresolvedFiber.stable_probability_concentrated_every_class| ; \path|StatusRecordsCoherence.ProbabilityAsStableUnresolvedFiber.stable_probability_iff_stable_and_charges_two| \\
    \Fref{F25} & faithful & \path|StatusRecordsCoherence.CausalityAsStableInterventionTransport.paper_causality_transport_separate_raw_witness_full| \\
    \Fref{F26} & faithful & \path|StatusRecordsCoherence.ModuliContingency.paper_moduli_contingency_split_full_general| \\
    \Fref{F27} & faithful & \path|StatusRecordsCoherence.ConservationAsOrbitDescent.partial_route_step_invariant_iff_orbit_descends| \\
    \Fref{F28} & faithful & \path|StatusRecordsCoherence.Composability.paper_composability_unique|; \path|StatusRecordsCoherence.Composability.paper_composability|; \path|StatusRecordsCoherence.Composability.assoc_identity_descent| \\
    \Fref{F29} & faithful & \path|StatusRecordsCoherence.PresentationInvariance.presentation_invariance_paper_conditional_strengthened|; \path|StatusRecordsCoherence.PresentationInvariance.presentation_relation_invariance_any_presentation|; \path|StatusRecordsCoherence.PresentationInvariance.presentation_invariance_general_relation| \\
    \Fref{F30} & faithful & \path|StatusRecordsCoherence.ExplanationAsCompressionWithDescent.explanation_as_compression_with_descent| \\
    \Fref{F31} & faithful & \path|StatusRecordsCoherence.CounterfactualLegitimacy.paper_counterfactual_legitimacy_without_surjectivity| \\
    \addlinespace
    \multicolumn{3}{@{}l}{\textbf{Self-reference / Limits}} \\
    \Fref{F33} & faithful & \path|SelfReferenceLimits.HorizonSufficiency.horizon_sufficiency_paper| \\
    \Fref{F34} & faithful & \path|SelfReferenceLimits.InformationLoss.information_loss_normal_form_paper| \\
    \Fref{F35} & faithful & \path|SelfReferenceLimits.ScaleDescentRenormalization.scale_descent_renormalization_paper_tower_strengthened|; \path|SelfReferenceLimits.ScaleDescentRenormalization.scale_descent_renormalized_unique| \\
    \Fref{F36} & faithful & \path|SelfReferenceLimits.SingularLimitCompletion.singular_limit_completion_paper_strengthened| \\
    \Fref{F37} & faithful & \path|SelfReferenceLimits.Complementarity.complementarity_split_criterion_general_unconditional|; \path|SelfReferenceLimits.Complementarity.complementarity_split_criterion_general| (companion; also assumes separate admissibility) \\
    \Fref{F38} & faithful & \path|SelfReferenceLimits.ArrowFromRecordMonotonicity.paper_arrow_from_record_monotonicity_fiber| \\
    \Fref{F39} & faithful & \path|SelfReferenceLimits.EntropyAsFiberVolume.paper_entropy_as_fiber_volume_conditional_mixing_at_pair|; \path|SelfReferenceLimits.EntropyAsFiberVolume.paper_shannon_entropy_as_fiber_volume|; \path|SelfReferenceLimits.EntropyAsFiberVolume.paper_shannon_fiber_entropy_instance|; \path|SelfReferenceLimits.EntropyAsFiberVolume.paper_entropy_as_fiber_volume_conditional_mixing| (companion; under global combine-monotonicity) \\
    \Fref{F42} & faithful & \path|SelfReferenceLimits.Irreducibility.paper_irreducibility_admissible|; \path|SelfReferenceLimits.Irreducibility.point_irreducible_all_maps_collapses|; \path|SelfReferenceLimits.Irreducibility.point_irreducible_isolating_maps_collapses|; \path|SelfReferenceLimits.Irreducibility.irreducibility_no_short_closed_description|; \path|SelfReferenceLimits.Irreducibility.admissible_closed_compression_has_irreducible_minimum| \\
    \Fref{F43} & faithful & \path|SelfReferenceLimits.ContinuumEmergence.continuum_limit_readout_dependent_tower_full|; \path|SelfReferenceLimits.ContinuumEmergence.continuum_limit_readout_dependent_tower_general| \\
    \Fref{F44} & faithful & \path|SelfReferenceLimits.Criticality.criticality_phase_boundary_descent_paper_full| \\
    \Fref{F45} & faithful & \path|SelfReferenceLimits.Locality.locality_domain_of_dependence_paper_strict| \\
    \Fref{F46} & faithful & \path|SelfReferenceLimits.LawStateSplit.law_state_split_paper_full| \\
    \Fref{F47} & faithful & \path|SelfReferenceLimits.FineTuning.fine_tuning_small_selector_region_paper|; \path|SelfReferenceLimits.FineTuning.complementary_selector_regions_not_both_small_real| \\
    \Fref{F48} & faithful & \path|SelfReferenceLimits.Topology.topology_global_gluing_invariant_paper| \\
    \Fref{F49} & faithful & \path|SelfReferenceLimits.CommonSourceNonlocal.common_source_nonlocal_correlation_repaired_strengthened_of_nonempty| \\
    \Fref{F50} & faithful & \path|SelfReferenceLimits.VacuumBackgroundSelection.paper_vacuum_background_selection_cardinality|; \path|SelfReferenceLimits.VacuumBackgroundSelection.vacuum_background_link_split|; \path|SelfReferenceLimits.VacuumBackgroundSelection.vacuum_background_link| (companion; also assumes the descent equation for the background clause) \\
    \Fref{F51} & faithful & \path|SelfReferenceLimits.Unification.unification_as_common_refinement_paper_exact_kernel|; \path|SelfReferenceLimits.Unification.common_refinement_square_and_kernel_of_projection_equations| \\
    \Fref{F52} & faithful & \path|SelfReferenceLimits.ClosureCompletenessBoundary.paper_closure_completeness_boundary_general_sharp| \\
    \addlinespace
    \multicolumn{3}{@{}l}{\textbf{Closure--Reality Boundary}} \\
    \Fref{F53} & faithful; clause (b) imported & \path|ClosureRealityBoundary.cer_direct_reduct_witness|, \path|ClosureRealityBoundary.cer_converse_closed_witness|, \path|ClosureRealityBoundary.cer_converse_failure_iff_closed_not_bare|, \path|ClosureRealityBoundary.cer_no_go_of_nonidentity_closure| (companion; assumes a closure operator), \path|ClosureRealityBoundary.finite_cer_paper_bundle|, \path|ClosureRealityBoundary.clSB_fixed_iff_structural_defect_zero|, \path|ClosureRealityBoundary.cer_no_go_of_nonidentity_extensive_operator|, \path|ClosureRealityBoundary.cer_no_go_iff_nonidentity_extensive_operator|, \path|ClosureRealityBoundary.cer_direct_no_go_of_nonidentity_operator|, \path|ClosureRealityBoundary.cer_direct_failure_iff_nonidentity_operator|, \path|ClosureRealityBoundary.closed_iff_aor|, \path|ClosureRealityBoundary.aor_saturation_cer_both_fail|, \path|ClosureRealityBoundary.finite_cer_direct_instance|, \path|ClosureRealityBoundary.finite_cer_converse_instance|, \path|ClosureRealityBoundary.cer_both_implications_iff_fixed_point_predicate|, \path|ClosureRealityBoundary.aor_realisability_unique_fixed_point_predicate|, \path|ClosureRealityBoundary.current_aor_accounted_completion_iff|, \path|ClosureRealityBoundary.current_aor_inference_closed_not_fully_discharged| \\
\end{longtable}
\endgroup

   % app:formalization
\section{Further layer instantiations}\label{app:further-instantiations}

This appendix collects the analogies that accompany the laws of the catalog, beyond those kept beside each law. They are grouped by law, in the order of the chapters, and all tagged as analogies; every special case is kept in the body.

\subsection*{\Fref{F10} Quotientality}
\begingroup\small
\begin{description}
\item[\emph{Trubetzkoy phonemic principle (structural linguistics).}] \textit{Analogy.}
For a language \(L\), the phoneme inventory \(\mathcal P_L\) is a
partition of phonetic realisations preserving every
morpheme-distinguishing minimal-pair contrast.  Two phones that
produce a minimal pair belong to distinct phonemes; phones that never
do are grouped into one phoneme only when they are also phonetically
similar (English \emph{h} and the final nasal of \emph{sing} never
contrast yet are distinct phonemes), so non-contrast alone is not an
equivalence.  Every morpheme-discriminating phonological observable
factors through the phoneme map, and merging two phonemes that form
a minimal pair collapses a morpheme distinction \(L\) tracks
\citep{Trubetzkoy1939GrundzugePhonologie}.

\item[\emph{Categorical perception of voice onset time (cognitive
psychophysics).}] \textit{Analogy.}
For a synthesised voice-onset-time continuum spanning \(/{\!b}/\) to
\(/{\!p}/\), human listeners exhibit \emph{categorical perception}:
within-category {ABX} discrimination is weak, falling toward but
generally not to chance, while between-category discrimination peaks
at the boundary near \(+25\)~ms.  Liberman and colleagues first
showed this pattern on a place-of-articulation continuum
(\(/{\!b}/\)--\(/{\!d}/\)--\(/{\!g}/\)).  The phonetic-category
partition is the coarsest equivalence on the acoustic continuum
preserving the identification probe, which factors through the
phonetic-category map; the ABX probe factors through it up to the
residual within-category discrimination
\citep{LibermanHarrisHoffmanGriffith1957Discrimination}.
\end{description}
\endgroup

\subsection*{\Fref{F11} Adequacy / No-Overread}
\begingroup\small
\begin{description}
\item[\emph{Post-processing invariance of differential privacy
(cryptography / privacy).}] \textit{Analogy.}
For an \((\varepsilon,\delta)\)-differentially-private mechanism
\(M:\mathcal D\to\mathcal R\), the post-processing property of differential privacy (proved for \(\varepsilon\)-DP by Dwork--McSherry--Nissim--Smith, and holding equally for \((\varepsilon,\delta)\)-DP) states that for any (possibly
randomised) \(g:\mathcal R\to\mathcal W\), the composition
\(g\circ M\) is itself \((\varepsilon,\delta)\)-DP.  Every post-processing
of \(M\) is therefore again \((\varepsilon,\delta)\)-DP; a published quantity whose
neighboring-database distinguishability exceeds the
\(e^{\varepsilon}\!\cdot\!\!\ (\cdot)+\delta\) bound is the
overread failure --- evidence of access outside \(M\)'s
randomisation
\citep{DworkMcSherryNissimSmith2006CalibratingNoise}.

\item[\emph{Registered Reports / preregistered analyses (scientific
methodology).}] \textit{Analogy.}
In the Registered Reports format the access quotient is the
\emph{Stage~1 preregistered protocol together with the locked
dataset}: confirmatory claims must be computable from those two
inputs alone, and any two execution histories sharing them must
produce the same confirmatory verdict.  HARKing, p-hacking, and
outcome switching are overread failures: a published claim
presented as confirmatory but computed from post-hoc analytic
choices fails to factor through the declared protocol-plus-data
quotient
\citep{Chambers2013RegisteredReports}.

\item[\emph{Federal Rules of Evidence: admissibility (legal
procedure).}] \textit{Analogy.}
Under the U.S. Federal Rules of Evidence (Rules 401--403 on
relevance, Rule 802 on hearsay), a finding of fact must factor
through the admitted-evidence quotient of the case record: two
factually distinct records with identical admitted-evidence content
must yield the same verdict.  A verdict resting on inadmissible
evidence is reversible error unless harmless, instantiating the
\Fref{F11} overread failure
\citep{FRE2024FederalRulesEvidence}.

\item[\emph{Halpern--Moses common knowledge (distributed-systems
epistemic discipline).}] \textit{Analogy.}
For a process \(i\) in a distributed system with local-state map
\(r_i\) and induced indistinguishability \(\sim_i\), the Kripke
semantics for \(K_i\varphi\) says that the truth of
\(K_i\varphi\) at a global state factors through \(r_i\): two
global states with the same local state for \(i\) make
\(K_i\varphi\) hold or fail together.  Accepted knowledge claims
by process \(i\) are therefore exactly those determined by its
view; a claim that depends on global state beyond \(r_i\) is the
overread failure.  Common knowledge \(C_G\varphi\) for a group \(G\)
is the group-level case: its truth factors through the finest
partition coarser than every local view \(r_i\), \(i\in G\), so it
holds or fails together on each connected component of the union of
the relations \(\sim_i\)
\citep{HalpernMoses1990KnowledgeCommonKnowledge}.
\end{description}
\endgroup

\subsection*{\Fref{F12} No-Free-Distinction}
\begingroup\small
\begin{description}
\item[\emph{Wolpert no-free-lunch theorem (statistical learning).}] \textit{Analogy.}
For supervised learning on a finite input space, averaged uniformly
over all possible target functions, every learning algorithm
achieves the same expected off-training-set error.  Off-sample
generalization is not a distinction obtainable from training data
alone; it requires the declared apparatus extension of an
inductive bias --- a hypothesis class restriction, a prior, or a
regulariser --- without which Wolpert's averaging identity collapses
every learner to chance
\citep{Wolpert1996NoFreeLunch}.

\item[\emph{Wootters--Zurek no-cloning theorem (quantum mechanics).}] \textit{Analogy.}
No unitary on \(\mathcal H_S\otimes\mathcal H_A\) sends
\(|\psi\rangle|0\rangle\) to \(|\psi\rangle|\psi\rangle\) for every
unknown \(|\psi\rangle\); the linearity of unitary evolution
forbids perfect copying of an unknown quantum state.  Acquiring a
distinguishing measurement of \(|\psi\rangle\) requires either
declaring a measurement basis (with collapse) or accepting the
fidelity bound of an approximate cloner --- the \Fref{F12} apparatus
cost in either case
\citep{WoottersZurek1982NoCloning}.

\item[\emph{Goldwasser--Micali semantic security (cryptography).}] \textit{Analogy.}
For an IND-CPA secure probabilistic public-key encryption scheme,
no probabilistic polynomial-time adversary holding only the public
key and a challenge ciphertext can distinguish encryptions of two
adversarially chosen messages with non-negligible advantage.
Producing such a distinction requires apparatus extension --- the
secret key, a decryption oracle, or super-polynomial computation
breaking the underlying trapdoor (quadratic residuosity, factoring)
\citep{GoldwasserMicali1984ProbabilisticEncryption}.

\item[\emph{Myerson revenue equivalence (mechanism design).}] \textit{Analogy.}
In a symmetric independent-private-values auction, Myerson's envelope-formula characterisation shows that in every
incentive-compatible mechanism the allocation rule fixes each
bidder's interim expected payment up to a constant, the payment of
the lowest type, so mechanisms with the same allocation rule and
lowest-type payments raise the same expected revenue; any
incentive-compatible mechanism whose allocation responds to private
values leaves the bidders an information rent.  An allocation-only mechanism with no transfers collapses
to type-pooling
\citep{Myerson1981OptimalAuctionDesign}.
\end{description}
\endgroup

\subsection*{\Fref{F13a} Hiddenness Normal Form}
\begingroup\small
\begin{description}
\item[\emph{Pearl back-door criterion and \(do\)-calculus (causal
inference).}] \textit{Analogy.}
In a causal DAG with observed variables \(V\) and unobserved
confounders \(U\), an unblocked back-door path from \(X\) to \(Y\)
through a common ancestor in \(U\) is the canonical witness of a
hidden upstream determining state.  The observational/interventional
gap \(P(Y\mid X)\ne P(Y\mid \mathrm{do}(X))\) is then the named
predictive surplus inside a \(P(V)\)-fiber, and \(do\)-calculus
gives rules for deciding when causal effects remain nonparametrically identifiable, for instance via a back-door adjustment set or a front-door mediator (an instrumental variable alone generally does not identify the causal effect; additional identifying assumptions are needed, and these need not be parametric)
\citep{Pearl1995CausalDiagrams}.

\item[\emph{Le Verrier's prediction of Neptune (celestial
mechanics).}] \textit{Analogy.}
In 1845--1846, Le Verrier and (independently, from 1843) Adams analysed the residual
perturbations of Uranus's orbit after accounting for all
gravitational influences then catalogued.  The unexplained
predictive surplus in Uranus's observed position was interpreted as
the gravitational signature of an undiscovered upstream determining
state, and Le Verrier computed an ecliptic-longitude prediction
which Galle confirmed telescopically on 1846 September 23 within
\(1^\circ\) of Le Verrier's position: the lawful exposure of the
hidden planet Neptune
\citep{LeVerrier1846RecherchesUranus}.

\item[\emph{{\AA}str{\"o}m belief-MDP reduction for {POMDPs}
(reinforcement learning).}] \textit{Analogy.}
For a partially observable Markov decision process, two histories
sharing the current observation can yield different distributions
over future observation--reward sequences only when their belief
states (posteriors over the hidden Markov state) differ; distinct
belief states need not yield distinct distributions.
{\AA}str{\"o}m's 1965 reduction theorem realises the lawful exposure:
the POMDP is equivalent to a fully observed MDP on the
continuous-belief simplex \(\Delta(S)\), with the belief state as a sufficient upstream-quotient refinement,
not in general the minimal one, that resolves the
observation-fiber surplus
\citep{Astrom1965OptimalControlIncomplete}.

\item[\emph{Fisher's epistatic variance component (quantitative
genetics).}] \textit{Analogy.}
Under the Cockerham--Kempthorne orthogonal decomposition,
phenotypic variance partitions as
\(\operatorname{Var}(Y)=V_A+V_D+V_I+V_E\) with
\(V_I\) the epistatic interaction component.  A non-zero \(V_I\) is
the named witness of an upstream multi-locus determining state:
the phenotypic variability within a single-locus marginal fiber
that additive single-locus regression cannot explain corresponds
under the declared decomposition to hidden gene--gene interaction
structure
\citep{Fisher1918CorrelationRelatives}.
\end{description}
\endgroup

\subsection*{\Fref{F13b} CSL Formed-Closure Law}
In each item the determining readout \(d\) is the named hidden cause (assignable cause, implementation defect, orbiting body, resonance); the analogy asserts neither that \(d\) is determining nor that the cited source establishes it.
\begingroup\small
\begin{description}
\item[\emph{Shewhart statistical process control (industrial
quality engineering).}] \textit{Analogy.}
In a formed manufacturing closure with declared in-control mean
\(\mu\), standard deviation \(\sigma\), and \(3\sigma\) control
limits calibrated from a baseline period, an out-of-control signal
(point outside the control limits or a non-random run pattern) is
the named predictive surplus that triggers root-cause investigation
into an assignable cause --- tool wear, raw-material change,
operator effect, environmental drift --- under declared sampling
assumptions (a sampling rule, baseline independence, and identical
distribution).  The control-chart signal is a witness, not
a proof: false alarms exist under the in-control distribution at the
declared rate
\citep{Shewhart1931EconomicControl}.

\item[\emph{Model checking against a temporal-logic specification
(formal verification).}] \textit{Analogy.}
For a finite-state model \(M\) and a temporal-logic
specification \(\varphi\) (branching-time CTL in the cited algorithm,
linear-time LTL in SPIN), an explicit-state or symbolic
model-checker produces a counterexample trace \(\tau\) when
\(M\not\models\varphi\) for a linear-time or universal property.  The trace is the named predictive surplus
witnessing a hidden upstream defect in the implementation that the
declared specification was meant to forbid, with lawful exposure via
the automated checker (SPIN, NuSMV, TLC for TLA+ specifications)
under declared model-checking assumptions (a closed system, finite
state, a fair scheduler, and sound temporal-logic semantics)
\citep{ClarkeEmersonSistla1986AutomaticVerification}.

\item[\emph{Exoplanet detection by radial velocity and transit
(observational astronomy).}] \textit{Analogy.}
For a formed exoplanet-detection closure with declared instrument
(spectrograph + Doppler pipeline; or photometric transit pipeline)
and declared sensitivity (Doppler precision in m\,s\(^{-1}\);
photometric precision in ppm), a periodic radial-velocity
oscillation or photometric dimming inconsistent with stellar
activity is the named predictive surplus witnessing a hidden
orbiting body.  Mayor and Queloz realised this lawful exposure with
the 1995 detection of 51 Pegasi b, the first confirmed exoplanet
around a Sun-like star
\citep{MayorQueloz199551Pegasi}.

\item[\emph{{LHC} Higgs-boson discovery as resonance detection
(high-energy physics).}] \textit{Analogy.}
For a formed LHC search closure with declared integrated luminosity,
detector acceptance \(\times\) efficiency, and background-only
prediction (Standard-Model processes excluding the hypothesised
Higgs), an observed event yield exceeding the background-only model
by \(\ge 5\sigma\) in a signal region is the named predictive
surplus witnessing a hidden upstream resonance / particle predicted
by the broader theory.  The ATLAS 2012 observation of a 125 GeV
resonance in declared diphoton and four-lepton mass windows realised
the lawful exposure, with follow-up spin-parity and branching-ratio
measurements identifying the Standard-Model Higgs boson
\citep{ATLAS2012HiggsObservation}.
\end{description}
\endgroup

\subsection*{\Fref{F24} Layer Multiplicity}
\begingroup\small
\begin{description}
\item[\emph{{ANSI} {SQL} transaction isolation levels (database
concurrency).}] \textit{Analogy.}
For concurrent database transactions, the isolation-level quotient
classifies the visibility of in-flight operations.  A transaction's read or write role can split that quotient,
producing a dirty read, a non-repeatable read, or a phantom.  Each
transaction declares in advance one of the four ANSI isolation levels
(read-uncommitted, read-committed, repeatable-read, serializable),
and each level prices a different set of these role-split
obstructions as excluded or tolerated; Berenson et al.\ show that the
ANSI anomaly definitions are ambiguous and miss further anomalies
such as lost updates and write skew
\citep{BerensonBernsteinGrayMeltonONeilONeil1995SQLIsolation}.

\item[\emph{Grammatical case-marking (structural linguistics).}] \textit{Analogy.}
In a case-marking language the morphological word-form quotient
classifies nouns by inflection.  When a syntactic role (subject,
object, oblique, possessor) splits that morphological quotient, the
language's case system resolves the obstruction by assigning a
named case --- nominative, accusative, ergative, absolutive,
dative, genitive, locative, instrumental --- from Blake's
typological inventory.  Different languages instantiate different
subsets of the inventory, but each role-split is closed by a named
case-formation category
\citep{Blake2001Case}.

\item[\emph{{OSI} seven-layer networking reference model (computer
networking).}] \textit{Analogy.}
In a networking stack the protocol-layer quotient classifies
communication functions.  When a new communication role splits the
current layer quotient, the ISO/IEC 7498-1 model resolves the
obstruction by placing the function at one of seven named layers
(physical, data-link, network, transport, session, presentation,
application), each licensing specific function types and obstructing
others
\citep{ISOIEC7498OSI}.

\item[\emph{Polymorphism taxonomy (programming-language theory).}] \textit{Analogy.}
In a typed language the type quotient classifies expressions by
their monomorphic type.  When a function's required role is
generic across multiple types --- a role-split obstruction in the
monomorphic quotient --- the Cardelli--Wegner taxonomy resolves
the obstruction by assigning the function to a named polymorphism
category: parametric (universal quantification), ad-hoc (overloading /
coercion), or inclusion (subtyping).  Each category
licenses a distinct layer-formation mechanism for the generic role
\citep{CardelliWegner1985UnderstandingTypes}.
\end{description}
\endgroup

\subsection*{\Fref{F6} No-Needles Stability}
\begingroup\small
\begin{description}
\item[\emph{Quantum relative-entropy monotonicity.}] \textit{Analogy.}
For any completely positive trace-preserving (CPTP) quantum channel
\(T\) and density operators \(\rho,\sigma\), the relative entropy
satisfies \(S(T(\rho)\,\|\,T(\sigma))\le S(\rho\,\|\,\sigma)\).  The
decomposition is the Stinespring dilation
\(T(\rho)=\operatorname{Tr}_E[V\rho V^{*}]\); the monotonicity
combines joint convexity of relative entropy (Lindblad's lemma) with
isometric embedding and partial trace.  The preorder is scalar
\(\le\) on entropy functionals, not the Loewner block ordering
\citep{Lindblad1975CompletelyPositiveMaps}.

\item[\emph{Csiszár \(f\)-divergence data-processing inequality
(information theory).}] \textit{Analogy.}
For any convex \(f:(0,\infty)\to\mathbb R\) with \(f(1)=0\), the
\(f\)-divergence
\(D_f(P\,\|\,Q)=\mathbb E_Q\!\bigl[f(\mathrm dP/\mathrm dQ)\bigr]\) is
monotone under every Markov kernel \(K\):
\(D_f(K\!\cdot\!P\,\|\,K\!\cdot\!Q)\le D_f(P\,\|\,Q)\).  The
monotonicity comes from the joint convexity of
\((p,q)\mapsto qf(p/q)\) plus Jensen on the kernel mixture; a general
\(f\)-divergence has no chain rule.  The preorder is scalar \(\le\) on a
divergence value rather than the Loewner cone
\citep{LieseVajda2006Divergences}.

\item[\emph{Strassen monotone-coupling stochastic dominance
(probability theory).}] \textit{Analogy.}
By Strassen's theorem, \(X\le_{\mathrm{st}}Y\) iff there is a
coupling \(\pi\) of \(\mu_X\) and \(\mu_Y\) supported on
\(\{(x,y):x\le y\}\); applying a stochastically monotone Markov
kernel \(K\) to both coordinates preserves the order, so
\(K\!\cdot\!\mu_X\le_{\mathrm{st}}K\!\cdot\!\mu_Y\).  The
decomposition is the coupling existence (a measure-theoretic object
on a product space); the preorder is stochastic dominance defined by
tail-probability inequalities, with no positive-semidefinite cone
involved \citep{Strassen1965ExistenceProbabilityMeasures}.

\item[\emph{Lyapunov's direct method (dynamical systems).}] \textit{Analogy.}
For a dissipative dynamical system \(\dot x=f(x)\) with a
positive-definite Lyapunov functional \(V\) and a non-negative
dissipation
residual \(W\), the Lie-derivative inequality
\(\dot V=\nabla V\!\cdot\!f(x)\le-W(x)\le 0\) along trajectories
propagates \(V(x(t))\le V(x_0)\) for all \(t\ge 0\).  The
decomposition is the pair \((V,W)\) chosen against the vector field's
structure; transitivity composes through the flow semigroup
\(\varphi_t\), giving the \Fref{F6} endpoint along a continuous-time
flow rather than via a matrix Schur complement
\citep{Lyapunov1892StabilityMotion}.
\end{description}
\endgroup

\subsection*{\Fref{F8} Strict Extension}
\begingroup\small
\begin{description}
\item[\emph{Hematopoietic CMP \(\to\) MEP/GMP bifurcation (stem-cell
biology).}] \textit{Analogy.}
The common myeloid progenitor (CMP) population, prospectively
isolated by Lin\(^{-}\) Sca-1\(^{-}\) c-Kit\(^{+}\) CD34\(^{+}\)
Fc\(\gamma\)RII/III\(^{\mathrm{lo}}\) surface markers, splits into
megakaryocyte--erythroid (MEP) and granulocyte--macrophage (GMP)
progenitors distinguished by CD34 and Fc\(\gamma\)RII/III.  Both
daughters retain the Lin\(^{-}\) Sca-1\(^{-}\) c-Kit\(^{+}\) part of
the CMP gate and leave it on the CD34 and Fc\(\gamma\)RII/III axes
(MEP are CD34\(^{-}\) Fc\(\gamma\)RII/III\(^{\mathrm{lo}}\), GMP are
CD34\(^{+}\) Fc\(\gamma\)RII/III\(^{\mathrm{hi}}\)), and the witness
split is two
clonally related siblings of one CMP division that commit to MEP and
GMP respectively \citep{AkashiTraverMiyamotoWeissman2000CMPClonogenic}.

\item[\emph{Parisi replica-symmetry-breaking hierarchy (statistical
mechanics).}] \textit{Analogy.}
In Parisi's mean-field solution of the Sherrington--Kirkpatrick spin
glass, the one-step replica-symmetry-breaking ansatz partitions
overlaps between pure states into two values \(\{q_0,q_1\}\); full RSB
refines this to a continuous overlap function \(q(x)\) for
\(x\in[0,1]\).  Cutting the full-RSB ultrametric tree at one
overlap level groups pure states into clusters with the two-level
overlap pattern of the one-step ansatz; this cluster partition, not
the one-step ansatz solution itself, is the retention.  The
ultrametric witness split is a triple of pure states in one cluster,
all of whose pairwise overlaps lie above the cut, with
\(q_{\alpha\beta}\ne q_{\alpha\gamma}\) under full RSB
\citep{Parisi1983OrderParameterSpinGlasses}.

\item[\emph{Verner's Law refining Grimm's Law (historical
linguistics).}] \textit{Analogy.}
Grimm's Law sends PIE \({*}p,{*}t,{*}k,{*}k^w\) (outside clusters after \({*}s\) or, for \({*}t\), after another obstruent) to
Proto-Germanic \({*}f,{*}\theta,{*}h,{*}h^w\); on the reflexes of the PIE voiceless stops, Verner's Law refines
this output by voicing a non-initial fricative when the preceding PIE
syllable was unaccented.  The refinement holds only on that domain:
the voiced Verner reflexes merge with the reflexes of the PIE voiced
aspirates.  The
grammatischer-Wechsel alternation within Germanic strong-verb
paradigms is the canonical witness: forms that differed only in PIE accent position remain merged under
Grimm alone but split by voicing under Grimm + Verner, even though the conditioning accent is itself
lost in attested Germanic
\citep{Verner1877AusnahmeLautverschiebung}.
\end{description}
\endgroup

\subsection*{\Fref{F9} No Unstatused Residual}
\begingroup\small
\begin{description}
\item[\emph{Progress and preservation.}] \textit{Analogy.}
The Wright--Felleisen syntactic type-soundness theorem states that
every well-typed reachable configuration of a source term retains
its typing judgment (preservation) and is either a value or admits
a reduction step (progress).  The progress class
\(\{\text{value},\,\text{can-step}\}\) together with the retained
typing judgment is the \Fref{F9} typed status.  The orphan residual is a
\emph{stuck term}: a closed non-value with no reduction rule and no
type, forbidden by progress together with preservation
\citep{WrightFelleisen1994SyntacticTypeSoundness}.

\item[\emph{Genetic code via ribosomal translation (molecular
biology).}] \textit{Analogy.}
The ribosome decodes each mRNA codon into a unique amino acid (the
61 sense codons) or a termination signal (the 3 stop codons),
following the standard genetic code table; codon--anticodon pairing
at the A-site enforces single-valuedness, with context-dependent
recoding of UGA to selenocysteine and of UAG to pyrrolysine at
specific sites as the known exceptions.  Ribosomes stalled at the 3\('\) end of mRNAs lacking an in-frame stop codon constitute the translation-failure residual; tmRNA trans-translation in bacteria, and Pelota--Hbs1 splitting followed by ribosome-associated quality-control ubiquitination in eukaryotes, rescue the ribosome and mark the peptide for degradation\footnote{\url{https://doi.org/10.1126/science.271.5251.990}; \url{https://doi.org/10.1016/j.molcel.2014.07.006}.}
\citep{Crick1968GeneticCode}.

\item[\emph{Wick--Wightman--Wigner superselection (foundational
quantum physics).}] \textit{Analogy.}
The WWW superselection rules require that every physically
admissible pure state lie in a unique joint eigenspace of the
superselection observables --- univalence \((-1)^{2J}\) and electric
charge \(Q\).
Coherent superpositions across distinct sectors are excluded from
the physical state space because the relative phase is unobservable
(the superselection observables commute with every observable), so
the would-be superposition reduces to a sector-mixture rather than a
pure state
\citep{WickWightmanWigner1952IntrinsicParity}.

\item[\emph{Gale--Shapley deferred-acceptance matching
(mechanism design).}] \textit{Analogy.}
In a two-sided matching market with complete strict preferences and
equal market sides, the Gale--Shapley deferred-acceptance algorithm
terminates with a stable matching \(\mu\) that is a function: every agent is assigned exactly one partner
\citep{GaleShapley1962CollegeAdmissions}; later work showed that
truth-telling is a dominant strategy for the proposing side
\citep{DubinsFreedman1981Machiavelli,Roth1982EconomicsMatching}.  An
orphan residual, an agent with no partner \(\mu(x)\), cannot occur: that
every agent is matched plays the role of the \Fref{F9} coverage rule,
and that \(\mu\) is a function plays the role of agreement of statuses
at a common locus.
\end{description}
\endgroup

\subsection*{\Fref{F4} Local-Global Obstruction Normal Form}
\begingroup\small
\begin{description}
\item[\emph{Schema mapping in data integration (computer science).}] \textit{Analogy.}
In data integration, source instances over local schemas are
glued to a mediated global schema through declared
tuple-generating and equality-generating dependencies (tgds, egds).
Local-overlap compatibility is the egd-and-tgd satisfaction check
at the chase; egd failure in the chase signals that no solution exists; under weakly acyclic tgds the chase terminates and, when no egd fails, yields a universal solution, unique up to homomorphic equivalence (its core up to isomorphism), from which certain answers are computed \citep{Fagin2005DataExchange}.

\item[\emph{Genome assembly from overlapping reads (biology).}] \textit{Analogy.}
Sequencing reads are local sequence data on intervals of an
unknown genome; overlap-layout-consensus and de Bruijn graph
assemblers test pairwise overlaps for sequence agreement and
attempt to stitch compatible reads into a single contig.  The
four statuses appear as overlap conflicts, fragmented contigs
without a single traversal, repeat-induced chimeric-vs-unique
ambiguity, and unique reference assembly
\citep{PevznerTangWaterman2001Eulerian}.

\item[\emph{Pose-graph SLAM with loop closures (engineering).}] \textit{Analogy.}
A robot trajectory and map are estimated from local sensor
observations whose loop closures impose pairwise pose constraints.
The four statuses appear as chi-squared loop-closure rejection
(perceptual aliasing), nonzero cocycle drift around a cycle (no
global \(SE(3)\) embedding), gauge ambiguity (global rigid
transform, monocular scale), and converged g2o\,/\,GTSAM solution
with positive-definite information matrix
\citep{Cadena2016SLAMSurvey}.
\end{description}
\endgroup

\subsection*{\Fref{F14} Duality Fixity / Self-Dual Trace Confinement}
\begingroup\small
\begin{description}
\item[\emph{Bender--Boettcher {PT}-symmetric Hamiltonians
(foundational physics).}] \textit{Analogy.}
For a spin-\(0\) quantum system, the composite \(\mathcal{PT}\)
operator is antilinear with \((\mathcal{PT})^2=\mathrm{id}\), an
involution on the Hilbert space.  Bender and Boettcher conjectured, on numerical and WKB evidence (reality for \(\varepsilon\ge0\) was later proved by Dorey, Dunning and Tateo), that
\(\mathcal{PT}\)-symmetric non-Hermitian Hamiltonians (such as
\(H=p^2+x^2(ix)^{\varepsilon}\), including \(H=p^2+ix^3\) at
\(\varepsilon=1\)) have entirely real spectra precisely when the
\(\mathcal{PT}\) symmetry is unbroken: each physical eigenstate is
a \(\mathcal{PT}\)-fixed eigenstate; a nonzero imaginary part of an eigenvalue,
the \(\mathcal{PT}\)-skew spectral component, occurs only when the symmetry is
broken. The analogy is with confinement to a fixed locus; no domination ladder
occurs in the cited work
\citep{BenderBoettcher1998PTSymmetric}.

\item[\emph{Real signals as the Hermitian-symmetric fixed locus (signal
processing).}] \textit{Analogy.}
For a finite-energy signal \(x[n]\in\mathbb C\) over \(\mathbb Z\), the involution \(J(X)(\omega)=\overline{X(-\omega)}\) on its discrete-time Fourier transform has as fixed locus the Hermitian-symmetric spectra, those of real signals.  The
\(J\)-anti-invariant spectral component is the \Fref{F14} anti-invariant
readout, and vanishing of its integrated power forces the spectrum
into that fixed locus, that is, \(x\) real (by Parseval its power is the energy of \(\operatorname{Im}x\))
\citep{OppenheimSchafer2010DiscreteTimeSP}.

\item[\emph{Galois conjugation and the fixed-field theorem
(algebraic number theory).}] \textit{Analogy.}
For a Galois extension \(L/K\) with cyclic order-2 Galois group
\(\langle\sigma\rangle\), the fundamental theorem of Galois theory
identifies the base field \(K\) with the fixed locus \(L^{\sigma}\)
of the involution \(\sigma\).  The anti-invariant component
\(z-\sigma(z)\) is the \Fref{F14} anti-invariant readout, and \(z-\sigma(z)=0\) forces
\(z\in K\) --- confinement to the fixed field
\citep{Galois1846Memoire}.

\item[\emph{Smith theory: \(\mathbb Z/2\)-cohomology and
fixed-point sets (algebraic topology).}] \textit{Analogy.}
P.~A. Smith's theorem relates the
\(\mathbb Z/2\)-cohomology of a topological space \(X\) under a
continuous involution \(J:X\to X\) to its fixed-point set
\(X^J\).  For a finite \(\mathbb Z/2\)-CW complex \(X\),
\(\dim H^*(X^J;\mathbb Z/2)\le\dim H^*(X;\mathbb Z/2)\) (the Smith--Floyd inequality), with
equality when \(X\) is totally non-homologous to zero.  With \(\mathbb Z/2\) coefficients the sign change of \(J\) is
invisible, so the \Fref{F14} anti-invariant readout has only a
counterpart here: Borel localisation, under which the equivariant
cohomology \(H^*_{\mathbb Z/2}(X;\mathbb Z/2)\), after inverting the
generator of \(H^1(B\mathbb Z/2;\mathbb Z/2)\), is carried by the
fixed-point set
\citep{Smith1938TransformationsFinitePeriod}.
\end{description}
\endgroup

\subsection*{\Fref{F15a} Selection Obstruction / Repair Normal Form}
\begingroup\small
\begin{description}
\item[\emph{Hairy ball theorem / Poincar{\'e}--Hopf (differential
topology).}] \textit{Analogy.}
On the 2-sphere \(S^2\), no continuous nonvanishing tangent vector
field exists.  The Poincar{\'e}--Hopf theorem identifies the
obstruction with the Euler characteristic
\(\chi(S^2)=2\neq 0\), and it rules out every nonvanishing field,
equivariant or not.  \(\mathrm{SO}(3)\)-equivariance is more
restrictive still: the rotations fixing a point fix no nonzero
tangent vector there, so the zero field is the only equivariant
section of \(TS^2\).  A declared breaker --- a preferred polar axis
or a stochastic perturbation --- selects a non-equivariant field,
which still vanishes somewhere; a nonvanishing direction field exists
only after points are removed, such as the poles of the declared
axis
\citep{Poincare1885CourbesEquationsDifferentielles}.

\item[\emph{Schelling focal points in coordination games (game
theory).}] \textit{Analogy.}
In a pure-coordination game with multiple symmetric equilibria, no
symmetric pure strategy can pin down a unique equilibrium.  Schelling
identified \emph{focal points} --- salient landmarks, precedents,
or conventions distinguishing one equilibrium from its symmetric
partners --- as the \Fref{F15a} declared breaker; mixed-strategy
randomisation provides an equivariant but non-selective
alternative
\citep{Schelling1960StrategyConflict}.
\end{description}
\endgroup

\subsection*{\Fref{F15b} Spontaneous Symmetry Breaking / Branch-Selection Formation}
\begingroup\small
\begin{description}
\item[\emph{Rayleigh--B{\'e}nard convection (fluid dynamics).}] \textit{Analogy.}
A horizontal fluid layer heated from below transitions, at a
critical Rayleigh number (\(\mathrm{Ra}_c\approx 1708\) for the
standard rigid no-slip boundary case and lower for stress-free
boundaries), from a translation- and rotation-invariant
conduction state to a patterned state of convection rolls or
hexagonal cells.  Each experimental realisation fixes a particular
roll axis or hexagonal sublattice, selecting a branch in the
continuous symmetry orbit
\citep{Chandrasekhar1961HydrodynamicStability}.

\item[\emph{Soai chiral autocatalysis (chemistry).}] \textit{Analogy.}
In the Soai reaction the product catalyses its own formation with
the same chirality; starting from a slightly enantioenriched mixture, autocatalysis amplifies the initial bias to nearly enantiopure product; started without chiral bias, either enantiomer results stochastically from run to run (reported by Soai and coworkers in 2003\footnote{\url{https://doi.org/10.1016/S0957-4166(02)00791-7}.}).  The underlying kinetics are chirally symmetric, but the
autocatalytic formation operator picks a single enantiomer per
realisation
\citep{SoaiShibataMoriokaChoji1995AsymmetricAutocatalysis}.

\item[\emph{Biological homochirality of life (origin-of-life
biology).}] \textit{Analogy.}
Terrestrial biology employs L-amino acids and D-sugars almost
exclusively, despite a mirror-symmetric prebiotic chemistry.
Candidate \Fref{F15b} formation-mechanism classes include
Frank-style autocatalytic amplification, photochemical bias from
circularly polarised light, and evolutionary locking after an
early enantiomeric bias; each names a declared branch-selection
source by which the biosphere settled into one chirality branch
\citep{Bonner1991OriginAmplificationBiomolecularChirality}.

\item[\emph{Crystallisation from supercooled liquid (condensed
matter).}] \textit{Analogy.}
A homogeneous, isotropic liquid below its freezing point is
thermodynamically metastable; nucleation of a crystalline phase
breaks continuous translations and \(\mathrm{SO}(3)\) rotations to
a discrete crystallographic subgroup, with each nucleation event
selecting a particular lattice orientation and translation phase
\citep{Turnbull1950FormationCrystalNuclei}.
\end{description}
\endgroup

\subsection*{\Fref{F40} Anomaly / Symmetry Obstruction}
\begingroup\small
\begin{description}
\item[\emph{Gimbal lock in Euler-angle parametrisations
(3D graphics).}] \textit{Analogy.}
The Lie group \(\mathrm{SO}(3)\) is a smooth 3-manifold, but its
Euler-angle parametrisation has coordinate singularities --- gimbal
lock --- where two rotation axes align.  At those configurations
DomainDescends fails: the parametrisation domain breaks down and
the orientation readout loses a degree of freedom.  Quaternion or
rotation-matrix representations are the declared apparatus repair
\citep{Kuipers1999QuaternionsRotationSequences}.

\item[\emph{Discretized translation symmetry.}] \textit{Analogy.}
A continuum PDE such as Navier--Stokes has continuous translational
symmetry, and its momentum equation is in flux form, so total momentum is conserved when net boundary momentum flux and net external force vanish (for the inviscid Euler equations this is the Noether charge of translation invariance).  A finite-difference discretisation on a fixed spatial
lattice respects only a discrete translation subgroup: both
SymmetryDescends and DomainDescends fail for the continuous group.  A
conservative flux-form discretisation still conserves discrete total
momentum exactly; a non-conservative one does not preserve the
downstream momentum readout.  Flux-form spatial discretisations,
paired in time with symplectic or energy-momentum-preserving
variational integrators, are the apparatus repair
\citep{HairerLubichWanner2006GeometricNumerical}.

\item[\emph{Triangular arbitrage in currency exchange (finance).}] \textit{Analogy.}
In a frictionless ideal market, exchange rates would compose
multiplicatively to a triangular product equal to 1
(e.g.\ \(\mathrm{USD/EUR}\cdot\mathrm{EUR/JPY}\cdot
\mathrm{JPY/USD}=1\)).  Quotes that update asynchronously across currency pairs break
GroupLawPreserved on the round-trip composition, and a product
different from unity offers a triangular-arbitrage opportunity in one
direction around the triangle; bid--ask spreads and transaction fees
lower the round-trip product in both directions and so absorb small
inconsistencies.  Competitive-arbitrage pressure, which restores
consistent quotes, and explicit modelling of friction are the
declared apparatus repairs
\citep{Aliber1973InterestRateParity}.

\item[\emph{Galilean vs special-relativistic velocity addition
(foundational physics).}] \textit{Analogy.}
The would-be Galilean velocity-addition group \(v_{AB}=v_{AC}+v_{CB}\)
fails to be closed on the physical velocity domain \(|v|<c\):
DomainDescends fails when Galilean composition of sublight
velocities can exit the declared domain, and the velocity readout
fails to be preserved.  Einstein's relativistic addition for collinear velocities,
\(v_{AB}=(v_{AC}+v_{CB})/(1+v_{AC}\cdot v_{CB}/c^2)\) --- equivalently
rapidity additivity --- is the declared apparatus repair
\citep{Einstein1905Elektrodynamik}.
\end{description}
\endgroup

\subsection*{\Fref{F41} Duality Equivalence}
\begingroup\small
\begin{description}
\item[\emph{Stone duality (logic and topology).}] \textit{Analogy.}
Stone's representation theorem gives a contravariant equivalence
between Boolean algebras and Stone spaces (compact totally
disconnected Hausdorff spaces): \(\mathrm{Spec}\) sends a Boolean
algebra to its ultrafilter space, \(\mathrm{Clop}\) sends a Stone
space to its Boolean algebra of clopen subsets, and
\(A\cong\mathrm{Clop}(\mathrm{Spec}(A))\),
\(X\cong\mathrm{Spec}(\mathrm{Clop}(X))\).  Role readouts (Boolean
operations \(\leftrightarrow\) topological operations on clopens),
routes (homomorphisms \(\leftrightarrow\) continuous maps), and
status (the Boolean-algebra / Stone-space structures) are
preserved on both sides
\citep{Stone1936RepresentationsBooleanAlgebras}.

\item[\emph{Gelfand duality (operator algebras and topology).}] \textit{Analogy.}
The Gelfand--Naimark theorem gives a contravariant equivalence
between commutative unital \(C^*\)-algebras and compact Hausdorff
spaces: \(\mathrm{Spec}\) returns the space of characters and
\(C(\cdot,\mathbb C)\) returns the algebra of continuous functions,
with the Gelfand transform an isometric \(*\)-isomorphism
\(A\cong C(\mathrm{Spec}(A))\) and homeomorphism
\(X\cong\mathrm{Spec}(C(X))\).  Role readouts (\(*\)-algebra
operations \(\leftrightarrow\) continuous-function operations),
routes (\(*\)-homomorphisms \(\leftrightarrow\) continuous maps),
and status (commutative-\(C^*\) versus compact-Hausdorff) are
preserved
\citep{GelfandNaimark1943NormedRings}.

\item[\emph{Curry--Howard correspondence (logic and type theory).}] \textit{Analogy.}
For the implicational fragment of intuitionistic natural deduction
and the simply typed \(\lambda\)-calculus (with product, sum and
empty types added for full propositional logic), Curry--Howard gives an exact syntactic
equivalence between propositions and types, proofs and
\(\lambda\)-terms: introduction/elimination rules correspond to
abstraction/application, and detour
elimination (proof normalisation) corresponds to
\(\beta\)-normalisation.  Roles (formula \(\leftrightarrow\) type,
proof \(\leftrightarrow\) program), routes (proof rules
\(\leftrightarrow\) typing rules), and status (provability
\(\leftrightarrow\) inhabitation, normalisation) are preserved by
mutually inverse translations between the two formal systems
\citep{Howard1980FormulaeAsTypes}.

\item[\emph{{AdS}/{CFT} holographic correspondence (string
theory; conjectural).}] \textit{Analogy.}
Maldacena's AdS/CFT correspondence conjectures that quantum gravity
on a \((d+1)\)-dimensional anti-de Sitter spacetime is dual to a
conformal field theory on its \(d\)-dimensional asymptotic
boundary, with the canonical example type IIB string theory on
\(\mathrm{AdS}_5\times S^5\) and \(\mathcal N=4\) super-Yang--Mills
on the boundary.  Within the accepted holographic dictionary, bulk
gravitational operators map to boundary CFT operators, partition
functions match, and AdS isometries map to CFT conformal
symmetries --- the proposed \Fref{F41} role/route/status preservation
across the duality
\citep{Maldacena1997LargeN}.
\end{description}
\endgroup

\subsection*{\Fref{F17} Objectivity as Common Quotient}
\begingroup\small
\begin{description}
\item[\emph{State-machine replication via Paxos (distributed
systems).}] \textit{Analogy.}
In a Paxos-style replicated state machine, each replica observes
its own log of accepted operations while the consensus protocol
commits a single sequence that all non-faulty replicas converge
on.  The per-replica log views are \Fref{F17} access quotients, the
committed sequence is the public quotient, and the safety property
of consensus is the common-across-accesses condition; split-brain
configurations are the \Fref{F17} access-objectivity defect
\citep{Lamport1998PartTimeParliament}.

\item[\emph{The {WGS-84} geodetic reference frame (geodesy).}] \textit{Analogy.}
National and regional geodetic datums (NAD-83, ETRS-89, JGD2011)
each provide a local access quotient on geodetic position, and
each maps to WGS-84 via published transformations.  WGS-84 is the
\Fref{F17} public quotient through which the local datums factor, and
agreement under WGS-84 reduction is the common-across-accesses
property
\citep{NGA2014WGS84}.

\item[\emph{The {SI} as public measurement quotient (metrology).}] \textit{Analogy.}
National metrology institutes (NIST, NPL, PTB, etc.) maintain
their own primary standards realising the seven SI base units
through the seven defining constants
\((c,h,e,k,N_A,K_{\mathrm{cd}},\Delta\nu_{\mathrm{Cs}})\).  The SI
is the public quotient through which all national standards
factor, and SI-traceable measurement agreement across laboratories
under the CIPM Mutual Recognition Arrangement is the \Fref{F17}
common-across-accesses property; unresolved Key-Comparison
discrepancies are the access-objectivity defect
\citep{BIPM2019SIBrochure}.

\item[\emph{Thermodynamic ensemble equivalence (statistical
mechanics).}] \textit{Analogy.}
In the thermodynamic limit and for short-range systems away from
phase transitions, the microcanonical, canonical, and grand-canonical
ensembles yield equivalent equations of state for macroscopic
intensive quantities.  Each ensemble is an access quotient on
phase space, and the thermodynamic equation of state is the \Fref{F17}
public quotient through which all three factor.  Ensemble
inequivalence at long-range systems / phase-transition
pathologies is the access-objectivity defect
\citep{Ruelle1969StatisticalMechanics}.
\end{description}
\endgroup

\subsection*{\Fref{F18} Lawhood Descent}
\begingroup\small
\begin{description}
\item[\emph{Stone--{\v C}ech compactification and continuous
extensions (topology).}] \textit{Analogy.}
For a Tychonoff space \(X\), the Stone--{\v C}ech
compactification \(\beta X\) is characterised by the universal
property that every continuous bounded real-valued function on
\(X\) extends uniquely to a continuous function on \(\beta X\).
A predicate declared on a dense subset is lawful on the closure
iff it admits a continuous extension --- the \Fref{F18} descent
criterion for topological laws
\citep{Cech1937BicompactSpaces}.
\end{description}
\endgroup

\subsection*{\Fref{F19} Object Persistence}
\begingroup\small
\begin{description}
\item[\emph{Comoving galaxy identity in FLRW cosmology
(cosmology).}] \textit{Analogy.}
In a Friedmann--Lema{\^\i}tre--Robertson--Walker universe, galaxies
moving with the Hubble flow keep their comoving coordinates as the
scale factor evolves.  Comoving-coordinate identity persists under
the cosmic-time transport; peculiar velocities and gravitational
mergers are the persistence obstruction
\citep{Friedmann1922KrummungRaumes}.

\item[\emph{Round-trip identity under {JSON} serialisation
(computer science).}] \textit{Analogy.}
For a data value \(D\), the round-trip
\(D\equiv\mathrm{deserialise}(\mathrm{serialise}(D))\) is the \Fref{F19}
persistence condition for the declared canonical-form identity
quotient.  Locale-dependent number formatting, time-zone-aware
datetimes, hash-randomised dictionary ordering, and floating-point
precision loss are common persistence obstructions
\citep{Crockford2006JSON}.

\item[\emph{Corporate personhood across statutory merger
(corporate law).}] \textit{Analogy.}
Under the Model Business Corporation Act \S\,11.07, a statutory
merger transfers the corporate personhood of constituent
corporations to a single surviving (or new) corporation, which
inherits the rights, properties, and obligations of the
constituents.  The legal-identity transport descends through the
corporate-personhood quotient under declared statutory procedures;
defective merger filings or undocumented dissolution without
successor are the persistence obstruction
\citep{MBCA2016Section1107}.
\end{description}
\endgroup

\subsection*{\Fref{F20} Fact / Record Formation}
\begingroup\small
\begin{description}
\item[\emph{Blockchain ledger and confirmation depth (distributed
computer science).}] \textit{Analogy.}
In a blockchain, each transaction is recorded by cryptographic
signature, included in a block, and chained to prior blocks via
hash pointers; ``stable fact'' is properly indexed by
confirmation depth.  Signature-verification failure is the
event-recording obstruction, hash tampering or chain
reorganisations and forks are the record-stability obstruction,
and a double-spend constitutes the fact-value obstruction
\citep{Nakamoto2008Bitcoin}.

\item[\emph{Forensic chain of custody (criminal procedure).}] \textit{Analogy.}
Physical evidence is admissible only when its chain of custody
is fully documented from collection through analysis: each
transfer between custodians is the recorded event, an unbroken
custodial log is the stable record, and laboratory results
matching the seized item are the consistent fact-value.  Gaps,
mishandling, or tampering are \Fref{F20} stable-fact failures
\citep{Rice2008ElectronicEvidence}.

\item[\emph{Pacioli double-entry bookkeeping (accounting).}] \textit{Analogy.}
Pacioli's double-entry method requires every transaction to be
posted as two equal opposing entries, preserving the accounting
identity \(\text{Assets}=\text{Liabilities}+\text{Equity}\).  An
unrecorded transaction is the event-recording obstruction, a
posting error is the record-stability obstruction, and a
trial-balance discrepancy is the fact-value obstruction
\citep{Pacioli1494Summa}.

\item[\emph{Paleontological taphonomy and the fossil record
(natural science).}] \textit{Analogy.}
Fossil-record formation is the natural-history analogue of \Fref{F20}:
preserved burial events, geological persistence, and consistent
inferential analyses yield robust paleontological facts.  Pre-burial
scavenging and decay are the event-recording obstruction;
metamorphism, erosion, and exhumation damage are the
record-stability obstruction; conflicting species-identification
or dating analyses are the fact-value obstruction.  The fossil
record is therefore selective rather than complete, with stable
facts indexed by preservation-and-inferential quality
\citep{KidwellHolland2002QualityFossilRecord}.
\end{description}
\endgroup

\subsection*{\Fref{F21} Reflexive Nonclosure}
\begingroup\small
\begin{description}
\item[\emph{Sarbanes-Oxley independent-audit mandate
(corporate-finance regulation).}] \textit{Analogy.}
Since the Securities Act of 1933 and the Securities Exchange Act of
1934, public-company financial statements have been audited by an
independent public accountant rather than by the issuer's own
accounting staff; the Sarbanes-Oxley Act of 2002 tightened
auditor-independence rules and added internal-control reporting.
Management's internal-controls assessment under \S404(a) is the
scoped lower audit; the independent auditor's attestation under
\S404(b), required of accelerated filers, is the meta-layer
certificate, and the auditor is itself overseen by the Public Company
Accounting Oversight Board that the Act created, producing a tower of
meta-audits.  Complete
self-attestation by the issuer is outside the scope of
the regime \citep{SarbanesOxley2002Act}.

\item[\emph{X.509 PKI external root-of-trust
(cryptographic protocol design).}] \textit{Analogy.}
In the X.509 public-key infrastructure, an endpoint
certificate cannot self-certify its trustworthiness to a
relying party: validation requires a chain of signatures
terminating at a root certificate authority in the
client's externally maintained trust store, which acts
as the meta-layer.  Checks on the leaf alone (syntax, validity period, key
usage) are the scoped lower audit; the leaf carries its issuer's
signature, not its own, and self-signed leaves are rejected by default by
mainstream TLS clients.  Complete in-certificate
self-trust is outside the scope of the protocol
\citep{Cooper2008RFC5280}.

\item[\emph{Rice's theorem on undecidability of semantic
properties (computability theory).}] \textit{Analogy.}
Rice's theorem rules out any same-layer universal
semantic decider for non-trivial behavioural properties
of partial recursive functions; specialised to
self-application, no program can act as a complete
same-layer
semantic self-certifier for a non-trivial behavioural
property.  Bounded-resource testing on a finite input
set within a step bound is a scoped lower audit, accepted
but insufficient for the full semantic claim.  Legitimate
reflexive certification of a particular instance must
invoke a strictly stronger meta-formalism (proof-assistant
formalisation against a specification, abstract
interpretation sound for the target property, or
model-checking of a decidable abstraction), which
certifies particular claims rather than serving as a
same-layer universal decider \citep{Rice1953Classes}.

\item[\emph{M\"unchhausen trilemma in epistemology
(analytic epistemology).}] \textit{Analogy.}
Hans Albert's M\"unchhausen trilemma argues that any
attempted self-grounding of a knowledge claim faces
three exhaustive options, none of which constitutes
complete same-layer self-certification: infinite
regress, logical circularity, or arbitrary termination
in unargued axioms.  A locally consistent dialectical
chain of reasons is the scoped lower audit; a legitimate
epistemic certificate must invoke a meta-layer
commitment (foundationalist intuition, coherentist
network, externalist reliability, or transcendental
argument), each of which exits the same-layer chain of
reasons \citep{Albert1985TreatiseCriticalReason}.
\end{description}
\endgroup

\subsection*{\Fref{F22} Interface Mediation}
\begingroup\small
\begin{description}
\item[\emph{Pharmacophore / structure-activity relationship
(medicinal chemistry).}] \textit{Analogy.}
Within the declared pharmacophore model of a target
receptor, the source quotient collapses candidate ligands
to the spatial / electronic arrangement of features
required for binding (hydrogen-bond donors and acceptors,
hydrophobic regions, charged groups); the boundary is the
receptor binding-site geometry; the update is binding
and downstream signalling; the target quotient is the
pharmacological response class.  Two molecules sharing
the same pharmacophore are predicted, within the model,
to produce the same response class at the same receptor;
deviations attributable to off-pharmacophore substituents
(allosteric effects, distinct ADMET profiles) are
diagnosed as boundary-leakage corrections outside the
declared interface \citep{Wermuth1998IUPACGlossary}.

\item[\emph{Vienna Convention on Diplomatic Relations
(international law).}] \textit{Analogy.}
The 1961 Vienna Convention formalises the diplomatic
interface between sovereign states.  The carrier is the
sending state's internal political and administrative
state; the source quotient collapses internal states
producing the same officially endorsed position; the
boundary is the recognised diplomatic channel (note
verbale, official meeting, protected correspondence)
delivered via accredited diplomatic agents; the update
is the receiving state's policy response through its
foreign ministry; the target quotient is the receiving
state's diplomatic posture.  Article 41 explicitly
forbids diplomats from interfering in the receiving
state's internal affairs, codifying the prohibition on
source-internal leakage past the diplomatic boundary
\citep{ViennaConventionDiplomatic1961}.

\item[\emph{Thermodynamic state functions and boundary
protocol (classical thermodynamics).}] \textit{Analogy.}
For a thermodynamic system following a declared
boundary protocol (e.g., quasi-static, adiabatic, or
isothermal exchange under specified constraints on
volume, pressure, and particle exchange), the carrier
is the microstate space; the source quotient is the
macroscopic state (energy, volume, particle number);
the boundary is the protocol-fixed heat and work
exchange together with its constraint structure; the
update is microscopic dynamical evolution; the target
quotient reads the post-process macrostate.  Within the
declared protocol, the macrostate change depends only
on the initial macrostate and the protocol-fixed
boundary data, not on the microstate trajectory; this is macroscopic determinism,
the premise of a thermodynamic description, and the state functions
\(U\) and \(S\) are the readouts whose changes depend only on the
initial and final equilibrium macrostates, not on the process path \citep{Callen1985Thermodynamics}.
\end{description}
\endgroup

\subsection*{\Fref{F25} Causality as Stable Intervention Transport}
\begingroup\small
\begin{description}
\item[\emph{Mendelian randomisation
(instrumental-variable epidemiology).}] \textit{Analogy.}
A germline genetic variant associated with a modifiable
exposure is used as a natural instrumental intervention
to estimate the causal effect of the exposure on a
disease outcome.  The carrier is the study population;
the context quotient collapses subjects sharing the
same confounder profile; the intervention space ranges
over genotypes at the chosen instrument locus; the
abstraction collapses genotypes to their associated
exposure level; the transport is the lifelong exposure
process; the readout is the disease outcome.
Admissibility rests on the conjunction of relevance,
independence (incorporating population-stratification
control), and the exclusion restriction, together with monotonicity (or effect homogeneity); the descended
residue is then the local-average exposure-outcome causal
effect \citep{DaveySmith2003Mendelian}.

\item[\emph{CRISPR-Cas9 directed gene knockout
(molecular biology).}] \textit{Analogy.}
An engineered guide-RNA-directed nuclease introduces a
targeted double-strand break at a specified genomic
locus, producing loss-of-function alleles in an
otherwise isogenic background.  The carrier is the
population of treated cells or organisms; the context
quotient collapses cells to their genetic background and
cell type; the intervention space ranges over candidate
guide RNAs and delivery conditions; the abstraction
collapses guide RNAs to the target locus they edit; the
transport is the editing event plus downstream
biological development; the target readout is the
phenotypic measurement.  Within the declared
editing-efficiency and off-target-control regime, the
phenotypic readout descends through
(background-and-cell-type, target-locus)
\citep{Jinek2012CRISPR}.

\item[\emph{Particle-scattering cross-section
measurement (high-energy physics).}] \textit{Analogy.}
A scattering experiment fires a beam of particles of
defined type and energy at a target of defined
composition; the carrier is the ensemble of individual
beam-target interaction events; the context quotient
collapses events to the target-and-beam apparatus
configuration; the intervention space parameterises
beam-energy and target choices; the abstraction
collapses operational details to the beam-target class;
the transport is the scattering process; the target
readout of a single event is the realised scattering
angle and energy.  Individual events are stochastic; the
\Fref{F25} reading is its probabilistic form, with the
differential cross-section \(d\sigma/d\Omega\) as the
stable kernel residue over repeated events.
The alpha-particle scattering off thin foils measured by Geiger and
Marsden from 1909 to 1913, with Rutherford's 1911 analysis, is the
founding instance
\citep{GeigerMarsden1913Deflexion,Rutherford1911Scattering}.

\item[\emph{Regression-discontinuity natural experiment
(empirical econometrics).}] \textit{Analogy.}
A regression-discontinuity design uses a known
deterministic cutoff in an assignment variable (test
score for scholarship eligibility, age for retirement
benefit, vote share for incumbency) to identify the
local average treatment effect at the threshold.  The
carrier is the population of units near the threshold;
the context quotient collapses units sharing the same
local cutoff neighbourhood of the running variable; the
intervention space is the treatment-assignment
indicator; the abstraction collapses delivery-channel
details to the just-above / just-below assignment; the
transport is the policy delivery process; the target
readout is the outcome of interest.  Admissibility
holds in the limit at the threshold; residue stability
is established by continuity of potential outcomes and
no-manipulation of the running variable around the
cutoff \citep{ThistlethwaiteCampbell1960Regression}.
\end{description}
\endgroup

\subsection*{\Fref{F26} Moduli / Contingency}
\begingroup\small
\begin{description}
\item[\emph{Instruction set architecture versus
micro-architecture (computer architecture).}] \textit{Analogy.}
The carrier is the space of computer-system
implementations; the lawfulness predicate requires
correct execution of the declared instruction set
architecture (ISA); the equivalence is observational
equivalence at the ISA boundary (same program, same
architectural-state update); the moduli quotient sends
an implementation to its ISA-conformant behaviour
class.  The ISA contract (instruction encodings,
register semantics, addressing modes, observable side
effects) is a necessary readout shared by all
conforming implementations.  The micro-architecture
(pipeline depth, cache hierarchy, branch-prediction
strategy, clock frequency) is presentation-dependent: it
varies among implementations of one behaviour class.  The
System/360 family, which fixed an ISA while shipping
many distinct micro-architectures, is the founding
instance \citep{AmdahlBlaauwBrooks1964System360}.

\item[\emph{Evolutionary contingency in macroevolution
(paleobiology).}] \textit{Analogy.}
The carrier is the space of possible biotic
trajectories on Earth; the lawfulness predicate
requires consistency with natural selection,
developmental constraints, and the physical
environment; the equivalence identifies trajectories
with identical sequences of macroevolutionary
bottleneck outcomes; the moduli quotient sends a
trajectory to its bottleneck-survival-pattern class.
Selection-driven adaptation, the existence of
developmental constraints, and thermodynamically
consistent matter/energy flow are necessary readouts.
The specific assemblage of surviving phyla on Earth
(e.g., the contingent survival of chordates after the
Cambrian extinction events) is a contingent readout
that varies across moduli classes; ``replaying the
tape of life'' would traverse a different moduli class
\citep{Gould1989WonderfulLife}.
\end{description}
\endgroup

\subsection*{\Fref{F28} Composability}
\begingroup\small
\begin{description}
\item[\emph{Retrosynthetic analysis and synthon-based
total synthesis (organic chemistry).}] \textit{Analogy.}
The carriers are organic compounds; the source quotients
collapse compounds to their synthon class (connectivity
and functional groups); the interface readouts are
reactive sites available for bond formation; the
composition is a declared coupling reaction within a
specified reaction family of declared functional-group
tolerance.  The residual records yield, regio- and
stereo-byproduct distribution, and
reaction-mass-efficiency residual.  Composite-outcome
descent holds
when the same target connectivity is obtained from
distinct yet retrosynthetically-equivalent precursors;
cross-residual descent holds when the residual
byproduct distribution is invariant under
representative-level choices within the same reaction
family \citep{Corey1989LogicSynthesis}.

\item[\emph{Intermodal maritime containerisation
(industrial logistics).}] \textit{Analogy.}
The ISO 668 standard fixes the external dimensions,
corner-fitting geometry, and rated load of intermodal
freight containers.  The carriers are containers loaded
by different shippers; the source quotients collapse
them to their ISO 668-conformant external-interface
class; the interface readouts are the ISO corner
fittings and the rated gross mass; the composition is
stacking on a vessel or chaining across
truck-train-vessel handoffs; the composite quotient
classifies the consolidated load by its
external-interface footprint.  The residual captures
internal load distribution, weight class, and
damage / customs / handling residuals.  Admissibility
restricts to rated compatible cargo (weight under
ratings, stacking-strength conformance, hazmat
compatibility) \citep{ISO668_2020}.

\item[\emph{{OSI} 7-layer protocol stack composition
(network engineering).}] \textit{Analogy.}
The OSI Basic Reference Model partitions network
communication into seven layers, each defined by
service primitives at Service Access Points (SAPs).
The carriers are layer-\(N\) entity implementations;
the source quotients collapse implementations to their
SAP-service-primitive contract; the interface readouts
record the service-primitive sequence delivered at the
SAP; the composition stacks a layer-\((N{+}1)\) entity
on top of a layer-\(N\) service by invoking layer-\(N\)
service primitives; the composite quotient classifies
the composite by its end-to-end service.  The residual
records PDU header overhead, fragmentation, and
retransmission residual.  Composite-outcome descent
holds because the layer-\((N{+}1)\) service depends
only on layer-\(N\) service primitives, not on the
specific implementation of layer \(N\)
\citep{ISOIEC7498OSI}.

\item[\emph{Golden Gate modular {DNA} cloning
(synthetic biology).}] \textit{Analogy.}
Golden Gate cloning uses Type IIS restriction enzymes
that cut DNA outside their recognition site, leaving
scarless four-base overhangs designed arbitrarily.
The carriers are donor plasmids carrying DNA parts
flanked by Type IIS recognition sequences; the source
quotients collapse plasmids to their part-class plus
overhang specification; the interface readouts are the
four-base overhang fingerprints; the composition is the
one-pot enzymatic digest-and-ligate reaction; the
composite quotient classifies the assembled multi-part
construct.  The residual records misligation products,
unligated linearised backbone, and off-target
re-religation events.  Composite-outcome descent holds
because overhang complementarity determines assembly
order independent of donor-plasmid backbone
\citep{Engler2008GoldenGate}.
\end{description}
\endgroup

\subsection*{\Fref{F30} Explanation as Compression with Descent}
\begingroup\small
\begin{description}
\item[\emph{Mendeleev periodic table explaining
chemical regularities (chemistry).}] \textit{Analogy.}
The history space is the set of chemical elements
with their observed properties (atomic weight,
valence, melting point, oxidation states, reaction
products).  Mendeleev's 1869 table is the
explanation: each element is assigned a
(period, group) position based on atomic weight and
chemical-similarity grouping (later refined by atomic
number).  The descended map sends a property profile
to its table position; the target map predicts the
existence and properties of then-unknown elements
(gallium, germanium, scandium were subsequently
discovered with properties anticipated from table
gaps).  The compression is strict: the table position
summarises chemical-family behaviour without
recovering all numerical property values
\citep{Mendeleev1869Periodic}.

\item[\emph{Boltzmann {H}-theorem and
statistical-mechanical irreversibility (statistical
physics).}] \textit{Analogy.}
The history space is the microscopic phase-space
trajectory set of a dilute gas; the base records the
full microstate.  The explanation is the one-particle
distribution function \(f(\mathbf r,\mathbf v,t)\) and
the H-functional
\(H[f]=\int f\log f\,d^3v\,d^3r\).  Under the
molecular-chaos hypothesis, \(H\) is non-increasing
along Boltzmann-equation evolution; identifying
macroscopic entropy with \(S=-k_B H\), the target is
the macroscopic thermodynamic readout.  The
compression is strict: distinct microstates with the
same one-particle distribution share the same
macroscopic readout \citep{Boltzmann1872Weitere}.

\item[\emph{Hart's \emph{The Concept of Law} as
jurisprudential summary (jurisprudence).}] \textit{Analogy.}
The history space is the recorded legal practice of a
jurisdiction (statutes, decided cases, official
enactments); the base records the full legal corpus.
Hart's explanation distinguishes primary rules of
obligation from secondary rules of recognition,
change, and adjudication, together with the
identified rule of recognition.  The target, in
Hart's construction, is the legal-validity verdict on
candidate rules: a candidate counts as valid law iff
it satisfies the rule of recognition.  The target map
is exact because the rule of recognition is itself
the criterion of validity; the compression is strict
because many distinct legal corpora share the same
rule-system structure \citep{Hart1961Concept}.
\end{description}
\endgroup

\subsection*{\Fref{F31} Counterfactual Legitimacy}
\begingroup\small
\begin{description}
\item[\emph{Fogel counterfactual economic history of US
growth without railroads (cliometrics).}] \textit{Analogy.}
The context quotient collapses economic trajectories
to their factor-endowment class (capital, labour,
natural resources); the intervention abstraction
collapses transportation mixes to (railroad available,
wagon and canal substitute) categories; the route is
the social-savings simulation that re-prices freight
using non-rail substitutes for each route; the target
truth records whether the counterfactual 1890 GNP
exceeds a specified threshold (Fogel's calculation
gives a deviation of less than 5\% of actual).  The
legitimacy of the counterfactual is the \Fref{F31} statement:
the social-savings estimate descends through
factor-endowment and transportation-mix classes
\citep{Fogel1964RailroadsAmericanGrowth}.

\item[\emph{Legal ``but-for'' causation in tort law
(jurisprudence).}] \textit{Analogy.}
In common-law tort doctrine, but-for causation asks
whether the plaintiff's harm would have occurred but
for the defendant's act.  The context quotient
collapses factual histories to their legally cognisable
class; the intervention abstraction collapses
candidate acts to their legal characterisation;
admissibility encodes proximate-cause limits; the
route is the counterfactual reconstruction of ``what
would have happened'' without the defendant's act; the
target truth records whether the harm would still have
occurred.  The legitimacy of the but-for
counterfactual is the \Fref{F31} statement: the verdict
descends through legally cognisable
(context, act) classes, and overreaches (e.g.,
multiple-sufficient-cause cases where but-for fails)
trigger
doctrinal corrections such as the substantial-factor
test \citep{HartHonore1985Causation}.

\item[\emph{Off-policy evaluation in reinforcement
learning (machine learning).}] \textit{Analogy.}
An agent estimates the expected return of a target
policy \(\pi_e\) using trajectories collected under a
different behaviour policy \(\pi_b\).  The context
quotient collapses environments to their (state-action
transition kernel, reward function) class; the
intervention abstraction collapses target policies to
their induced visitation; admissibility requires the
support condition \(\pi_b>0\) wherever \(\pi_e>0\);
the route is the importance-sampling re-weighting; the
target truth records whether the estimated return of
\(\pi_e\) exceeds a specified threshold.  The
Precup-Sutton-Singh result gives conditions under
which the off-policy estimator descends from
behaviour-policy trajectories to a legitimate
counterfactual verdict on the target policy's value
\citep{PrecupSuttonSingh2000Eligibility}.

\item[\emph{Climate-model counterfactual emission
scenarios (climate science).}] \textit{Analogy.}
In atmospheric general-circulation modelling, the
context quotient collapses climate trajectories to
their baseline-forcing class (pre-industrial CO\(_2\),
solar and orbital parameters); the intervention
abstraction collapses scenarios to an idealised
greenhouse-gas forcing class (e.g., doubled
atmospheric CO\(_2\) concentration); the route is the
numerical integration of the coupled atmosphere
model under the chosen forcing; the target truth
records whether a specified climate diagnostic (e.g.,
equilibrium global-mean surface temperature change)
exceeds a threshold.  Manabe \& Wetherald 1967 is the
founding instance, with its calculation of equilibrium
climate sensitivity to doubled CO\(_2\) in a one-dimensional
radiative-convective model, a calculation the same authors repeated
in a general-circulation model in 1975; legitimacy
is the \Fref{F31} statement that the threshold-crossing
verdict descends through (baseline, idealised forcing)
classes \citep{ManabeWetherald1967Thermal,ManabeWetherald1975Doubling}.
\end{description}
\endgroup

\subsection*{\Fref{F33} Horizon Sufficiency / Holographic Compression}
\begingroup\small
\begin{description}
\item[\emph{{TCP/IP} encapsulation (network protocols).}] \textit{Analogy.}
A network message is encapsulated as a sequence of
nested headers wrapping a payload.  The interior is the
full packet state (headers and payload together); the
boundary readout records the IP and TCP header fields
(source/destination addresses, ports, sequence numbers,
flags, checksum).  The exterior classifier records the
IP protocol-number class (TCP, UDP, ICMP) and the
link-layer class.  Routing, flow-control, retransmission,
and demultiplexing decisions descend through the product
quotient: two packets sharing the same header fields and
the same protocol stack class receive the same routing
verdict, regardless of payload bytes
\citep{Postel1981RFC791,Postel1981RFC793}.

\item[\emph{Fisher sufficient statistic (mathematical
statistics).}] \textit{Analogy.}
For an exponential-family parametric model
\(p(x\mid\theta)\), a sufficient statistic
\(T(X_1,\dots,X_n)\) summarises a sample so that the
conditional distribution of the sample given \(T\) is
independent of \(\theta\).  The interior is the raw
sample space; the boundary readout is the sufficient
statistic; the exterior factor is a one-point set (an ancillary statistic is a function of the same sample and cannot fill \Fref{F33}'s separate exterior slot).  By the Fisher--Neyman factorization theorem,
likelihood-based inference functionals --- likelihood
ratio, score where defined, maximum-likelihood maximizing set (and any estimate selected solely from \(T\)), and
likelihood-equivalent functionals --- descend through \(T\), hence through
the product quotient \citep{Fisher1922Mathematical}.

\item[\emph{Cauchy's integral formula in complex analysis
(complex analysis).}] \textit{Analogy.}
For a function \(f\) holomorphic on a simply connected
domain \(D\subseteq\mathbb C\) and a point \(z_0\) in
the interior of a positively oriented simple closed
contour \(\gamma\subset D\),
\[
  f(z_0)=\frac{1}{2\pi i}\oint_\gamma \frac{f(z)}{z-z_0}\,
  dz.
\]
The interior parameterises the choice of contour
\(\gamma\) together with the holomorphic-function trace
\(f|_\gamma\) on it; the boundary readout records that
trace; the exterior classifier records the holomorphy
class of \(f\).  The target value \(f(z_0)\) at any
interior point descends exactly through the product
quotient: it is given by the Cauchy kernel integral
applied to the boundary trace
\citep{Ahlfors1979ComplexAnalysis}.
\end{description}
\endgroup

\subsection*{\Fref{F35} Scale Descent / Renormalization}
\begingroup\small
\begin{description}
\item[\emph{Multigrid methods on a nested mesh
hierarchy (numerical analysis).}] \textit{Analogy.}
For a linear elliptic PDE discretised on a nested
sequence of meshes
\(V_0\subset V_1\subset V_2\subset\cdots\) with mesh
sizes \(h_0>h_1>h_2>\cdots\), the carrier is the
finest-grid function space \(V_N\).  The scale quotient
is the composite restriction \(R_{n+1}\cdots R_N:V_N\to V_n\); the coarsening
is the one-step restriction \(R_n:V_n\to V_{n-1}\).  The discrete Galerkin variational
form \(A_n u_n=f_n\) at each level is the target law;
Galerkin coarsening \(A_{n-1}=R_n A_n P_n\) (with
prolongation \(P_n\) and \(R_n=P_n^{\!\top}\)) makes the coarse
operator the variational restriction of the fine one: for a symmetric
positive-definite \(A_n\) the coarse solution is the energy-norm
projection of the fine solution onto the coarse space, not its
restriction, so the coarse law is a descended approximation rather
than an exact descended law.  The multigrid V-cycle corrects it with
fine-level smoothing for asymptotically optimal solution costs \citep{Hackbusch1985Multigrid}.

\item[\emph{Coarse-grained molecular dynamics
(computational physics and materials).}] \textit{Analogy.}
For a biomolecular system, the carrier is the
all-atom phase-space trajectory.  The scale quotient
groups atoms into coarse-grained beads (e.g.,
MARTINI mapping of heavy atoms to lipid beads); the
coarsening iterates the grouping to even coarser
representations.  The target law at each level is a
structural-statistical observable matched between
the all-atom and coarse-grained models (radial
distribution functions, free-energy landscapes,
thermodynamic averages); the renormalised law is the
effective bead Hamiltonian or potential of mean force
chosen to reproduce those observables.  Scale descent
succeeds in the equilibrium structural regime and
becomes path-dependent for dynamical observables
sensitive to integrated-out fast atomic motions
\citep{Voth2008CoarseGraining}.

\item[\emph{Lucas critique on policy-invariant micro
parameters (macroeconomics).}] \textit{Analogy.}
The carrier is the set of individual agent decision
histories under a declared policy regime.  The scale
quotient collapses individual decisions to
macroeconomic aggregates (output, inflation,
employment); the coarsening sends aggregates under
one policy regime to aggregates under another via
macro-econometric extrapolation.  The Lucas critique
observes that reduced-form macro laws do not descend
through this coarsening when the policy change
itself alters agents' decision rules: the obstruction
set is non-empty.  Legitimate scale descent requires
identifying policy-invariant deep-structural micro
parameters and aggregating through them
\citep{Lucas1976Critique}.
\end{description}
\endgroup

\subsection*{\Fref{F36} Singular Limit Completion}
\begingroup\small
\begin{description}
\item[\emph{Prandtl boundary-layer matched asymptotic
expansion (fluid dynamics).}] \textit{Analogy.}
For the incompressible Navier-Stokes equations with
small viscosity \(\nu\), the set of limiting
processes is the family of solutions parameterised by
\(\nu\).  The limit quotient sends each solution to
its formal \(\nu\to 0\) limit; the target readout
records the velocity field.  The naive Euler limit
fails to satisfy the no-slip boundary condition;
Prandtl's boundary-layer theory introduces an inner
expansion in the rescaled wall-normal coordinate
\(y/\sqrt{\nu}\) and matches it to the outer Euler
solution.  The matched-asymptotic structure is the
canonical path-memory repair, carrying the
boundary-layer thickness as an additional quotient
parameter that completes the singular limit
\citep{Prandtl1904BoundaryLayer}.

\item[\emph{Critical-point scaling at second-order
phase transitions (statistical mechanics).}] \textit{Analogy.}
For a magnetic system near its critical temperature
\(T_c\), the set of limiting processes is the family
of approach paths to \((T,h)=(T_c,0)\) in the
temperature-field plane.  Order parameter,
susceptibility, and correlation length diverge or
vanish with universal exponents that depend on the approach path (\(|m|\propto(-t)^\beta\) on a selected ordered branch as \(t\uparrow0\) along \(h=0\), and \(|m|\propto|h|^{1/\delta}\) along \(t=0\)), with non-universal amplitudes.
The Widom scaling hypothesis is the path-memory
repair: it records the universal scaling functions
parameterised by \(h/|t|^{\beta\delta}\) (with
\(t=(T-T_c)/T_c\)), completing the critical-point
readout by labelling each approach direction
\citep{Stanley1971PhaseTransitions}.
\end{description}
\endgroup

\subsection*{\Fref{F37} Complementarity as Non-Joint Access}
\begingroup\small
\begin{description}
\item[\emph{Joint diagonalisation of non-commuting
Hermitians (linear algebra).}] \textit{Analogy.}
For a finite-dimensional Hilbert space \(H\) and two
Hermitian operators \(A,B\) on \(H\), the access
quotient \(q_A\) records the spectral labelling of
\(A\) in a basis and \(q_B\) records that of \(B\).
A joint admissible quotient is a single orthonormal
basis simultaneously diagonalising both operators.
By the spectral theorem, such a basis exists iff the
commutator \([A,B]=AB-BA\) vanishes.  When
\([A,B]\ne 0\), no admissible joint quotient exists;
the access quotients are \Fref{F37}-complementary within
the orthonormal-basis class
\citep{Halmos1974FiniteDimensional}.

\item[\emph{Intuitionistic vs classical logic
(mathematical logic).}] \textit{Analogy.}
For a propositional language \(L\) with atomic
propositions, the access quotient \(q_A\) sends a
formula \(\varphi\) to its intuitionistic-validity
class under the Brouwer-Heyting-Kolmogorov (BHK)
interpretation; \(q_B\) sends \(\varphi\) to its
classical-proof class requiring the law of excluded
middle (LEM) and proof by contradiction.  Within the
admissibility class of BHK-respecting proof systems,
no joint quotient validates LEM for undecidable
propositions \(P\): classical proof of
\(P\vee\neg P\) is admissible classically but lacks
a constructive witness.  The two proof-system
quotients are \Fref{F37}-complementary, formalised by
Heyting's intuitionistic logic
\citep{Heyting1930Logik}.

\item[\emph{{CAP} theorem in distributed systems
(distributed systems).}] \textit{Analogy.}
For a replicated distributed data store, \(q_A\)
records linearisable-consistency observations and
\(q_B\) records always-available reads and writes.
The admissibility constraint is partition tolerance:
the implementation must continue under arbitrary
network partitions.  Gilbert-Lynch's formalisation of
Brewer's CAP conjecture shows that, under network
partitions, no admissible implementation jointly
realises both \(q_A\) and \(q_B\); the access
quotients are \Fref{F37}-complementary within the
partition-tolerant implementation class
\citep{GilbertLynch2002CAP}.

\item[\emph{Russell's paradox under naive
comprehension (set-theoretic foundations).}] \textit{Analogy.}
For a candidate set \(R\), the access quotient
\(q_A\) imposes the self-membership predicate
\(x\in R\iff x\notin x\); the access quotient
\(q_B\) requires \(R\) to be a set in the admissible
class.  Substituting \(R\) for \(x\) yields
\(R\in R\iff R\notin R\), so no joint admissible
quotient exists under unrestricted comprehension.
The standard remedies (ZFC's separation schema, which forms \(\{x\in A:x\notin x\}\) only inside a given set \(A\), and Russell--Whitehead type
stratification) refine admissibility to block this
joint quotient \citep{Russell1902Letter}.
\end{description}
\endgroup

\subsection*{\Fref{F38} Arrow from Record Monotonicity}
\begingroup\small
\begin{description}
\item[\emph{Steno's law of superposition (geology).}] \textit{Analogy.}
The source history set is the configuration of
sediment-deposition processes; \(\kappa\) is the
geological-time evolution.  The record quotient sends
each deposition event to its stratum class (named bed,
geological epoch); the status readout is the
stratigraphic depth.  Steno's law of superposition
asserts that, in undisturbed sequences, older strata
lie below younger strata.  Folding, faulting, and
overturning are \Fref{F38}-obstructions that require
unfolding to recover the arrow \citep{Steno1669Prodromus}.

\item[\emph{Dendrochronological tree-ring sequence
(paleoclimate and archaeology).}] \textit{Analogy.}
The source history set is the sequence of annual
tree-ring formation events; \(\kappa\) records one
annual growth cycle.  The record quotient sends a
tree-ring to its ring-pattern signature; the status
readout is the calendrical year identified by master
tree-ring chronology cross-dating.  Each successive
ring is added outside the previous ones, producing a
strictly monotone chronological sequence.  Douglass's
master chronology cross-dates archaeological wood to
its year of growth via pattern matching against the
reference sequence \citep{Douglass1919ClimaticCycles}.

\item[\emph{Libby radiocarbon dating (archaeology and
nuclear chemistry).}] \textit{Analogy.}
The source history set is the moment of organism death
(after which exchange of \(^{14}\)C with the atmosphere
ceases); \(\kappa\) records radioactive decay over
time.  The record quotient sends a sample to its
measured \(^{14}\)C\(/^{12}\)C ratio; the status readout
is the radiocarbon age inferred from the exponential decay law;
conventional radiocarbon ages use the Libby half-life of 5568 years
(the measured value is about 5730 years) and are converted to
calendar ages by calibration against tree-ring and varve
chronologies.  The ratio decreases monotonically, providing a
chronological arrow up to about 50{,}000 years before present \citep{Libby1955RadiocarbonDating}.

\item[\emph{Historical phonological reconstruction.}] \textit{Analogy.}
The source history set is the lexicon of a
reconstructed proto-language; the target is the
lexicon of an attested daughter language; \(\kappa\)
records the regular sound-change rules carrying
proto-forms forward to daughter reflexes.  The record
quotient sends a proto-form to its phonological class;
the status readout is the chronological layer in the
family tree, increasing from the proto-language toward
each attested daughter.  Brugmann's regularity
principle asserts that sound changes are exceptionless
within a defined phonological environment, providing a
monotone forward arrow along the family tree
\citep{Brugmann1886Grundriss}.
\end{description}
\endgroup

\subsection*{\Fref{F39} Entropy as Fiber Volume}
\begingroup\small
\begin{description}
\item[\emph{Herfindahl-Hirschman market-concentration
index (industrial economics).}] \textit{Analogy.}
For a market with firms having shares \(s_i\)
(\(\sum_i s_i=1\)), the access quotient sends each
customer transaction to the firm serving it.  The
Herfindahl-Hirschman index
\(\mathrm{HHI}=\sum_i s_i^2\) is the probability that
two random transactions are served by the same firm.
Under firm grouping (mergers, parent-company
aggregation), nonnegativity of the cross terms in the square of a grouped share
forces the sum of squared shares to non-decrease, so HHI is monotone under
coarsening.  The US Department of Justice merger
guidelines use HHI thresholds to identify market
concentration concerns \citep{Hirschman1964HHI}.

\item[\emph{Logarithm of the multinomial coefficient
(combinatorics).}] \textit{Analogy.}
For \(N\) distinguishable items distributed into \(k\)
labelled cells, the fine quotient sends an assignment
to its cell-content profile
\((n_1,\dots,n_k)\); the fine fiber has cardinality
\(\binom{N}{n_1,\dots,n_k}=N!/\prod_i n_i!\).
Coarsening merges two cells \(i,j\) into a single
labelled cell of size \(n_{\{i,j\}}=n_i+n_j\); the
coarse fiber is the disjoint union of fine fibers
with \(n_i'+n_j'=n_{\{i,j\}}\), so its cardinality is
\(\sum_{n_i'+n_j'=n_{\{i,j\}}}
\binom{N}{n_i',n_j',\dots}\), at least any single
summand.  The combinatorial fiber-volume entropy is
thus monotone non-decreasing under cell coarsening
\citep{Knuth1997TAOCP}.

\item[\emph{Kingman coalescent effective population
size (population genetics).}] \textit{Analogy.}
For a haploid neutral genealogical model, the access
quotient sends a genealogy of \(n\) sampled lineages
to its lineage-merger pattern.  The Kingman coalescent
gives pairwise coalescence rate \(1/N_e\), where
\(N_e\) is the effective population size; with \(n\) lineages the expected waiting time to the next merger is \(N_e/\binom n2\).  Changing \(N_e\) rescales coalescent time rather than coarsening a quotient, so this is only a loose analogue of the \Fref{F39} inequality; coarsening a finite carrier of merger patterns with fiber-volume entropy gives an instance \citep{Kingman1982Coalescent}.
\end{description}
\endgroup

\subsection*{\Fref{F42} Irreducibility / No Short Closed Description}
\begingroup\small
\begin{description}
\item[\emph{Quine-McCluskey minimum sum-of-products
form (digital-logic design).}] \textit{Analogy.}
For a Boolean function \(f:\{0,1\}^n\to\{0,1\}\), the
admissible compression class is sum-of-products
expressions; routes are Boolean-algebra equivalences
(De Morgan, absorption, distribution).  Quine-McCluskey
yields a minimum-literal expression by generating all prime implicants and then solving a minimum-weight cover problem, weighting each prime implicant by its number of literals (NP-hard in general) for the minterms not covered by essential prime implicants; the
function is \Fref{F42}-irreducible at the minimum-cover
expression \citep{Quine1952Boolean}.

\item[\emph{Information-theoretic
\(\Omega(n\log n)\) sorting lower bound (analysis of
algorithms).}] \textit{Analogy.}
For comparison-sorting on \(n\) distinct keys, the
admissible compression class is comparison-based
decision trees; routes are two-key comparisons; the
sorted permutation is the target.  A correct decision
tree must have at least \(n!\) leaves, so its
worst-case depth is \(\Omega(\log n!) = \Omega(n\log
n)\).  No asymptotically shallower comparison-based
decision tree exists; the \Fref{F42} irreducibility is the
unconditional \(\Omega(n\log n)\) lower bound
\citep{Knuth1998TAOCPv3}.

\item[\emph{Minimal genome for self-replicating cell
(synthetic biology).}] \textit{Analogy.}
Hutchison and colleagues designed and synthesized
JCVI-syn3.0, a \emph{Mycoplasma} genome of 531
kilobases encoding 473 genes that supports
self-replication under declared laboratory growth
conditions.  The design retained genes classified as
essential and also quasi-essential genes needed for
robust growth \citep{Hutchison2016MinimalGenome}.
The analogy is with the search for a short closed
description within a declared class of designs.  The
experiment tests only part of that class, so it does
not establish the minimality that \Fref{F42}
requires.

\item[\emph{Crystallographic irreducible Brillouin
zone (solid-state physics).}] \textit{Analogy.}
For a periodic crystalline solid with space group
\(G\), the admissible compression is the
restriction of crystal momentum \(\mathbf k\) to the
irreducible Brillouin zone (IBZ), the fundamental
domain of the reciprocal-lattice point-group action;
routes are point-group operations.  The
Bouckaert-Smoluchowski-Wigner analysis shows the
IBZ is the minimal subset of the full Brillouin
zone that, together with symmetry operations,
generates all distinct electronic states.  The IBZ
is \Fref{F42}-irreducible by construction
\citep{BouckaertSmoluchowskiWigner1936Theory}.
\end{description}
\endgroup

\subsection*{\Fref{F43} Continuum Emergence}
\begingroup\small
\begin{description}
\item[\emph{p-adic completion in number theory
(algebraic number theory).}] \textit{Analogy.}
For a prime \(p\), the discrete approximations are
the finite quotient rings \(\mathbb Z/p^i\mathbb Z\)
with reduction maps
\(\phi_{ij}:\mathbb Z/p^j\to\mathbb Z/p^i\).  The
candidate limit is the p-adic integers
\(\mathbb Z_p=\varprojlim\mathbb Z/p^i\mathbb Z\),
with projections to each finite quotient.  Tower
coherence is congruence compatibility, cofinality
is directedness of p-power indices, and vanishing
residuals correspond to convergence in the p-adic
metric.  \(\mathbb Z_p\) is the p-adic completion of
\(\mathbb Z\); the non-archimedean continuum
\(\mathbb Q_p=\mathbb Z_p[1/p]\) is the p-adic completion of
\(\mathbb Q\), distinct from the archimedean reals
\citep{Hensel1908Theorie}.

\item[\emph{Thermodynamic limit in statistical
mechanics (statistical physics).}] \textit{Analogy.}
For a classical-spin lattice model on a finite
volume \(\Lambda\subset\mathbb Z^d\), the discrete
approximations are the finite-volume Gibbs kernels at inverse
temperature \(\beta\), which condition on the configuration outside
the volume and are consistent under nesting of volumes; finite-volume
Gibbs measures with one fixed boundary condition are not consistent
under restriction.  The candidate limit is
the set of Dobrushin-Lanford-Ruelle (DLR) states
on \(\{+1,-1\}^{\mathbb Z^d}\), with projections
to finite-volume marginals.  Tower coherence is
the DLR consistency condition, cofinality is
exhaustion by finite volumes, and vanishing
residuals correspond to absence of unstable
boundary contributions.  The set of infinite-volume
Gibbs states emerges as the \Fref{F43} stable limit quotient; it need
not be a singleton, and at low temperature in dimension \(d\ge 2\)
the Ising model has several (phase coexistence)
\citep{Ruelle1969StatisticalMechanics}.

\item[\emph{Krull absolute Galois group as inverse
limit (field theory).}] \textit{Analogy.}
For a field \(K\) with separable closure
\(K^{\mathrm{sep}}\), the discrete approximations
are finite Galois groups
\(\mathrm{Gal}(L_i/K)\) for finite Galois
extensions \(L_i\), with restriction maps for
\(L_i\subseteq L_j\).  The candidate limit is the
absolute Galois group
\(\mathrm{Gal}(K^{\mathrm{sep}}/K)
=\varprojlim_i\mathrm{Gal}(L_i/K)\), with
projections by restriction.  Tower coherence is
compatibility of restrictions, cofinality is
directedness of finite Galois extensions, and
vanishing residuals correspond to profinite
separation by the finite quotients.  Krull's
infinite Galois theory establishes the absolute
Galois group with its profinite (Krull) topology
as the \Fref{F43} stable limit quotient
\citep{Krull1928Galois}.

\item[\emph{Born-Huang continuum elasticity from
atomic lattice (solid-state physics).}] \textit{Analogy.}
For a crystalline solid, the discrete
approximations are atomic-lattice models at
finer-and-finer lattice spacings; comparison maps
are inter-lattice coarse-graining operations.
The candidate limit is the continuum elasticity
theory governing displacement fields on
\(\mathbb R^3\); projections restrict continuum
solutions to lattice points.  Tower coherence is
compatibility of inter-lattice strain
representations, cofinality is the existence of
sufficiently fine lattices for any continuum
scale, and vanishing residuals correspond to the
long-wavelength approximation being valid.  The
Born-Huang derivation of elastic constants from
interatomic potentials is the canonical
solid-state \Fref{F43} continuum emergence
\citep{BornHuang1954Dynamical}.
\end{description}
\endgroup

\subsection*{\Fref{F44} Criticality / Phase Boundary}
\begingroup\small
\begin{description}
\item[\emph{Reynolds laminar-turbulent transition
(fluid dynamics).}] \textit{Analogy.}
For incompressible smooth-pipe flow, the control
space is the Reynolds number
\(\mathrm{Re}=\rho v D/\mu\); the phase readout
records laminar vs turbulent.  Pipe-flow
experiments since Reynolds's identify a critical
\(\mathrm{Re}_c\approx 2300\) as the phase boundary at which laminar
flow loses empirical robustness to finite-amplitude perturbations at
ordinary inflow disturbance levels; with carefully quieted inflow,
laminar flow persists far above this value, to about 13{,}000 in
Reynolds's own experiments.  The
empirical robustness margin vanishes, the
intermittent-puff response is non-uniform, and the
residual status is described by transition
intermittency statistics
\citep{Reynolds1883Turbulent}.
\end{description}
\endgroup

\subsection*{\Fref{F47} Fine-Tuning as Small Selector Region}
\begingroup\small
\begin{description}
\item[\emph{Pharmaceutical therapeutic window
(clinical pharmacology).}] \textit{Analogy.}
The moduli space is the set of dose-and-timing
protocols; the readout is the steady-state plasma
concentration; the target predicate restricts to
concentrations between the minimum effective
concentration (MEC) and the maximum tolerated
concentration (MTC).  The therapeutic-window selector
region is small for drugs with narrow therapeutic
indices (warfarin, lithium, digoxin); realized
fine-tuning is the clinically attained dose protocol
placing plasma concentration within the window
\citep{GoodmanGilman2018Pharmacology}.

\item[\emph{Mechanical tolerance stack-up gap
(manufacturing engineering).}] \textit{Analogy.}
The moduli space is the set of component dimension
tuples within their tolerance bands; the readout is
the resulting critical assembly dimension; the
target predicate restricts to assemblies meeting the
critical-dimension specification.  The selector
region shrinks as the assembly chain lengthens;
realized fine-tuning is the actually manufactured
component-dimension tuple landing inside the
allowable region \citep{Bjorke1989Tolerance}.

\item[\emph{Hutchinson ecological niche (ecology).}] \textit{Analogy.}
The moduli space is the multi-dimensional
environmental parameter space (temperature,
humidity, pH, salinity, prey abundance); the readout
is the species' population growth rate; the target
predicate restricts to non-negative growth.  Within
the declared environmental envelope for the species,
the niche occupies a small fraction of total moduli
volume; realized fine-tuning is the species' actual
occupancy of the niche
\citep{Hutchinson1957Concluding}.

\item[\emph{Nyquist closed-loop stability region
(control theory).}] \textit{Analogy.}
The moduli space is the set of (plant,
controller-parameter) combinations; the readout is
the closed-loop pole location in the complex plane;
the target predicate restricts to all poles having
negative real part.  The Nyquist criterion
characterises the stability selector region: the
closed-loop system is stable iff the number of
counter-clockwise encirclements of \(-1+0i\) equals
the number of open-loop right-half-plane poles.  For
systems with large gain, delay, or open-loop
instability, the stability region can be small,
instantiating \Fref{F47} fine-tuning of controller
parameters \citep{Nyquist1932Regeneration}.
\end{description}
\endgroup

\subsection*{\Fref{F49} Common-Source / Nonlocal Correlation Normal Form}
\begingroup\small
\begin{description}
\item[\emph{Doolittle lateral gene transfer
(microbial phylogenetics).}] \textit{Analogy.}
For prokaryotic phylogenomics, the source quotient
sends a gene-distribution configuration to its
vertical-inheritance most-recent-common-ancestor
class; the readouts are per-gene tree topologies
and presence patterns.  Vertical inheritance gives
CommonSource correlations across co-inherited
genes; Doolittle observed that horizontal gene
transfer events mediated by phages, plasmids, and
conjugation break the common-source pattern,
constituting an interface channel.  A pattern is LocallyExplainable when it factors through vertical descent, a transfer interface, or a direct route; \Fref{F49} gives this as a disjunction, not an additive decomposition \citep{Doolittle1999Phylogenetic}.

\item[\emph{Plagiarism detection via shared-source
fingerprinting (digital forensics).}] \textit{Analogy.}
For pairs of documents, the source quotient sends
a pair to its shared-source text together with
declared edit / paraphrase / template
transformations; each document is fingerprinted via
n-gram or hashed-shingle representations; the
correlation readout is an overlap metric (Jaccard
similarity, MinHash).  CommonSource holds when the
overlap descends through the
shared-source-plus-declared-transformations
quotient; the interface
captures legitimate citation, common literature,
and formulaic templates not classified as
plagiarism \citep{Heintze1996Scalable}.

\item[\emph{Mosteller-Wallace stylometric authorship
attribution (statistical humanities).}] \textit{Analogy.}
For the disputed Federalist Papers, the source
quotient assigns a paper to a candidate author; the
readouts are author-conditioned function-word
frequency distributions
\(P(\mathrm{wordfreq}\mid\sigma)\) estimated from
each candidate's known corpus.  The Bayesian
posterior \(P(\sigma\mid\mathrm{wordfreq})\) is a probabilistic analogue of CommonSource, which \Fref{F49} covers only as a deterministic factorization; joint
authorship or formulaic shared templates
instantiate the interface, and residual anomalies
index unmodeled influences
\citep{MostellerWallace1964Federalist}.

\item[\emph{Common-cause failure analysis in
fault trees (safety engineering).}] \textit{Analogy.}
For a safety-critical system, the source quotient
sends a failure event to its underlying
common-cause basic event (loss of off-site power,
fire,
earthquake).  A joint failure fixed by a shared basic event is the deterministic CommonSource pattern; statistical common-cause correlation lies outside \Fref{F49}.  Fault-tree analysis
identifies the basic events that constitute the
source, codifying methodology for nuclear-plant
reliability and risk assessment
\citep{Vesely1981FaultTree}.
\end{description}
\endgroup

\subsection*{\Fref{F50} Vacuum / Background Selection}
\begingroup\small
\begin{description}
\item[\emph{Ostwald rule of stages: polymorph
selection (crystallography).}] \textit{Analogy.}
For a supersaturated solution, the moduli quotient
collapses crystallisation trajectories sharing the
same crystalline polymorph; the descended
background records the polymorph type.  Ostwald's
rule of stages predicts kinetic selection of a
metastable polymorph close in free energy to the
supersaturated state.  The vacuum predicate
identifies the ultimately stable polymorph as the
singleton target-formed modulus
\citep{Ostwald1897Studien}.

\item[\emph{Zermelo axiomatisation (foundations
of mathematics).}] \textit{Analogy.}
The closure space is the set of formal theories
purporting to be foundations of mathematics; the
moduli quotient sends a theory to its axiom-system
class (Zermelo 1908 set theory,
Fraenkel-Skolem-extended ZF, ZFC, NBG, intuitionistic
ZF,
Martin-Löf type theory).  Zermelo's 1908
axiomatisation selected a specific moduli value —
Zermelo set theory with separation and choice —
as a singleton background; subsequent extensions
(Fraenkel Replacement, Foundation, large cardinals)
refine that selection within the same \Fref{F50}
background-moduli framework
\citep{Zermelo1908Untersuchungen}.

\item[\emph{Hund's rules for atomic ground state
(atomic physics).}] \textit{Analogy.}
For a multi-electron atom, the moduli quotient
sends Pauli-allowed configurations to their
\((S,L,J)\) term-symbol class; the descended
background is the ground-state term symbol.
Hund's rules select the ground state: (i) maximize
total \(S\); (ii) for given \(S\), maximize \(L\);
(iii) for given \((S,L)\), minimize \(J\) if the
subshell is less than half-filled, otherwise
maximize \(J\).  Within the LS-coupling regime, the
vacuum predicate is the Hund-selected ground term
symbol \citep{Hund1925Linienspektren}.

\item[\emph{Pittendrigh circadian phase entrainment
(chronobiology).}] \textit{Analogy.}
For a circadian oscillator under a periodic
zeitgeber, the moduli quotient sends oscillator
trajectories sharing the same entrained phase to a
class; the descended background records the
entrained phase.  Pittendrigh's phase-response-curve
analysis explains how zeitgebers select a definite
entrained phase from the continuous phase moduli
space; the vacuum predicate is the biologically
realised entrained phase
\citep{Pittendrigh1960Circadian}.
\end{description}
\endgroup

\subsection*{\Fref{F51} Unification as Common Refinement}
\begingroup\small
\begin{description}
\item[\emph{Cartesian analytic geometry
(mathematics).}] \textit{Analogy.}
The parent closure is Cartesian analytic geometry;
the children are Euclidean geometry (synthetic
proofs) and algebra (symbolic equations).  Cartesian
coordinates supply interpretations in both
directions: every coordinate-representable
algebraic curve maps to a geometric object and
vice versa.  Euclidean theorems become algebraic
identities on coordinates, and algebraic equations
acquire geometric interpretation as curves and
surfaces \citep{Descartes1637Geometrie}.

\item[\emph{Banach functional analysis (functional
analysis).}] \textit{Analogy.}
The parent closure is Banach's framework of complete
normed vector spaces with bounded linear operators;
the children are finite-dimensional linear algebra
and classical function-space analysis.  Both embed
as concrete Banach spaces, and theorems
(Hahn-Banach, open mapping, closed graph, uniform
boundedness) supply unified results across both
children within the framework's declared scope
\citep{Banach1932Theorie}.

\item[\emph{L\'evi-Strauss structural anthropology
(anthropology).}] \textit{Analogy.}
The parent closure is L\'evi-Strauss's structural
analysis of myth via mythemes (binary oppositions,
mediating figures); the children are
ethnographically distinct mythological corpora.
Each corpus is interpreted in terms of its
structural elements, and the parent identifies
mythematic isomorphisms across geographically
isolated cultures \citep{LeviStrauss1958Anthropologie}.

\item[\emph{Huxley modern synthesis (evolutionary
biology).}] \textit{Analogy.}
The parent closure is the modern synthesis (Huxley
1942), unifying Mendelian genetics and Darwinian
natural selection together with cytology,
paleontology, and biogeography.  Population
genetics supplies the common refinement:
allele-frequency dynamics under Mendelian
inheritance extend, with selection coefficients, to a
unified
predictive framework recovering both Mendelian
ratios and Darwinian adaptive change as special
cases \citep{Huxley1942Evolution}.
\end{description}
\endgroup

\subsection*{\Fref{F52} Closure Completeness Boundary}
\begingroup\small
\begin{description}
\item[\emph{{PMBOK} project-scope management
(project management).}] \textit{Analogy.}
The residuals are deliverables and acceptance
criteria; directions are project activities and
stakeholder requests.  The PMBOK scope statement
declares deliverables (\emph{inScope}), exclusions
(\emph{namedOutside}), assumptions, constraints,
and acceptance criteria; scope creep is a
stakeholder request that is neither.  The
meta-audit is project-sponsor sign-off
\citep{PMI2017PMBOK}.

\item[\emph{International Building Code
jurisdictional scope (construction codes).}] \textit{Analogy.}
The residuals are code-mandated safety,
structural, fire-protection, and accessibility
requirements; directions are declarable
construction features.  The IBC declares scope by
occupancy classification, construction type, and
jurisdictional adoption, with detached one- and two-family dwellings and townhouses, each not more than three stories above grade plane and, for townhouses, with a separate means of egress, referred to the International Residential Code (IBC 101.2) \emph{namedOutside}.  The meta-audit is the
building inspector's certificate of occupancy
\citep{ICC2021IBC}.

\item[\emph{Patent-claim scope under {M}arkman
hearings (intellectual-property law).}] \textit{Analogy.}
The residuals are patent-claim terms requiring
construction; directions are declarable claim
scopes and accused infringing embodiments.  Under
\emph{Markman v. Westview Instruments}, claim
construction is a matter of law: each disputed term
is construed (\emph{inScope}) or held indefinite
under \S112 (\emph{namedOutside} the valid claim).
The meta-audit is the Federal Circuit affirmance of
district-court claim construction
\citep{Markman1996Westview}.

\item[\emph{{FDA} drug-label indication scope
(regulatory affairs).}] \textit{Analogy.}
The residuals are safety and efficacy items required
on the FDA-approved label; directions are medical
uses.  21 CFR 201.57 prescribes the format and
content of the ``Indications and Usage'' section,
with FDA-approved indications \emph{inScope} and
off-label uses \emph{namedOutside}.  The meta-audit
is FDA approval of the (Supplemental) New Drug
Application \citep{FDA21CFR20157}.
\end{description}
\endgroup

\subsection*{Direct CER no-go (\Cref{thm:CER-direct-no-go})}
\begingroup\small
\begin{description}
\item[\emph{Preclinical reproducibility record
(research methodology).}] \textit{Analogy.}
The carrier \(C\) is the reported experimental
result; \(\mathfrak S\) imposes the
reproducibility audit (lab notebook, raw data,
randomisation log, blinded scoring).  The
closure-relevant datum \(d\) is the
replication-enabling record.  Begley and Ellis
report that Amgen scientists confirmed only 6 of 53 ``landmark''
preclinical findings, and attribute the failures in part to missing
controls, unblinded scoring and selective reporting.  Read through
the theorem, the bench experiments were bare-realised yet the record
required to close the published claim was inadequate, so the
reduct satisfies \(\operatorname{BareReal}
_{\mathfrak S}\) but the reproducibility
residual has no satisfying record
\citep{Begley2012Drug}.

\item[\emph{Software bill of materials (software
supply chain).}] \textit{Analogy.}
The carrier \(C\) is the compiled binary in
production; \(\mathfrak S\) imposes the
supply-chain audit (component inventory, licence
terms, vulnerability mapping, build provenance).
The closure-relevant datum \(d\) is the SBOM
listing every linked component with supplier, component name, version, other unique identifiers and dependency relationships, per the NTIA minimum elements.
The binary runs without the SBOM, so the reduct
satisfies \(\operatorname{BareReal}_{\mathfrak S}\)
yet the retained claim ``free of
known-vulnerable components'' has a residual
whose audit datum (the SBOM) was never produced
\citep{NTIA2021SBOM}.

\item[\emph{Photometric flux without calibration
log (observational astronomy).}] \textit{Analogy.}
The carrier \(C\) is the recorded photon count on
a CCD; \(\mathfrak S\) imposes the photometric
audit (flat field, dark current, atmospheric
extinction, zero-point).  The closure-relevant
datum \(d\) is the per-exposure calibration log
binding ADU counts to absolute flux, as in the
SDSS \"ubercal solution.  The raw integration is
bare-realised even when \(d\) is absent, so the
reduct satisfies \(\operatorname{BareReal}
_{\mathfrak S}\) yet the retained claim ``object
has AB magnitude \(m\)'' has a residual whose
audit datum is the missing calibration
\citep{Padmanabhan2008Improved}.

\item[\emph{Chain of title for real property
(real-property law).}] \textit{Analogy.}
The carrier \(C\) is the parcel together with
the occupant's actual possession; \(\mathfrak S\)
imposes the marketable-title audit (recorded
deeds, mortgages, easements, liens).  The
closure-relevant datum \(d\) is the unbroken
chain of recorded conveyances from a sovereign
source.  Adverse occupancy bare-realises the
parcel even when an unrecorded prior deed or
undischarged mortgage breaks the chain, so the
reduct satisfies \(\operatorname{BareReal}
_{\mathfrak S}\) but the retained claim
``marketable title'' has a residual whose
required record is the missing deed
\citep{Palomar2003Patton}.
\end{description}
\endgroup

\subsection*{Converse CER no-go (\Cref{thm:CER-converse-no-go})}
\begingroup\small
\begin{description}
\item[\emph{Ratified RFC without deployed
implementation (Internet standards).}] \textit{Analogy.}
The carrier \(C\) is the standards-track RFC
considered as a closure-of-record: working-group
consensus, IETF Last Call, IESG approval, RFC
Editor publication, IANA registry allocation per
RFC 2026.  Under \(\mathfrak S\) the standards
process is closed at publication.  When no
production implementation deploys the protocol
(orphan or stillborn RFCs), no Internet host
realises \(C\), so \(C=\mathrm{Cl}_{\mathfrak S}
(C)\) yet \(\neg\operatorname{BareReal}_{\mathfrak
S}(C)\) \citep{Bradner1996RFC2026}.

\item[\emph{Statute fallen into desuetude
(statutory law).}] \textit{Analogy.}
The carrier \(C\) is the enacted statute
considered as a closure-of-record: bill passed by
both chambers, presented and signed (or enacted
over veto), codified in the official compilation.
Under \(\mathfrak S\) the legislative procedure is
closed at codification.  When prosecutors decline
to charge, courts decline to apply, and citizens
disregard the statute over sustained time, no
enforcement instance realises \(C\), so
\(C=\mathrm{Cl}_{\mathfrak S}(C)\) yet \(\neg
\operatorname{BareReal}_{\mathfrak S}(C)\); desuetude is not a recognised doctrine in
most United States jurisdictions (West Virginia is the usual
exception), and Calabresi proposes that courts be allowed to treat
such obsolete statutes as open to judicial revision \citep{Calabresi1982Common}.

\item[\emph{Superconducting Super Collider cancellation.}] \textit{Analogy.}
The carrier \(C\) is the proposed 40\,TeV
(centre-of-mass) proton--proton collider considered as a
closure-of-record: Conceptual Design Report,
Site-Specific
Conceptual Design Report, technical design
reports, DOE design reviews, environmental impact
statement, congressional authorisations.  Under
\(\mathfrak S\) the megaproject design-review
audit is closed at the review-completion
boundary.  Congressional cancellation in 1993
terminated construction before operations, so no
running collider instantiates \(C\).  Hence
\(C=\mathrm{Cl}_{\mathfrak S}(C)\) yet
\(\neg\operatorname{BareReal}_{\mathfrak S}(C)\)
\citep{Riordan2015Tunnel}.

\item[\emph{Standard {ML} signature without a
structure (programming-language type systems).}] \textit{Analogy.}
The carrier \(C\) is a Standard ML signature
considered as a closure-of-record: abstract type
declarations, value specifications with types,
exception declarations, all type-checked against
the SML definition.  Under \(\mathfrak S\) the
type-system audit is closed at signature
elaboration.  Without a matching structure binding
no runtime value inhabits the abstract types and
no function body computes the declared values, so
\(C=\mathrm{Cl}_{\mathfrak S}(C)\) yet \(\neg
\operatorname{BareReal}_{\mathfrak S}(C)\)
\citep{MacQueen1984Modules}.
\end{description}
\endgroup

\subsection*{\Fref{F53} Closure--Reality Boundary resolution}
\begingroup\small
\begin{description}
\item[\emph{Audited financial statements (financial reporting).}]
\textit{Analogy.} Closing a ledger's inference rules does not supply missing
transaction evidence or an auditor's attestation. Accounting needs actual
witnesses satisfying the declared obligations; acceptance of an opinion is
an additional consumer requirement \citep{FASB2009ASC}.

\item[\emph{ACM Artifact Review and Badging (computational reproducibility).}]
\textit{Analogy.} Deriving consequences of an artifact manifest does not
create a missing executable or reproduction log. Those witnesses must be
supplied and checked before the relevant policy can award a badge
\citep{ACM2020Artifact}.

\item[\emph{CompCert formally verified compiler (programming-language verification).}]
\textit{Analogy.} Inference closure under proof rules does not create a
missing semantic-preservation certificate. The actual proof must satisfy
the original source-to-target specification \citep{Leroy2009Formal}.

\item[\emph{End-to-end verifiable elections (election integrity).}]
\textit{Analogy.} Closing the declared public record under inference does
not create missing shuffle or decryption proofs. Actual supplied certificates
and the verification policy determine whether the original tally obligation
is delivered \citep{Adida2008Helios}.
\end{description}
\endgroup
   % app:further-instantiations
\section{Version history}\label{app:version-history}

This appendix records the differences between successive versions.  No
statement or proof of the paper depends on it.

\subsection{Version 3}\label{subsec:version-3}

Version 3 keeps every numbered result of Version 2, with its name,
number and place.

\paragraph{Results.}
\begin{itemize}
\item \Cref{thm:F13a} takes its source hypothesis \((\mathsf S)\)
  directly as a hypothesis on the interface.  Version~2 stated it for an
  arbitrary source predicate with a certificate and an obligation
  supplying \((\mathsf S)\); the theorem used no property of that
  predicate, so the two forms are equivalent.  The reading paragraph
  gives an interface satisfying \((\mathsf S)\) with a predictive
  surplus.
\item The formal development no longer contains an axiom for
  \Fref{F13a}: the property that Version~2 recorded as a guarded
  trust-base axiom is now the explicit hypothesis \((\mathsf S)\).  The
  opacity-guard subsection and the tagged-axiom register of
  Version~2 are removed accordingly.
\item \Cref{thm:F6} is formalized for real matrices, and over any
  ordered field; a remark after its proof records that the theorem holds
  over any ordered field.  Version~2 formalized it over the rationals
  only, so its formalization is now faithful.
\item The summaries of \Fref{F14} carry the measurability and
  square-integrability of the readout, and the summary of \Fref{F36}
  carries the surjectivity of the limit quotient.
\end{itemize}

\paragraph{Provenance and analogies.}
\begin{itemize}
\item The reading paragraph of \Cref{thm:F14} identifies the odd
  readout of the self-dual trace confinement paper with the case
  \(U=\mathrm{id}\), and distinguishes it from that paper's response
  involution.
\item The description of the Riemann-hypothesis application of
  \Fref{F14} follows the current version of that paper: its conditional
  theorem uses an auxiliary ledger of represented zeros, and is
  equivalent to the Riemann hypothesis on the bare shell of all
  nontrivial zeros; its concrete version uses finite windows of zeros
  closed under \(s\mapsto1-\bar s\), with unit weights, that eventually
  contain every zero, and is likewise equivalent to the Riemann
  hypothesis.
\item The description of the substrate of \Fref{F6} follows the current
  version of the Needle Killer paper, which uses Moore--Penrose
  pseudoinverses of positive semidefinite audit forms.
\item The hidden Markov model analogy of \Fref{F13a} states the
  generic identifiability theorem of Allman, Matias and Rhodes with its
  hypotheses and separates identifiability from estimation; Version~2
  stated a sufficient condition that does not suffice.
\item The Landauer analogy of \Fref{F12} is restricted to the erasure of
  an unbiased bit with a thermal reservoir and states the
  entropy-dependent bound.
\item The matching analogy of \Fref{F9} cites the incentive results of
  Dubins--Freedman and Roth, and the minimal-genome analogy of
  \Fref{F42} no longer asserts the irreducibility of the experimental
  genome.
\item Citations of the companion papers refer to their current versions.
\end{itemize}

\paragraph{Presentation.}
\begin{itemize}
\item The introduction, \Cref{sec:apparatus}, the chapter openings of
  the Access, Stability, Transport and Symmetry chapters, and
  \Cref{sec:anchor:closure-reality-boundary,sec:substrate-instances,sec:sketches,sec:nonclaims}
  are rewritten for readability; \Cref{sec:anchor:closure-reality-boundary}
  is retitled.
\item The \emph{Interaction status} paragraphs, whose first sentence held
  for every law by construction, are replaced by nonclaims paragraphs,
  and \Cref{def:birdint,def:normal-form} and
  \Cref{fig:normal-form-template,fig:axis-topology} are revised
  accordingly.
\item Notes on the formalization are collected in
  \Cref{app:formalization}: the body no longer carries formalization
  footnotes, and the summary tables no longer carry a coverage column.
\item In \Cref{sec:status-records-coherence,sec:self-reference-limits},
  proofs that repeated the descent argument cite the descent criterion
  of \Cref{sec:apparatus}, and the reading paragraphs no longer repeat
  the statements they precede; the split set of \Cref{thm:F19} is named
  \(\Delta_{\mathrm{split}}\).
\item The introduction states that most laws are elementary
  consequences of the descent criterion.
\item The paragraphs on notation and displayed equivalences in
  \Cref{sec:apparatus} are shortened, and the descent criterion is
  stated there with its proof.
\end{itemize}

\subsection{Version 2}\label{subsec:version-2}

Version 2 keeps every numbered result of Version 1, with its name, number
and place. This subsection lists the differences between the two versions:
first the principal change to each result, then further changes to results,
Lean coverage, summary rows and instantiations in the order in which they were
made, and then summaries for the Lean artifact, the apparatus and the
instantiations. The body statements are authoritative; each entry below summarizes a change.

\paragraph{Changes to results.}
\begin{itemize}
\item \Cref{thm:F2} states both repairs by their universal properties: the
  source repair \(\langle q,r\circ F\rangle\) is the coarsest refinement of
  \(q\) through which \(r\circ F\) descends, and a target map \(c\) makes
  \(c\circ r\circ F\) descend exactly when it factors through \(c_F\).
\item \Cref{thm:F3} defines when memory is forced and states the memory
  repair \(\langle\currQ,\predM\rangle\) by its universal property, as an
  instance of \Cref{thm:F2}.
\item \Cref{thm:F4} states its status partition as ``exactly one of four'',
  with each status defined, and states that a globalizing family is
  compatible; its proof gives the exclusivity step that uses this.
\item \Cref{thm:F6} reads its needle clause in the language of
  \Cref{sec:apparatus}. Its proof derives the layer-dissolving decomposition from the projection identity and the monotonicity of the Loewner order under congruence, and the sizes of the bounds and of \(S\) are stated. It defines a needle as a transported dissolving direction exceeding its budget and states that no needle exists.
\item \Cref{thm:F7} defines sufficiency and adds the split-pair form and the
  partition into exactly one of four sufficiency statuses.
\item \Cref{lem:F8} defines retention, definability from the old quotient and
  strictness, assumes the old quotient surjective, proves the split-fiber step
  from non-definability, and adds the partition into exactly one of four
  comparisons. Without surjectivity the equivalence fails: an empty history
  set with a nonempty old quotient and an empty new quotient gives a strict
  extension with no split fiber.
\item \Cref{thm:F9} is stated for the residual records of a package, with the
  inherited rule (RC) and agreement of statuses at a common locus as
  hypotheses, and adds that both hypotheses are sharp.
\item \Cref{thm:F11} names the inherited rule (NO) among its hypotheses,
  states that its factorization is unique, and says which hypothesis gives
  which conclusion. It adds that constancy on the access fibers, emptiness of the overread obstruction and unique factorization are equivalent, and that (NO) holds exactly when every accepted exact current claim factors uniquely through the access quotient.
\item \Cref{thm:F12} defines current visibility, states the properties of the
  financed refinement \(\langle\accessQ,d\rangle\), including its
  minimality, and states its last conclusion under the no-overread rule.
\item \Cref{thm:F13a} names its source predicates in paper terms, adds the
  upstream readout of the determining state, defines quotient exposure,
  predictive collapse and predictive surplus from the maps, and proves its
  fifth conjunct from the fourth; an adopted, guarded source obligation supplies
  the first four conjuncts and surjectivity of the current quotient, and no
  concrete PvNP adapter is constructed. It states the unconditional equivalence of predictive surplus with the failure of predictive collapse for a surjective current quotient, and marks the first four conjuncts as supplied by that obligation.
\item \Cref{thm:F13b} is a direct theorem for a closure with a declared
  determining-state readout: predictive surplus and hidden-state surplus are equivalent, and every surplus pair is separated by the
  determining state. Version~1 stated it under an obligation axiom.
\item \Cref{thm:F14} is a direct theorem: it states the two properties of
  the positive cone and the trace identity its proof uses, with the
  measurability, integrability and equivariance of the readout and proves that a positive anti-invariant
  ledger dominated by bounds of vanishing trace is zero, and that a
  separating readout confines visible mass to the fixed locus. Version~1
  stated it under a source certificate. Its readout is measurable, square-integrable and anti-invariant, \(\psi^-(\Jinv x)=-U\psi^-(x)\), for a declared unitary involution \(U\), and its separation condition is a null-set condition.
\item \Cref{thm:F15a} no longer carries a breaker readout as an unused hypothesis
  (the breaker readout now enters only through the breaker-extension clause), and its
  proof does not use equivariance of the orbit map.
\item \Cref{thm:F15b} lets the formation operator select a branch in
  \(\mathcal C\) itself, with \(q\circ F=q\); in Version~1 it took values in
  the branch quotient, on which the group acts trivially, so the required
  inequality \(g\cdot F(c)\ne F(c)\) could not hold. It is a direct theorem:
  a branch-selection formation record exists exactly when some ground state
  is not invariant under the symmetry, and every such record selects a branch
  in the orbit of the ground state that is not fixed by the group and is not
  the value of any equivariant selector. Its formation record is constant on orbits, so the identity map is not a record, and for a symmetric ground-state predicate the record exhibits two distinct ground states in one orbit.
\item \Cref{thm:F17} states its universal property as ``any quotient that
  factors through every \(q_i\) also factors through \(p\)''; Version~1
  stated the reverse direction, which forces \(p\) to be injective. Part (7)
  of the proof uses descent of \(r\) through \(p\) and of \(p\) through \(r\).
\item \Cref{thm:F18} defines lawfulness of a predicate as invariance under
  the bijections that preserve the access classes, assumes the access quotient
  surjective, and proves that lawfulness is descent.
\item The proof of part (5) of \Cref{thm:F19} uses the obstruction on pairs
  of classes in \(Q_i\times Q_j\). \Cref{thm:F19} declares its target set \(Q_k\).
\item The proof of \Cref{thm:F20} names the constructed descended transport
  \(\theta^{\flat}\) instead of redefining the given map \(\bar\theta\), and
  the theorem adds that, for the descended transport, fact-value stability is
  the raw fact-value equation. Its packaging clause characterizes idempotence by the equality of image and fixed locus.
\item \Cref{thm:F21} states the rules of the audit policy and a conclusion
  about the certifying meta-claim, assumes its three statuses distinct, and
  adds a two-claim model showing that its fourth rule cannot be dropped.
\item \Cref{thm:F23} defines stability as fiber-supported ledgers with the
  declared transport equation, and adds the dichotomy between descent and
  unresolvedness, the value of \(\operatorname{Pr}_{q,r}\) for a descending
  readout, and that the readout distributes the weight of each fiber. It adds that a ledger concentrated on its fiber gives the unrestricted mass of each readout value.
\item \Cref{thm:F22} states the equivalence of an empty mediation obstruction
  with the absence of boundary leakage and of internal insufficiency; Version~1
  stated one direction.
\item \Cref{thm:F24} states that a role split implies the strict layer
  extension \(q_s\); Version~1 asserted that exactly one of eight resolution
  cases holds, and the other cases are now described in the reading paragraph
  as declared statuses. It adds that \(q_s\) is coarsest in the factorization preorder among maps through which \(q\) and \(s\) both factor. Its corestriction to its image is the coarsest surjective quotient with that property.
\item \Cref{thm:F25} defines a causal effect as a stable nonzero contrast
  between two admissible interventions, and adds that, for the descended
  transports, every stable causal effect has raw witnesses.
\item \Cref{thm:F27} defines conservation as invariance under every declared
  step and the orbits as the classes of the generated equivalence.
\item \Cref{thm:F26} defines moduli readouts, presentation dependence,
  necessity and contingency, and proves the descent criterion, the exclusive
  split of a moduli readout into necessary and contingent, and the transport
  criterion; the converse step uses a pair with the same modulus.
\item \Cref{thm:F28} defines composability as descent through parts together
  with interfaces, with associativity and identity as separate conditions,
  and names \(\ell\) and \(\varrho\) as the two bracketings of a triple
  composite.
\item \Cref{thm:F29} defines compatibility of \(\sim\) with \(\pi\), states that
  presentation invariance implies invariance under \(\sim\); when
  \(\sim\) is an equivalence relation, the converse holds for every
  two-valued readout exactly when \(\sim\) is the kernel of \(\pi\).
  It also justifies \(\bar T\circ\pi=T\).
\item \Cref{thm:F30} adds the obstruction criterion for strict compression
  and the continuation clause that uses \(\kappa\), \(e'\) and \(t'\).
\item \Cref{thm:F31} assumes \(q\) and \(\alpha\) surjective, and its
  factorization form requires admissibility to descend; its reading paragraph
  notes that no obstruction uses the target quotient.
\item \Cref{thm:F34} assumes \(\rho\circ\tau\) surjective onto \(R\); in
  Version~1 surjectivity onto its image imposed no condition, and the
  equivalence failed for an empty source with a nonempty recovery quotient.
\item \Cref{thm:F35} assumes the finer scale quotient surjective, which makes
  the descended law unique, and proves both directions of the renormalization
  criterion.
\item The last line of \Cref{thm:F36} states the obstruction criterion for the
  renormalized readout; in Version~1 it was a tautology.
\item \Cref{thm:F37} quantifies over a declared class of admissible joint
  carriers and states separate admissibility; in Version~1 the product map
  was always a joint quotient, so the right-hand side could not hold. It adds
  the split-pair criterion.
\item \Cref{thm:F39} assumes a declared zero of the entropy values, attained
  exactly by quotients with singleton fibers, and adds that, when the mixing
  functional vanishes exactly on pairs with the same fibers, zero target
  entropy is zero mixing \(m(\langle q,t\rangle,q)\), as for conditional
  Shannon entropy. Its mixing clause assumes only that \(m(\langle q,t\rangle,q)\) vanishes exactly when the two quotients have the same fibers, which conditional Shannon entropy satisfies.
\item \Cref{thm:F40} concerns a symmetry that is claimed and not recorded as
  explicitly broken, and adds that the group law holds for the descended
  action of a group action. It defines the anomaly as a nonempty action or domain obstruction or a failure of readout preservation, and the induced action when the action obstruction is empty; with the induced action, an anomaly is exactly the failure of one of the four conditions.
\item \Cref{thm:F41} writes the route actions as \(\mathrm{act}_A\),
  \(\mathrm{act}_B\), states the status-value map between the packages, says
  that role-preserving duality equivalence is defined by the displayed
  conditions, and adds that the inverse map also intertwines the routes. Its role readouts live on the quotients, and the inverse map carries roles back.
\item \Cref{thm:F42} writes the cost of the competing description at \(h\)
  as \(\operatorname{cost}(e'(h))\), with the compression as a subscript of
  the irreducibility predicate, and adds that an irreducible closed
  compression exists when costs are well founded and some compression is
  closed.
\item \Cref{thm:F43} defines the stable limit quotient by its four conditions
  and proves the implication; no limit is constructed.
\item \Cref{thm:F44} writes its defining line with \(:\Longleftrightarrow\),
  states the surjectivity its descent clause uses, and adds that phase
  boundaries and criticality descend when their data do.
\item \Cref{thm:F45} declares the order on local quotients used in its
  minimality clause and defines its locality predicates as factorizations.
\item \Cref{thm:F46} writes its defining line with \(:\Longleftrightarrow\)
  and adds the selected-modulus and transport clauses. It no longer assumes a
  descended readout: descent is fiber constancy, the descended readout is
  unique, and a law-readout is a constant readout.
\item \Cref{thm:F47} writes its defining line with \(:\Longleftrightarrow\)
  and adds the consequences of realized fine-tuning and the monotonicity of
  smallness under a stricter target.
\item \Cref{thm:F49} states the tests for the joint readout together with the
  local access data, adds a declared direct-route explanation, the
  surjectivity hypotheses and a nonempty history space, and classifies an
  unexplained present correlation as a residual, nonlocal exactly when the
  declared locality package fails. Its explanation predicates are defined by
  fiber constancy, so the displayed factorizations are proved equivalences.
\item \Cref{thm:F38} and \Cref{thm:F50} open with the fiber-constancy
  criterion for the transported record and the raw background, and
  \Cref{thm:F48} defines a topological readout as a factorization through the
  gluing quotient and adds that every topological readout is exactly one of a
  complete and a coarse invariant.
\item \Cref{thm:F50} declares its set of formed targets and adds that, for the
  descended background, necessity is constancy of the raw background.
\item \Cref{thm:F51} defines exact unification by the conjunction of
  parent formation, role preservation and closure of the cross residuals, and
  declares the route systems and status translations the display uses; the
  reading paragraph says that exact unification is defined by the displayed
  conditions, and the theorem adds the common-refinement square and the
  inclusion of the parent kernel in both child kernels.
\item \Cref{thm:F52} writes its defining line with \(:\Longleftrightarrow\)
  and adds its obstruction criterion. \Cref{sec:sketches} describes it as a characterization for declared residuals and directions.
\item \Cref{prop:F53} states that its positive equivalence combines the
  closure fixed-point criterion with the definition of the retained operational predicate, with
  the structural defect named. The
  hypotheses of \Cref{thm:CER-direct-no-go,thm:CER-converse-no-go} are
  spelled out before them, the direct one through a declared reduct map, and a
  finite model with a closure operator follows each and satisfies its
  hypothesis. Both theorems state the finite model as their last clause, so
  neither implication follows from the axioms of a closure operator with
  reduct stability, and \Cref{thm:CER-converse-no-go} adds that the converse
  fails exactly when some fixed point of the closure is not bare-realised.
\item \Cref{thm:F10} and \Cref{thm:F5} assume their quotient maps surjective,
  which the coarsest-quotient clause and the recombination criterion need
  when a carrier is empty. \Cref{thm:F13a} states that its certificate is
  used only as the hypothesis of conjuncts 1--4 and that its
  certificate-free content is the last clause. \Cref{thm:F11} and
  \Cref{thm:F9} say which conclusions restate their inherited rules.
  \Cref{thm:F6} states the scope of its substrate Lean anchor, and the
  closure--reality row lists its Lean theorems and says that the positive
  equivalence holds in Lean by unfolding definitions. Summary rows for
  \Fref{F4}, \Fref{F26}, \Fref{F40} and \Fref{F52} match their theorems.
\item \Cref{thm:F13a} describes its source certificate as the guard of an
  adopted source obligation, with no identification with an interface of the
  source paper constructed, and \Cref{thm:F11} is stated
  in three parts: the factorization criterion for every claim, its
  consequence under (NO), and the characterization of (NO). Defining lines of
  \Cref{thm:F29}, \Cref{thm:F36} and \Cref{thm:F40} use
  \(:\Longleftrightarrow\), \Cref{thm:F25} writes its context classes as \(c\),
  and \Cref{thm:F20} and \Cref{thm:F50} name their record-level transport and
  background as declared maps. \Cref{def:diagnosis-source} states the determining hypothesis to which
  \Cref{thm:F13b} applies. \Cref{thm:F15a}
  adds that for the orbit quotient with its induced action, the action on
  orbits is trivial and every equivariant section takes values in \(X^G\);
  when \(O\ne\emptyset\), such a section can exist only if
  \(X^G\ne\emptyset\). The Lean theorems for \Fref{F20}, \Fref{F25} and
  \Fref{F50} assume the descent equation only for the clause that uses it.
\item \Cref{thm:F29} adds section independence, the descent criterion for
  relations on presentation tuples and the independence of equivariance;
  \Cref{thm:F38} defines the arrow at the level of histories and proves its
  record-level form; \Cref{thm:F24} and \Cref{thm:F31} define the strict layer
  extension and counterfactual legitimacy. In \Cref{thm:F36},
  \Cref{thm:F40}, \Cref{thm:F42} and \Cref{thm:F45} definitions precede the
  proved lines. The Lean theorems for \Fref{F11}, \Fref{F23}, \Fref{F24},
  \Fref{F29} and \Fref{F39} carry each hypothesis only where the paper uses
  it. The summary of \Fref{F14} says that the sources take their domination
  bounds as input, and the catalog-boundary definition is linked to
  \Cref{thm:F52}.
\item \Cref{thm:F5} adds the reverse comparison and its four statuses,
  \Cref{thm:F7} states the bundled form of sufficiency, \Cref{thm:F15a} states
  the breaker-extension clauses, and \Cref{thm:F40} lists its four conditions.
  The Lean theorems for \Fref{F7} and \Fref{F17} take the paper's hypotheses:
  the per-probe family and the factorization and minimality of the common
  quotient. \Cref{thm:F13a} says what its certificate identifies and that
  the formal predicate is opaque, the summaries of \Fref{F9} state its
  hypotheses, the displayed-equivalence convention treats every line without
  \(:\Longleftrightarrow\) as an assertion, and Lean names in the body are
  collected in \Cref{app:formalization}.
\item Interpretive sentences that closed the proofs of
  \Cref{thm:F27,thm:F28,thm:F29,thm:F43} follow them.
\item Surjectivity is assumed only where a statement says so
  (\Cref{sec:apparatus}). \Cref{thm:F5,thm:F25,thm:F48,thm:F49,thm:F50} define
  their predicates before the conclusion, and \Cref{thm:F43} defines its test
  set, cofinality and vanishing. The footnotes to \Cref{thm:F13a} and
  \Cref{prop:F53} say which clauses are checked in Lean and which are taken
  from a source paper, the Lean ledger for \Fref{F14} carries the declared
  involution \(U\), and the Access introduction and summary state what
  \Cref{thm:F12} proves. The coverage table follows the chapter order, and
  three instantiations are retagged or reworded to match their content.
\item \Cref{thm:F50} classifies the vacuum criterion as empty, singleton or
  multiple and states existence and uniqueness of a vacuum;
  \Cref{thm:F25} defines its descent predicates and gives each raw-witness
  clause only its own descent equation; \Cref{thm:F38} assumes surjectivity of
  the source record quotient only. \Cref{def:closure-formation} defines the
  formed closure of a rule set. The Lean theorems for \Fref{F9},
  \Fref{F25}, \Fref{F38} and \Fref{F50} state these forms, with an arbitrary
  status type for \Fref{F9}; the coverage row of \Fref{F53} lists first the
  theorems that take the paper's hypotheses. The formalization appendix
  states that the trust-base axiom can refute its guard, the answer to the
  anchor-row objection says what each anchor row adds, and the catalog
  application of \Cref{thm:F52} states which of its clauses are met.
\item \Cref{thm:F2} states the factorization of the source repair through
  every refinement, \Cref{thm:F45} adds that the local repair carries the
  target, and \Cref{thm:F51} derives its square and kernel inclusions from
  the projection equations alone. \Cref{thm:F14} states its conclusion up to
  null sets, and the readings of \Cref{thm:F20,thm:F23} say what their
  predicates join. The Lean theorems for \Fref{F2}, \Fref{F9}, \Fref{F12},
  \Fref{F36}, \Fref{F43}, \Fref{F44}, \Fref{F45} and \Fref{F51} state every
  displayed clause, \Fref{F43} with index-dependent stage and residual
  types; the coverage row of \Fref{F53} lists only applicable theorems, and
  the formalization appendix states what the opaque guard of \Fref{F13a}
  permits.
\item \Cref{thm:F13a} lists surjectivity of the current quotient among the
  imported data, \Cref{thm:F33} states the family clause as an equivalence
  for the bundled readout, and \Cref{thm:F6,thm:F37} say which setting words
  add no condition. The abstract and introduction count the laws proved in
  full and the laws with direct Lean theorems separately. The Lean
  development instantiates the \Fref{F14} laws for positive semidefinite
  rational matrices, proves the \Fref{F33} family clause, and proves the
  fixed-point form of the \Fref{F53} criterion for the structural defect of
  the retained five-flag encoding. The coverage vocabulary is Theorem, Anchor and Partial
  with alignments faithful, restricted and anchored, and the build
  instructions are replaced by a statement of the trust base.
\item The readings of \Cref{thm:F17,thm:F19,thm:F20} name the case of each
  status, the comparisons with G\"odel's and Noether's theorems and with the
  information paradox state what each classical result derives, the framing
  of \Cref{thm:F37} states that non-joint access is relative to the declared
  class of joint carriers, and the prose around \Cref{thm:F6} states what the
  theorem proves. The Lean theorem for \Fref{F52} allows an arbitrary verdict
  set, and the coverage rows of \Fref{F8} and \Fref{F42} list the theorems
  that carry their hypothesis-free clauses.
\item \Cref{thm:F49} states when a correlation residual is present including
  the admissibility and stability of the route, the \Fref{F21} nonclaim names
  the rules its conclusion uses, and the reading of \Cref{thm:F19} gives its
  obstruction sets in ordinary notation. Surjectivity is assumed only where a
  statement says so, a faithful Lean row may state more than the paper, and
  the build statement names its scope. The abstract counts the laws in one
  split, and the presentation figure describes where each part of a law
  appears.
\item \Cref{thm:F35} states the next-scale descent criterion that its
  linking equation gives, \Cref{thm:F23} characterizes stable probability as
  stability together with failure of descent, \Cref{thm:F40} names the
  induced action separately from the candidate action, and the residual of
  the closure--reality apparatus has a stated codomain. Two instantiations
  whose data do not meet their law's hypotheses are tagged as analogies, and
  the abstract states the scope of the formalized matrix inequality.
\item The Myhill--Nerode, crystallographic, Gr\"obner-basis and general
  relativity instantiations are restated so that they meet the hypotheses of
  their laws; \Cref{thm:F47} keeps only the predicates its conclusion uses;
  gates are described as maps on carriers; the introduction describes the
  eight themes as chapter labels; the abstract counts the laws proved in full
  separately from those with direct Lean theorems; and two displays are set
  within the margins.
\item \Cref{thm:F42} measures irreducibility against a declared class of
  admissible compressions and states what it reduces to when every map is
  admissible, and \Cref{thm:F21} shows that rule~3 cannot be dropped either;
  the Lean development proves both. The residual of the closure--reality
  apparatus is the set of its five structural defect atoms, the hiddenness
  interface defines \(\operatorname{HiddenFromCurrent}\) with its other
  predicates, and a paragraph on the pairing refinement precedes the
  catalog. The readings of \Cref{thm:F24,thm:F46} relate them to
  \Cref{thm:F12,thm:F26}, the similarity instantiation of \Cref{thm:F29}
  states that its listed invariants do not separate similarity classes, and
  the build statement names Lean's standard axioms.
\item \Cref{thm:F23} adds a clause for finite weights, in which descent is
  equivalent to every fiber probability being a point mass, and its
  concentration clause is proved from additivity alone; \Cref{thm:F43}
  characterizes the stable limit quotient for every index;
  \Cref{thm:F45} characterizes minimal domains under strict factorization;
  \Cref{thm:F51} states when the parent carries exactly the children's
  distinctions; and \Cref{thm:F52} shows that each of its four conditions is
  independent of the others. The Lean development proves these clauses. The
  catalog boundary, horizon-bounded support, record formation and diagnosis
  source data are defined by formulas, \Cref{thm:F14} assumes only that the
  traces have infimum zero, and the quantum instantiation of
  \Cref{thm:F19} places decoherence in the dissolution obstruction.
\item \Cref{thm:F27} allows partial steps, so that moves applying only where
  they apply are instances, and \Cref{thm:CER-converse-no-go} shows that both
  closure--reality implications fail for every nontrivial closure operator on
  subsets; the Lean development proves both. The Reed--Solomon instantiation
  places the error vector in the source carrier, the introduction says how the
  eight themes are assigned to chapters, the coverage row of \Fref{F23} names
  the theorem for each clause, the formalization appendix gives each restricted
  or anchored row its own paragraph, and the pairing paragraph lists every law
  that uses the pairing.
\item \Cref{thm:F23} defines its transported ledger and target mixture
  through a ledger transport, a target mixture and the class-level
  continuation, \Cref{thm:F49} defines its correlation record from the joint
  readout and the local access classes, and the duality ledger of
  \Cref{thm:F14} is strongly measurable. The crystallographic, Burnside and
  quantum instantiations state the predicate or functional they use, the
  abstract states the rational-matrix scope of \Fref{F6}, the formalization
  appendix distinguishes the inherited diagonal-support proof from the
  Foundations IV endpoint proof for arbitrary rational blocks, the \Fref{F13a} note names its
  axiom, guard and axiom-free clause, the introduction relates its coverage
  counts to those of the abstract, and two summary-table rows state what their
  theorems prove.
\item The closure--reality anchor defines the retained operational predicate by the vanishing
  of the residual in the body and in the proof, \Cref{thm:F20} says when its
  three stable-fact conditions concern one recorded fact, and
  \Cref{thm:F33} states horizon-boundedness and defines horizon
  sufficiency. The inherited-rules paragraph counts the laws that use them,
  the source of \Fref{F13b} is called the CSL paper, access is stated for any
  observation signature, the Access chapter points new readers to the
  Transport chapter, and three summary-table rows state what their theorems
  prove.
\item The Lehmann--Scheff\'e instantiation uses the likelihood-ratio function,
  whose fibers are those of the minimal sufficient statistic, and the Markov
  instantiation uses the transition kernel of a countable chain. The
  closure--reality chapter states what the no-go theorems show and that
  the retained operational predicate is not compared with bare realisability, the
  formalization appendix states the restriction of the two anchored rows,
  the abstract's coverage sentences are regrouped, the chapter groupings are
  called axes throughout, the covering condition of \Cref{thm:F43} is
  explained, and the predictive-state and residual-coverage statements are
  corrected.
\item \Cref{thm:F44} shows that a phase transition along a path of
  neighbouring formed controls crosses a phase boundary, \Cref{thm:F39} adds
  the fiber-volume instance, and \Cref{thm:F13a} states the exposure,
  collapse and surplus equivalences under the certificate; the Lean
  development proves these clauses, proves the \Fref{F6} endpoint for
  arbitrary rational blocks, and defines closedness in the closure--reality
  anchor as a fixed point of the saturation operator. The \Fref{F13b}
  statement says that surplus pairs are separated by the determining state,
  the proof of \Cref{thm:F11} cites the no-overread rule, the reading path
  names the one cross-chapter use, and each Lean footnote names its row of
  the coverage table.
\item \Cref{thm:F44} adds the clause for continuous paths in a topological
  control space, and \Cref{thm:F42} states its collapse clause for any
  admissible class that isolates \(h\) with every value; the Lean
  development proves both. \(\BirdInt\)-compatibility is defined by listing in
  the law's setup, \Cref{thm:F40} declares the record of explicit breaking,
  \Cref{thm:F15a} assumes the orbit-quotient property only where it is used,
  the statements of \Cref{thm:F17,thm:F19,thm:F36} gloss their predicates,
  the reading paragraph of \Cref{thm:F43} names every conjunct unused by the
  construction, and the summary rows of \Fref{F15a} and \Fref{F39} follow
  their statements.
\item \Cref{prop:F53} adds its clause on the saturation of the retained five-flag encoding,
  and the closure--reality chapter states the general content of its two no-go
  theorems. The notation of \Cref{thm:F6} allows any symmetric \(C^{+}\) and any
  symmetric \(K_{LL}^{\dagger}\) satisfying the projection identity, the reading
  paragraphs of \Cref{thm:F13a,thm:F39,thm:F42} state what their predicates do
  and do not require, the introduction and abstract state which laws import
  results and how Lean covers them, the Access chapter can be read first, and
  the interaction-status convention adds no premise.
\item \Cref{thm:F13b} characterizes determining readouts by factorization
  and splits predictive surplus for an arbitrary readout, \Cref{thm:F29}
  characterizes presentation invariance by a generating relation of the
  kernel, and \Cref{thm:F36} relates the completion of a renormalized target
  obtained by post-processing to that of its raw target; the Lean development
  proves the three clauses. The encoding of the reflective closure used by the
  saturation clause is described with the definition, every
  interaction-status paragraph names its setup data, the instantiations of
  \Fref{F46} state that the classical law is an input, and the reading
  paragraphs of \Cref{thm:F43,thm:F44} state what their conclusions add.
\item \Cref{thm:F42} adds its global form over a base description,
  \Cref{thm:F19} shows that holonomy is coherent when the three class maps are
  the descended ones, \Cref{thm:F14} requires only that the traces approach
  zero from above, the no-overread rule names its acceptance predicate, the
  fiber-volume instance of \Cref{thm:F39} states where its mixing term is
  evaluated, the formalization appendix names the declarations of
  \Fref{F13a}, and every law is noted as \(\BirdInt\)-compatible by
  construction.
\item \Cref{thm:F42} adds the existence of a globally irreducible description
  of minimal cost, \Cref{thm:F47} adds a quantitative clause for
  complementary selector regions, \Cref{thm:F39} requires combine-monotonicity
  only at the pair it uses, so that the fiber-volume instance satisfies the
  theorem, the closure--reality no-go theorems hold for every extensive
  operator, and the Lean development proves these clauses and the full
  descent equivalence of \Cref{thm:F35}. \Cref{thm:F43} states its unique
  compatible readout as the conclusion, the direct no-go theorem characterizes
  failure of its implication, the vocabulary defines a closure space, the
  reading path lists the arguments that are not descent arguments, and the
  Reflexive Nonclosure law states which of its conclusions are premises.
\item The point-mass clause of \Cref{thm:F23} requires weights only on each
  fiber, so distinct classes may share a ledger, and the Lean development proves
  it in that form. \Cref{thm:F6} is stated for arbitrary block matrices with the
  symmetry and projection identities it uses, \Cref{prop:F53} displays the
  fixed-point equivalence it asserts, two instantiations are tagged as
  analogies, and \Cref{thm:F17} writes its defects out.
\item \Cref{thm:F13b} identifies the pairing with the predictive quotient as
  the coarsest determining one, \Cref{thm:F49} shows that a correlation
  residual is exactly the absence of a local explanation, and \Cref{thm:F41}
  shows that the inverse class map and the role-value map are determined; the
  Lean development proves these clauses. The definition of the current and
  predictive quotients states what a quotient and refinement are, the
  interaction-status convention states that it adds no hypothesis, the
  closure--reality proof states what the no-go theorems show, and the coverage
  rows of \Fref{F23} and \Fref{F39} name the declarations of their clauses.
\item The closure--reality no-go theorems characterize the extensive operators
  for which a reduct-stable predicate makes both implications fail, and the Lean
  development proves the characterization. \Cref{thm:F18} states its invariance
  condition in the statement, \Cref{thm:F39} defines its entropy predicates in
  the statement, \Cref{thm:F44} states what its boundary means for continuous
  paths, the settings of \Cref{thm:F9,thm:F11,thm:F41} are stated as intended
  readings, and the nonclaims name the \(\BirdInt\) judgment as background.
\item \Cref{thm:F44} asks along a finite path only that each control have
  neighbours of the next control's phase, which in the topological case means
  that it lies in the closure of that phase class, and \Cref{prop:F53} shows
  that both closure--reality implications hold exactly for the fixed-point
  predicate of the closure, so that the retained operational predicate is the
  unique predicate for which closure equals reality holds for
  \(\operatorname{Sat}_5\); the Lean development
  proves both. The reading of \Cref{thm:F17} is stated in ordinary terms, the
  footnote to \Cref{thm:F14} describes its Lean instance exactly, and the
  Vitali and sheaf-locality instantiations name the readout and the minimal
  domain.
\item The introduction identifies the inherited no-overread rule (NO)
  as a premise in \Fref{F11} and \Fref{F12} and summarizes the coverage
  in two sentences; the direct
  closure--reality no-go names the carrier at which the implication fails;
  clause (d) of \Cref{prop:F53} is stated for every operator; the reading of
  \Cref{thm:F44} explains its neighbourhood relation; \Cref{thm:F43} glosses
  its covering condition in the statement; and the summary rows of
  \Fref{F8}, \Fref{F23} and \Fref{F44} state the results rather than their
  definitions.
\item The Lean development proves the characterization of the phase boundary
  in clause (iv) of \Cref{thm:F44} and the closure form of the step hypothesis
  of clause (iii). \Cref{thm:F42} defines its obstruction sets in the
  statement, \Cref{thm:F27} states its route structure in the statement, the
  reading of \Cref{thm:F49} explains its surjectivity hypothesis, the anomaly
  clause of \Cref{thm:F40} is stated for unbroken symmetries, the interaction-status
  convention says that the paragraphs may be skipped, and the direct
  closure--reality no-go states that its first clause holds by definition.
\item The Lean development proves the two factorizations of the canonical
  pairing in \Cref{thm:F13b}. The reading of \Cref{thm:F15b} states when a
  non-invariant ground state is a spontaneous breaking, \Cref{thm:F26} explains
  its obstructed clause, \Cref{thm:F19} gains a summary in ordinary terms, the
  Access chapter cites the apparatus pairing refinement in place of
  \Cref{thm:F3}, and the reading groups
  of the Status, Records and Coherence chapter are listed.
\item The reading of \Cref{thm:F23} states which data its stability condition
  constrains; \Cref{prop:F53} describes its right-hand side as the defect
  characterization of the closure's fixed points; the pointwise clause of
  \Cref{thm:F42} states what it depends on for large admissible classes; the
  opening of the Self-Reference and Limits chapter states the horizon law as
  proved and locates the self-audit law; the certificate of \Cref{thm:F13a} is
  described as fixing one interface; and the introduction explains the
  F-codes.
\item \Cref{prop:F53} adds clause (e): both closure--reality implications fail
  for the retained saturation \(\operatorname{Sat}_5\) itself, for a
  reduct-stable predicate; \Cref{thm:F37}
  states its factorization form and the product corollary; and the Lean
  development proves these and the empty-moduli clause of \Cref{thm:F26}. The
  reading of \Cref{thm:F40} states what the recorded flag affects, the
  correlation predicate of \Cref{thm:F49} is described as defined from the
  data, and one instantiation of the Self-Reference and Limits chapter is
  relabelled an analogy.
\item The Lean footnotes point to the per-row map of the formalization
  appendix; the dependency hints of the axis figure state the proof-level uses;
  the convention on displayed lines names the lines that unfold a definition;
  \Cref{thm:F28} defines its associativity and identity predicates; the
  comparison statuses of \Cref{thm:F19} are stated with the theorem; the
  reading of \Cref{thm:F23} states which clauses use the continuation data; and
  the coverage row of \Fref{F53} lists the finite instances.
\item \Cref{thm:F4} adds the characterization of the sheaf condition for a
  cover, and \Cref{thm:F14} adds, for \(U=\mathrm{id}\), the equivalence between
  a vanishing ledger and full measure of the fixed locus; the Lean development
  proves both (the second for the abstract ledger under explicit laws). The
  converse closure--reality no-go identifies its closure-instantiation-gap
  hypothesis with a witnessing carrier, \Cref{thm:F23} names the operation and zero of clause (v),
  the global clause of \Cref{thm:F42} drops an unused surjectivity, and the
  reading order states that the Transport chapter can be read first.
\item Clause (e) of \Cref{prop:F53} names its predicate and requires only a
  carrier over the scope; \Cref{thm:F19} defines its persistence predicates in
  the statement; the last line of \Cref{thm:F22} holds for every \(\bar U\); the
  reading paragraphs of \Cref{thm:F15a,thm:F44} name their symbols in order;
  the anchor status describes what the Lean proves for \Fref{F6} and
  \Fref{F53}; the interaction-status convention points to the nonclaims it
  carries; and the reading of \Cref{thm:F49} states when the source map is
  injective.
\item The convention on displayed lines adds examples of definitions and of
  lines combining definitions with descent criteria; \Cref{prop:F53} states at its head which clauses are imported,
  definitional and proved; the interaction-status paragraphs of the Access,
  Stability and Symmetry chapters are titled with their nonclaims; the proof
  of \Cref{thm:F6} calls its identity the projection identity; and the general
  clause of \Cref{thm:F43} names its hypotheses.
\item The special-case readings of the catalog rows (query determinacy, the
  gas and Coulomb models, connected genus, the Riemann and removable-limit
  quotients, the pyramid carrier) state the hypotheses under which they
  instantiate their rows; \Cref{thm:F43} separates its hypotheses from its
  conclusion; and \Cref{thm:F25} states the uniqueness of the descended
  outcome maps, with the full Lean statement cited in its coverage row.
\item The final clause of \Cref{thm:F14} holds whenever \(U\) has no
  \(-1\)-eigenvector, with a one-point ledger showing the condition is needed;
  \Cref{prop:F53} states the uniqueness of the retained predicate in clause (c)
  and defines the flag set and reducts of clause (e); the proof of
  \Cref{thm:CER-direct-no-go} points to where its extensive-operator clause is
  proved; the substrate table describes what \Fref{F53} imports; and the
  Berry-phase, cryptographic-hash and Voronoi special cases and the weight
  hypothesis of \Cref{thm:F23}(vi) are stated exactly.
\item The obstructed-closure part of \Cref{thm:F26}(iii) holds without the
  standing assumption, which then forces the closure space to be empty;
  \Cref{thm:F20} adds that a stable fact transports the event it records;
  the closure--reality chapter says which operator ``retained'' and
  ``current'' refer to; and the meiosis, abstract-data-type, de Bruijn,
  Vitali, martingale and fundamental-group special cases are corrected.
\item Special cases state the data of their laws and the hypotheses they
  need: the elliptic boundary-value problem (bounded Lipschitz domain,
  coefficient and source classes, Neumann compatibility), the holonomy,
  Berry-phase and Preisach cases of \Fref{F3}, integrability of the
  martingale, the minimal domains of \Fref{F45}, the Bayes closure space,
  the Euler-characteristic transport and the path-integral readout.
\item \Cref{thm:F46} characterizes state-dependent readouts by pairs of
  distinct moduli classes; \Cref{thm:F41} adds that the dualities preserve
  quotient fibers; \Cref{thm:F22} states uniqueness of the mediated update;
  \Cref{thm:F37} drops the surjectivity of its access quotients;
  \Cref{thm:F39} states the laws its fiber-volume instance satisfies; the
  terms of clause (b) of \Cref{prop:F53} are defined; and the proofs of
  \Cref{thm:F15a}, \Cref{thm:F17}, \Cref{thm:F31}, \Cref{thm:F42} and
  \Cref{thm:F44} state only what they prove.
\end{itemize}

\paragraph{Lean artifact and coverage.}
The catalog has 49 direct theorem rows, two anchor rows, and one conditional
row. Each direct row has a Lean theorem, and 48 direct rows have faithful Lean
coverage of the displayed result; \Cref{app:formalization} gives the alignment
of every row. Version~1 described all direct rows as faithful and
listed the declarations for \Fref{F13b} and \Fref{F15b} as theorems, although
they were obligation axioms. Version~2 removes both obligation axioms, proves
\Fref{F13b}, \Fref{F14} and \Fref{F15b} without a guarded axiom, and adds Lean
theorems for the clauses added above. The Lean theorems for \Fref{F11},
\Fref{F17} and \Fref{F21} take the paper's hypotheses and rules explicitly
instead of reading the conclusions off record fields, and the closure--reality
no-go theorems are proved over an abstract carrier and instantiated by the
finite model. One conditional law (\Cref{thm:F13a})
and one trust-base axiom remain; the axiom supplies only the four
source-specific conjuncts of \Cref{thm:F13a} and surjectivity of its current
quotient.

\paragraph{Apparatus and instantiations.} The current and predictive quotient maps are written \(\currQ\) and \(\predM\), distinct from their target sets, and defining lines in several laws use \(:\Longleftrightarrow\). Three instantiations whose data are
empirical are tagged as analogies, and the meiosis instantiation counts
crossovers by parity.
\Cref{sec:apparatus} gives a reading guide: the chapters and axes, the status
vocabulary, the acronyms, the meaning of the layer instantiations, a reading
path and a map of results with the section of each result. It adds paragraphs
on the conditional laws, the two rules inherited from earlier papers,
displayed as formulas, the words that name a law's setting, the emptiness
predicates, obstruction sets, reused letters and displayed equivalences. A
small running example is carried through the introduction and the chapter
openings. A paragraph \emph{Reading the statement} precedes each result whose
statement uses predicates that need explanation, and a paragraph \emph{What the
source supplies} follows each conditional law. The chapter titles name their
subject, and the conditional laws are marked in their subsection titles. The
summary rows paraphrase the statements, symbols that named two objects were
renamed, and the result codes link to their results. The interaction status
and nonclaims of each row carry run-in labels. Each layer instantiation is
tagged as a special case or an analogy; seventy-three instantiations were
corrected where they misstated a condition, a formula, a direction, a date or
an attribution, and two references were added. Each law keeps its special cases beside it, and every law with analogies keeps at least one beside it (\Fref{F29}, \Fref{F45} and \Fref{F46} have special cases only); the other analogies are collected in
\Cref{app:further-instantiations}. The substrate map is given once, in
\Cref{tab:substrate-instances}.

\paragraph{Development of the Lean artifact.}
Before Version 1, an axis-level review of the Lean development led to three
strengthened statements. For \Fref{F2}, the canonical target repair is built
from the equivalence generated by identifying \(rF(x)\) and \(rF(x')\)
whenever \(q(x)=q(x')\), matching the minimal-repair universal property of
the statement, in place of a fallback to a one-point target. For \Fref{F10},
the access quotient \(q_I\) is constructed from a lawful observation
signature rather than assumed, and the universal property covers observable
factorization together with the coarsest-quotient property. For \Fref{F53},
the two closure--reality no-go statements are derived from guarded
structural assumptions (reduct stability on the direct side and a closure
instantiation gap on the converse side) rather than assumed as raw
counterexamples.

\paragraph{Closure operators of the closure--reality chapter.}
The current revision distinguishes the retained five-flag saturation
\(\operatorname{Sat}_5\) from AOR's current evidence-preserving inference
reflector. Earlier descriptions identifying those operators are superseded
by \Cref{def:reflective-closure,prop:F53}. The retained saturation theorem
and CER controls remain valid. Current AOR full accounting has its separate
source-witness criterion, and a Lean counterexample establishes that inference
closure alone does not pay the obligation \(5\le2\).

\paragraph{Source scope of F13a and matrix scope of F6.}
\Fref{F13a} remains conditional on an adopted guarded source obligation.
The current PvNP Lean does not construct the required catalogue adapter
or its opaque certificate; earlier descriptions of its four clauses as
proved source imports are superseded by
\Cref{thm:F13a,subsec:narrowing-disclosure}.
For \Fref{F6}, arbitrary symmetric blocks support the algebraic budget
theorem, while the optimization and positive-residual interpretation
requires the positive-semidefinite setting \(C^{+}\succeq0\) of the
Notation paragraph before \Cref{thm:F6}.
   % app:version-history

\clearpage
\bibliographystyle{plainnat}
\bibliography{references}

\end{document}
